% !TeX program = lualatex
\documentclass[11pt,leqno]{article}
\usepackage{amsmath,amssymb}
\usepackage{fontspec}
\usepackage{unicode-math}
\setmainfont{STIXTwoText}[Extension=.otf,UprightFont=*-Regular,ItalicFont=*-Italic,BoldFont=*-Bold,BoldItalicFont=*-BoldItalic]
\setmathfont{STIXTwoMath-Regular.otf}
\usepackage[main=french]{babel}
\usepackage{tikz-cd}
\usetikzlibrary{babel}
\usepackage[a4paper,inner=27mm,outer=23mm,top=25mm,bottom=27mm]{geometry}
\usepackage{xcolor}
\usepackage[hidelinks,unicode]{hyperref}
\setlength{\parindent}{0pt}
\setlength{\parskip}{0.55\baselineskip}
\setlength{\emergencystretch}{3em}
\newcommand{\Page}[1]{\ifvmode\noindent\else\unskip\ \fi{\footnotesize\textcolor{gray}{[p.\,#1]}}\ \ignorespaces}
\title{La conjecture de Weil. II}
\author{Pierre Deligne}
\date{}
\begin{document}
\begin{center}\small\itshape Version corrigée : les 92 coquilles de l'imprimé relevées dans corrections/Deligne\_1980\_La\_conjecture\_de\_Weil\_II\_misprints\_corrected.csv y sont corrigées. L'édition principale est le texte tel qu'imprimé.\end{center}

\maketitle
% ---- 137-144 ----
\Page{137}
\section*{\centering\normalfont LA CONJECTURE DE WEIL.II}

\begin{center}
\textit{par} \textsc{Pierre} DELIGNE
\end{center}

\section*{\centering\normalfont\normalsize SOMMAIRE}

\begingroup
\small
\noindent\textsc{Introduction}\dotfill\quad 138\par
\medskip
\noindent\textsc{Notations et conventions}\dotfill\quad 144\par
\medskip
\noindent\textbf{1. Pureté}\dotfill\quad 146\par
\noindent\hspace*{1em}\makebox[3.3em][l]{(1.1)}Faisceaux \(\ell\)-adiques\dotfill\quad 146\par
\noindent\hspace*{1em}\makebox[3.3em][l]{(1.2)}Poids\dotfill\quad 153\par
\noindent\hspace*{1em}\makebox[3.3em][l]{(1.3)}Poids déterminantiels\dotfill\quad 156\par
\noindent\hspace*{1em}\makebox[3.3em][l]{(1.4)}Cohomologie des courbes et fonctions \(\mathrm{L}\) : rappels\dotfill\quad 162\par
\noindent\hspace*{1em}\makebox[3.3em][l]{(1.5)}Un critère de pureté\dotfill\quad 164\par
\noindent\hspace*{1em}\makebox[3.3em][l]{(1.6)}Autour de Jacobson-Morosov\dotfill\quad 165\par
\noindent\hspace*{1em}\makebox[3.3em][l]{(1.7)}Monodromie locale\dotfill\quad 170\par
\noindent\hspace*{1em}\makebox[3.3em][l]{(1.8)}Monodromie locale des faisceaux purs\dotfill\quad 175\par
\noindent\hspace*{1em}\makebox[3.3em][l]{(1.9)}Monodromie locale modérée des faisceaux mixtes\dotfill\quad 179\par
\noindent\hspace*{1em}\makebox[3.3em][l]{(1.10)}Valeurs absolues non archimédiennes\dotfill\quad 182\par
\noindent\hspace*{1em}\makebox[3.3em][l]{(1.11)}Spécialisation de la monodromie\dotfill\quad 184\par
\medskip
\noindent\textbf{2. La méthode de Hadamard-de la Vallée-Poussin}\dotfill\quad 187\par
\noindent\hspace*{1em}\makebox[3.3em][l]{(2.1)}La méthode\dotfill\quad 187\par
\noindent\hspace*{1em}\makebox[3.3em][l]{(2.2)}Pureté et compacité\dotfill\quad 192\par
\medskip
\noindent\textbf{3. Le théorème fondamental}\dotfill\quad 197\par
\noindent\hspace*{1em}\makebox[3.3em][l]{(3.1)}Un calcul de cycles évanescents\dotfill\quad 197\par
\noindent\hspace*{1em}\makebox[3.3em][l]{(3.2)}Dimension 1\dotfill\quad 200\par
\noindent\hspace*{1em}\makebox[3.3em][l]{(3.3)}Le cas général\dotfill\quad 204\par
\noindent\hspace*{1em}\makebox[3.3em][l]{(3.4)}Application : la structure des faisceaux mixtes\dotfill\quad 207\par
\noindent\hspace*{1em}\makebox[3.3em][l]{(3.5)}Application : théorèmes d’équidistribution\dotfill\quad 210\par
\noindent\hspace*{1em}\makebox[3.3em][l]{(3.6)}Application : le théorème local des cycles invariants\dotfill\quad 212\par
\noindent\hspace*{1em}\makebox[3.3em][l]{(3.7)}Retour à I.8\dotfill\quad 215\par
\medskip
\noindent\textbf{4. Pinceaux de Lefschetz}\dotfill\quad 217\par
\noindent\hspace*{1em}\makebox[3.3em][l]{(4.1)}Le théorème de Lefschetz difficile\dotfill\quad 217\par
\noindent\hspace*{1em}\makebox[3.3em][l]{(4.2)}Rappels sur les pinceaux de Lefschetz\dotfill\quad 219\par
\noindent\hspace*{1em}\makebox[3.3em][l]{(4.3)}Complément à SGA 7 XIX 4\dotfill\quad 221\par
\noindent\hspace*{1em}\makebox[3.3em][l]{(4.4)}La monodromie des pinceaux de Lefschetz\dotfill\quad 227\par
\noindent\hspace*{1em}\makebox[3.3em][l]{(4.5)}Le théorème du pgcd\dotfill\quad 231\par
\endgroup

\Page{138}
\begingroup
\small
\noindent\textbf{5. Application au \(\mathbf{Q}_\ell\)-type d’homotopie}\dotfill\quad 234\par
\noindent\hspace*{1em}\makebox[3.3em][l]{(5.1)}Le \(\mathbf{Z}\{t\}\)-complexe de De Rham, d’après Grothendieck et Miller\dotfill\quad 234\par
\noindent\hspace*{1em}\makebox[3.3em][l]{(5.2)}Le \(\mathbf{Q}_\ell\)-type d’homotopie\dotfill\quad 237\par
\noindent\hspace*{1em}\makebox[3.3em][l]{(5.3)}Graduations par le poids\dotfill\quad 240\par
\medskip
\noindent\textbf{6. Le formalisme des faisceaux mixtes}\dotfill\quad 243\par
\noindent\hspace*{1em}\makebox[3.3em][l]{(6.1)}Stabilité\dotfill\quad 243\par
\noindent\hspace*{1em}\makebox[3.3em][l]{(6.2)}Complexes purs\dotfill\quad 247\par
\endgroup

\section*{Introduction}

Dans [1], cité I par la suite, nous avons démontré la conjecture de Weil donnant la valeur absolue complexe des valeurs propres de Frobenius sur la cohomologie d’une variété projective et lisse définie sur un corps fini. Nous étudions ici la cohomologie à valeur dans un faisceau; il s’agit de passer des propriétés ponctuelles du faisceau à celles de sa cohomologie.

Soient donc \(\mathrm{X}_0\) un schéma de type fini sur \(\mathbf{F}_q\) et \(\mathscr{F}_0\) un \(\overline{\mathbf{Q}}_\ell\)-faisceau sur \(\mathrm{X}_0\). On suppose choisie une clôture algébrique \(\mathbf{F}\) de \(\mathbf{F}_q\), et on indique par la suppression de l’indice 0 l’extension du corps de base de \(\mathbf{F}_q\) à \(\mathbf{F}\) (cf. (0.7)). Pour \(x_0\in|\mathrm{X}_0|\) un point fermé de \(\mathrm{X}_0\) et \(x\in\mathrm{X}(\mathbf{F})\) au-dessus, on dispose de l’automorphisme de Frobenius \(\mathrm{F}_{x_0}^{*}\) de la fibre \(\mathscr{F}_x\) de \(\mathscr{F}\) en \(x\) (I, (1.11) à (1.13), ou (1.18)). On appellera ses valeurs propres les \textit{valeurs propres} de \(\mathrm{F}_{x_0}\) \textit{sur} \(\mathscr{F}_0\). On dit que \(\mathscr{F}_0\) est \textit{ponctuellement pur} de poids \(n\) si, pour tout \(x_0\in|\mathrm{X}_0|\), les valeurs propres de \(\mathrm{F}_{x_0}\) sur \(\mathscr{F}_0\) sont des nombres algébriques dont tous les conjugués complexes sont de valeur absolue \(\mathrm{N}(x_0)^{n/2}\). On dit que \(\mathscr{F}_0\) est \textit{mixte} s’il est extension itérée de faisceaux ponctuellement purs; les poids de ceux-ci sont les poids de \(\mathscr{F}_0\). Notre résultat essentiel est le

\textit{Théorème 1} (3.3.1). — \textit{Soient \(f:\mathrm{X}_0\to\mathrm{S}_0\) un morphisme de schémas de type fini sur \(\mathbf{F}_q\), et \(\mathscr{F}_0\) un faisceau mixte de poids \(\leq n\) sur \(\mathrm{X}_0\). Alors, pour chaque \(i\), le faisceau \(\mathrm{R}^i f_!\mathscr{F}_0\) sur \(\mathrm{S}_0\) est mixte de poids \(\leq n+i\).}

Pour \(\mathrm{S}_0=\operatorname{Spec}(\mathbf{F}_q)\), ce théorème dit que, pour chaque valeur propre \(\alpha\) de Frobenius sur \(\mathrm{H}_c^i(\mathrm{X},\mathscr{F})\), il existe un entier \(m\leq n+i\) (le poids de \(\alpha\)) tel que les conjugués complexes de \(\alpha\) soient tous de valeur absolue \(q^{m/2}\).

La dualité de Poincaré permet parfois de compléter ces majorations par des minorations (3.3.5). Par exemple, si \(\mathrm{X}_0\) est propre et lisse, et que le faisceau \(\mathscr{F}_0\) est lisse (\(=\)constant tordu, dans une autre terminologie) et ponctuellement pur de poids \(n\), alors les valeurs propres de Frobenius sur \(\mathrm{H}^i(\mathrm{X},\mathscr{F})\) sont toutes de poids \(n+i\) ce que nous exprimerons encore en disant que \(\mathrm{H}^i(\mathrm{X},\mathscr{F})\) est pur de poids \(n+i\).

Pour \(\mathscr{F}_0=\mathbf{Q}_\ell\) (de poids 0), on retrouve le résultat principal de I (avec l’hypothèse projectif et lisse remplacée par propre et lisse).

\Page{139}
Une réduction facile, parallèle à la preuve du théorème de finitude pour les \(\mathrm{R}^i f_!\) (cf. SGA 4, XIV, 1) ramène le théorème 1 au théorème suivant, et à une étude locale à l’infini des faisceaux lisses ponctuellement purs sur une courbe (C ci-dessous).

\textit{Théorème 2} (cf. (3.2.3)). — \textit{Soient \(\mathrm{X}_0\) une courbe propre et lisse sur \(\mathbf{F}_q\), \(j:\mathrm{U}_0\to\mathrm{X}_0\) l’inclusion d’un ouvert dense, et \(\mathscr{F}_0\) un faisceau lisse ponctuellement pur de poids \(n\) sur \(\mathrm{U}_0\). Alors, \(\mathrm{H}^i(\mathrm{X},j_*\mathscr{F})\) est pur de poids \(n+i\).}

Voici les grandes lignes de la démonstration.

A) \textit{Nettoyages.} — (i) Si \(u:\mathrm{X}'_0\to\mathrm{X}_0\) est un morphisme fini surjectif de source une courbe propre et lisse \(\mathrm{X}'_0\), et qu’on désigne par \({}'\) le changement de base par \(u\), les \(\mathrm{H}^i(\mathrm{X},j_*\mathscr{F})\) sont des facteurs directs des \(\mathrm{H}^i(\mathrm{X}',j'_*\mathscr{F}')\). Cet argument permet de se ramener au cas où la monodromie locale de \(\mathscr{F}\) en les points de \(\mathrm{X}-\mathrm{U}\) est unipotente.

(ii) Un théorème de dualité assure que \(\mathrm{H}^i(\mathrm{X},j_*\mathscr{F})\) et \(\mathrm{H}^{2-i}(\mathrm{X},j_*(\mathscr{F}^{\vee}))\) sont en dualité parfaite, à valeur dans \(\overline{\mathbf{Q}}_\ell(-1)\). Ceci ramène à vérifier que les conjugués complexes \(\alpha'\) des valeurs propres \(\alpha\) de Frobenius sur \(\mathrm{H}^i(\mathrm{X},j_*\mathscr{F})\) sont de valeur absolue \(|\alpha'|\leq q^{(n+i)/2}\). Le cas difficile est celui du \(\mathrm{H}^1\).

B) \textit{Plongements complexes.} — Soit \(\alpha\in\overline{\mathbf{Q}}_\ell\) une valeur propre de Frobenius sur \(\mathrm{H}^i(\mathrm{X},j_*\mathscr{F})\). Il est commode de reformuler comme suit les estimations à vérifier : pour chaque isomorphisme \(\iota:\overline{\mathbf{Q}}_\ell\to\mathbf{C}\), on a \(|\iota\alpha|\leq q^{(n+i)/2}\).

Dans la démonstration, chaque isomorphisme \(\iota\) sera traité séparément; ceci conduit à introduire les notions de faisceaux ponctuellement \(\iota\)-pur ou \(\iota\)-mixte (il s’agit de ne pas regarder tous les conjugués complexes des valeurs propres \(\alpha\), mais seulement \(\iota\alpha\)). Il est par ailleurs commode de permettre aux poids d’être des nombres réels quelconques. Voir (1.2.6). On prouvera le théorème 2, avec pur remplacé par \(\iota\)-pur. Je renvoie à (1.2.11) le lecteur qui, comme moi, répugne à l’axiome du choix, implicite dans l’usage d’isomorphismes entre \(\overline{\mathbf{Q}}_\ell\) et \(\mathbf{C}\).

C) \textit{Monodromie locale des faisceaux \(\iota\)-purs.} — Posons \(\mathrm{S}_0=\mathrm{X}_0-\mathrm{U}_0\). On commence par montrer que, si \(\mathscr{F}_0\) est lisse et ponctuellement \(\iota\)-pur de poids \(\beta\in\mathbf{R}\), le poids \(w_{\mathrm{N}(x_0)}(\alpha)=2\log_{\mathrm{N}(x_0)}|\iota\alpha|\) d’une valeur propre \(\alpha\) de \(\mathrm{F}_{x_0}\) sur \(j_*\mathscr{F}_0\), pour \(x_0\in\mathrm{S}_0\), est de la forme \(\beta-m\), avec \(m\) entier \(\geq 0\), et on détermine \(m\) en terme de la monodromie locale de \(\mathscr{F}_0\) en \(x_0\) (1.8.4). Plus généralement, on détermine \(w_{\mathrm{N}(x_0)}(\alpha)\) pour \(\alpha\) valeur propre de \(\mathrm{F}_{x_0}\) sur \(\mathscr{F}_0\), au sens (1.10.2). Première étape : montrer que \(w_{\mathrm{N}(x_0)}(\alpha)\leq\beta+2\). On le fait en exploitant la formule de Grothendieck pour la fonction \(\mathrm{Z}(\mathrm{U}_0,\mathscr{F}_0,t)\) : appliquant \(\iota\), on trouve à gauche un produit infini convergeant pour \(w_q(t)<-\beta-2\), à droite une fraction rationnelle de numérateur \(\iota\det(\mathrm{I}-\mathrm{F}t,\mathrm{H}_c^1(\mathrm{U},\mathscr{F}))\), et on utilise le fait que les \((j_*\mathscr{F})_x\) pour \(x\in\mathrm{S}\) contribuent à \(\mathrm{H}_c^1(\mathrm{U},\mathscr{F})\). Deuxième étape : appliquer ce résultat aux puissances tensorielles de \(\mathscr{F}_0\) et de son dual, en tenant compte de la structure connue de la monodromie locale.

\Page{140}
Ceci acquis, il est loisible, et commode, d’étudier \(j_!\mathscr{F}_0\) plutôt que \(j_*\mathscr{F}_0\) : si \(\mathscr{F}_0\) est lisse et ponctuellement \(\iota\)-pur de poids \(\beta\in\mathbf{R}\), il s’agit de montrer que les valeurs propres \(\alpha\) de Frobenius sur \(\mathrm{H}_c^1(\mathrm{U},\mathscr{F})=\mathrm{H}^1(\mathrm{X},j_!\mathscr{F})\) sont de poids \(w_q(\alpha)\leq\beta+1\). On supposera pour simplifier l’écriture que \(\beta=0\). On se ramène à ce cas par torsion (1.2.7).

D) \textit{Passage à \(\mathrm{X}_0\times\mathrm{X}_0\).} — La principale idée géométrique est de passer de \((\mathrm{U}_0,\mathscr{F}_0)\) à \((\mathrm{U}_0\times\mathrm{U}_0,\mathscr{F}_0\boxtimes\mathscr{F}_0)\) (où \(\mathscr{F}_0\boxtimes\mathscr{F}_0=\mathit{pr}_1^*\mathscr{F}_0\otimes\mathit{pr}_2^*\mathscr{F}_0\)), et d’analyser
\[
\mathrm{H}_c^*(\mathrm{U}\times\mathrm{U},\mathscr{F}\boxtimes\mathscr{F})
=
\mathrm{H}_c^*(\mathrm{X}\times\mathrm{X},j_!\mathscr{F}\boxtimes j_!\mathscr{F})
\]
à l’aide d’un pinceau de sections hyperplanes de \(\mathrm{X}_0\times\mathrm{X}_0\), supposé en position générale. Ceci pour un plongement projectif convenable de \(\mathrm{X}_0\times\mathrm{X}_0\). On suppose effectué le nettoyage A(i).

Le but est de prouver que le poids d’une valeur propre \(\alpha\) de Frobenius sur \(\mathrm{H}^2(\mathrm{X}\times\mathrm{X},j_!\mathscr{F}\boxtimes j_!\mathscr{F})\) vérifie \(w_q(\alpha)\leq 3=2+1\). Ceci acquis, si \(\alpha\) est maintenant une valeur propre de Frobenius sur \(\mathrm{H}^1(\mathrm{X},j_!\mathscr{F})\), la formule de Künneth assure que \(\alpha^2\) est valeur propre de Frobenius sur \(\mathrm{H}^2(\mathrm{X}\times\mathrm{X},j_!\mathscr{F}\boxtimes j_!\mathscr{F})\), et on obtient \(w_q(\alpha)\leq 1+\frac{1}{2}\). Plus généralement, supposons le théorème déjà vérifié à \(\delta\) près, au sens suivant : pour toute courbe, et tout faisceau (encore de poids zéro) les valeurs propres de Frobenius sur \(\mathrm{H}^1(\mathrm{X},j_!\mathscr{F})\) sont de poids \(w_q(\alpha)\leq 1+\delta\). Le but est alors de prouver que le poids d’une valeur propre \(\alpha\) de Frobenius sur \(\mathrm{H}^2(\mathrm{X}\times\mathrm{X},j_!\mathscr{F}\boxtimes j_!\mathscr{F})\) vérifie \(w_q(\alpha)\leq 2+\delta\). Procédant comme plus haut, ceci améliore l’estimation originale en \(w_q(\alpha)\leq 1+\delta/2\) pour \(\mathrm{H}^1(\mathrm{X},j_!\mathscr{F})\), et on itère le procédé. Cette itération remplace l’usage de grandes puissances cartésiennes en (I.7.3).

Prenons un pinceau (assez général) de sections hyperplanes de \(\mathrm{X}_0\times\mathrm{X}_0\). Les sections hyperplanes du pinceau sont les fibres d’un morphisme \(f:\mathrm{Y}_0\to\mathbf{P}_0^1\) où \(\mathrm{Y}_0\) se déduit de \(\mathrm{X}_0\times\mathrm{X}_0\) en éclatant un nombre fini de points. Soit \(\mathscr{G}_0\) l’image inverse de \(j_!\mathscr{F}_0\boxtimes j_!\mathscr{F}_0\) sur \(\mathrm{Y}_0\). La suite spectrale de Leray pour \(f\) réduit, pour l’essentiel, l’étude de \(\mathrm{H}^2(\mathrm{X}\times\mathrm{X},j_!\mathscr{F}\boxtimes j_!\mathscr{F})\) à celle de \(\mathrm{H}^1(\mathbf{P}^1,\mathrm{R}^1 f_*\mathscr{G})\).

La généralisation (1.5.1) de la majoration fondamentale (I.3.2) permet de montrer que, sur un certain ouvert \(\mathrm{V}_0\) de \(\mathbf{P}_0^1\) où le faisceau \(\mathrm{R}^1 f_*\mathscr{G}_0\) est lisse, il est extension successive de faisceaux lisses ponctuellement \(\iota\)-purs. Des hypothèses (I.3.2) seule subsiste l’exigence que, sur \(\mathrm{V}_0\), \(\mathrm{R}^1 f_*\mathscr{G}_0\) soit facteur direct d’un faisceau lisse \(\iota\)-réel (un faisceau \(\mathscr{H}_0\) tel que les polynômes \(\iota\det(\mathrm{I}-\mathrm{F}_x t,\mathscr{H}_0)\) soient à coefficients réels). Les résultats du § 2 (voir E) ci-dessous) permettent de le déduire de ce que \(\mathscr{F}_0\) lui-même est facteur direct d’un faisceau \(\iota\)-réel, à savoir \(\mathscr{F}_0\oplus\mathscr{F}_0^\vee\) : parce que \(\mathscr{F}_0\) est de poids 0, son dual joue le rôle d’un complexe conjugué.

Comment calculer les poids de \(\mathrm{R}^1 f_*\mathscr{G}_0|\mathrm{V}_0\) ? Dans (I.3.2), l’hypothèse (ii) assurait que le faisceau considéré n’ait pas de sous-faisceau lisse non trivial — donc soit ponctuellement \(\iota\)-pur d’après (1.5.1) — et l’hypothèse (i) permettait de calculer son poids. Dans l’application (I.6.3) de (I.3.2), ces hypothèses résultaient de la théorie des
\Page{141}
pinceaux de Lefschetz, et plus particulièrement du théorème de conjugaison des cycles évanescents. Ici, \(\mathrm{R}^1 f_*\mathscr{G}_0\) admet des points de ramification de trois types géométriquement distincts, et il faut un autre argument. La théorie des cycles évanescents permet de contrôler le saut de \(\mathrm{R}^1 f_*\mathscr{G}_0\) en un point de \(\mathbf{P}_0^1-\mathrm{V}_0\) à partir d’informations locales sur \((\mathrm{X}_0\times\mathrm{X}_0,j_!\mathscr{F}_0\boxtimes j_!\mathscr{F}_0)\), à savoir, les groupes de cycles évanescents. Un calcul local et les résultats de C, appliqués à \(\mathscr{F}_0\), permettent de déterminer les poids des valeurs propres de Frobenius sur les groupes de cycles évanescents. On trouve que ces poids sont des entiers. Appliquant alors les résultats de C aux quotients successifs de \(\mathrm{R}^1 f_*\mathscr{G}_0|\mathrm{V}_0\), pour une filtration convenable, on parvient à montrer que \(\mathrm{R}^1 f_*\mathscr{G}_0\) est extension successive de faisceaux lisses sur \(\mathrm{V}_0\) et soit \((a)\) lisse sur \(\mathbf{P}^1\) tout entier (de poids non détectés par les cycles évanescents), soit \((b)\) dont la restriction à \(\mathrm{V}_0\) est ponctuellement \(\iota\)-pure de poids entier. Heureusement les faisceaux \((a)\) ne contribuent pas au groupe \(\mathrm{H}^1(\mathbf{P}^1,\mathrm{R}^1 f_*\mathscr{G})\) qui nous intéresse puisque \(\mathbf{P}^1\) est simplement connexe et que \(\mathrm{H}^1(\mathbf{P}^1,\mathbf{Q}_\ell)=0\).

Quels poids entiers \(n\) apparaissent ? Si on suppose déjà connu le théorème, avec une erreur \(\delta<1\), et qu’on l’applique aux fibres de \(f\), il fournit l’estimation \(n\leq 1+\delta\). Puisque \(n\) est entier, on a même \(n\leq 1\). Appliquant le théorème, avec l’erreur \(\delta\), à \(\mathbf{P}^1\), et aux quotients successifs de \(\mathrm{R}^1 f_*\mathscr{G}_0\) pour la filtration introduite précédemment, on trouve que les valeurs propres \(\alpha\) de Frobenius sur \(\mathrm{H}^1(\mathbf{P}^1,\mathrm{R}^1 f_*\mathscr{G})\) sont de poids \(w_q(\alpha)\leq 2+\delta\) — notre but. Au départ toutefois, on a seulement \(\delta=1\), et permettre \(n=2\) ruine l’argument.

E) \textit{Le poids d’un \(\mathrm{H}^1\) est \(<2\).} — Nous prouvons directement que, pour \(\mathscr{F}_0\) lisse et ponctuellement \(\iota\)-pur de poids 0, les valeurs propres \(\alpha\) de Frobenius sur \(\mathrm{H}^1(\mathrm{X},j_!\mathscr{F})\) sont de poids \(w_q(\alpha)<2\). Ceci appliqué à la fibre de \(f\) en un point de \(\mathrm{V}_0\), suffit à faire démarrer la démonstration. Ce résultat permet aussi de séparer les valeurs propres \(\alpha\) de Frobenius sur \(\mathrm{H}^1(\mathrm{X},j_!\mathscr{F})\), de poids \(w_q(\alpha)<2\), de celles sur \(\mathrm{H}^2(\mathrm{X},j_!\mathscr{F})\), de poids 2. Appliqué aux fibres de \(f\), il est utilisé dans la démonstration du fait que \(\mathrm{R}^1 f_*\mathscr{G}_0\) est, sur \(\mathrm{V}_0\), facteur direct d’un faisceau lisse \(\iota\)-réel. La démonstration est inspirée de la démonstration de Mertens du théorème de Hadamard-de la Vallée-Poussin (non nullité de \(\zeta(s)\) pour \(\mathscr{R}s=1\)).

On peut déduire de E) que si \(\mathscr{F}_0\) est un faisceau lisse ponctuellement \(\iota\)-pur sur une courbe lisse \(\mathrm{U}_0\), la fibre \(\mathscr{F}_u\) de \(\mathscr{F}\) en \(u\in|\mathrm{U}|\) est une représentation complètement réductible du groupe de monodromie géométrique (3.4.1) (iii); cf. (3.4.14). Via le dictionnaire heuristique de [2], I, ceci correspond au théorème de semi-simplicité [2], II (4.2.6) en théorie des variations de structures de Hodge. Un cas particulier du théorème montre que si \(\mathrm{X}\) est une variété projective non singulière sur \(k\) algébriquement clos, que \((\mathrm{H}_t)_{t\in\mathbf{P}^1}\) est un pinceau de Lefschetz de sections hyperplanes de \(\mathrm{X}\), avec \(\mathrm{H}_t\) singulier pour \(t\in\mathrm{S}\), et que \(u\in\mathbf{P}^1-\mathrm{S}\), les \(\mathrm{H}^i(\mathrm{H}_u,\mathbf{Q}_\ell)\) sont des représentations complètement réductibles de \(\pi_1(\mathbf{P}^1-\mathrm{S},u)\) : un argument de spécialisation ramène à supposer que \(k=\mathbf{F}\), \(\mathrm{X}\) et le pinceau sont alors définis sur un corps fini, et l’hypothèse de pureté requise résulte de la conjecture de Weil pour les \(\mathrm{H}_t\). Comme il est bien connu, les arguments
\Page{142}
originaux de Lefschetz permettent de déduire de cette complète réductibilité le théorème de Lefschetz difficile sur le cup-produit itéré par la classe de cohomologie d’une section hyperplane (4.1.1). On en donne aussi une variante pour la cohomologie à valeurs dans un faisceau (voire un complexe) (6.2.13).

D’autres énoncés purement géométriques résultent de l’existence d’un formalisme des poids : citons le théorème local des cycles invariants (3.6.1), et les résultats annoncés dans [3] quant à l’existence de graduation par le poids sur des \(\mathbf{Q}_\ell\)-types d’homotopie (5.3.4). On dispose de résultats parallèles en théorie de Hodge, laquelle permet aussi de parler de poids : [12], [9].

Le théorème 1 s’insère dans un formalisme souvent beaucoup plus facile à manier que I. En (3.7), on montre comment il permet de simplifier les preuves des applications données en (I.8). Pour le mode d’emploi de son application aux majorations de sommes trigonométriques, je renvoie à SGA \(4\frac{1}{2}\) [Sommes trigonométriques] (et spécialement aux n\textsuperscript{os} 1, 2, 3, 7).

Les résultats essentiels de l’article concernent les schémas de type fini sur un corps fini, avec deux « mais » : \textit{a)} Ils ont des conséquences géométriques, pour un schéma de type fini sur la clôture algébrique d’un corps fini, et des arguments de spécialisation permettent de passer de là au cas d’un corps de base algébriquement clos quelconque. \textit{b)} Certains résultats valent pour des schémas de type fini sur \(\mathbf{Z}\). Ils se déduisent immédiatement de la démonstration des résultats analogues sur un corps fini, mais non de leur énoncé. Nous les omettrons de la revue ci-dessous, section par section, de l’article.

En (1.1), on rappelle ce qu’est un \(\overline{\mathbf{Q}}_\ell\)-faisceau constructible, avec quelques variantes. On donne aussi un formalisme pour une catégorie dérivée correspondante. En (1.2), on définit les faisceaux ponctuellement purs ou \(\iota\)-purs, mixtes, ..., et on énonce, sous une forme plus précise, une conjecture selon laquelle sur un schéma de type fini sur \(\mathbf{F}_q\), tout \(\overline{\mathbf{Q}}_\ell\)-faisceau constructible est \(\iota\)-mixte. Les n\textsuperscript{os} (1.3) à (1.5) sont consacrés à la généralisation annoncée de (I.3.2). En (1.3), on exploite le théorème de Grothendieck (1.3.8) pour définir les « poids déterminantiels » d’un faisceau lisse \(\mathscr{F}_0\) sur une courbe lisse \(\mathrm{U}_0\). Ils ne dépendent que des puissances extérieures maximales des constituants de \(\mathscr{F}_0\), et contrôlent le \(\mathrm{H}^0\) (et donc le \(\mathrm{H}_c^2\)) de tout faisceau déduit de \(\mathscr{F}_0\) par passage à un espace de tenseurs. La proposition (1.3.13), d’apparence anodine, jouera le rôle joué dans (I.3.2) par la théorie classique des invariants pour le groupe symplectique. Les n\textsuperscript{os} (1.6) à (1.8) sont consacrés au point C) ci-dessus de la démonstration. En (1.6) des préliminaires algébriques : la théorie d’un endomorphisme nilpotent d’un espace vectoriel; en (1.7) des préliminaires géométriques : comment la monodromie locale fournit des endomorphismes nilpotents. La fin, multi-dimensionnelle, de (1.7) ne servira qu’en (1.9), lui-même inutile au reste de l’article.

Le n\textsuperscript{o} (1.9) est une application de (1.8) à la monodromie locale d’un faisceau mixte lisse sur le complément d’un diviseur à croisements normaux, le long duquel il est modérément ramifié. Un résultat analogue, pour les variations de structures de Hodge, vient d’être obtenu par Cattani et Kaplan.

\Page{143}
Dans le n\textsuperscript{o} (1.10), on applique les méthodes de (1.6) à (1.8) à des valeurs absolues non archimédiennes. Le résultat de semi-continuité (1.10.7) a permis de montrer que, pour une intersection complète générale de caractéristique \(p\), le polygone de Newton (défini par les pentes de la cohomologie cristalline) coïncide avec le polygone de Hodge. Que ce problème, de nature \(p\)-adique, n’ait pu être résolu que par des méthodes \(\ell\)-adiques tient à l’absence jusqu’à présent d’une bonne théorie des cycles évanescents en cohomologie \(p\)-adique.

Le n\textsuperscript{o} (1.11) (spécialisation de la monodromie) est technique et sans surprise.

Le § 2 contient la partie E ci-dessus de la démonstration. En (3.5), par des arguments connus, on en déduit un théorème d’équidistribution, dont un cas particulier est l’analogue pour un corps de fonctions de la conjecture de Sato-Tate (sur la distribution des angles de Frobenius, pour une courbe elliptique sur \(\mathrm{K}\)).

Le n\textsuperscript{o} (3.3) contient la preuve du théorème 1, et quelques corollaires. Le n\textsuperscript{o} (3.2) donne celle du théorème 2, et (3.1) donne le calcul de cycles évanescents utilisé (cf. D ci-dessus). En (3.4), le théorème 1 est appliqué à l’étude des extensions successives qui définissent un faisceau mixte. On prouve en particulier qu’un faisceau mixte lisse admet une « filtration par le poids » par des sous-faisceaux lisses, et que la monodromie géométrique d’un faisceau lisse ponctuellement pur sur un schéma normal est semi-simple. En (3.6), le théorème local des cycles invariants; en (3.7), des simplifications aux preuves de (I.8).

Le n\textsuperscript{o} (4.1) contient la preuve du théorème de Lefschetz difficile. Les n\textsuperscript{os} (4.2) et (4.3) apportent quelques compléments à SGA 7. Dans (4.3), l’usage d’un topos \textit{ad hoc} permet de transposer les arguments d’aspect « théorie de Morse » utilisés par Lefschetz. Dans (4.4), on montre que le groupe de monodromie géométrique de la partie évanescente de la cohomologie des sections hyperplanes d’un pinceau de Lefschetz (supposé assez général dans le cas sauvage de caractéristique 2) est Zariski dense dans un groupe orthogonal ou symplectique, ou est un groupe de Weyl. Ce dernier cas est étudié en détail. En (4.5), on en déduit le « théorème du pgcd » utilisé dans [7] pour comparer les cohomologies \(\ell\)-adiques et cristallines des variétés projectives et lisses.

Au § 5, on applique le yoga des poids au \(\mathbf{Q}_\ell\)-type d’homotopie. La définition de ce dernier utilise de façon essentielle la construction par Miller [8] et Grothendieck d’une algèbre différentielle graduée (anti-) commutative attachée à un ensemble simplicial qui permette le calcul de sa cohomologie entière.

Le § 6 enfin développe le formalisme des faisceaux mixtes. On montre, avec des démonstrations inspirées de SGA \(4\frac{1}{2}\) [Th. finitude], que leur catégorie est remarquablement stable. On définit aussi une notion de complexe pur, généralisant celle de faisceau lisse ponctuellement pur sur un schéma lisse (6.2.4), (6.2.5). L’image directe \(\mathrm{R}f_*\) par un morphisme propre transforme complexes purs en complexes purs. On leur généralise le théorème local des cycles invariants et, sous des hypothèses de dimension, le théorème global des cycles invariants et le théorème de Lefschetz difficile.

\Page{144}
\section*{Notations et conventions.}

\textbf{(0.1)} Dans tout cet article, sauf mention expresse du contraire, on fixe un nombre premier \(\ell\), et on ne considère que des espaces algébriques noethériens séparés sur lesquels \(\ell\) est inversible. On les appelle simplement \textit{schémas}.

Le lecteur qui répugne aux espaces algébriques remplacera dans (0.1) « espace algébrique » par « schéma ».

\textbf{(0.2)} On fixe une clôture algébrique \(\overline{\mathbf{Q}}_\ell\) du corps \(\mathbf{Q}_\ell\) des nombres \(\ell\)-adiques. On désigne par \(\iota\) un isomorphisme de corps de \(\overline{\mathbf{Q}}_\ell\) avec \(\mathbf{C}\) (cf. (1.2.11)). Exception : au n\textsuperscript{o} (1.10), on désignera encore par \(\iota\) un isomorphisme de corps de \(\overline{\mathbf{Q}}_\ell\) avec une clôture algébrique \(\overline{\mathbf{Q}}_{\ell'}\) du corps des nombres \(\ell'\)-adiques. Le mot \textit{faisceau} signifiera selon le numéro « faisceau pour la topologie étale », « \(\overline{\mathbf{Q}}_\ell\)-faisceau constructible » ou « faisceau de Weil ». Voir (1.1.5), (1.3.2).

Si \(f:\mathrm{X}\to\mathrm{Y}\) est un morphisme, et \(\mathscr{F}\) un faisceau sur \(\mathrm{Y}\), on écrira souvent, lorsqu’il n’y a pas d’ambiguïté, \(\mathrm{H}^*(\mathrm{X},\mathscr{F})\) pour \(\mathrm{H}^*(\mathrm{X},f^*\mathscr{F})\). De même pour l’hypercohomologie, avec ou sans supports.

\textbf{(0.3)} \textit{Un point géométrique} \(\bar{x}\) de \(\mathrm{X}\) est un morphisme du spectre d’un corps algébriquement clos — noté \(k(\bar{x})\) — dans \(\mathrm{X}\). Il est dit \textit{localisé en} \(x\in\mathrm{X}\) si \(x\) est son image dans \(\mathrm{X}\). Il est dit \textit{algébrique} si \(k(\bar{x})\) est une clôture algébrique de \(k(x)\).

Chaque point géométrique définit un point géométrique algébrique : remplacer \(k(\bar{x})\) par la clôture algébrique de \(k(x)\) dans \(k(\bar{x})\). On utilisera parfois cette construction pour tacitement étendre aux points géométriques une terminologie introduite pour les points géométriques algébriques.

Chaque point \(x\) de \(\mathrm{X}\) à corps résiduel séparablement clos définit un point géométrique \(\bar{x}\) : prendre la clôture parfaite de \(k(x)\). On identifiera souvent \(x\) et \(\bar{x}\).

Si \(f:\mathrm{X}\to\mathrm{Y}\) est un morphisme, tout point géométrique \(\bar{x}\) de \(\mathrm{X}\) définit par composition un point géométrique \(f(\bar{x})\) de \(\mathrm{Y}\) qu’on notera parfois simplement \(\bar{x}\).

\textbf{(0.4)} Sauf mention expresse du contraire, \textit{localement} signifie localement pour la topologie étale.

Si \(x\) est un point d’un schéma \(\mathrm{X}\), nous noterons \(\mathrm{X}_{(x)}\) l’hensélisé de \(\mathrm{X}\) en \(x\) (le spectre de l’hensélisé de l’anneau local de \(\mathrm{X}\) en \(x\)). Si \(\bar{x}\) est un point géométrique (0.3) de \(\mathrm{X}\), nous noterons \(\mathrm{X}_{(\bar{x})}\) le \textit{localisé strict} (\(=\)\textit{hensélisé strict}) de \(\mathrm{X}\) en \(\bar{x}\). Son corps résiduel est la clôture séparable de \(k(x)\) dans \(k(\bar{x})\).

\textbf{(0.5)} Une \textit{courbe lisse} sur un corps \(k\) est un schéma lisse purement de dimension 1 sur \(k\).

Soit \(\mathrm{X}\) un schéma de type fini sur un corps \(k\). Si \(\mathrm{X}\) est connexe, le morphisme structural \(\mathrm{X}\to\operatorname{Spec}(k)\) admet une unique factorisation \(\mathrm{X}\to\operatorname{Spec}(k')\to\operatorname{Spec}(k)\) avec \(k'\)

% ---- 145-152 ----
\Page{145}
extension finie séparable de \(k\) et \(\mathrm{X}/k'\) géométriquement connexe. En général, chaque composante connexe de \(\mathrm{X}\) admet une telle factorisation.

On a là, en particulier, un procédé mécanique pour ramener les propriétés des courbes lisses sur les corps finis à celles des courbes lisses absolument irréductibles sur les corps finis. Nous l’emploierons souvent tacitement, pour passer des unes aux autres.

\textbf{(0.6)} Un \textit{trait} est le spectre d’un anneau de valuation discrète \(\mathrm{V}\). Il est dit hensélien si \(\mathrm{V}\) l’est, et \textit{strictement local} s’il est hensélien et que son corps résiduel est séparablement clos. L’expression « un \textit{trait} \((\mathrm{S},\eta,s)\) » signifie : un trait \(\mathrm{S}\) de point fermé \(s\) et de point générique \(\eta\). L’expression « un trait \((\mathrm{S},\eta,s,\bar\eta,\bar s)\) » signifie : un trait \((\mathrm{S},\eta,s)\), muni d’un point géométrique \(\bar s\) localisé en \(s\), et dont le localisé strict \(\mathrm{S}_{(\bar s)}\) est muni d’un point géométrique générique \(\bar\eta\).

\textbf{(0.7)} On désignera toujours par \(p\) un nombre premier \(\ne\ell\), par \(q\) une puissance de \(p\), par \(\mathbf{F}_q\) un corps fini à \(q\) éléments et par \(\mathbf{F}\) une clôture algébrique de \(\mathbf{F}_q\) (resp. de \(\mathbf{F}_p\), si \(\mathbf{F}_q\) n’a pas été introduit).

Les conventions suivantes seront souvent en vigueur : \(\mathbf{F}_q\) et \(\mathbf{F}\) sont fixés ; les objets sur \(\mathbf{F}_q\) (schémas, ou faisceaux sur des \(\mathbf{F}_q\)-schémas) sont notés avec un indice \({}_0\). La suppression de cet indice indique l’extension des scalaires de \(\mathbf{F}_q\) à \(\mathbf{F}\). Par exemple, si \(\mathscr{F}_0\) est un faisceau sur le \(\mathbf{F}_q\)-schéma \(\mathrm{X}_0\), on note \(\mathscr{F}\) son image inverse sur \(\mathrm{X}=\mathrm{X}_0\otimes_{\mathbf{F}_q}\mathbf{F}\).

Les points géométriques \(\operatorname{Spec}(\overline{\mathrm{K}})\to\mathrm{X}_0\) de \(\mathrm{X}_0\) seront toujours supposés définis par un point géométrique de \(\mathrm{X}\), i.e. \(\overline{\mathrm{K}}\) sera supposé muni d’un \(\mathbf{F}_q\)-morphisme de \(\mathbf{F}\) dans \(\overline{\mathrm{K}}\).

\textbf{(0.8)} Si \(\mathrm{X}\) est un schéma de type fini sur un corps ou sur \(\mathbf{Z}\), on notera \(|\mathrm{X}|\) l’ensemble de ses points fermés. Avec les notations de (0.7), si \(\mathrm{X}_0\) est de type fini sur \(\mathbf{F}_q\), tout point \(\bar x\in|\mathrm{X}|\) définit un point géométrique encore noté \(\bar x:\operatorname{Spec}(\mathbf{F})\to\mathrm{X}_0\). On parlera simplement du point géométrique \(\bar x\in|\mathrm{X}|\).

\textbf{(0.9)} Dans cet article, le point de vue sera plus galoisien que dans I ; ceci nous amènera souvent à écrire \(\mathrm{F}\) là où dans I eut figuré \(\mathrm{F}^{*}\).

\textbf{(0.10)} Une filtration \(\mathrm{F}\) d’un objet \(\mathrm{V}\) d’une catégorie abélienne est dite \textit{finie} s’il existe \(n\) et \(m\) tels que \(\mathrm{F}^n\mathrm{V}=\mathrm{V}\) et \(\mathrm{F}^m=0\).

\textbf{(0.11)} On écrit lim ind et lim proj pour les limites inductives et projectives. Des guillemets indiquent une limite prise dans une catégorie de ind-objets ou de pro-objets ; ainsi « lim ind » \(\mathrm{X}_i\) désigne le ind-objet défini par le système inductif des \(\mathrm{X}_i\) (supposé filtrant).

On écrit \(:=\) pour une égalité dont le second membre est la définition du premier.

\Page{146}
\section*{\normalfont I. — PURETÉ}

\section*{\textbf{(1.1) Faisceaux \(\ell\)-adiques.}}

\textbf{(1.1.1)} Rappelons la définition d’un \(\overline{\mathbf{Q}}_\ell\)-faisceau constructible sur un schéma (0.1) \(\mathrm{X}\) (cf. SGA 5, VI (1.4.2)).

\textit{a)} Soit \(\mathrm{A}\) un anneau noethérien de torsion. Un faisceau de \(\mathrm{A}\)-modules \(\mathscr{F}\) est dit \textit{constructible} (SGA 4, IX, § 2) s’il existe une partition finie de \(\mathrm{X}\) en parties localement fermées \(\mathrm{X}_i\) telle que les restrictions \(\mathscr{F}|\mathrm{X}_i\) soient localement constantes, de fibres des \(\mathrm{A}\)-modules de type fini.

\textit{b)} Soient \(\mathrm{R}\) un anneau local noethérien de caractéristique résiduelle \(\ell\), et \(m\) son idéal maximal. On suppose \(\mathrm{R}\) complet pour la topologie \(m\)-adique. La catégorie des \(\mathrm{R}\)-\textit{faisceaux constructibles} est la catégorie 2-limite projective des catégories de faisceaux de \(\mathrm{R}/\mathrm{I}\)-modules, pour \(\mathrm{I}\) un idéal ouvert de \(\mathrm{R}\). Pour la notion générale de limite projective (resp. inductive) 2-catégorique, le lecteur peut consulter SGA 4, VI (6.10) (resp. (6.3)). Nous l’expliciterons à chaque usage, pour qu’il puisse se dispenser de le faire. Ici : Un \(\mathrm{R}\)-faisceau constructible \(\mathscr{F}\) sur \(\mathrm{X}\) est un système projectif de faisceaux de \(\mathrm{R}\)-modules, indexé par les idéaux ouverts \(\mathrm{I}\) de \(\mathrm{R}\), tel que :

\noindent \(\alpha)\) \(\mathscr{F}_{\mathrm{I}}\) est annulé par \(\mathrm{I}\), et est constructible en tant que faisceau de \(\mathrm{R}/\mathrm{I}\)-modules ;

\noindent \(\beta)\) pour \(\mathrm{I}\supset\mathrm{J}\), le morphisme de transition de \(\mathscr{F}_{\mathrm{J}}\) dans \(\mathscr{F}_{\mathrm{I}}\) induit un isomorphisme de \(\mathscr{F}_{\mathrm{J}}\otimes\mathrm{R}/\mathrm{I}\) avec \(\mathscr{F}_{\mathrm{I}}\).

L’application \(n\mapsto m^n\), des entiers \(>0\) dans l’ensemble des idéaux ouverts de \(\mathrm{R}\) est cofinale. Ceci permet, dans les définitions précédentes, de remplacer la limite projective sur \(\mathrm{I}\) par une limite projective sur \(n\), avec \(\mathrm{R}/\mathrm{I}\) remplacé par \(\mathrm{R}/m^n\).

Le foncteur \(\mathscr{F}\mapsto\text{« lim proj »}\,\mathscr{F}_{\mathrm{J}}\) est pleinement fidèle, de la catégorie des \(\mathrm{R}\)-faisceaux constructibles dans celle des pro-faisceaux de \(\mathrm{R}\)-modules (pro-objets de la catégorie des faisceaux de \(\mathrm{R}\)-modules), et on appellera encore \(\mathrm{R}\)-faisceau constructible tout pro-faisceau dans l’image essentielle. Ceci conduit à abandonner la notation \(\mathscr{F}_{\mathrm{I}}\). On écrira plutôt \(\mathscr{F}\otimes\mathrm{R}/\mathrm{I}\). Justification de cette notation : pour \((\mathscr{F}_{\mathrm{I}})\) comme avant, le système projectif en \(\mathrm{J}\) des \(\mathscr{F}_{\mathrm{J}}\otimes\mathrm{R}/\mathrm{I}\) est essentiellement constant, de valeur \(\mathscr{F}_{\mathrm{I}}\).

On dit que \(\mathscr{F}\) est \textit{lisse} (dans une autre terminologie, \textit{constant tordu}) si les \(\mathscr{F}\otimes\mathrm{R}/\mathrm{I}\) sont localement constants.

\textit{c)} Soient \(\mathrm{E}\) une extension finie de \(\mathbf{Q}_\ell\) et \(\mathrm{R}\) la clôture intégrale de \(\mathbf{Z}_\ell\) dans \(\mathrm{E}\). La catégorie des \(\mathrm{E}\)-\textit{faisceaux constructibles} est le quotient de la catégorie des \(\mathrm{R}\)-faisceaux \Page{147}
constructibles par la sous-catégorie épaisse des faisceaux de torsion. En d’autres termes :

\noindent \(\alpha)\) On dispose d’un foncteur essentiellement surjectif \(\mathscr{F}\to\mathscr{F}\otimes\mathrm{E}\) des \(\mathrm{R}\)-faisceaux constructibles dans les \(\mathrm{E}\)-faisceaux constructibles.

\noindent \(\beta)\) On a \(\operatorname{Hom}(\mathscr{F}\otimes\mathrm{E},\mathscr{G}\otimes\mathrm{E})=\operatorname{Hom}(\mathscr{F},\mathscr{G})\otimes_{\mathrm{R}}\mathrm{E}\).

Un \(\mathrm{E}\)-faisceau constructible est \textit{lisse} s’il est localement de la forme \(\mathscr{F}\otimes_{\mathrm{R}}\mathrm{E}\), avec \(\mathscr{F}\) lisse.

\textit{d)} Pour \(\mathrm{F}\) une extension finie de \(\mathrm{E}\), on dispose d’un foncteur \(\mathscr{F}\mapsto\mathscr{F}\otimes_{\mathrm{E}}\mathrm{F}\) des \(\mathrm{E}\)-faisceaux constructibles dans les \(\mathrm{F}\)-faisceaux constructibles. On a :
\[
\operatorname{Hom}(\mathscr{F}\otimes_{\mathrm{E}}\mathrm{F},\mathscr{G}\otimes_{\mathrm{E}}\mathrm{F})
=\operatorname{Hom}(\mathscr{F},\mathscr{G})\otimes_{\mathrm{E}}\mathrm{F}.
\]
Pour une extension itérée, on a un isomorphisme canonique :
\[
(\mathscr{F}\otimes_{\mathrm{E}}\mathrm{F})\otimes_{\mathrm{F}}\mathrm{G}
=\mathscr{F}\otimes_{\mathrm{E}}\mathrm{G}.
\tag*{(1)}
\]

La catégorie des \(\overline{\mathbf{Q}}_\ell\)-\textit{faisceaux constructibles} est la catégorie \(\overline{\mathbf{Q}}_\ell\)-linéaire 2-limite inductive des catégories des \(\mathrm{E}\)-faisceaux constructibles, pour \(\mathrm{E}\subset\overline{\mathbf{Q}}_\ell\) une extension finie de plus en plus grande de \(\mathbf{Q}_\ell\). En d’autres termes :

\noindent \(\alpha)\) Pour \(\mathrm{E}\subset\overline{\mathbf{Q}}_\ell\) une extension finie de \(\mathbf{Q}_\ell\), on dispose d’un foncteur \(\mathrm{E}\)-linéaire \(\mathscr{F}\mapsto\mathscr{F}\otimes_{\mathrm{E}}\overline{\mathbf{Q}}_\ell\) des \(\mathrm{E}\)-faisceaux constructibles dans les \(\overline{\mathbf{Q}}_\ell\)-faisceaux constructibles. Il induit des isomorphismes :
\[
\operatorname{Hom}(\mathscr{F},\mathscr{G})\otimes_{\mathrm{E}}\overline{\mathbf{Q}}_\ell
\overset{\sim}{\to}
\operatorname{Hom}(\mathscr{F}\otimes\overline{\mathbf{Q}}_\ell,\mathscr{G}\otimes\overline{\mathbf{Q}}_\ell).
\]

\noindent \(\beta)\) On a des isomorphismes canoniques \((\mathscr{F}\otimes_{\mathrm{E}}\mathrm{F})\otimes_{\mathrm{F}}\overline{\mathbf{Q}}_\ell=\mathscr{F}\otimes_{\mathrm{E}}\overline{\mathbf{Q}}_\ell\), compatibles à (1).

\noindent \(\gamma)\) Tout \(\overline{\mathbf{Q}}_\ell\)-faisceau constructible est de la forme \(\mathscr{F}\otimes_{\mathrm{E}}\overline{\mathbf{Q}}_\ell\), pour \(\mathrm{E}\) et \(\mathscr{F}\) convenables.

Un \(\overline{\mathbf{Q}}_\ell\)-faisceau constructible est \textit{lisse} s’il est localement de la forme \(\mathscr{F}\otimes_{\mathrm{E}}\overline{\mathbf{Q}}_\ell\), avec \(\mathscr{F}\) lisse.

Pour que ces définitions soient utilisables, il faut disposer d’un formalisme du type habituel pour ces faisceaux. Un tel formalisme est partiellement construit dans SGA 5, V, VI. Dans ces exposés, Jouanolou montre que les \(\mathbf{Z}_\ell\)-faisceaux constructibles sur \(\mathrm{X}\) forment une catégorie abélienne noethérienne (VI (1.1.3)), donne la relation entre \(\mathbf{Z}_\ell\)-faisceaux constructibles lisses et représentations du groupe fondamental (VI (1.2.5)), montre que les \(\mathbf{Z}_\ell\)-faisceaux constructibles sont lisses sur les strates d’une stratification de \(\mathrm{X}\) (VI (1.2.6)), définit produits tensoriels et images réciproques, et définit et étudie les foncteurs images directes à support propre \(\mathrm{R}^i f_!\) (pour \(f:\mathrm{X}\to\mathrm{Y}\) un morphisme de type fini de schémas noethériens) (VI (2.2)). Grâce à SGA \(4\frac{1}{2}\) [finitude], la théorie qu’il développe s’applique aussi bien aux foncteurs \(\mathrm{R}^i f_*\) et \(\mathrm{R}^i f^!\), lorsque \(f\) est un morphisme de type fini entre schémas de type fini sur un schéma régulier excellent de dimension \(\leq1\). Elle s’applique de même aux foncteurs « faisceaux de cycles évanescents » \(\mathrm{R}^i\Psi\) et \(\mathrm{R}^i\Phi\).

\Page{148}
\textbf{(1.1.2)} Au § 6, nous travaillerons de façon essentielle dans une catégorie dérivée d’une catégorie de faisceaux \(\ell\)-adiques. Il s’agit d’une théorie qui n’est pas au point. Nous nous réfugierons dans le cadre suivant.

Soient \(\mathrm{R}\) l’anneau des entiers d’une extension finie \(\mathrm{E}_\lambda\) de \(\mathbf{Q}_\ell\), et \(m\) son idéal maximal. Si \(\mathrm{X}\) est un schéma, on note \(\mathrm{D}^-(\mathrm{X},\mathrm{R})\) la catégorie 2-limite projective sur \(n\).
\[
\mathrm{D}^-(\mathrm{X},\mathrm{R})
:=2\text{-}\operatorname{lim\,proj}\mathrm{D}^-(\mathrm{X},\mathrm{R}/m^n).
\]
Un objet \(\mathrm{K}\) de \(\mathrm{D}^-(\mathrm{X},\mathrm{R})\) s’identifie à un système projectif défini par des complexes \(\mathrm{K}_n\in\operatorname{Ob}\mathrm{D}^-(\mathrm{X},\mathrm{R}/m^n)\) et par des isomorphismes :
\[
\mathrm{K}_{n+1}\overset{\mathrm{L}}{\otimes}_{\mathrm{R}/m^{n+1}}\mathrm{R}/m^n
\overset{\sim}{\to}\mathrm{K}_n
\quad
(\text{dans }\mathrm{D}^-(\mathrm{X},\mathrm{R}/m^n)).
\]
Le complexe \(\mathrm{K}_n\) se notera plutôt \(\mathrm{K}\overset{\mathrm{L}}{\otimes}\mathrm{R}/m^n\).

On s’intéressera surtout à la sous-catégorie \(\mathrm{D}_c^b(\mathrm{X},\mathrm{R})\) formée des \(\mathrm{K}\in\operatorname{Ob}\mathrm{D}^-(\mathrm{X},\mathrm{R})\) tels que \(\mathrm{K}\overset{\mathrm{L}}{\otimes}\mathrm{R}/m\) soit borné à cohomologie constructible. Pour tout \(n\), \(\mathrm{K}\overset{\mathrm{L}}{\otimes}\mathrm{R}/m^n\) est alors dans \(\mathrm{D}_{\mathit{ctf}}^b(\mathrm{X},\mathrm{R}/m^n)\) (complexes bornés, de Tor-dimension finie, dont les faisceaux de cohomologie sont constructibles ; cf. \textit{a)} ci-dessous).

Les remarques suivantes justifient l’usage que nous ferons de \(\mathrm{D}_c^b(\mathrm{X},\mathrm{R})\).

\textit{a)} Si \(\mathrm{K}\) est dans \(\mathrm{D}_c^b(\mathrm{X},\mathrm{R})\), les \(\mathscr{H}^i\mathrm{K}:=\text{« lim proj »}\,\mathscr{H}^i(\mathrm{K}\overset{\mathrm{L}}{\otimes}\mathrm{R}/m^n)\) sont des \(\mathrm{R}\)-faisceaux constructibles. Découpant \(\mathrm{X}\), on se ramène, pour le prouver, au cas où les \(\mathscr{H}^i(\mathrm{K}\overset{\mathrm{L}}{\otimes}\mathrm{R}/m)\) sont localement constants. Tout complexe à composantes plates, représentant \(\mathrm{K}\overset{\mathrm{L}}{\otimes}\mathrm{R}/m^n\), admet une filtration (la filtration \(m\)-adique) de quotients successifs quasi-isomorphes à \(\mathrm{K}\overset{\mathrm{L}}{\otimes}\mathrm{R}/m\). Les \(\mathscr{H}^i(\mathrm{K}\overset{\mathrm{L}}{\otimes}\mathrm{R}/m^n)\) sont donc localement constants, et il suffit d’étudier leur fibre en un point. Dans le cas ponctuel enfin, on applique l’argument de SGA 5, XV, p. 32.

\textit{b)} Réciproquement, si \(\mathscr{F}\) est un \(\mathrm{R}\)-faisceau constructible sans torsion, le système projectif des complexes réduits à \(\mathscr{F}\otimes\mathrm{R}/m^n\) en degré 0 est dans \(\mathrm{D}_c^b(\mathrm{X},\mathrm{R})\). Nous n’aurons pas besoin de considérer le cas où \(\mathscr{F}\) a de la torsion ; il faudrait remplacer \(\mathscr{F}\otimes\mathrm{R}/m^n\) par \(\tau_{\geq-1}(\mathscr{F}\otimes\mathrm{R}/m^k)\overset{\mathrm{L}}{\otimes}_{\mathrm{R}/m^k}\mathrm{R}/m^n\) (\(k\) assez grand pour que \(m^{k-n}\) tue la torsion de \(\mathscr{F}\) ; le complexe obtenu est alors indépendant de \(k\)).

\textit{c)} Sous des hypothèses de finitude convenable, par exemple lorsqu’on travaille dans la catégorie des schémas de type fini sur \(\mathrm{S}\) avec \(\mathrm{S}\) régulier de dimension \(\leq1\), les catégories \(\mathrm{D}_{\mathit{ctf}}^b(\mathrm{X},\mathrm{R}/m^n)\) sont stables par les quatre opérations \(\mathrm{R}f_*\), \(f^*\), \(\mathrm{R}f_!\), \(\mathrm{R}f^!\), ainsi que par les opérations internes \(\overset{\mathrm{L}}{\otimes}\) et \(\mathrm{R}\operatorname{Hom}\). De plus, toutes ces opérations commutent à la réduction mod \(m^n\). Elles s’étendent donc trivialement à \(\mathrm{D}_{\mathit{ctf}}^b(\mathrm{X},\mathrm{R})\). De même pour les foncteurs \(\mathrm{R}\Psi\) et \(\mathrm{R}\Phi\) de la théorie des cycles évanescents. Voir SGA \(4\frac{1}{2}\) [Th. finitude].

\Page{149}
\textit{Exemple.} — Soient \(\mathrm{X}\) de type fini sur \(\mathrm{S}\), comme ci-dessus, et \(\mathscr{F}\) et \(\mathscr{G}\) deux \(\mathrm{R}\)-faisceaux constructibles sans torsion sur \(\mathrm{X}\). Les faisceaux \(\mathscr{E}xt^i_{\mathrm{R}/m^n}(\mathscr{F}\otimes\mathrm{R}/m^n,\mathscr{G}\otimes\mathrm{R}/m^n)\) forment alors un système projectif, et on pose :
\[
\mathscr{E}xt^i(\mathscr{F},\mathscr{G})
:=\mathscr{H}^i\mathrm{R}\,\mathscr{H}om(\mathscr{F},\mathscr{G})
=\text{« lim proj »}\,\mathscr{E}xt^i_{\mathrm{R}/m^n}
(\mathscr{F}\otimes\mathrm{R}/m^n,\mathscr{G}\otimes\mathrm{R}/m^n).
\]
Cet \(\mathscr{E}xt^i\) est à nouveau un \(\mathrm{R}\)-faisceau constructible. Cette construction passe au quotient pour définir des \(\mathscr{E}xt^i\) de \(\mathrm{E}_\lambda\)-faisceaux constructibles, et s’étend par passage à la limite aux \(\overline{\mathbf{Q}}_\ell\)-faisceaux constructibles.

\textit{d)} Si, quels que soient \(\mathrm{K}\) et \(\mathrm{L}\) dans \(\mathrm{D}_c^b(\mathrm{X},\mathrm{R})\), les groupes \(\operatorname{Hom}(\mathrm{K}\overset{\mathrm{L}}{\otimes}\mathrm{R}/m^n,\mathrm{L}\overset{\mathrm{L}}{\otimes}\mathrm{R}/m^n)\) sont finis (tel est le cas pour \(\mathrm{X}\) de type fini sur un corps fini ou algébriquement clos), alors la catégorie \(\mathrm{D}_c^b(\mathrm{X},\mathrm{R})\) est triangulée, de triangles distingués les triangles dont la réduction modulo \(m^n\) est distinguée pour tout \(n\). En général, la considération de triangles est avantageusement remplacée par celle de complexes filtrés : on définit \(\mathrm{DF}_{\mathit{ctf}}^b(\mathrm{X},\mathrm{R}/m^n)\) comme la catégorie dérivée de la catégorie des complexes bornés filtrés (de filtration finie), dont le gradué est de tor-dimension finie (cf. L. Illusie, Complexe cotangent et déformations, I, LN 239 (Springer), chap. V, n\textsuperscript{os} 1-3). On pose :
\[
\mathrm{DF}_c^b(\mathrm{X},\mathrm{R})
:=2\text{-}\operatorname{lim\,proj}\mathrm{DF}_{\mathit{ctf}}^b(\mathrm{X},\mathrm{R}/m^n),
\]
et chaque fois qu’on aimerait considérer un triangle dans \(\mathrm{D}_c^b(\mathrm{X},\mathrm{R})\), on peut faire mieux : construire un objet de \(\mathrm{DF}_c^b(\mathrm{X},\mathrm{R})\), dont la filtration n’a que deux crans, les sommets du triangle cherché étant les deux composantes du gradué, et le complexe sous-jacent (oubli de la filtration).

Par abus de langage, nous appellerons parfois \textit{triangles} de tels objets de \(\mathrm{DF}_c^b(\mathrm{X},\mathrm{R})\). Nous les écrirons comme une suite exacte courte.

\textit{e)} La définition des foncteurs de troncature \(\tau_{\leq n}\) est délicate. Si \(\mathrm{K}\) est un complexe de faisceaux, le complexe \(\tau_{\leq n}\mathrm{K}\) est le complexe :
\[
\ldots\mathrm{K}^i\ldots\to\mathrm{K}^{n-1}\to\operatorname{Ker}d\to0\ldots
\]
Ce foncteur transforme quasi-isomorphismes en quasi-isomorphismes, donc passe à la catégorie dérivée. Mais il ne respecte pas la Tor-dimension finie. On contourne comme suit la difficulté. Si \(\mathrm{K}\) est dans \(\mathrm{D}_c^b(\mathrm{X},\mathrm{R})\), et que le complexe de faisceaux \(\mathrm{K}_k\) est un représentant de \(\mathrm{K}\overset{\mathrm{L}}{\otimes}\mathrm{R}/m^k\), soit \(\tau'_{\leq n}\mathrm{K}_k\) le sous-complexe de \(\tau_{\leq n}\mathrm{K}_k\) obtenu en remplaçant \(\operatorname{Ker}d\) par le sous-faisceau des cycles dont l’image dans \(\mathscr{H}^n(\mathrm{K}\overset{\mathrm{L}}{\otimes}\mathrm{R}/m^k)\) est dans \(\operatorname{Im}(\mathscr{H}^n\mathrm{K}\to\mathscr{H}^n(\mathrm{K}\overset{\mathrm{L}}{\otimes}\mathrm{R}/m^k))\). Le complexe \(\tau'_{\leq n}\mathrm{K}_k\) est dans \(\mathrm{D}_{\mathit{ctf}}^b(\mathrm{X},\mathrm{R}/m^k)\), et l’on pose :
\[
\tau_{\leq n}\mathrm{K}
:=\text{« lim proj »}\,\tau'_{\leq n}(\mathrm{K}\overset{\mathrm{L}}{\otimes}\mathrm{R}/m^k).
\]

Cette construction repose de façon essentielle sur le fait que \(\mathrm{R}\) est régulier de dimension 1.

\Page{150}
\textit{Variante} \textbf{(1.1.3)}. — On note \(\mathrm{D}_c^b(\mathrm{X},\mathrm{E}_\lambda)\) la catégorie déduite de \(\mathrm{D}_c^b(\mathrm{X},\mathrm{R})\) en étendant les scalaires de \(\mathrm{R}\) à \(\mathrm{E}_\lambda\) : on dispose d’un foncteur essentiellement surjectif \(\otimes_{\mathrm{R}}\mathrm{E}_\lambda:\mathrm{D}_c^b(\mathrm{X},\mathrm{R})\to\mathrm{D}_c^b(\mathrm{X},\mathrm{E}_\lambda)\), et :
\[
\operatorname{Hom}(\mathrm{K},\mathrm{L})\otimes_{\mathrm{R}}\mathrm{E}_\lambda
\overset{\sim}{\to}
\operatorname{Hom}(\mathrm{K}\otimes\mathrm{E}_\lambda,\mathrm{L}\otimes\mathrm{E}_\lambda).
\]
On pose ensuite :
\[
\mathrm{D}_c^b(\mathrm{X},\overline{\mathbf{Q}}_\ell)
:=2\text{-}\operatorname{lim\,ind}\mathrm{D}_c^b(\mathrm{X},\mathrm{E}_\lambda)
\]
(limite sur les extensions \(\mathrm{E}_\lambda\subset\overline{\mathbf{Q}}_\ell\) de \(\mathbf{Q}_\ell\)).

\textbf{(1.1.4)} Dans la suite de cet article, nous utiliserons librement pour les \(\mathbf{Z}_\ell\)-, \(\mathbf{Q}_\ell\)- ou \(\overline{\mathbf{Q}}_\ell\)-faisceaux les théorèmes connus pour les faisceaux constructibles de \(\mathrm{R}/m^n\)-modules. Les exposés SGA 5 V, VI, XV (pour la formule des traces, cf. aussi SGA \(4\frac{1}{2}\) [Rapport]) et les remarques ci-dessus justifient nos arguments. Là où une difficulté apparaîtrait, elle sera signalée.

\textbf{(1.1.5)} \textit{Dans ce numéro et le suivant, nous dirons simplement} faisceau \textit{pour} \(\overline{\mathbf{Q}}_\ell\)-faisceau constructible. Cette convention sera relayée par (1.3.2).

\textbf{(1.1.6)} Une \textit{représentation \(\ell\)-adique} d’un groupe profini \(\pi\) (ou, plus généralement, d’une extension d’un groupe discret de type fini par un groupe profini) sur un \(\overline{\mathbf{Q}}_\ell\)-espace vectoriel \(\mathrm{V}\) est un homomorphisme \(\sigma:\pi\to\operatorname{GL}(\mathrm{V})\) tel qu’il existe une extension finie \(\mathrm{E}\) de \(\mathbf{Q}_\ell\) dans \(\overline{\mathbf{Q}}_\ell\), et une \(\mathrm{E}\)-structure \(\mathrm{V}_{\mathrm{E}}\) sur \(\mathrm{V}\), telles que \(\sigma\) se factorise par un homomorphisme continu de \(\pi\) dans \(\operatorname{GL}(\mathrm{V}_{\mathrm{E}})\). On définit de même les \textit{représentations} de \(\pi\) dans un groupe algébrique sur \(\overline{\mathbf{Q}}_\ell\).

Pour \(\mathrm{X}\) connexe, pointé par un point géométrique \(\bar x\), on a la relation habituelle entre faisceaux lisses et représentations continues du groupe fondamental \(\pi_1(\mathrm{X},\bar x)\) (une équivalence de catégories \(\mathscr{F}\mapsto\text{le }\pi_1\text{-module }\mathscr{F}_{\bar x}\)). Nous transporterons aux faisceaux la terminologie familière pour les représentations. Que \(\mathrm{X}\) soit connexe ou non, nous dirons encore qu’un faisceau lisse \(\mathscr{F}\) sur \(\mathrm{X}\) est \textit{simple} ou \textit{irréductible} s’il est non nul et n’a pas de sous-faisceau lisse non trivial, et \textit{semi-simple} s’il est somme de faisceaux simples. Une \textit{suite de Jordan-Hölder} de \(\mathscr{F}\) est une filtration finie de \(\mathscr{F}\) par des sous-faisceaux lisses, les \(\operatorname{Gr}_{\mathrm{F}}^i(\mathscr{F})\) étant irréductibles, ou nuls. Les \textit{constituants} de \(\mathscr{F}\) sont les \(\operatorname{Gr}_{\mathrm{F}}^i(\mathscr{F})\) non nuls pour une suite de Jordan-Hölder de \(\mathscr{F}\) et le \textit{semi-simplifié} de \(\mathscr{F}\) est la somme directe de ses constituants. Constituants, et semi-simplifié, sont définis à isomorphisme (non unique) près.

\textbf{(1.1.7)} Soient \(k\) un corps fini à \(q\) éléments, et \(\bar k\) une clôture algébrique de \(k\). La \textit{substitution de Frobenius} \(\varphi\in\operatorname{Gal}(\bar k/k)\) est l’automorphisme \(x\mapsto x^q\) de \(\bar k\). Le \textit{Frobenius géométrique} \(\mathrm{F}\in\operatorname{Gal}(\bar k/k)\) est l’inverse de \(\varphi\). Le \textit{groupe de Weil} \(\mathrm{W}(\bar k/k)\) est le sous-groupe de \(\operatorname{Gal}(\bar k/k)\) formé des puissances entières de \(\mathrm{F}\). Il est isomorphe à \(\mathbf{Z}\), et \(\operatorname{Gal}(\bar k/k)\) est \Page{151}
son complété profini. On identifiera souvent \(\mathrm{W}(\bar k/k)\) à \(\mathbf{Z}\), et \(\operatorname{Gal}(\bar k/k)\) à \(\widehat{\mathbf{Z}}\), par l’isomorphisme qui envoie \(\mathrm{F}\) sur 1.

Spécialisons (1.1.6) à \(\mathrm{X}=\operatorname{Spec}(k)\). Si \(\mathscr{F}\) est un faisceau sur \(\operatorname{Spec}(k)\), la fibre \(\mathscr{F}_{\bar k}\) de \(\mathscr{F}\) en le point géométrique \(\operatorname{Spec}(\bar k)\) est un espace vectoriel sur \(\overline{\mathbf{Q}}_\ell\), sur lequel \(\operatorname{Gal}(\bar k/k)\) agit par transport de structure. Cette action est continue : l’application \(n\to\mathrm{F}^n\) se prolonge par continuité en un homomorphisme de \(\widehat{\mathbf{Z}}\) dans \(\operatorname{GL}(\mathscr{F}_{\bar x})\). En d’autres termes, les valeurs propres de \(\mathrm{F}\) sont des unités \(\ell\)-adiques. Le foncteur \(\mathscr{F}\to\mathscr{F}_{\bar x}\) est une équivalence de la catégorie des faisceaux sur \(\operatorname{Spec}(k)\) avec celle des espaces vectoriels sur \(\overline{\mathbf{Q}}_\ell\), munis d’un automorphisme \(\mathrm{F}\) de valeurs propres des unités \(\ell\)-adiques.

\textbf{(1.1.8)} Soient \(\mathrm{X}\) un schéma de type fini sur \(\mathbf{Z}\), \(\mathscr{F}\) un faisceau sur \(\mathrm{X}\), \(y\) un point de \(\mathrm{X}\) à valeur dans un corps fini \(k\) : \(y:\operatorname{Spec}(k)\to\mathrm{X}\) et \(\bar k\) une clôture algébrique de \(k\). L’image inverse de \(\mathscr{F}\) sur \(\operatorname{Spec}(k)\) est justifiable de (1.1.7) : \(\operatorname{Gal}(\bar k/k)\) agit sur \(\mathscr{F}_{\bar k}\). Par abus de notation, nous écrirons \(\det(1-\mathrm{F}_y t,\mathscr{F})\) pour \(\det(1-\mathrm{F}t,\mathscr{F}_{\bar k})\). Nous parlerons de même des valeurs propres de \(\mathrm{F}_y\) agissant sur \(\mathscr{F}\ldots\)

Supposons \(\mathrm{X}\) connexe, muni d’un point géométrique \(\bar x\). On dispose alors d’une classe de conjugaison d’applications \(\operatorname{Gal}(\bar k/k)=\pi_1(\operatorname{Spec}(k),\operatorname{Spec}(\bar k))\to\pi_1(\mathrm{X},\bar x)\). On note \(\mathrm{F}_y\) l’image du Frobenius géométrique par une quelconque de ces applications. Pour \(\mathscr{F}\) lisse défini par une représentation \(\mathrm{V}\) de \(\pi_1(\mathrm{X},\bar x)\), on a \(\det(1-\mathrm{F}_y t,\mathscr{F})=\det(1-\mathrm{F}_y t,\mathrm{V})\), ce qui justifie la notation.

Un cas particulier : \(y\) est un point fermé de \(\mathrm{X}\), et on considère le morphisme d’inclusion de \(\operatorname{Spec}(k(y))\) dans \(\mathrm{X}\). On peut souvent se borner à considérer les Frobenius ainsi obtenus. Si \(z\) est un point de \(\mathrm{X}\) à valeur dans un corps fini \(k\), d’image \(y\) dans \(\mathrm{X}\), et que \([k:k(y)]=n\), le Frobenius relatif à \(z\) est la puissance \(n\)-ième du Frobenius relatif à \(y\).

\textbf{(1.1.9)} Lorsqu’on étudie les schémas de type fini sur \(\mathbf{F}_p\), il est souvent commode de remplacer les groupes de Galois (resp. faisceaux) par les groupes de Weil (resp. faisceaux de Weil) définis ci-dessous. Le lecteur prendra garde que la définition des faisceaux de Weil est relative au choix d’une clôture algébrique \(\mathbf{F}\) de \(\mathbf{F}_p\).

\textit{Définition} \textbf{(1.1.10)}. — \textit{Avec les notations} (0.7) :

(i) \textit{Un \textup{faisceau de Weil} \(\mathscr{F}_0\) sur \(\mathrm{X}_0\) consiste en un faisceau \(\mathscr{F}\) sur \(\mathrm{X}\), muni d’une action de \(\mathrm{W}(\mathbf{F}/\mathbf{F}_q)\) \textup{(1.1.7)} sur \((\mathrm{X},\mathscr{F})\) qui induise sur \(\mathrm{X}=\mathrm{X}_0\otimes_{\mathbf{F}_q}\mathbf{F}\) l’action déduite de l’action de \(\mathrm{W}(\mathbf{F}/\mathbf{F}_q)\) sur \(\mathbf{F}\).}

(ii) \textit{Soit \(\bar x\) un point géométrique de \(\mathrm{X}\). Le \textup{groupe de Weil} \(\mathrm{W}(\mathrm{X}_0,\bar x)\) est l’image inverse de \(\mathrm{W}(\mathbf{F}/\mathbf{F}_q)\) par l’application naturelle de \(\pi_1(\mathrm{X}_0,\bar x)\) dans \(\pi_1(\operatorname{Spec}(\mathbf{F}_q),\bar x)=\operatorname{Gal}(\mathbf{F}/\mathbf{F}_q)\).}

Conformément aux définitions générales, un automorphisme de \((\mathrm{X},\mathscr{F})\) est un couple \((f,g)\) formé d’un automorphisme \(f\) de \(\mathrm{X}\) et d’un isomorphisme \(g:f_*\mathscr{F}\to\mathscr{F}\) (il peut être plus commode de considérer, plutôt que \(g\), l’isomorphisme \(f^*(g)^{-1}:f^*\mathscr{F}\to\mathscr{F}\)). Ceci pour expliquer le mot « action ».

\Page{152}
Précisons que la \textit{topologie} du groupe de Weil
\[
\mathrm{W}(\mathrm{X}_0,\bar x)
=\pi_1(\mathrm{X}_0,\bar x)\times_{\operatorname{Gal}(\mathbf{F}/\mathbf{F}_q)}
\mathrm{W}(\mathbf{F}/\mathbf{F}_q)
\]
est la topologie produit : le sous-groupe \(\pi_1(\mathrm{X},\bar x)=\operatorname{Ker}(\pi_1(\mathrm{X}_0,\bar x)\to\pi_1(\operatorname{Spec}(\mathbf{F}_q),\bar x))\) est ouvert et fermé.

Un schéma \(\mathrm{X}_0\) sur \(\mathbf{F}_q\) est aussi un schéma sur \(\mathbf{F}_p\) : \(\mathrm{X}_0\to\operatorname{Spec}(\mathbf{F}_q)\to\operatorname{Spec}(\mathbf{F}_p)\), et une clôture algébrique de \(\mathbf{F}_q\) est aussi une clôture algébrique de \(\mathbf{F}_p\). On vérifie que remplacer ainsi \(\mathbf{F}_q\) par \(\mathbf{F}_p\) ne change ni la catégorie des faisceaux de Weil, ni le groupe de Weil (utiliser que \(\mathrm{X}_0\otimes_{\mathbf{F}_p}\mathbf{F}=\coprod\limits_\sigma\mathrm{X}_0\otimes_{\mathbf{F}_q,\sigma}\mathbf{F}\), où \(\sigma\) parcourt \(\operatorname{Hom}(\mathbf{F}_q,\mathbf{F})\sim\operatorname{Gal}(\mathbf{F}_q/\mathbf{F}_p)\)).

\textit{Les notations \textup{(0.7)} sont en vigueur dans la fin de ce numéro.}

\textbf{(1.1.11)} L’image inverse sur \(\mathrm{X}\) d’un faisceau sur \(\mathrm{X}_0\) est munie d’une action par transport de structure de \(\operatorname{Gal}(\mathbf{F}/\mathbf{F}_q)\). Cette construction nous permet d’identifier les faisceaux sur \(\mathrm{X}_0\) à des faisceaux de Weil particuliers (en fait, ceux pour lesquels l’action de \(\mathrm{W}(\mathbf{F}/\mathbf{F}_q)\) se prolonge en une action continue de \(\operatorname{Gal}(\mathbf{F}/\mathbf{F}_q)\)).

Nous admettrons, pour les faisceaux de Weil, l’analogue des résultats de Grothendieck pour les faisceaux ordinaires. La généralisation est facile, à l’aide de (1.3.14).

\textbf{(1.1.12)} Si \(\widetilde{\mathrm{X}}\) est le revêtement universel de \(\mathrm{X}_0\), pointé par le point géométrique \(\bar x\), on dispose d’un morphisme équivariant (relativement au morphisme naturel de \(\mathrm{W}(\mathrm{X}_0,\bar x)\) dans \(\mathrm{W}(\mathbf{F}/\mathbf{F}_q)\)) :
\[
(\widetilde{\mathrm{X}},\ \text{avec action de }\mathrm{W}(\mathrm{X}_0,\bar x))
\to
(\mathrm{X},\ \text{avec action de }\mathrm{W}(\mathbf{F}/\mathbf{F}_q)).
\]
L’image inverse sur \(\widetilde{\mathrm{X}}\) d’un faisceau de Weil est donc munie d’une action de \(\mathrm{W}(\mathrm{X}_0,\bar x)\). Si le faisceau est lisse, sa fibre en \(\bar x\) s’identifie à ses sections globales sur \(\widetilde{\mathrm{X}}\), d’où un foncteur « fibre en \(\bar x\) » :
\[
(\text{faisceaux de Weil lisses})
\to
(\text{représentations de }\mathrm{W}(\mathrm{X}_0,\bar x)),
\]
où « représentation » est défini comme en (1.1.6). Des arguments standard montrent que, si \(\mathrm{X}_0\) est connexe, ce foncteur est une équivalence de catégorie. Il nous permet, comme en (1.1.6), de transposer aux faisceaux de Weil lisses le langage des représentations.

\textbf{(1.1.13)} Supposons \(\mathrm{X}_0\) connexe. Pour \(y\in|\mathrm{X}_0|\) (ou, plus généralement, pour \(y\) un point de \(\mathrm{X}_0\) à valeur dans un corps fini) les Frobenius \(\mathrm{F}_y\) sont dans \(\mathrm{W}(\mathrm{X}_0,\bar x)\), l’image de \(\mathrm{F}_y\) dans \(\operatorname{Gal}(\bar k/\mathbf{F}_q)\) étant la puissance entière \(\varphi^{-\deg(y)}\) (où \(\deg y=[k(y):\mathbf{F}_q]\)) de la substitution de Frobenius. L’élément \(\mathrm{F}_y\) de \(\mathrm{W}(\mathrm{X}_0,\bar x)\) est bien défini à \(\mathrm{W}(\mathrm{X}_0,\bar x)\)-conjugaison près, car son centralisateur dans \(\pi_1(\mathrm{X}_0,\bar x)\) s’envoie sur un sous-groupe ouvert de \(\operatorname{Gal}(\bar k/\mathbf{F}_q)\) (si \(\bar y\) est un point géométrique localisé en \(y\), ce centralisateur contient une copie de \(\operatorname{Gal}(\bar y/y)\), dont l’image dans \(\operatorname{Gal}(\bar k/\mathbf{F}_q)\), identifié à \(\widehat{\mathbf{Z}}\), est \(\deg(y).\widehat{\mathbf{Z}}\)).

% ---- 153-160 ----
\Page{153}

Nous appellerons \textit{degrés} les flèches à droite du diagramme :
\[
\begin{tikzcd}[column sep=1.4em,row sep=1.8em]
0 \arrow[r]
  & \pi_1(\mathrm{X},\bar{x}) \arrow[r] \arrow[d,equal]
  & \mathrm{W}(\mathrm{X}_0,\bar{x}) \arrow[r] \arrow[d,hook]
  & \mathbf{Z} \arrow[d,hook] \\
0 \arrow[r]
  & \pi_1(\mathrm{X},\bar{x}) \arrow[r]
  & \pi_1(\mathrm{X}_0,\bar{x}) \arrow[r]
  & \widehat{\mathbf{Z}}
\end{tikzcd}
\tag*{\textbf{(1.1.13.1)}}
\]
on a \(\deg \mathrm{F}_y=\deg y\). Les flèches \(\deg\) sont surjectives si \(\mathrm{X}_0\) est géométriquement connexe. En général, si le plus grand sous-corps fini de \(\Gamma(\mathrm{X}_0,\mathscr{O}_{\mathrm{X}_0})\) est de degré \(n\) sur \(\mathbf{F}_q\), leurs images sont \(n\mathbf{Z}\) et \(n\widehat{\mathbf{Z}}\).

\textbf{(1.1.14)} Soit \(\mathscr{F}\) un faisceau sur \(\mathrm{X}\). Si \(\mathrm{F}:\mathrm{X}\to\mathrm{X}\) est l’endomorphisme de Frobenius de \(\mathrm{X}\), déduit par extension des scalaires de l’endomorphisme de Frobenius \(x\mapsto x^q\) de \(\mathrm{X}_0\) (lequel est \(\mathbf{F}_q\)-linéaire), et que \(\operatorname{id}_{\mathrm{X}_0}\times\mathrm{F}:\mathrm{X}\to\mathrm{X}\) est l’action du Frobenius géométrique \(\mathrm{F}\in\operatorname{Gal}(\mathbf{F}/\mathbf{F}_q)\), on voit comme dans SGA 5 XV qu’il y a un isomorphisme canonique \((\operatorname{id}_{\mathrm{X}_0}\times\mathrm{F})^*\mathscr{F}\simeq\mathrm{F}^*\mathscr{F}\). Pour \(\mathscr{F}_0\) un faisceau de Weil sur \(\mathrm{X}_0\), on en déduit une correspondance de Frobenius \(\mathrm{F}^*\mathscr{F}\xrightarrow{\sim}\mathscr{F}\), et le foncteur \(\mathscr{F}_0\mapsto(\mathscr{F},\mathrm{F}^*)\) est une équivalence de la catégorie des faisceaux de Weil avec celle des faisceaux \(\mathscr{F}\) sur \(\mathrm{X}\), munis d’un isomorphisme \(\mathrm{F}^*:\mathrm{F}^*\mathscr{F}\xrightarrow{\sim}\mathscr{F}\). Pour les faisceaux constructibles de \(\mathbf{Z}/\ell^n\)-modules, ou pour les \(\mathbf{Z}_\ell\)-faisceaux constructibles, on dispose d’un résultat analogue. Le lecteur prendra garde qu’il n’en va pas de même pour les \(\mathbf{Q}_\ell\)- ou \(\overline{\mathbf{Q}}_\ell\)-faisceaux : prendre par exemple \(\mathscr{F}=\mathbf{Q}_\ell\), et \(\mathrm{F}^*=\) multiplication par \(u\in\mathbf{Q}_\ell^*\) qui ne soit pas une unité \(\ell\)-adique.

Le morphisme composé \(\mathrm{F}^*:\mathrm{H}^i(\mathrm{X},\mathscr{F})\xrightarrow{\sim}\mathrm{H}^i(\mathrm{X},\mathrm{F}^*\mathscr{F})\xrightarrow{\sim}\mathrm{H}^i(\mathrm{X},\mathscr{F})\) déduit de la correspondance de Frobenius coïncide avec l’action de \(\mathrm{F}\in\mathrm{W}(\bar{k}/\mathbf{F}_q)\), agissant sur \(\mathrm{H}^i(\mathrm{X},\mathscr{F})\) par transport de structure. De même pour les \(\mathrm{H}_c^i\)\ldots

\textbf{(1.1.15)} On appelle \(\pi_1(\mathrm{X}_0,\bar{x})\) le \textit{groupe fondamental arithmétique}, et \(\pi_1(\mathrm{X},\bar{x})\) le \textit{groupe fondamental géométrique}. Si \(\mathrm{X}_0\) est normal, ce sont des groupes profinis.

Si \(\mathscr{F}_0\) est un faisceau de Weil lisse sur \(\mathrm{X}_0\), son \textit{groupe de monodromie géométrique} est l’image de \(\pi_1(\mathrm{X},\bar{x})\) dans \(\mathrm{GL}(\mathscr{F}_{\bar{x}})\).

\section*{(1.2) Poids}

\textit{Définition} \textbf{(1.2.1)}. — \textit{Soient \(q\) une puissance d’un nombre premier, et \(n\in\mathbf{Z}\). Un nombre \(\alpha\) est dit \textup{pur de poids} \(n\), rel. \(q\), s’il est algébrique et que tous ses conjugués complexes sont de valeur absolue \(q^{n/2}\).}

Dans la suite, « nombre » signifiera « élément de \(\overline{\mathbf{Q}}_\ell\) ».

\textit{Définition} \textbf{(1.2.2)}. — \textit{Soient \(\mathrm{X}\) un schéma de type fini sur \(\mathbf{Z}\) et \(\mathscr{F}\) un faisceau sur \(\mathrm{X}\).}

(i) \textit{On dit que \(\mathscr{F}\) est \textup{ponctuellement pur} s’il existe un entier \(n\), le \textup{poids} de \(\mathscr{F}\), tel que, pour tout \(x\in|\mathrm{X}|\), les valeurs propres de \(\mathrm{F}_x\) soient pures de poids \(n\), rel. \(\mathrm{N}(x)\).}

(ii) \textit{\(\mathscr{F}\) est \textup{mixte} s’il admet une filtration finie de quotients successifs des faisceaux ponctuellement purs. Les poids de ceux de ces quotients qui sont non nuls sont les \textup{poids ponctuels} de \(\mathscr{F}\).}

\Page{154}

D’après ces définitions, le faisceau \(0\) est ponctuellement pur et, pour tout \(n\), est de poids \(n\). Il est mixte, d’ensemble de poids vide.

\textit{Variante} \textbf{(1.2.3)}. — Si \(k\) est un corps fini, les faisceaux sur \(\operatorname{Spec}(k)\) s’identifient aux représentations \(\ell\)-adiques de \(\operatorname{Gal}(\bar{k}/k)\). On leur transporte la terminologie (1.2.2).

\textit{Variante} \textbf{(1.2.4)}. — Une même terminologie s’applique aux faisceaux de Weil (1.1.10), ainsi qu’aux \(\overline{\mathbf{Q}}_\ell\)-vectoriels munis d’une action de Frobenius, en ce qui concerne la variante (1.2.3).

\textit{Stabilités} \textbf{(1.2.5)}. — (i) \textit{La catégorie des faisceaux ponctuellement purs de poids \(n\) est stable par les opérations de passage au quotient, à un sous-faisceau, par extension, image réciproque, image directe par un morphisme fini.}

(ii) \textit{Le produit tensoriel de deux faisceaux ponctuellement purs de poids \(n\) et \(m\) est ponctuellement pur de poids \(n+m\). Le dual d’un faisceau lisse ponctuellement pur de poids \(n\) est ponctuellement pur de poids \(-n\).}

(iii) \textit{La catégorie des faisceaux mixtes est stable par les opérations de \textup{(i)}, et par produit tensoriel.}

(iv) \textit{Notons enfin que \(\overline{\mathbf{Q}}_\ell(1)\) est ponctuellement pur de poids \(-2\).}

\textbf{(1.2.6)} Pour montrer que \(\alpha\in\overline{\mathbf{Q}}_\ell^*\) est pur de poids \(n\), il suffit de vérifier que pour tout isomorphisme \(\iota\) de \(\overline{\mathbf{Q}}_\ell\) avec \(\mathbf{C}\), le nombre complexe \(\iota\alpha\) est de valeur absolue \(q^{n/2}\). Le nombre \(\alpha\) est alors automatiquement algébrique, sans quoi \(\iota\alpha\) pourrait être n’importe quel nombre transcendant. Dans toute la suite, nos arguments concerneront séparément chaque isomorphisme \(\iota\). C’est pourquoi, bien que les concepts importants soient ceux de (1.2.2), la terminologie suivante nous sera utile.

Soit \(\iota\) un isomorphisme (de corps) de \(\overline{\mathbf{Q}}_\ell\) avec \(\mathbf{C}\). Si \(q\) est une puissance d’un nombre premier, et que \(\alpha\in\overline{\mathbf{Q}}_\ell^*\), on pose :
\[
w_q(\alpha):=2\log_q|\iota\alpha|;
\tag*{\textbf{(1.2.6.1)}}
\]
c’est le \textit{\(\iota\)-poids} de \(\alpha\), rel. \(q\). Pour tout nombre réel \(\beta\), un faisceau \(\mathscr{F}\) est dit \textit{ponctuellement \(\iota\)-pur de poids} \(\beta\) si, pour tout \(x\in|\mathrm{X}|\) et toute valeur propre \(\alpha\) de \(\mathrm{F}_x\) agissant sur \(\mathscr{F}\), on a \(w_{\mathrm{N}(x)}(\alpha)=\beta\). On dit que \(\mathscr{F}\) est \(\iota\)-mixte s’il admet une filtration finie de quotients successifs ponctuellement \(\iota\)-purs. Les poids de ceux de ces quotients qui sont non nuls sont les \(\iota\)-poids (ponctuels) de \(\mathscr{F}\). On transpose de même les variantes (1.2.3) et (1.2.4). L’analogue de (1.2.5) reste vrai. On notera que, pour \(k\) un corps fini, une représentation \(\mathrm{V}\) de \(\mathrm{W}(\bar{k}/k)\) est automatiquement \(\iota\)-mixte.

\textbf{(1.2.7)} \textit{Torsion.} Pour \(b\in\overline{\mathbf{Q}}_\ell^*\), nous appellerons \(\overline{\mathbf{Q}}_\ell^{(b)}\) un quelconque faisceau de Weil sur \(\operatorname{Spec}(\mathbf{F}_p)\), de rang un et pour lequel \(\mathrm{F}\) soit la multiplication par \(b\). Pour \(b\) une unité \(\ell\)-adique, c’est un faisceau ordinaire. Si un choix a été fait d’une clôture algébrique \(\mathbf{F}\)
\Page{155}
de \(\mathbf{F}_p\), on le normalisera par le choix d’un isomorphisme entre sa fibre géométrique en \(\mathbf{F}\) et \(\overline{\mathbf{Q}}_\ell\). Pour \(\mathrm{X}\) un schéma de caractéristique \(p\) (i.e. sur \(\mathbf{F}_p\)), \(\mathscr{F}\) un faisceau sur \(\mathrm{X}\) et \(b\) une unité \(\ell\)-adique (resp. \(\mathrm{X}\) de type fini sur \(\mathbf{F}_p\), \(\mathscr{F}\) un faisceau de Weil et \(b\in\overline{\mathbf{Q}}_\ell^*\)), on note \(\mathscr{F}^{(b)}\) le produit tensoriel de \(\mathscr{F}\) par l’image inverse de \(\overline{\mathbf{Q}}_\ell^{(b)}\); on dit que \(\mathscr{F}^{(b)}\) se déduit de \(\mathscr{F}\) par \textit{torsion}. Le faisceau \(\overline{\mathbf{Q}}_\ell^{(b)}\) est ponctuellement \(\iota\)-pur de poids \(w_p(b)\); la torsion \(\mathscr{F}\mapsto\mathscr{F}^{(b)}\) accroît donc les poids par \(w_p(b)\).

\textit{Remarque} \textbf{(1.2.8)}. — Dans les définitions des faisceaux ponctuellement purs, et des faisceaux mixtes, nous avons imposé aux poids d’être entiers. Dans celles des faisceaux ponctuellement \(\iota\)-purs, et des faisceaux \(\iota\)-mixtes, il est commode de leur permettre d’être des nombres réels quelconques. Le théorème profond (3.4.1) montre \textit{a posteriori} que pour \(\mathrm{X}\) de type fini sur \(\mathbf{F}_p\), cette généralité est largement illusoire : tout faisceau \(\iota\)-mixte est somme directe de faisceaux déduits par torsion (1.2.7) de faisceaux \(\iota\)-mixtes de poids entiers, et tout faisceau \(\iota\)-mixte pour tout \(\iota\) est somme directe de faisceaux déduits par torsion de faisceaux mixtes.

Tout faisceau lisse est extension successive de faisceaux lisses irréductibles, et nous verrons en (1.3.6) que tout faisceau lisse sur \(\mathrm{X}\) normal connexe de type fini sur \(\mathbf{F}_p\) se déduit par torsion d’un faisceau dont le déterminant (\(=\)puissance extérieure maximale) est défini par un caractère d’ordre fini du groupe fondamental de \(\mathrm{X}\). La

\textit{Conjecture} \textbf{(1.2.9)}. — \textit{Sur \(\mathrm{X}\) de type fini sur \(\mathbf{F}_p\), tout faisceau est \(\iota\)-mixte,}

\noindent est donc conséquence de la conjecture (1.2.10) plus précise ci-dessous.

Un de nos outils principaux, le théorème (1.5.1), est un résultat partiel en direction de (1.2.9).

\textit{Conjecture} \textbf{(1.2.10)}. — \textit{Soient \(\mathrm{X}\) normal connexe de type fini sur \(\mathbf{F}_p\), et \(\mathscr{F}\) un faisceau lisse irréductible dont le déterminant est défini par un caractère d’ordre fini du groupe fondamental.}

(i) \textit{\(\mathscr{F}\) est pur de poids \(0\).}

(ii) \textit{Il existe un corps de nombres \(\mathrm{E}\subset\overline{\mathbf{Q}}_\ell\) tel que le polynôme \(\det(1-\mathrm{F}_xt,\mathscr{F})\) pour \(x\in|\mathrm{X}|\), soit à coefficients dans \(\mathrm{E}\).}

(iii) \textit{Pour \(\lambda\) une place non archimédienne première à \(p\) de \(\mathrm{E}\), les racines inverses dans \(\overline{\mathrm{E}}_\lambda\) du polynôme \(\det(1-\mathrm{F}_xt,\mathscr{F})\) (valeurs propres de Frobenius) sont des unités \(\lambda\)-adiques.}

(iv) \textit{Pour \(\lambda\) divisant \(p\), la valuation des racines inverses vérifie :}
\[
|v(\alpha)/v(\mathrm{N}x)|\leq1/2\operatorname{rg}\mathscr{F}.
\]

(v) \textit{Pour \(\mathrm{E}\) convenable (peut-être plus grand qu’en \textup{(ii)}), et chaque place non archimédienne \(\lambda\) première à \(p\), il existe un \(\mathrm{E}_\lambda\)-faisceau compatible à \(\mathscr{F}\) (mêmes valeurs propres des Frobenius).}

(vi) \textit{Pour \(\lambda\) divisant \(p\), on espère des petits camarades cristallins.}

Les parties (i), (iii), (iv) de la conjecture résultent du cas particulier où \(\mathrm{X}\) est une courbe. Pour \(\mathscr{F}\) de rang \(2\) sur une courbe, les travaux de Drinfeld ramènent la
\Page{156}
conjecture ((vi) jusqu’à présent exclue) au contrôle de la constante de l’équation fonctionnelle de la fonction \(\mathrm{L}\) attachée à \(\mathscr{F}\) (il s’agit de \(\mathrm{L}(\mathscr{F},\chi)\) : \(\mathrm{L}\) et ses tordues). En particulier, si \(\mathscr{F}\) appartient à un système compatible infini de représentations \(\lambda\)-adiques, (i) à (v) sont vérifiés.

\textit{Remarque} \textbf{(1.2.11)}. — Je ne prétends pas croire à l’existence d’isomorphismes entre \(\overline{\mathbf{Q}}_\ell\) et \(\mathbf{C}\), et ceux-ci ne sont qu’une commodité d’exposition. Chaque fois que nous prouverons qu’un nombre est pur, un fragment facile de la démonstration suffirait à établir qu’il est algébrique. Pour le reste des arguments, il suffirait alors de considérer les plongements complexes du sous-corps de \(\overline{\mathbf{Q}}_\ell\) formé des nombres algébriques, et ceci ne requiert pas l’axiome du choix.

\textit{Définition} \textbf{(1.2.12)}. — \textit{Soient \(\mathrm{X}\) un schéma de type fini sur \(\mathbf{Z}\) et \(\mathscr{F}\) un faisceau lisse sur \(\mathrm{X}\). On dit que \(\mathscr{F}\) est \textup{totalement réel} si pour tout \(x\in|\mathrm{X}|\), les coefficients du polynôme caractéristique de \(\mathrm{F}_x\) agissant sur \(\mathscr{F}\) sont des nombres algébriques totalement réels.}

La variante suivante nous sera utile en égale caractéristique.

\textit{Variante} \textbf{(1.2.13)}. — \textit{On dit que \(\mathscr{F}\) est \(\iota\)-réel si, pour tout \(x\in|\mathrm{X}|\), le polynôme \(\iota\det(1-\mathrm{F}_xt,\mathscr{F})\) est à coefficients réels.}

\textit{Remarque} \textbf{(1.2.14)}. — Tout faisceau lisse ponctuellement pur \(\mathscr{F}\) est facteur direct d’un faisceau lisse ponctuellement pur totalement réel, à savoir, s’il est de poids \(n\), le faisceau \(\mathscr{F}\oplus\mathscr{F}^{\vee}(-n)\). De même, un faisceau lisse ponctuellement \(\iota\)-pur de poids entier est facteur direct d’un faisceau lisse ponctuellement \(\iota\)-pur \(\iota\)-réel.

Dans ces variantes, on peut prendre pour \(\mathscr{F}\) un faisceau de Weil; pour les faisceaux de Weil, un faisceau lisse ponctuellement \(\iota\)-pur de poids quelconque est facteur direct d’un faisceau lisse ponctuellement \(\iota\)-pur \(\iota\)-réel.

\section*{(1.3) Poids déterminantiels}

\textit{Les conventions \textup{(0.7)} sont en vigueur.}

Soient \(\mathrm{X}_0\) un schéma normal géométriquement connexe de type fini sur \(\mathbf{F}_q\) et \(\bar{x}\) un point géométrique de \(\mathrm{X}_0\). Dans cette section, nous donnons des conséquences du théorème suivant :

\textit{Théorème} \textbf{(1.3.1)}. — \textit{L’image du groupe fondamental géométrique \(\pi_1(\mathrm{X},\bar{x})\) dans le quotient de \(\mathrm{W}(\mathrm{X}_0,\bar{x})\) par l’adhérence de son groupe dérivé est le produit d’un groupe fini d’ordre premier à \(p\) par un pro-\(p\)-groupe.}

Ce théorème est bien connu, mais je n’ai pu le trouver dans la littérature. Pour nos résultats principaux, seul le cas où \(\mathrm{X}\) est une courbe nous importe. Ci-dessous une preuve dans ce cas. Le cas général sera traité en (1.11.4).

\Page{157}

Soient \(\overline{\mathrm{X}}_0\) la complétée projective et lisse de \(\mathrm{X}_0\), et \(\mathrm{S}=\overline{\mathrm{X}}_0-\mathrm{X}_0\). Notons \(k\) le corps des fonctions de \(\mathrm{X}_0\), \(k_v\) son complété en une place \(v\) (identifiée à un point fermé \(v\in|\overline{\mathrm{X}}_0|\)), \(r_v^*\) le groupe des unités de \(k_v^*\) et \(k_{\mathbf{A}}^*\) le groupe des idèles de \(k\). La théorie du corps de classe identifie \(\mathrm{W}(\mathrm{X}_0,\bar{x})^{ab}\) au groupe de classes d’idèles \((\prod\limits_{v\notin\mathrm{S}}r_v^*)\backslash k_{\mathbf{A}}^*/k^*\), la limite projective des groupes de classes de diviseurs :
\[
\mathrm{C}_m
=\{\text{diviseurs sur }\mathrm{X}_0\}/
 \{\text{ceux de la forme }(\varphi),\ \text{avec }\varphi\equiv1\,(m)\}
\]
pour \(m\) un « module » de plus en plus grand, concentré sur \(\mathrm{S}\). L’image du groupe fondamental géométrique s’obtient en remplaçant \(k_{\mathbf{A}}^*\) par le sous-groupe \(k_{\mathbf{A}}^1\) des idèles de norme \(1\), i.e. en remplaçant les \(\mathrm{C}_m\) par les sous-groupes \(\mathrm{C}_m^0\) des classes de diviseurs de degré \(0\). Le noyau de l’application :
\[
(\prod_{v\notin\mathrm{S}}r_v^*)\backslash k_{\mathbf{A}}^1/k^*
\longrightarrow
(\prod_v r_v^*)\backslash k_{\mathbf{A}}^1/k^*
\]
est un quotient de \(\prod\limits_{v\in\mathrm{S}}r_v^*\), donc est le produit d’un pro-\(p\)-groupe par un groupe fini, et le groupe image \(\prod\limits_v r_v^*\backslash k_{\mathbf{A}}^1/k^*\) des classes de diviseurs de degré \(0\), pour l’équivalence linéaire, est fini (Weil, BNT IV, Th. 7). La variété de Picard \(\operatorname{Pic}^0(\mathrm{X})\) est en effet de type fini, et ce groupe de classes de diviseurs est le groupe des points rationnels sur \(\mathbf{F}_q\) de \(\operatorname{Pic}^0(\mathrm{X})\).

\textbf{(1.3.2)} \textit{Dans la suite de cet article, nous appellerons simplement \textup{faisceaux} les faisceaux de Weil \textup{(1.1.10)}. Les \(\overline{\mathbf{Q}}_\ell\)-faisceaux constructibles seront appelés \textup{faisceaux étales}.}

Cette convention nous amènera parfois à appliquer à des \(\mathbf{Q}_\ell\)-faisceaux des résultats prouvés pour des \(\overline{\mathbf{Q}}_\ell\)-faisceaux. La justification sera toujours immédiate.

\textbf{(1.3.3)} Soit \(\mathscr{F}_0\) un faisceau de rang \(1\) sur \(\mathrm{X}_0\). Le groupe \(\mathrm{W}(\mathrm{X}_0,\bar{x})\) agit sur \(\mathscr{F}_{\bar{x}}\) par homothéties, via un caractère \(\chi:\mathrm{W}(\mathrm{X}_0,\bar{x})\to\overline{\mathbf{Q}}_\ell^*\). Réciproquement, pour tout tel caractère \(\chi\), l’équivalence (1.1.12) nous fournit un faisceau \(\mathscr{L}_0(\chi)\), muni d’une trivialisation \(\mathscr{L}(\chi)_{\bar{x}}\sim\overline{\mathbf{Q}}_\ell\).

Le caractère \(\chi\) se factorise par le plus grand quotient abélien de \(\mathrm{W}(\mathrm{X}_0,\bar{x})\). D’après (1.3.1), le groupe de monodromie géométrique de \(\mathscr{F}_0\), i.e. l’image par \(\chi\) de \(\pi_1(\mathrm{X},\bar{x})\), est donc le produit d’un groupe fini par un pro-\(p\)-groupe. Cette image est aussi un sous-groupe compact de \(\mathrm{E}^*\), pour \(\mathrm{E}\subset\overline{\mathbf{Q}}_\ell\) une extension finie de \(\mathbf{Q}_\ell\). C’est donc le produit d’un groupe fini par un pro-\(\ell\)-groupe et, puisque \(p\ne\ell\), c’est un groupe fini. En particulier, une puissance \(\chi^n\) \((n>0)\) de \(\chi\) est triviale sur \(\pi_1(\mathrm{X},\bar{x})\), i.e. de la forme \(w\mapsto b^{\deg(w)}\).
Si \(c\) est une racine \(n\)-ième de \(b\), le caractère \(\chi\) est le produit du caractère \(w\mapsto c^{\deg(w)}\), par un caractère \(\varepsilon\) d’ordre \(n\), i.e. à valeurs dans le groupe cyclique des racines \(n\)-ièmes de l’unité de \(\overline{\mathbf{Q}}_\ell\). En particulier, pour \(y\in|\mathrm{X}_0|\), on a \(|\iota\chi(\mathrm{F}_y)|=|\iota c|^{\deg(y)}\) : le faisceau \(\mathscr{F}_0\) est ponctuellement \(\iota\)-pur, de poids le \(\iota\)-poids de \(c\) rel. \(q\) (1.2.6.1). On a prouvé :

\textit{Proposition} \textbf{(1.3.4)}. — \textit{Soient \(\mathscr{F}_0\) un faisceau lisse de rang \(1\) sur \(\mathrm{X}_0\) et \(\chi:\mathrm{W}(\mathrm{X}_0,\bar{x})\to\overline{\mathbf{Q}}_\ell^*\) le caractère correspondant.}

\Page{158}

(i) \textit{Le groupe de monodromie géométrique de \(\mathscr{F}_0\) est fini. Plus précisément, le caractère \(\chi\) est le produit d’un caractère \(w\mapsto c^{\deg(w)}\) par un caractère d’ordre fini.}

(ii) \textit{\(\mathscr{F}_0\) est ponctuellement \(\iota\)-pur, de poids le poids de \(c\), rel. \(q\).}

\textit{Variante.} — Si \(\mathrm{X}_0\) est une courbe, la démonstration de (1.3.1) est complète — indépendante de (1.11) — d’où déjà (1.3.4). Voici, d’après N. Katz, comment déduire le cas général de (1.3.4) de ce cas particulier.

Soient \(\mathrm{E}\subset\overline{\mathbf{Q}}_\ell\) comme plus haut et \(\mathrm{M}\) le groupe des éléments d’ordre fini de \(\mathrm{E}^*\) (racines de l’unité). C’est un groupe fini. On commence par montrer que, quels que soient \(x\) et \(y\) dans \(|\mathrm{X}_0|\) on a \(\chi(\mathrm{F}_x)^{1/\deg(x)}=\chi(\mathrm{F}_y)^{1/\deg(y)}\) dans \(\mathrm{E}^*\otimes\mathbf{Q}\) (noté multiplicativement) : joignant \(x\) et \(y\) par une courbe, on se ramène au cas où \(\mathrm{X}_0\) est une courbe et on applique (1.3.4). Quel que soit le \(0\)-cycle de degré \(0\) \(\sum n(x)x\), on a donc \(\chi(\prod\mathrm{F}_x^{n(x)})\in\mathrm{M}\). Il résulte du théorème de Čebotarev que ces produits de Frobenius sont denses dans le \(\pi_1\) géométrique, que \(\chi\) envoie donc dans \(\mathrm{M}\) : la restriction de \(\chi\) au \(\pi_1\) géométrique est d’ordre fini, et on conclut comme précédemment.

\textit{Définition} \textbf{(1.3.5)}. — \textit{Soit \(\mathscr{F}_0\) un faisceau lisse sur \(\mathrm{X}_0\). Les \textup{poids déterminantiels (rel. \(\iota\))} de \(\mathscr{F}_0\) sont les nombres \(\frac{1}{d}\) \textup{(\(\iota\)-poids de \(\overset{d}{\wedge}\mathscr{G}_0\))} pour \(\mathscr{G}_0\) un constituant de \(\mathscr{F}_0\), et \(d\) son rang.}

Il est clair qu’un faisceau lisse ponctuellement \(\iota\)-pur de poids \(\beta\) est aussi purement de poids déterminantiel \(\beta\).

\textbf{(1.3.6)} Si \(\mathscr{F}_0\) est lisse, irréductible de rang \(n\) et que \(b\in\mathrm{Q}_\ell^*\) est tel que, pour un \(y\in|\mathrm{X}_0|\), on a \(\det(\mathrm{F}_y)=b^{n\deg(y)[\mathbf{F}_q:\mathbf{F}_p]}\), le faisceau \(\mathscr{F}_0^{(b^{-1})}\) (1.2.7) est de poids déterminantiel \(0\) (rel. \(\iota\)) pour tout \(\iota\) : sa puissance extérieure maximale est définie par un caractère d’ordre fini.

\textbf{(1.3.7)} Pour étudier le comportement des poids déterminantiels par produit tensoriel et puissance extérieure, il nous sera commode de passer du langage des faisceaux à celui des représentations de \textit{groupes algébriques de monodromie}. Le dictionnaire est inspiré de Serre [11], II (1.3).

Soit donc \(\mathscr{F}_0\) un faisceau lisse sur \(\mathrm{X}_0\). Soit \(\mathrm{G}^0\) le sous-groupe algébrique de \(\mathrm{GL}(\mathscr{F}_{\bar{x}})\) adhérence de Zariski du groupe de monodromie géométrique (1.1.15). Par \textit{push-out} à partir de la suite exacte \(0\to\pi_1(\mathrm{X},\bar{x})\to\mathrm{W}(\mathrm{X}_0,\bar{x})\to\mathbf{Z}\to0\), on définit un groupe algébrique (non de type fini) \(\mathrm{G}\), extension de \(\mathbf{Z}\) par \(\mathrm{G}^0\) :
\[
\begin{tikzcd}[column sep=1.4em,row sep=1.6em]
0 \arrow[r]
  & \pi_1(\mathrm{X},\bar{x}) \arrow[r] \arrow[d]
  & \mathrm{W}(\mathrm{X}_0,\bar{x}) \arrow[r] \arrow[d]
  & \mathbf{Z} \arrow[r] \arrow[d,equal]
  & 0 \\
0 \arrow[r]
  & \mathrm{G}^0 \arrow[r]
  & \mathrm{G} \arrow[r] \arrow[d]
  & \mathbf{Z} \arrow[r]
  & 0 \\
  & & \mathrm{GL}(\mathscr{F}_{\bar{x}})
\end{tikzcd}
\tag*{\textbf{(1.3.7.1)}}
\]

\Page{159}

Voici la définition exacte : si \(w\) est un élément de \(\mathrm{W}(\mathrm{X}_0,\bar{x})\) de degré \(1\), \(\mathrm{W}(\mathrm{X}_0,\bar{x})\) est le produit semi-direct du groupe \(w^{\mathbf{Z}}\) engendré par \(w\) par \(\pi_1(\mathrm{X},\bar{x})\). L’image \(\bar{w}\) de \(w\) dans \(\mathrm{GL}(\mathscr{F}_{\bar{x}})\) normalise celle de \(\pi_1(\mathrm{X},\bar{x})\) (le groupe de monodromie géométrique). Elle normalise donc son adhérence de Zariski \(\mathrm{G}^0\) et \(\mathrm{G}\) est le produit semi-direct de \(\mathbf{Z}\) par \(\mathrm{G}^0\), relativement à l’action \(\operatorname{int}(\bar{w})\) de \(\mathbf{Z}\) sur \(\mathrm{G}^0\). Le groupe \(\mathrm{G}\) obtenu s’insère dans un diagramme (1.3.7.1), et ceci le caractérise à isomorphisme unique près.

Chaque représentation linéaire du groupe algébrique \(\mathrm{G}\) définit une représentation \(\ell\)-adique de \(\mathrm{W}(\mathrm{X}_0,\bar{x})\), donc un faisceau de Weil sur \(\mathrm{X}_0\). Le faisceau \(\mathscr{F}_0\) de départ correspond à la représentation tautologique (1.3.7.1) de \(\mathrm{G}\) dans \(\mathrm{GL}(\mathscr{F}_{\bar{x}})\).

Si \(\mathscr{F}_0^{(i)}\) est une famille finie de faisceaux lisses, et qu’on applique la construction précédente à la somme \(\mathscr{F}_0\) des \(\mathscr{F}_0^{(i)}\), on peut dans (1.3.7.1) remplacer \(\mathrm{GL}(\mathscr{F}_{\bar{x}})\) par le sous-groupe \(\prod\mathrm{GL}(\mathscr{F}_{\bar{x}}^{(i)})\). Chaque faisceau \(\mathscr{F}_0^{(i)}\) est déterminé par la représentation de \(\mathrm{G}\) dans le facteur correspondant \(\mathrm{GL}(\mathscr{F}_{\bar{x}}^{(i)})\).

\textit{Théorème} \textbf{(1.3.8)} (Grothendieck). — \textit{Soient \(\mathscr{F}_0\) un faisceau lisse sur \(\mathrm{X}_0\), \(\mathrm{G}\) et \(\mathrm{G}^0\) comme ci-dessus et \(\mathrm{G}^{00}\) la composante neutre de \(\mathrm{G}^0\). Alors, le radical du groupe algébrique \(\mathrm{G}^{00}\) est unipotent.}

Nous prouverons d’abord le

\textit{Corollaire} \textbf{(1.3.9)}. — \textit{Si \(\mathscr{F}_0\) est semi-simple, l’adhérence de Zariski \(\mathrm{G}^0\) du groupe de monodromie géométrique de \(\mathscr{F}_0\) est extension d’un groupe fini par un groupe semi-simple.}

Si la représentation de \(\mathrm{W}(\mathrm{X}_0,x)\) sur \(\mathscr{F}_{0x}\) est semi-simple, il en va de même de sa restriction au sous-groupe distingué \(\pi_1(\mathrm{X},x)\) : la somme des \(\pi_1(\mathrm{X},x)\)-sous-modules simples est \(\mathrm{W}\)-stable, donc a un supplémentaire\ldots{} Nous prouverons (1.3.9) en supposant seulement que la représentation de \(\pi_1(\mathrm{X},x)\) est semi-simple. Sous cette hypothèse, \(\mathrm{G}^{00}\) est réductif, et il s’agit de prouver que son plus grand tore central \(\mathrm{T}_1\) est trivial. Il est isogène au plus grand tore \(\mathrm{T}\) quotient de \(\mathrm{G}^{00}\).

Supposons tout d’abord que \(\mathrm{G}^0=\mathrm{G}^{00}\) et que \(\mathrm{G}\) est un produit : \(\mathrm{G}\simeq\mathrm{G}^{00}\times\mathbf{Z}\). L’application \(\mathrm{W}(\mathrm{X}_0,x)\to\mathrm{G}\) fournit alors des morphismes \(\mathrm{W}(\mathrm{X}_0,x)\to\mathrm{T}\), tels que l’image de \(\pi_1(\mathrm{X},x)\) soit Zariski dense. Puisque \(\mathrm{T}\) est un produit de groupes multiplicatifs, il résulte de (1.3.4) que cette image est finie, et \(\mathrm{T}=\{1\}\).

Nous ramènerons le cas général à ce cas particulier en remplaçant \(\mathrm{X}_0\) par un revêtement fini, et \(\mathbf{F}_q\) par une extension finie, donc \(\mathrm{W}(\mathrm{X}_0,x)\) et \(\mathrm{G}\) par des sous-groupes d’indice fini. Le groupe \(\mathrm{W}(\mathrm{X}_0,x)\) agit sur \(\mathrm{T}_1\) en respectant l’ensemble fini \(\mathrm{F}\subset\mathrm{X}(\mathrm{T}_1)\) des caractères par lesquels \(\mathrm{T}_1\) agit sur \(\mathscr{F}_{0\bar{x}}\). Cet ensemble engendre \(\mathrm{X}(\mathrm{T}_1)\) car \(\mathrm{G}^0\), donc \(\mathrm{T}_1\), agit fidèlement; le groupe \(\mathrm{W}(\mathrm{X}_0,x)\) agit donc sur \(\mathrm{T}_1\) via un quotient fini, et on se ramène à supposer cette action triviale. Le groupe des automorphismes extérieurs de \(\mathrm{G}^{00}\), triviaux sur \(\mathrm{T}_1\), est fini : on peut supposer que \(\mathrm{W}(\mathrm{X}_0,x)\) agit sur \(\mathrm{G}^{00}\) par automorphismes intérieurs. On se ramène enfin à supposer que \(\mathrm{G}^0=\mathrm{G}^{00}\), et on a alors \(\mathrm{G}\simeq\mathrm{G}^{00}\times\mathbf{Z}\).

\Page{160}

Prouvons (1.3.8). Soient \(\mathrm{G}\) une suite de Jordan-Hölder pour \(\mathscr{F}_0\), \(\mathrm{P}\) le sous-groupe de \(\mathrm{GL}(\mathscr{F}_x)\) qui respecte la filtration \(\mathrm{G}\), et \(\mathrm{N}\) le sous-groupe de \(\mathrm{P}\) qui agit trivialement sur les quotients successifs (radical unipotent), et \(\mathrm{L}=\mathrm{P}/\mathrm{N}\). On a \(\mathrm{G}^0\subset\mathrm{P}\), et l’analogue \(\mathrm{G}_1^0\) de \(\mathrm{G}^0\) pour \(\operatorname{Gr}_{\mathrm{G}}(\mathscr{F}_0)\) est l’image de \(\mathrm{G}^0\) dans \(\mathrm{L}\) : le groupe algébrique \(\mathrm{G}^0\) est extension de \(\mathrm{G}_1^0\) par un sous-groupe de \(\mathrm{N}\), automatiquement unipotent, de sorte qu’il suffit de prouver (1.3.8) pour \(\operatorname{Gr}_{\mathrm{G}}(\mathscr{F}_0)\), justifiable de (1.3.9).

\textit{Lemme} \textbf{(1.3.10)}. — \textit{Soit \(\mathrm{G}\) un schéma en groupes extension de \(\mathbf{Z}\) par un groupe algébrique \(\mathrm{G}^0\), dont la composante connexe de l’identité \(\mathrm{G}^{00}\) est réductive. Les conditions suivantes sont équivalentes.}

(i) \textit{\(\mathrm{G}\) admet une représentation linéaire dont la restriction à \(\mathrm{G}^0\) est fidèle.}

(ii) \textit{\(\mathrm{G}\) admet une représentation linéaire dont la restriction à \(\mathrm{G}^0\) a un noyau fini.}

(iii) \textit{Soit \(\mathrm{Z}^{00}\) la composante connexe du centre de \(\mathrm{G}^{00}\). La restriction à \(\mathrm{Z}^{00}\) de tout automorphisme intérieur de \(\mathrm{G}\) est d’ordre fini.}

(iv) \textit{Le centre de \(\mathrm{G}\) s’envoie sur un sous-groupe d’indice fini de \(\mathbf{Z}\).}

L’implication \(\text{(i)}\Rightarrow\text{(ii)}\) est triviale. Prouvons que \(\text{(ii)}\Rightarrow\text{(iii)}\) : si \(\rho\) est une représentation linéaire comme en (ii), les poids de \(\mathrm{Z}^{00}\) dans cette représentation forment un ensemble fini de caractères, qui engendrent rationnellement le groupe des caractères de \(\mathrm{Z}^{00}\). Cet ensemble fini est respecté par l’action de \(\mathrm{G}\) sur \(\mathrm{Z}^{00}\) par automorphismes intérieurs. Un sous-groupe d’indice fini de \(\mathrm{G}\) agit donc trivialement sur cet ensemble, et sur \(\mathrm{Z}^{00}\).

Prouvons que \(\text{(iii)}\Rightarrow\text{(iv)}\). Soit \(g\) un élément de \(\mathrm{G}\) d’image \(1\) dans \(\mathbf{Z}\) et soit \(\gamma\) l’automorphisme \(x\mapsto gxg^{-1}\) de \(\mathrm{G}^{00}\). On sait que le groupe des automorphismes d’un groupe réductif est le produit semi-direct du groupe des automorphismes respectant un épinglage par le groupe adjoint (automorphismes intérieurs) (SGA 3 XXIV (1.3)) et qu’un automorphisme qui respecte un épinglage est d’ordre fini si et seulement si sa restriction au centre connexe l’est. L’hypothèse (iii) assure donc que \(\gamma\) est le produit d’un automorphisme d’ordre fini respectant un épinglage par un automorphisme intérieur. Modifiant \(g\) par un élément de \(\mathrm{G}^{00}\), on peut supposer \(\gamma\) d’ordre fini. La restriction de \(\operatorname{Int}(g)\) à \(\mathrm{G}^0\) est alors elle aussi d’ordre fini. Si elle est d’ordre \(n\), \(g^n\) est central (comme centralisant \(\mathrm{G}^0\) et \(g\)), et (iv) en résulte.

Enfin, sous l’hypothèse (iv), \(\mathrm{G}\) admet un sous-groupe central \(\mathrm{Z}\), engendré par un élément \(g\) d’image \(n\ne0\) dans \(\mathbf{Z}\); le groupe \(\mathrm{G}/\mathrm{Z}\) est algébrique et toute représentation fidèle de \(\mathrm{G}/\mathrm{Z}\) fournit une représentation de \(\mathrm{G}\) dont la restriction à \(\mathrm{G}^0\) est fidèle.

\textit{Corollaire} \textbf{(1.3.11)}. — \textit{Supposons \(\mathscr{F}_0\) semi-simple, et soit \(\mathrm{Z}\) le centre de \(\mathrm{G}\). Alors, les noyaux et conoyaux de \(\deg:\mathrm{Z}\to\mathbf{Z}\) sont finis.}

Le noyau \(\mathrm{Z}\cap\mathrm{G}^0\) est contenu dans le centre de \(\mathrm{G}^0\), qui est fini d’après (1.3.9). D’après (1.3.9) encore, la condition (1.3.10) (iii) est trivialement vérifiée (on a \(\mathrm{Z}^{00}=\{e\}\)) et le conoyau est fini d’après (1.3.10) (iv).

% ---- 161-168 ----
\Page{161}
\textit{Corollaire} \textbf{(1.3.12)}. — \textit{Supposons \(\mathcal{F}_0\) semi-simple, et soit \(g\) un élément central de degré \(m\ne0\) de \(\mathrm{G}\). Pour que le faisceau défini par une représentation \(\mathrm{V}\) de \(\mathrm{G}\) soit purement de \(\iota\)-poids déterminantiel \(\beta\), il faut et il suffit que les valeurs propres \(\alpha\) de \(g\) agissant sur \(\mathrm{V}\) vérifient toutes \(|\iota\alpha|=q^{m\beta/2}\).}

On se ramène à supposer \(\mathrm{V}\) simple, auquel cas \(g\) est scalaire. On remplace alors \(\mathrm{V}\) par sa puissance extérieure maximale, pour se ramener au cas où \(\mathrm{V}\) est de rang \(1\). Ce cas est laissé au lecteur (cf. (1.3.4)).

Cette interprétation du poids déterminantiel fournit aussitôt la proposition suivante, qui est le résultat principal de ce numéro.

\textit{Proposition} \textbf{(1.3.13)}. — (i) \textit{Soit \(f:\mathrm{X}'_0\to\mathrm{X}_0\) un morphisme dominant de schémas normaux connexes de type fini sur \(\mathbf{F}_q\). Pour qu’un faisceau lisse sur \(\mathrm{X}_0\) soit purement de \(\iota\)-poids déterminantiel \(\beta\), il faut et il suffit que son image inverse sur \(\mathrm{X}'_0\) le soit.}

(ii) \textit{Si des faisceaux lisses \(\mathcal{F}_0\) et \(\mathcal{G}_0\) sur \(\mathrm{X}_0\) sont purement de \(\iota\)-poids déterminantiel \(\beta\) et \(\gamma\), alors \(\mathcal{F}_0\otimes\mathcal{G}_0\) est purement de \(\iota\)-poids déterminantiel \(\beta+\gamma\).}

(iii) \textit{Soit \(\mathcal{F}_0\) un faisceau lisse sur \(\mathrm{X}_0\), et soit \(n(\beta)\) la somme des rangs des constituants de \(\mathcal{F}_0\) de \(\iota\)-poids déterminantiel \(\beta\). Alors, les poids déterminantiels de \(\bigwedge\limits^a\mathcal{F}_0\) sont les \(\sum a(\beta)\beta\) avec \(\sum a(\beta)=a\) et \(a(\beta)\le n(\beta)\).}

Pour (ii), on prendra \(\mathrm{G}\) défini par \(\mathcal{F}_0\) et \(\mathcal{G}_0\) ((1.3.7), dernier alinéa).

\textit{Proposition} \textbf{(1.3.14)} (voir (1.3.2)). — \textit{Pour qu’un faisceau lisse irréductible \(\mathcal{F}_0\) de rang \(r\) sur \(\mathrm{X}_0\) soit un faisceau étale, il faut et il suffit que \(\bigwedge\limits^r\mathcal{F}_0\) en soit un. Tout faisceau lisse irréductible est donc obtenu par torsion \textup{(1.2.7)} à partir d’un faisceau étale.}

Soit \(g\) central dans \(\mathrm{G}\) comme en (1.3.12), il agit sur \(\mathcal{F}_{0\bar{x}}\) par multiplication par un scalaire \(u\). Le groupe \(\mathrm{G}\) agit sur \(\bigwedge\limits^r\mathcal{F}_{0\bar{x}}\) via un quotient discret, qui est aussi un quotient de \(\mathrm{W}(\mathrm{X}_0,x)\), et pour que \(\bigwedge\limits^r\mathcal{F}_0\) soit un faisceau étale, il faut et il suffit que \(u\) soit une unité \(\ell\)-adique. Supposons-le, et soit \(\mathrm{E}\) une extension finie de \(\mathbf{Q}_\ell\) telle que l’action de \(\mathrm{W}(\mathrm{X}_0,x)\) se factorise par \(\mathrm{GL}(r,\mathrm{E})\). Il s’agit de montrer que l’image de \(\mathrm{W}(\mathrm{X}_0,x)\) dans \(\mathrm{GL}(r,\mathrm{E})\) est bornée : l’image inverse d’un sous-groupe ouvert sera alors ouverte d’indice fini, et la représentation se prolongera au complété \(\pi_1(\mathrm{X}_0,x)\) de \(\mathrm{W}(\mathrm{X}_0,x)\).

Soient \(\mathrm{W}_1\) le sous-groupe d’indice fini de \(\mathrm{W}(\mathrm{X}_0,x)\) image inverse de \(\mathrm{G}^{00}\times g^{\mathbf{Z}}\), \(\mathrm{W}_1^0\) l’image inverse de \(\mathrm{G}^{00}\), et \(\rho\) le morphisme composé \(\mathrm{W}_1\to\mathrm{G}^{00}\times g^{\mathbf{Z}}\xrightarrow{\mathrm{pr}_1}\mathrm{G}^{00}\). Puisque les \(u^n\) (\(n\in\mathbf{Z}\)) forment un ensemble borné, il suffit de montrer que \(\rho(\mathrm{W}_1)\subset\mathrm{G}^{00}(\mathrm{E})\) est borné. Ce sous-groupe normalise \(\rho(\mathrm{W}_1^0)\), qui est compact et Zariski dense. On conclut par le

\textit{Lemme} \textbf{(1.3.15)}. — \textit{Dans un groupe semi-simple \(\mathrm{H}\) sur un corps local \(\mathrm{E}\), le normalisateur d’un sous-groupe compact et Zariski dense \(\mathrm{K}\subset\mathrm{H}(\mathrm{E})\) est compact.}

Nous ne traiterons que le cas où \(\mathrm{E}\) est de caractéristique \(0\). L’application logarithme, d’abord définie sur un voisinage assez petit de \(e\), est étendue à la réunion des sous-groupes \Page{162}compacts de \(\mathrm{H}(\mathrm{E})\) par la formule \(\log(x)=\frac{1}{n}\log(x^n)\). La densité de \(\mathrm{K}\), et le fait que \(\mathrm{H}\) soit semi-simple, impliquent que \(\log(\mathrm{K})\) engendre \(\operatorname{Lie}\mathrm{H}\) comme espace vectoriel sur \(\mathrm{E}\). Soient \(\mathcal{O}_{\mathrm{E}}\) l’anneau de la valuation de \(\mathrm{E}\), et \(\mathrm{L}^0\) le \(\mathcal{O}_{\mathrm{E}}\)-module engendré par le compact \(\log\mathrm{K}\subset\operatorname{Lie}\mathrm{H}\). Il est stable sous le normalisateur de \(\mathrm{K}\). La représentation adjointe étant à noyau fini, et \(\operatorname{ad}:\mathrm{H}(\mathrm{E})\to\mathrm{GL}(\operatorname{Lie}\mathrm{H})\) propre, il en résulte que ce normalisateur est compact.

\section*{(1.4) Cohomologie des courbes et fonctions L : rappels}

Ce numéro rassemble quelques résultats bien connus dont nous ferons un usage constant.

\textbf{(1.4.1)} Si \(\mathrm{X}\) est une courbe sur un corps algébriquement clos \(k\), et \(\mathcal{F}\) un faisceau sur \(\mathrm{X}\), les groupes \(\mathrm{H}^0(\mathrm{X},\mathcal{F})\), \(\mathrm{H}_c^0(\mathrm{X},\mathcal{F})\), \(\mathrm{H}^2(\mathrm{X},\mathcal{F})\), \(\mathrm{H}_c^2(\mathrm{X},\mathcal{F})\) sont de nature élémentaire :

\textit{a)} \(\mathrm{H}^0(\mathrm{X},\mathcal{F})\) est le groupe des sections globales, \(\mathrm{H}_c^0(\mathrm{X},\mathcal{F})\) le sous-groupe de celles à support compact (à support fini, si \(\mathrm{X}\) n’a pas de composante irréductible complète). Si \(\mathrm{X}\) est connexe, de point base \(x\), et que \(\mathcal{F}\) est lisse, défini par une représentation \(\mathrm{V}\) de \(\pi_1(\mathrm{X},x)\), on a :
\[
\mathrm{H}^0(\mathrm{X},\mathcal{F})=\mathrm{V}^{\pi_1(\mathrm{X},x)};
\tag*{\textbf{(1.4.1.1)}}
\]
\[
\mathrm{H}_c^0(\mathrm{X},\mathcal{F})=
\begin{cases}
\mathrm{V}^{\pi_1(\mathrm{X},x)} & \text{pour \(\mathrm{X}\) complet}\\
0 & \text{sinon.}
\end{cases}
\tag*{\textbf{(1.4.1.2)}}
\]

\textit{b)} Si \(\mathrm{U}\) est le complément d’une partie finie de \(\mathrm{X}\), on a \(\mathrm{H}_c^2(\mathrm{U},\mathcal{F})\xrightarrow{\sim}\mathrm{H}_c^2(\mathrm{X},\mathcal{F})\). Prenant \(\mathrm{U}\) assez petit pour que \(\mathrm{U}_{\mathrm{red}}\) soit lisse et \(\mathcal{F}\) lisse sur \(\mathrm{U}\), ceci permet de calculer le \(\mathrm{H}_c^2\) par dualité de Poincaré à partir du calcul \textit{a)} du \(\mathrm{H}^0\) : si \(\mathrm{U}_{\mathrm{red}}\) est lisse connexe, de point base \(x\), et que, sur \(\mathrm{U}\), \(\mathcal{F}\) est défini par une représentation \(\mathrm{V}\) de \(\pi_1(\mathrm{U},x)\), on a :
\[
\mathrm{H}_c^2(\mathrm{X},\mathcal{F})=\mathrm{V}_{\pi_1(\mathrm{U},x)}(-1).
\tag*{\textbf{(1.4.1.3)}}
\]

\textbf{(1.4.2)} Avec les conventions (0.7), soient \(\mathrm{X}_0\) une courbe lisse absolument irréductible sur \(\mathbf{F}_q\), et \(\mathcal{F}_0\) un faisceau lisse sur \(\mathrm{X}_0\). Soit \(\mathcal{F}'_0\) (resp. \(\mathcal{F}''_0\)) le plus grand sous-faisceau (resp. faisceau quotient) de \(\mathcal{F}_0\) qui devienne constant sur \(\mathrm{X}\). C’est l’image inverse d’un faisceau \(\mathrm{F}'_0\) (resp. \(\mathrm{F}''_0\)) sur \(\operatorname{Spec}(\mathbf{F}_q)\), et la discussion (1.4.1) donne :
\[
\mathrm{H}^0(\mathrm{X},\mathcal{F})=\mathrm{F}'
\tag*{\textbf{(1.4.2.1)}}
\]
\[
\mathrm{H}_c^0(\mathrm{X},\mathcal{F})=\mathrm{F}'\text{ ou }0
\text{ selon que }\mathrm{X}\text{ est ou n’est pas complet}
\tag*{\textbf{(1.4.2.2)}}
\]
\[
\mathrm{H}_c^2(\mathrm{X},\mathcal{F})=\mathrm{F}''(-1).
\tag*{\textbf{(1.4.2.3)}}
\]

Les faisceaux \(\mathcal{F}'_0\) et \(\mathrm{F}'_0\) (resp. \(\mathcal{F}''_0\) et \(\mathrm{F}''_0\)) ont les mêmes poids déterminantiels, et ceux-ci sont parmi les poids déterminantiels de \(\mathcal{F}_0\). D’après (1.4.2), on a donc le

\Page{163}
\textit{Corollaire} \textbf{(1.4.3)}. — \textit{Si \(\alpha\) est une valeur propre de \(\mathrm{F}^*\) sur \(\mathrm{H}^0(\mathrm{X},\mathcal{F})\) ou \(\mathrm{H}_c^0(\mathrm{X},\mathcal{F})\) (resp. \(\mathrm{H}_c^2(\mathrm{X},\mathcal{F})\)), il existe un poids déterminantiel \(\beta\) de \(\mathcal{F}_0\) tel que \(w_q(\alpha)=\beta\) (resp. \(w_q(\alpha)=\beta+2\)).}

Plus trivialement, on a aussi

\textit{Corollaire} \textbf{(1.4.4)}. — \textit{Si \(\alpha\) est une valeur propre de \(\mathrm{F}^*\) sur \(\mathrm{H}^0(\mathrm{X},\mathcal{F})\) ou \(\mathrm{H}_c^0(\mathrm{X},\mathcal{F})\) (resp. \(\mathrm{H}_c^2(\mathrm{X},\mathcal{F})\)), pour tout \(x\in|\mathrm{X}_0|\), \(\alpha^{\deg(x)}\) (resp. \((q^{-1}\alpha)^{\deg(x)}\)) est une valeur propre de \(\mathrm{F}_x\) sur \(\mathcal{F}_0\).}

\textbf{(1.4.5)} Si \(\mathrm{X}_0\) est un schéma de type fini sur \(\mathbf{F}_q\), et que \(\mathcal{F}_0\) est un faisceau sur \(\mathrm{X}_0\), on a d’après Grothendieck une égalité de séries formelles \(\ell\)-adiques :
\[
\prod_{x\in|\mathrm{X}_0|}
\det(1-\mathrm{F}_x t^{\deg x},\mathcal{F})^{-1}
=
\prod_i\det(1-\mathrm{F}t,\mathrm{H}_c^i(\mathrm{X},\mathcal{F}))^{(-1)^{i+1}}.
\tag*{\textbf{(1.4.5.1)}}
\]
Appliquons-lui \(\iota\) : on obtient une égalité entre séries formelles complexes, qui exprime que le membre de gauche est le développement en série de Taylor de la fonction rationnelle au membre de droite. Si le produit au membre de gauche converge pour \(|t|<\mathrm{R}\), on obtient dans ce disque une égalité entre fonctions analytiques. Voici un critère de convergence :

\textit{Proposition} \textbf{(1.4.6)}. — \textit{Si, pour \(x\in|\mathrm{X}_0|\), les valeurs propres \(\alpha\) de \(\mathrm{F}_x\) sur \(\mathcal{F}\) vérifient \(w_{\mathrm{N}(x)}(\alpha)\le\beta\), le produit \(\prod_x\iota\det(1-\mathrm{F}_x t^{\deg x},\mathcal{F})^{-1}\) converge pour \(|t|<q^{-\beta/2-\dim\mathrm{X}_0}\), donc n’a ni zéro ni pôle dans ce disque.}

Posons \(d=\dim\mathrm{X}_0\), et découpons \(\mathrm{X}_0\) en morceaux dont chacun est quasi-fini sur un espace affine \(\mathbf{A}_0^d\) : on trouve une majoration :
\[
(\text{nombre de points }x\in|\mathrm{X}_0|\text{ de degré }n)
\le\mathrm{C}^{\mathrm{te}}\cdot q^{dn}.
\]
Ceci ramène la convergence du produit à celle de la série géométrique :
\[
\sum_n q^{nd}\cdot q^{n\beta/2}\cdot|t|^n.
\]

\textbf{(1.4.7)} Si \(\dim\mathrm{X}_0\le1\), les \(\mathrm{H}_c^i\) sont nuls pour \(i\ne0,1,2\). La formule déduite de (1.4.5.1) par application de \(\iota\) se réduit donc à :
\[
\prod_{x\in|\mathrm{X}_0|}
\iota\det(1-\mathrm{F}_x t^{\deg x},\mathcal{F})^{-1}
=
\frac{\iota\det(1-\mathrm{F}t,\mathrm{H}_c^1(\mathrm{X},\mathcal{F}))}
{\iota\det(1-\mathrm{F}t,\mathrm{H}_c^0(\mathrm{X},\mathcal{F}))
 \iota\det(1-\mathrm{F}t,\mathrm{H}_c^2(\mathrm{X},\mathcal{F}))}.
\tag*{\textbf{(1.4.7.1)}}
\]
Si \(\mathrm{X}_0\) est une courbe lisse affine et que \(\mathcal{F}_0\) est lisse, le \(\mathrm{H}_c^0\) est nul et la formule se simplifie en :
\[
\prod_{x\in|\mathrm{X}_0|}
\iota\det(1-\mathrm{F}_x t^{\deg x},\mathcal{F})^{-1}
=
\frac{\iota\det(1-\mathrm{F}t,\mathrm{H}_c^1(\mathrm{X},\mathcal{F}))}
{\iota\det(1-\mathrm{F}t,\mathrm{H}_c^2(\mathrm{X},\mathcal{F}))}.
\tag*{\textbf{(1.4.7.2)}}
\]

Dans l’usage que nous ferons de ces formules, le membre de gauche sera contrôlé par (1.4.6), et les pôles au membre de droite par (1.4.3).

\Page{164}
\section*{(1.5) Un critère de pureté}

\textit{Les conventions \textup{(0.7)} sont en vigueur.} Soit \(\mathrm{X}_0\) une courbe lisse absolument irréductible sur \(\mathbf{F}_q\). La démonstration du théorème suivant est parallèle à celle du théorème (I.3.2), qu’il généralise.

\textit{Théorème} \textbf{(1.5.1)}. — \textit{Les constituants de tout faisceau lisse \(\iota\)-réel sur \(\mathrm{X}_0\) sont \(\iota\)-purs.}

\textit{Lemme} \textbf{(1.5.2)}. — \textit{Soient \(\mathcal{F}_0\) un faisceau lisse \(\iota\)-réel sur \(\mathrm{X}_0\), et \(r\) le plus grand de ses poids déterminantiels (rel. \(\iota\)) \textup{(1.3.5)}. Alors, pour tout \(x\in|\mathrm{X}_0|\), et toute valeur propre \(\alpha\) de \(\mathrm{F}_x\) sur \(\mathcal{F}\), on a \(w_{\mathrm{N}(x)}(\alpha)\le r\).}

Quitte à ôter de \(\mathrm{X}_0\) un point (autre que celui auquel on s’intéresse), on se ramène à supposer \(\mathrm{X}_0\) affine. Pour chaque entier positif \(k\), on a alors une égalité (1.4.7.2) :
\[
\prod_{x\in|\mathrm{X}_0|}
\iota\det\left(1-\mathrm{F}_x t^{\deg x},
 \mathop{\otimes}\limits^{2k}\mathcal{F}\right)^{-1}
=
\frac{\iota\det\left(1-\mathrm{F}t,
 \mathrm{H}_c^1\left(\mathrm{X},\mathop{\otimes}\limits^{2k}\mathcal{F}\right)\right)}
{\iota\det\left(1-\mathrm{F}t,
 \mathrm{H}_c^2\left(\mathrm{X},\mathop{\otimes}\limits^{2k}\mathcal{F}\right)\right)}.
\tag*{\textbf{(1.5.2.1)}}
\]

Dans cette formule :

\textit{a)} Les facteurs \(\iota\det(1-\mathrm{F}_x t^{\deg x},\mathop{\otimes}\limits^{2k}\mathcal{F}_0)^{-1}\) au membre de gauche sont des séries formelles de terme constant \(1\) à coefficients réels \(\ge0\). La démonstration est comme en (I.3.3), (3.4), avec « rationnel » remplacé par « réel ».

\textit{b)} D’après (1.3.13) (ii), les poids déterminantiels des constituants de \(\mathop{\otimes}\limits^{2k}\mathcal{F}_0\) sont \(\le2kr\). D’après (1.4.3), le membre de droite n’a donc pas de pôle pour \(|t|<q^{-(2kr+2)/2}\).

D’après (I.3.5), (3.6), les facteurs \(\iota\det(1-\mathrm{F}_x t^{\deg x},\mathop{\otimes}\limits^{2k}\mathcal{F}_0)^{-1}\) n’ont pas de pôle pour \(|t|<q^{-(2kr+2)/2}\). En particulier, si \(\alpha\) est valeur propre de \(\mathrm{F}_x\) sur \(\mathcal{F}_0\), \(\iota\alpha^{-2k/\deg x}\) est un pôle, et :
\begin{align*}
&|\iota\alpha|^{2k/\deg x}\le q^{(2kr+2)/2},\\
\text{soit}\qquad &w_{\mathrm{N}(x)}(\alpha)\le r+1/k.
\end{align*}
Le lemme s’obtient en faisant tendre \(k\) vers l’infini.

\textbf{(1.5.3)} \textit{Preuve de} (1.5.1). — Soit \(\mathcal{F}_0\) un faisceau lisse \(\iota\)-réel sur \(\mathrm{X}_0\). Pour \(\beta\in\mathbf{R}\), soient \(\mathcal{F}_0(\beta)\) la somme des constituants de \(\mathcal{F}_0\) de poids déterminantiels \(\beta\) (rel. \(\iota\)), et \(n(\beta)\) le rang de \(\mathcal{F}_0(\beta)\). Soient \(x\in|\mathrm{X}_0|\), et \(\alpha_i^\beta\) (\(1\le i\le n(\beta)\)) les valeurs propres de \(\mathrm{F}_x\) sur \(\mathcal{F}(\beta)\), comptées avec leur multiplicité. Il nous faut prouver que \(w_{\mathrm{N}(x)}(\alpha_i^\beta)=\beta\).

Par définition du poids déterminantiel, on a pour tout \(\gamma\) :
\[
\sum_i w_{\mathrm{N}(x)}(\alpha_i^\gamma)=n(\gamma)\cdot\gamma.
\tag*{\textbf{(1.5.3.1)}}
\]

Supposons, ce qui est loisible, que \(\mathcal{F}_0(\beta)\ne0\), et soit \(\mathrm{N}\) la somme des \(n(\gamma)\) pour \(\gamma>\beta\). D’après (1.3.13) (iii), les poids déterminantiels de la puissance exté-\Page{165}rieure \((\mathrm{N}+1)\)-ième de \(\mathcal{F}_0\) sont \(\le\beta+\sum\limits_{\gamma>\beta}n(\gamma)\cdot\gamma\). Puisque chaque \(\alpha_i^\beta\cdot\prod\limits_{\gamma>\beta}\prod\limits_i \alpha_i^\gamma\) est valeur propre de \(\mathrm{F}_x\) sur \(\bigwedge\limits^{\mathrm{N}+1}\mathcal{F}_0\), on a par (1.5.2) :
\[
w_{\mathrm{N}(x)}(\alpha_i^\beta)
+\sum_{\gamma>\beta}\sum_i w_{\mathrm{N}(x)}(\alpha_i^\gamma)
\le\beta+\sum_{\gamma>\beta}n(\gamma)\cdot\gamma.
\]
Soustrayant de cette inégalité les égalités (1.5.3.1) pour \(\gamma>\beta\), on trouve :
\[
w_{\mathrm{N}(x)}(\alpha_i^\beta)\le\beta.
\tag*{\textbf{(1.5.3.2)}}
\]

Pour déduire de cette inégalité une égalité, on peut soit passer au faisceau dual, soit observer que si on somme sur \(i\) les inégalités (1.5.3.2), on trouve l’égalité (1.5.3.1) pour \(\beta\).

\section*{(1.6) Autour de Jacobson-Morosov}

\textit{Proposition} \textbf{(1.6.1)}. — \textit{Soit \(\mathrm{N}\) un endomorphisme nilpotent d’un objet \(\mathrm{V}\) d’une catégorie abélienne. Il existe alors une et une seule filtration finie \textup{(0.10)} \(\mathrm{M}\) de \(\mathrm{V}\), croissante, telle que \(\mathrm{N}\mathrm{M}_i\subset\mathrm{M}_{i-2}\) et que \(\mathrm{N}^k\) induise un isomorphisme \(\operatorname{Gr}_k^{\mathrm{M}}\mathrm{V}\xrightarrow{\sim}\operatorname{Gr}_{-k}^{\mathrm{M}}\mathrm{V}\).}

Prouvons l’existence (l’unicité se prouve de même). On procède par récurrence sur un entier \(d\) tel que \(\mathrm{N}^{d+1}=0\). Pour \(d=0\) (\(\mathrm{N}=0\)), on prend une filtration \(\mathrm{M}\) triviale (\(\operatorname{Gr}_i^{\mathrm{M}}\mathrm{V}=0\) pour \(i\ne0\)). Ensuite, on prend \(\mathrm{M}_d=\mathrm{V}\), \(\mathrm{M}_{d-1}=\operatorname{Ker}\mathrm{N}^d\), \(\mathrm{M}_{-d}=\operatorname{Im}\mathrm{N}^d\) et \(\mathrm{M}_{-d-1}=0\), de sorte que \(\operatorname{Gr}_d^{\mathrm{M}}\mathrm{V}\) et \(\operatorname{Gr}_{-d}^{\mathrm{M}}\mathrm{V}\) soient les co-images et images de \(\mathrm{N}^d\). Sur \(\operatorname{Ker}\mathrm{N}^d/\operatorname{Im}\mathrm{N}^d\), la puissance \(d\)-ième de \(\mathrm{N}\) induit \(0\). Ceci permet de définir \(\mathrm{M}\) sur \(\operatorname{Ker}\mathrm{N}^d/\operatorname{Im}\mathrm{N}^d\) par récurrence, et on définit \(\mathrm{M}_i\) sur \(\mathrm{V}\) (\(d-1\ge i\ge-d\)) comme l’image inverse dans \(\operatorname{Ker}\mathrm{N}^d\) du sous-objet \(\mathrm{M}_i\) pour \(\operatorname{Ker}\mathrm{N}^d/\operatorname{Im}\mathrm{N}^d\).

\textit{Remarque} \textbf{(1.6.2)}. — La caractérisation (1.6.1) de \(\mathrm{M}\) montre sa compatibilité au passage à la catégorie duale.

\textit{Définition} \textbf{(1.6.3)}. — \textit{La \textup{partie primitive} \(\mathrm{P}_i(\mathrm{V})\) (ou simplement \(\mathrm{P}_i\)) de \(\operatorname{Gr}_i^{\mathrm{M}}\mathrm{V}\) est le noyau du morphisme induit par \(\mathrm{N}\) de \(\operatorname{Gr}_i^{\mathrm{M}}\mathrm{V}\) dans \(\operatorname{Gr}_{i-2}^{\mathrm{M}}\mathrm{V}\).}

\textbf{(1.6.4)} Si \(i>0\), \(\mathrm{N}:\operatorname{Gr}_i^{\mathrm{M}}\mathrm{V}\to\operatorname{Gr}_{i-2}^{\mathrm{M}}\mathrm{V}\) est injectif, puisque \(\mathrm{N}^{i-1}\circ\mathrm{N}\) est un isomorphisme, et \(\mathrm{P}_i=0\). Si \(i\ge0\), et si l’on écrit que \(\mathrm{N}\circ\mathrm{N}^{i+1}:\operatorname{Gr}_{i+2}^{\mathrm{M}}\mathrm{V}\to\operatorname{Gr}_{-i}^{\mathrm{M}}\mathrm{V}\to\operatorname{Gr}_{-i-2}^{\mathrm{M}}\mathrm{V}\) est un isomorphisme, on trouve que \(\operatorname{Gr}_{-i}^{\mathrm{M}}\mathrm{V}\) est somme directe de \(\mathrm{P}_{-i}\) et de l’image de \(\operatorname{Gr}_{i+2}^{\mathrm{M}}\) par \(\mathrm{N}^{i+1}\); le morphisme \(\mathrm{N}\) induit un isomorphisme de cette image avec \(\operatorname{Gr}_{-i-2}^{\mathrm{M}}\mathrm{V}\). Répétant cette construction, on construit par récurrence décroissante sur \(i\ge0\) un isomorphisme de \(\operatorname{Gr}_{-i}^{\mathrm{M}}\mathrm{V}\) avec la somme des \(\mathrm{P}_{-j}\) pour \(j\ge i\) et \(j\) congru à \(i\) mod \(2\). Utilisant de plus les isomorphismes \(\mathrm{N}^i:\operatorname{Gr}_i^{\mathrm{M}}\mathrm{V}\to\operatorname{Gr}_{-i}^{\mathrm{M}}\mathrm{V}\), on obtient des isomorphismes :
\[
\operatorname{Gr}_i^{\mathrm{M}}\mathrm{V}
\simeq
\bigoplus_{\substack{j\ge|i|\\j\equiv i(2)}}\mathrm{P}_{-j}.
\tag*{\textbf{(1.6.4.1)}}
\]

\Page{166}
Via ces isomorphismes, le morphisme \(\mathrm{N}\) de \(\operatorname{Gr}_i^{\mathrm{M}}\mathrm{V}\) dans \(\operatorname{Gr}_{i-2}^{\mathrm{M}}\mathrm{V}\) est l’inclusion évidente pour \(i\ge1\) (\(i-2\ge-1\)), la projection évidente pour \(i-2\le-1\).
\[
\begin{tikzcd}[
  ampersand replacement=\&,
  row sep=0.2em,
  column sep=0.4em,
  cells={nodes={anchor=west,inner sep=1pt}}
]
\phantom{\operatorname{Gr}_4}
  \arrow[dd,bend right=35,start anchor=west,end anchor=west]
  \&[0.8em] \&[2em] \& \& \& \&
  \cdots\arrow[dd,start anchor=south west,end anchor=north west] \& \\
\operatorname{Gr}_3
  \arrow[dd,bend left=35,start anchor=east,end anchor=east]
  \& \colon \& \& \& \&
  \mathrm{P}_{-3}\arrow[dd,start anchor=south west,end anchor=north west]
  \& \& \cdots \\
\operatorname{Gr}_2
  \arrow[dd,bend right=35,start anchor=west,end anchor=west]
  \& \colon \& \& \&
  \mathrm{P}_{-2}\arrow[dd,start anchor=south west,end anchor=north west]
  \& \&
  \mathrm{P}_{-4}\arrow[dd,start anchor=south west,end anchor=north west]
  \& \cdots \\
\operatorname{Gr}_1
  \arrow[dd,bend left=35,start anchor=east,end anchor=east]
  \& \colon \& \&
  \mathrm{P}_{-1}\arrow[dd,start anchor=south west,end anchor=north west]
  \& \&
  \mathrm{P}_{-3}\arrow[dd,start anchor=south west,end anchor=north west]
  \& \& \cdots \\
\operatorname{Gr}_0
  \arrow[dd,bend right=35,start anchor=west,end anchor=west]
  \& \colon \& \mathrm{P}_0 \& \&
  \mathrm{P}_{-2}\arrow[dd,start anchor=south west,end anchor=north west]
  \& \&
  \mathrm{P}_{-4}\arrow[dd,start anchor=south west,end anchor=north west]
  \& \cdots \\
\operatorname{Gr}_{-1}
  \arrow[dd,bend left=35,start anchor=east,end anchor=east]
  \& \colon \& \& \mathrm{P}_{-1} \& \&
  \mathrm{P}_{-3}\arrow[dd,start anchor=south west,end anchor=north west]
  \& \& \cdots \\
\operatorname{Gr}_{-2}
  \arrow[dd,bend right=35,start anchor=west,end anchor=west]
  \& \colon \& \& \& \mathrm{P}_{-2} \& \&
  \mathrm{P}_{-4}\arrow[dd,start anchor=south west,end anchor=north west]
  \& \cdots \\
\operatorname{Gr}_{-3}
  \& \colon \& \& \& \& \mathrm{P}_{-3} \& \& \cdots \\
\phantom{\operatorname{Gr}_{-4}}
  \& \& \& \& \& \& \cdots \&
\end{tikzcd}
\]

\textit{Lemme} \textbf{(1.6.5)}. — \textit{Le morphisme} \(\mathrm{N}:(\mathrm{V},\mathrm{M})\to(\mathrm{V},\mathrm{M}\text{ décalé de }2)\) \textit{est strictement compatible aux filtrations.}

Il faut prouver que \(\mathrm{N}\mathrm{M}_{i+2}\supset\operatorname{Im}(\mathrm{N})\cap\mathrm{M}_i\). Distinguons deux cas. Si \(i<0\), le morphisme \(\mathrm{N}:\mathrm{M}_{i+2}\to\mathrm{M}_i\) a un gradué surjectif, donc est surjectif et \(\mathrm{N}\mathrm{M}_{i+2}=\mathrm{M}_i\), d’où \textit{a fortiori} l’assertion. Si \(i+2>0\), le morphisme \(\mathrm{N}:\mathrm{V}/\mathrm{M}_{i+2}\to\mathrm{V}/\mathrm{M}_i\) a un gradué injectif, donc est injectif, et \(\mathrm{N}^{-1}\mathrm{M}_i\subset\mathrm{M}_{i+2}\), d’où \textit{a fortiori} l’assertion.

Ceci permet d’échanger les constructions \(\operatorname{Ker}\mathrm{N}\) et \(\operatorname{Gr}^{\mathrm{M}}\) :

\textit{Corollaire} \textbf{(1.6.6)}. — \textit{L’inclusion de \(\operatorname{Ker}\mathrm{N}\) dans \(\mathrm{V}\) induit des isomorphismes :}
\[
\operatorname{Gr}_i^{\mathrm{M}}(\operatorname{Ker}\mathrm{N})
\xrightarrow{\sim}\mathrm{P}_i.
\]

\textbf{(1.6.7)} Supposons maintenant que \(\mathrm{V}\) soit un espace vectoriel de dimension finie sur un corps \(k\). Nous commencerons par décrire \(\mathrm{M}\) lorsque, dans une base \(\mathbf{e}\) de \(\mathrm{V}\), la matrice de \(\mathrm{N}\) est
\[
\begin{pmatrix}
0 & 0 & 0 & \cdot & \cdot & \cdot & 0 & 0\\
1 & 0 & 0 & \cdot & \cdot & \cdot & 0 & 0\\
0 & 1 & 0 & \cdot & \cdot & \cdot & 0 & 0\\
\cdot & \cdot & \cdot & \cdot & \cdot & \cdot & \cdot & \cdot\\
0 & 0 & 0 & \cdot & \cdot & \cdot & 1 & 0
\end{pmatrix}.
\tag*{\textbf{(1.6.7.1)}}
\]
Posons \(\dim\mathrm{V}=d+1\) et indexons les vecteurs de base par les entiers allant de \(2\) en \(2\) de \(d\) à \(-d\), de sorte que \(\mathrm{N}e_i=e_{i-2}\) pour \(i\ne -d\) et que \(\mathrm{N}e_{-d}=0\). Alors, \(\mathrm{M}_i\) est engendré par les \(e_j\) pour \(j\le i\).

En général, \(\mathrm{V}\) est somme de sous-espaces \(\mathrm{V}_\alpha\) stables par \(\mathrm{N}\), avec \(\mathrm{N}\) du type (1.6.7.1) sur \(\mathrm{V}_\alpha\), et la filtration \(\mathrm{M}\) de \(\mathrm{V}\) est somme des filtrations ci-dessus des \(\mathrm{V}_\alpha\).

\Page{167}
\textbf{(1.6.8)} Si \(k\) est de caractéristique \(0\), on peut interpréter la filtration \(\mathrm{M}\) en terme du théorème de Jacobson-Morosov : si \(u:\mathrm{SL}(2)\to\mathrm{GL}(\mathrm{V})\) satisfait :
\[
du\begin{pmatrix}0&0\\1&0\end{pmatrix}=\mathrm{N},
\]
et que \(\mathrm{V}_j\) est le sous-espace de \(\mathrm{V}\) formé des vecteurs \(v\) tels que :
\[
u\begin{pmatrix}\lambda&0\\0&\lambda^{-1}\end{pmatrix}(v)=\lambda^jv,
\]
alors \(\mathrm{M}_i\) est la somme des \(\mathrm{V}_j\) pour \(j\le i\). Pour le voir, on se ramène à supposer que \(\mathrm{V}\) est une représentation irréductible de \(\mathrm{SL}(2)\) (i.e. une puissance symétrique de la représentation évidente). On est alors dans la situation de (1.6.7).

Cette interprétation de \(\mathrm{M}\) donne son comportement par produit tensoriel et passage au dual (pour lequel on peut plutôt invoquer (1.6.2)).

\textit{Proposition} \textbf{(1.6.9)}. — \textit{Définissons le produit tensoriel \((\mathrm{V},\mathrm{N})\) de \((\mathrm{V}',\mathrm{N}')\) et \((\mathrm{V}'',\mathrm{N}'')\) par \(\mathrm{V}=\mathrm{V}'\otimes\mathrm{V}''\), \(\mathrm{N}=\mathrm{N}'\otimes1+1\otimes\mathrm{N}''\), et le dual de \((\mathrm{V}',\mathrm{N}')\) comme étant \((\mathrm{V}^{\prime *},-{}^t\mathrm{N}')\).}

(i) \textit{Si \(k\) est de caractéristique \(0\), la filtration \(\mathrm{M}\) d’un produit tensoriel est le produit tensoriel des filtrations \(\mathrm{M}\) des facteurs (\(\mathrm{M}_i(\mathrm{V}'\otimes\mathrm{V}'')=\sum\limits_{i'+i''=i}\mathrm{M}_{i'}(\mathrm{V}')\otimes\mathrm{M}_{i''}(\mathrm{V}'')\)).}

(ii) \textit{La filtration \(\mathrm{M}\) d’un dual est le dual de la filtration \(\mathrm{M}\) de l’espace de départ (\(\mathrm{M}_i(\mathrm{V}^*)=\mathrm{M}_{-i-1}(\mathrm{V})^\perp\)).}

(iii) \textit{Sous ces hypothèses, la formation de \(\operatorname{Gr}^{\mathrm{M}}\) est donc compatible au produit tensoriel et au passage au dual.}

\textbf{(1.6.10)} On suppose dorénavant que \(\mathrm{V}\) est de dimension finie sur \(k\) de caractéristique \(0\). Sur \(\operatorname{Gr}^{\mathrm{M}}(\mathrm{V})\), il existe alors une unique représentation \(v\) de \(\mathrm{SL}(2)\) telle que \(dv\begin{pmatrix}0&0\\1&0\end{pmatrix}\) soit \(\mathrm{N}:\operatorname{Gr}_*^{\mathrm{M}}\mathrm{V}\to\operatorname{Gr}_{*-2}^{\mathrm{M}}\mathrm{V}\) et que \(\begin{pmatrix}\lambda&0\\0&\lambda^{-1}\end{pmatrix}\) agisse par multiplication par \(\lambda^j\) sur \(\operatorname{Gr}_j^{\mathrm{M}}(\mathrm{V})\) : l’existence se voit en choisissant \(u\) comme en (1.6.8), en identifiant \(\operatorname{Gr}_j^{\mathrm{M}}(\mathrm{V})\) à \(\mathrm{V}_j\) à l’aide du scindage \(\mathrm{V}_j\) de la filtration et en transportant \(u\) au gradué par cette identification. L’unicité résulte de Bourbaki, Lie VIII, § 11, n\textsuperscript{o} 1, lemme 1. La formation de \(v\) est compatible au produit tensoriel et au passage au dual pris comme en (1.6.9). Ceci va nous permettre de déterminer le comportement de \(\mathrm{P}_j\) par produit tensoriel et passage au dual.

\textbf{(1.6.11)} Choisissons une fois pour toutes, pour chaque entier \(d\ge0\), une représentation irréductible \(\mathrm{S}_d\) de \(\mathrm{SL}(2)\) de dimension \(d+1\), et un vecteur de plus bas poids \(e_{-d}\in\mathrm{S}_d\) : \(e_{-d}\ne0\), et \(\begin{pmatrix}\lambda&0\\0&\lambda^{-1}\end{pmatrix}\) transforme \(e_{-d}\) en \(\lambda^{-d}e_{-d}\). Posons :
\[
\mathrm{P}(d,d')=\{j\mid |d-d'|\le j\le d+d'\text{ et }j\equiv d\pm d'(2)\},
\tag*{\textbf{(1.6.11.1)}}
\]
\Page{168}et choisissons des isomorphismes de représentations :
\[
\mathrm{S}_d\otimes\mathrm{S}_{d'}
\xrightarrow{\sim}
\bigoplus_{j\in\mathrm{P}(d,d')}\mathrm{S}_j,
\tag*{\textbf{(1.6.11.2)}}
\]
\[
\mathrm{S}_d^*\xrightarrow{\sim}\mathrm{S}_d.
\tag*{\textbf{(1.6.11.3)}}
\]

Pour \(\mathrm{V}\), \(\mathrm{N}\), \(\mathrm{M}\) et \(\mathrm{P}\) comme précédemment, il existe un unique isomorphisme de représentations de \(\mathrm{SL}(2)\) :
\[
\alpha:\bigoplus_j\mathrm{S}_j\otimes\mathrm{P}_{-j}
\xrightarrow{\sim}\operatorname{Gr}^{\mathrm{M}}(\mathrm{V}),
\tag*{\textbf{(1.6.11.4)}}
\]
où \(\mathrm{SL}(2)\) agit trivialement sur \(\mathrm{P}_j\) à la source, et par (1.6.10) au but, tel que le composé :
\[
\begin{tikzcd}[ampersand replacement=\&,column sep=large]
\mathrm{P}_{-j}
  \arrow[r,"{e_{-j}\otimes}"]
  \& \mathrm{S}_j\otimes\mathrm{P}_{-j}
  \arrow[r,"\alpha"]
  \& \operatorname{Gr}^{\mathrm{M}}(\mathrm{V})
\end{tikzcd}
\]
soit l’inclusion identique de \(\mathrm{P}_{-j}\). Ceci fournit des isomorphismes :
\[
\mathrm{P}_{-j}\simeq
\operatorname{Hom}_{\mathrm{SL}(2)}
(\mathrm{S}_j,\operatorname{Gr}^{\mathrm{M}}(\mathrm{V})).
\tag*{\textbf{(1.6.11.5)}}
\]

\textbf{(1.6.12)} Soient maintenant \(\mathrm{V}'\) et \(\mathrm{V}''\) comme en (1.6.10), \(\mathrm{V}=\mathrm{V}'\otimes\mathrm{V}''\) et écrivons \(\mathrm{P}_j\), \(\mathrm{P}'_j\) et \(\mathrm{P}''_j\) pour \(\mathrm{P}_j(\mathrm{V})\), \(\mathrm{P}_j(\mathrm{V}')\) et \(\mathrm{P}_j(\mathrm{V}'')\). D’après (1.6.10), le produit tensoriel des isomorphismes (1.6.11.4) pour \(\mathrm{V}'\) et \(\mathrm{V}''\) nous fournit un isomorphisme :
\[
\bigoplus_{j',j''}
(\mathrm{S}_{j'}\otimes\mathrm{S}_{j''})
\otimes(\mathrm{P}'_{-j'}\otimes\mathrm{P}''_{-j''})
\xrightarrow{\sim}\operatorname{Gr}^{\mathrm{M}}(\mathrm{V}).
\]
Le composant avec les isomorphismes (1.6.11.2) on obtient :
\[
\bigoplus
\left(\mathrm{S}_j\otimes
\bigoplus_{j\in\mathrm{P}(j',j'')}
\mathrm{P}'_{-j'}\otimes\mathrm{P}''_{-j''}\right)
\xrightarrow{\sim}\operatorname{Gr}^{\mathrm{M}}(\mathrm{V}),
\]
d’où enfin un isomorphisme, ne dépendant que des choix (1.6.11.2) :
\[
\mathrm{P}_{-j}\xrightarrow{\sim}
\bigoplus_{j\in\mathrm{P}(j',j'')}
\mathrm{P}'_{-j'}\otimes\mathrm{P}''_{-j''}.
\tag*{\textbf{(1.6.12.1)}}
\]

Pour le passage au dual, on trouve de même des isomorphismes, ne dépendant que des choix (1.6.11.3) :
\[
\mathrm{P}_{-j}(\mathrm{V}^*)\xrightarrow{\sim}\mathrm{P}_{-j}(\mathrm{V})^*.
\tag*{\textbf{(1.6.12.2)}}
\]

\textit{Proposition} \textbf{(1.6.13)}. — \textit{Soient \(\mathrm{V}\) un objet d’une catégorie abélienne, muni d’une filtration finie croissante \(\mathrm{W}\), et \(\mathrm{N}\) un endomorphisme nilpotent de \(\mathrm{V}\) qui respecte \(\mathrm{W}\). Il existe alors au plus une filtration finie croissante \(\mathrm{M}\) de \(\mathrm{V}\) telle que \(\mathrm{N}\mathrm{M}_i\subset\mathrm{M}_{i-2}\) et que \(\mathrm{N}^k\) induise des isomorphismes \(\operatorname{Gr}_{i+k}^{\mathrm{M}}\operatorname{Gr}_i^{\mathrm{W}}\mathrm{V}\xrightarrow{\sim}\operatorname{Gr}_{i-k}^{\mathrm{M}}\operatorname{Gr}_i^{\mathrm{W}}\mathrm{V}\).}

On procède par récurrence sur la longueur de \(\mathrm{W}\). Ceci nous permet de supposer que, pour un entier \(a\), on a \(\mathrm{W}_a=\mathrm{V}\), et que (1.6.13) est vrai pour \(\mathrm{W}_{a-1}\), muni de la filtration induite par \(\mathrm{W}\) et de l’endomorphisme induit par \(\mathrm{N}\). Un décalage simultané, de même amplitude, sur les filtrations \(\mathrm{W}\) et \(\mathrm{M}\) respecte la condition de (1.6.13). Ceci nous permet de simplifier la notation en supposant que \(a=0\).

% ---- 169-176 ----
\Page{169}
Soit \(\mathrm{M}\) une filtration vérifiant les conditions de (1.6.13). La filtration induite sur \(\mathrm{W}_{-1}\) vérifie les mêmes conditions, donc est uniquement déterminée. La filtration quotient est la filtration (1.6.1) de \(\operatorname{Gr}_0^{\mathrm{W}}\mathrm{V}\). Soit \(c\) un entier tel que \(\operatorname{Gr}_i^{\mathrm{M}}\operatorname{Gr}_0^{\mathrm{W}}\mathrm{V}\) ne soit non nul que pour \(-c\leq i\leq c\). Vérifions les identités suivantes, pour \(i\geq0\) :

\noindent 1) pour \(i>c\), on a \(\mathrm{M}_{-i}=\mathrm{M}_{-i}(\mathrm{W}_{-1})\);

\noindent 2) \(\mathrm{M}_{-i}=\mathrm{M}_{-i}(\mathrm{W}_{-1})+\mathrm{N}^i(\mathrm{M}_i)\);

\noindent 3) \(\mathrm{M}_i=\operatorname{Ker}(\mathrm{N}^{i+1}:\mathrm{V}\to\mathrm{V}/\mathrm{M}_{-i-2})\).

\noindent L’assertion 1) exprime que \(\mathrm{M}_{-i}\operatorname{Gr}_0^{\mathrm{W}}\mathrm{V}=0\) pour \(i>c\). Dans 2), l’inclusion \(\supset\) résulte de l’hypothèse \(\mathrm{N}\mathrm{M}_j\subset\mathrm{M}_{j-2}\). L’inclusion \(\subset\) résulte de ce que l’image de \(\mathrm{N}^i\mathrm{M}_i\) dans \(\operatorname{Gr}_0^{\mathrm{W}}(\mathrm{V})\) est \(\mathrm{N}^i\mathrm{M}_i(\operatorname{Gr}_0^{\mathrm{W}}\mathrm{V})=\mathrm{M}_{-i}(\operatorname{Gr}_0^{\mathrm{W}}\mathrm{V})\). Dans 3), l’inclusion \(\subset\) résulte de l’hypothèse \(\mathrm{N}\mathrm{M}_j\subset\mathrm{M}_{j-2}\). Il suffit donc de prouver l’égalité après passage aux \(\operatorname{Gr}_j^{\mathrm{W}}\), avec \(j\leq0\), et pour cela de prouver l’inclusion \(\supset\) après passage aux \(\operatorname{Gr}_j^{\mathrm{W}}\). Posons \(\mathrm{G}=\operatorname{Gr}_j^{\mathrm{W}}\mathrm{V}\), et notons encore \(\mathrm{M}\) la filtration induite par \(\mathrm{M}\) sur \(\mathrm{G}\). On a :
\[
\operatorname{Gr}_j^{\mathrm{W}}\operatorname{Ker}
(\mathrm{N}^{i+1}:\mathrm{V}\to\mathrm{V}/\mathrm{M}_{-i-2})
\subset
\operatorname{Ker}
(\mathrm{N}^{i+1}:\mathrm{G}\to\mathrm{G}/\mathrm{M}_{-i-2}),
\]
et ce dernier noyau coïncide avec \(\mathrm{M}_i\mathrm{G}=\operatorname{Gr}_j^{\mathrm{W}}(\mathrm{M}_i)\) : l’application \(\mathrm{N}^{i+1}\) de \(\mathrm{G}/\mathrm{M}_i\) dans \(\mathrm{G}/\mathrm{M}_{-i-2}\) est injective, car son \(\operatorname{Gr}^{\mathrm{M}}\) l’est, i.e. les applications \(\mathrm{N}^{i+1}:\operatorname{Gr}_k^{\mathrm{M}}\mathrm{G}\to\operatorname{Gr}_{k-2i-2}^{\mathrm{M}}\mathrm{G}\) pour \(k>i\) le sont : l’application itérée de \(\mathrm{N}\) de \(\operatorname{Gr}_k^{\mathrm{M}}\) dans \(\operatorname{Gr}_{2j-k}^{\mathrm{M}}\) est par hypothèse un isomorphisme, et \(k-2i-2\geq -k\geq2j-k\).

Enfin, une récurrence descendante sur \(|i|\) montre que les identités 1), 2), 3) déterminent uniquement \(\mathrm{M}\), en terme de sa restriction à \(\mathrm{W}_{-1}\) et de \(\mathrm{N}\).

\textbf{(1.6.14)} Dans les applications, la catégorie abélienne considérée sera celle des espaces vectoriels sur un corps \(k\) de caractéristique \(0\), et nous disposerons non pas d’un endomorphisme nilpotent \(\mathrm{N}\) de \(\mathrm{V}\), mais d’une action nilpotente d’une algèbre de Lie \(\mathfrak{N}\) de dimension \(1\). La théorie précédente s’applique dès qu’on choisit \(\mathrm{N}\neq0\) dans \(\mathfrak{N}\). Dans (1.6.13), l’existence de la filtration \(\mathrm{M}\), et cette filtration même, si elle existe, sont indépendantes du choix de \(\mathrm{N}\). L’isomorphisme \(\mathrm{N}^k\) de \(\operatorname{Gr}_{i+k}^{\mathrm{M}}\operatorname{Gr}_i^{\mathrm{W}}\mathrm{V}\) avec \(\operatorname{Gr}_{i-k}^{\mathrm{M}}\operatorname{Gr}_i^{\mathrm{W}}\mathrm{V}\) a pour avatar un isomorphisme canonique de \(\operatorname{Gr}_{i+k}^{\mathrm{M}}\operatorname{Gr}_i^{\mathrm{W}}\mathrm{V}\otimes\mathfrak{N}^{\otimes k}\) avec \(\operatorname{Gr}_{i-k}^{\mathrm{M}}\operatorname{Gr}_i^{\mathrm{W}}\mathrm{V}\). Pour tout espace vectoriel \(\mathrm{Y}\), écrivons \(\mathrm{Y}(k)\) pour \(\mathrm{Y}\otimes\mathfrak{N}^{\otimes k}\). L’isomorphisme s’écrit :
\[
\operatorname{Gr}_{i+k}^{\mathrm{M}}\operatorname{Gr}_i^{\mathrm{W}}\mathrm{V}(k)
\xrightarrow{\sim}
\operatorname{Gr}_{i-k}^{\mathrm{M}}\operatorname{Gr}_i^{\mathrm{W}}\mathrm{V}.
\tag*{\textbf{(1.6.14.1)}}
\]

Plaçons-nous dans la situation étudiée de (1.6.1) à (1.6.12). Ici, la filtration \(\mathrm{M}\) et \(\mathrm{P}_j\) ne dépendent pas du choix de \(\mathrm{N}\), mais les isomorphismes (1.6.4.1) et (1.6.12.1 et 2) en dépendent. Si \(\mathrm{N}\) est remplacé par \(\lambda\mathrm{N}\), chacune de leurs composantes, notée génériquement \(\alpha:\mathrm{X}\to\mathrm{Y}\), est multipliée par une puissance convenable \(\lambda^{-k}\) de \(\lambda\). Elle définit alors une application \(\alpha\otimes\mathrm{N}^{\otimes k}:\mathrm{X}\to\mathrm{Y}(k)\) qui, elle, est indépendante du choix de \(\mathrm{N}\). Les isomorphismes définis précédemment deviennent des isomorphismes indépendants du choix de \(\mathrm{N}\) :
\Page{170}
\[
\operatorname{Gr}_i^{\mathrm{M}}(\operatorname{Ker}\mathrm{N})
\simeq\mathrm{P}_i
\qquad\text{(pour (1.6.6));}
\tag*{\textbf{(1.6.14.2)}}
\]
\[
\operatorname{Gr}_i^{\mathrm{M}}\mathrm{V}
\simeq
\bigoplus_{\substack{j\geq|i|\\j\equiv i(2)}}
\mathrm{P}_{-j}\left(\frac{i+j}{2}\right)
\qquad\text{(pour (1.6.4.1));}
\tag*{\textbf{(1.6.14.3)}}
\]
\begin{align*}
\mathrm{P}_{-j}
&\simeq
\bigoplus_{j\in\mathrm{P}(j'j'')}
\mathrm{P}'_{-j'}\otimes\mathrm{P}''_{-j''}
\left(\frac{j-j'-j''}{2}\right)
\tag*{\textbf{(1.6.14.4)}}\\
&\qquad\text{(pour (1.6.12.1); }\mathrm{P}(j',j'')
\text{ défini par (1.6.11.1));}
\end{align*}
\[
\mathrm{P}_{-j}(\mathrm{V}^*)
\simeq\mathrm{P}_{-j}(\mathrm{V})^*(j)
\qquad\text{(pour (1.6.12.2)).}
\tag*{\textbf{(1.6.14.5)}}
\]

Pour se rappeler ces formules : si on regarde \(\operatorname{Gr}_i\) et \(\mathrm{P}_i\) comme de poids \(i\), et \(\mathfrak{N}\) comme de poids \(-2\), les deux membres sont toujours isobares du même poids.

Nous utiliserons aussi une variante faisceautique de (1.6.14).

\section*{(1.7) Monodromie locale}

\textbf{(1.7.1)} Soient \(\mathrm{R}\) un anneau de valuation discrète hensélien, \(\mathrm{K}\) son corps des fractions, \(k=\mathrm{R}/\mathfrak{m}\) son corps résiduel, \(\overline{\mathrm{K}}\) une clôture algébrique de \(\mathrm{K}\), et \(\bar{k}\) le corps résiduel du normalisé de \(\mathrm{R}\) dans \(\overline{\mathrm{K}}\). L’extension \(\bar{k}\) de \(k\) en est une clôture algébrique. Il est souvent commode de choisir d’abord \(\bar{k}\), et ensuite \(\overline{\mathrm{K}}\), comme suit : on prend une clôture algébrique \(\bar{k}\) de \(k\), puis une clôture algébrique \(\overline{\mathrm{K}}\) du corps des fractions de l’hensélisé strict de \(\mathrm{R}\), rel. \(\bar{k}\) (cf. (0.6)).

Le groupe d’inertie est défini comme un noyau :
\[
0\to\mathrm{I}\to
\operatorname{Gal}(\overline{\mathrm{K}}/\mathrm{K})
\to\operatorname{Gal}(\bar{k}/k)\to0.
\tag*{\textbf{(1.7.1.1)}}
\]
Lui-même, ou son image dans une représentation, s’appellent encore groupe de monodromie locale. Rappelons que si \(k\) est d’exposant caractéristique \(p\), on dispose d’une suite exacte ([10], [5], \S~2) :
\[
0\to\mathrm{P}\to\mathrm{I}
\xrightarrow{t}\widehat{\mathbf{Z}}_{(p')}(1)\to0
\]
où \(\mathrm{P}\) est un pro-\(p\)-groupe. Le quotient \(\widehat{\mathbf{Z}}_{(p')}(1)=\prod\limits_{\ell\neq p}\mathbf{Z}_\ell(1)\) est le groupe d’inertie modéré. Prenant la \(\ell\)-composante de \(t\), on trouve :
\[
t_\ell:\mathrm{I}\to\mathbf{Z}_\ell(1)
\]
de noyau un groupe pro-fini d’ordre premier à \(\ell\).

\textbf{(1.7.2)} Soit \(\rho:\mathrm{I}\to\operatorname{GL}(\mathrm{V})\) une représentation \(\ell\)-adique de \(\mathrm{I}\). Grothendieck a montré que l’hypothèse suivante est souvent vérifiée en pratique (voir SGA 7 I).

\((*)\) La restriction de \(\rho\) à un sous-groupe d’indice fini \(\mathrm{I}_1\) de \(\mathrm{I}\) est unipotente.

Si elle l’est, il existe une unique représentation nilpotente \(\bar{\rho}\) de \(\overline{\mathbf{Q}}_\ell(1)\), vue comme algèbre de Lie de dimension \(1\), dans \(\mathrm{V}\), telle qu’on ait :
\[
\rho(\sigma)=\exp(\bar{\rho}t_\ell(\sigma))
\qquad\text{pour}\qquad\sigma\in\mathrm{I}_1\subset\mathrm{I}.
\tag*{\textbf{(1.7.2.1)}}
\]
\Page{171}
On appelle \(\bar{\rho}\) le \textit{logarithme de la partie unipotente de la monodromie locale}. On peut identifier \(\bar{\rho}\) à un morphisme nilpotent \(\mathrm{N}:\mathrm{V}(1)\to\mathrm{V}\). La formule de définition s’écrit alors :
\[
\rho(\sigma)=\exp(\mathrm{N}.t_\ell(\sigma))
\qquad\text{pour}\qquad\sigma\in\mathrm{I}_1\subset\mathrm{I}.
\tag*{\textbf{(1.7.2.2)}}
\]

La théorie (1.6.14) s’applique à cette situation. Le twist \((k)\) de loc. cit. n’est autre qu’un twist à la Tate. Les filtrations (1.6.1) et (1.6.13) de \(\mathrm{V}\) s’appelleront respectivement \textit{filtration de monodromie locale de} \(\mathrm{V}\) et \textit{filtration de monodromie locale de} \(\mathrm{V}\), \textit{rel.} \(\mathrm{W}\).

\textbf{(1.7.3)} Supposons le corps résiduel fini à \(q\) éléments. On définit le groupe de Weil \(\mathrm{W}(\overline{\mathrm{K}}/\mathrm{K})\) comme le sous-groupe de \(\operatorname{Gal}(\overline{\mathrm{K}}/\mathrm{K})\) image inverse de \(\mathrm{W}(\bar{k}/k)\).

Soit \(\mathrm{V}\) une représentation \(\ell\)-adique de \(\mathrm{W}(\overline{\mathrm{K}}/\mathrm{K})\). D’après Grothendieck (SGA 7 I), la condition \((*)\) de (1.7.2) est automatiquement vérifiée. Le logarithme de la partie unipotente de la monodromie \(\mathrm{N}:\mathrm{V}(1)\to\mathrm{V}\) est donc défini. Il commute à l’action de \(\mathrm{W}(\overline{\mathrm{K}}/\mathrm{K})\).

\textit{Lemme} \textbf{(1.7.4)}. — \textit{Si \(\mathrm{F}'\) et \(\mathrm{F}''\) sont deux relèvements dans \(\mathrm{W}(\overline{\mathrm{K}}/\mathrm{K})\) de \(\mathrm{F}\in\mathrm{W}(\bar{k}/k)\), les valeurs propres de \(\mathrm{F}'\) et \(\mathrm{F}''\) coïncident, à multiplication près par des racines de l’unité.}

Pour le vérifier, on peut remplacer \(\mathrm{V}\) par sa semi-simplifiée, donc supposer que l’action du groupe d’inertie \(\mathrm{I}\) se factorise par un quotient fini. Dans ce cas, il existe \(n>0\) tel que \(\mathrm{F}^{\prime n}\) et \(\mathrm{F}^{\prime\prime n}\) aient même image dans \(\operatorname{GL}(\mathrm{V})\), et le lemme en résulte.

Quel que soit \(\iota\), les \(\iota\)-poids des valeurs propres de \(\mathrm{F}'\), ou de \(\mathrm{F}''\) (rel. \(q\)) sont donc les mêmes. A l’imitation de (1.2), on définit alors ce qu’est une représentation pure, mixte, \(\iota\)-pure, et les poids d’une représentation.

\textit{Proposition-définition} \textbf{(1.7.5)}. — \textit{Supposons que les \(\iota\)-poids de \(\mathrm{V}\) soient entiers. Il existe alors une et une seule filtration finie croissante \(\mathrm{M}\) de \(\mathrm{V}\), stable par \(\mathrm{W}(\overline{\mathrm{K}}/\mathrm{K})\), telle que \(\operatorname{Gr}_i^{\mathrm{M}}(\mathrm{V})\) soit \(\iota\)-pur de poids \(i\). On l’appelle la \textup{filtration par le poids}. On a \(\mathrm{N}\mathrm{M}_i(1)\subset\mathrm{M}_{i-2}\).}

Choisissons un relèvement \(\mathrm{F}'\) de Frobenius dans \(\mathrm{W}(\overline{\mathrm{K}}/\mathrm{K})\). Soit \(\mathrm{V}'_i\) la somme des sous-espaces propres généralisés de \(\mathrm{F}'\) dans \(\mathrm{V}\) correspondant aux valeurs propres \(\alpha\) telles que \(w_q(\alpha)=i\), et posons \(\mathrm{M}'_i=\prod\limits_{j\leq i}\mathrm{V}'_j\). On a \(\mathrm{N}\mathrm{M}'_i\subset\mathrm{M}'_{i-2}\). Il s’agit de prouver que \(\mathrm{M}'\) est indépendant de \(\mathrm{F}'\) : cela assurera que \(\mathrm{M}'\) est stable par \(\mathrm{W}(\overline{\mathrm{K}}/\mathrm{K})\), et on aura \(\mathrm{M}=\mathrm{M}'\). Soit donc \(\mathrm{M}''\) la filtration définie par un relèvement \(\mathrm{F}''\). Soit \(\mathrm{I}_1\) comme dans (1.7.2.1). Puisque \(\mathrm{I}/\mathrm{I}_1\) est fini, il existe \(n>0\) tel que \(\mathrm{F}^{\prime n}=\mathrm{F}^{\prime\prime n}\bmod\mathrm{I}_1\). Pour \(\lambda\) convenable dans \(\mathbf{Q}_\ell(1)\), on a alors dans \(\operatorname{GL}(\mathrm{V})\) :
\[
\mathrm{F}^{\prime\prime n}=\exp(\lambda\mathrm{N})\mathrm{F}^{\prime n},
\]
donc \(\mathrm{F}^{\prime\prime n}=\exp(\mu\mathrm{N})\mathrm{F}^{\prime n}\exp(\mu\mathrm{N})^{-1}\) si \(\mu=\lambda/(1-q^{-n})\).

Puisque \(\mathrm{V}'_i\) est encore la somme des sous-espaces propres généralisés de \(\mathrm{F}^{\prime n}\), relatifs aux valeurs propres \(\alpha\) telles que \(w_{q^n}(\alpha)=i\), l’automorphisme \(\exp(\mu\mathrm{N})\) transforme \(\mathrm{M}'\) en \(\mathrm{M}''\). Puisque \(\mathrm{N}\) respecte \(\mathrm{M}'\), on a gagné.

\Page{172}
\textit{Corollaire} \textbf{(1.7.6)}. — \textit{Si \(\mathrm{V}\) est \(\iota\)-pur, l’action de \(\mathrm{I}\) se factorise par un quotient fini.}

Par torsion, on se ramène à supposer \(\mathrm{V}\) de poids entier. La filtration \(\mathrm{M}\) est par hypothèse réduite à un seul cran. Puisque \(\mathrm{N}\mathrm{M}_i\subset\mathrm{M}_{i-2}\), on a donc \(\mathrm{N}=0\), et (1.7.6) en résulte.

\textit{Remarque} \textbf{(1.7.7)}. — La preuve de (1.7.5) montre que n’importe quelle représentation \(\mathrm{V}\) admet une unique décomposition \(\mathrm{V}=\sum\mathrm{V}^{\alpha}\) \((\alpha\in\overline{\mathbf{Q}}_\ell^*/(q^{\mathbf{Z}}.\text{racines de }1))\), telle que les valeurs propres des relèvements de Frobenius dans \(\mathrm{V}^{\alpha}\) soient dans \(\alpha\). Chaque \(\mathrm{V}^{\alpha}\) est obtenu par torsion à partir d’une représentation mixte.

\textit{Variante} \textbf{(1.7.8)}. — Si \(\mathrm{D}\) est un diviseur lisse dans un schéma régulier \(\mathrm{X}\), les constructions qui précèdent ont un analogue pour les faisceaux sur \(\mathrm{X}-\mathrm{D}\), \textit{modérément ramifiés} le long de \(\mathrm{D}\). Un peu plus généralement, soit \(\mathrm{D}\) un diviseur à croisements normaux de \(\mathrm{X}\), réunion transverse de diviseurs lisses \(\mathrm{D}_i\) \((i\in\mathrm{I})\). On se localise, pour que chaque \(\mathrm{D}_i\) admette une équation globale \(t_i=0\). Soit \(\mathrm{E}\) l’intersection des \(\mathrm{D}_i\). Pour \(n\) inversible sur \(\mathrm{X}\), posons \(\mathrm{X}_n=\mathrm{X}[t_i^{1/n}]_{i\in\mathrm{I}}\) (le sous-schéma de l’espace affine sur \(\mathrm{X}\), de coordonnées \(\mathrm{T}_i\), défini par les équations \(\mathrm{T}_i^n=t_i\)). Le revêtement ramifié \(\pi:\mathrm{X}_n\to\mathrm{X}\) de \(\mathrm{X}\) est étale au-dessus de \(\mathrm{X}-\mathrm{D}\), et totalement ramifié au-dessus de \(\mathrm{E}\); au-dessus de \(\mathrm{E}\), \(\pi\) admet la section \(t_i^{1/n}=0\); on appelle restriction à \(\mathrm{E}\) d’un faisceau sur \(\mathrm{X}_n\) son image inverse par cette section. On définit une action de \((\mu_n)^{\mathrm{I}}\) sur \(\mathrm{X}_n\) en faisant agir \((r_i)_{i\in\mathrm{I}}\) sur \(t_i^{1/n}\) par multiplication par \(r_i\).

Soient \(\mathbf{L}\) un ensemble de nombres premiers inversibles sur \(\mathrm{X}\), et \(\mathscr{F}\) un faisceau d’ensembles sur \(\mathrm{X}-\mathrm{D}\), localement constant, modérément ramifié le long de \(\mathrm{D}\), et même \(\mathbf{L}\)-ramifié le long de \(\mathrm{D}\), au sens suivant : pour \(d\) le point générique d’une composante irréductible de \(\mathrm{D}\), l’hensélisé \(\mathrm{X}_{(d)}\) est un trait hensélien; si \(\bar{\eta}\) en est un point générique géométrique, on demande que le groupe d’inertie agisse sur \(\mathscr{F}_{\bar{\eta}}\) via un \(\mathbf{L}\)-groupe. Le lemme d’Abhyankar assure alors l’existence d’un entier \(n\), dont tous les facteurs premiers sont dans \(\mathbf{L}\), tel que l’image inverse de \(\mathscr{F}\) sur \(\pi^{-1}(\mathrm{X}-\mathrm{D})\subset\mathrm{X}_n\) soit la restriction à \(\pi^{-1}(\mathrm{X}-\mathrm{D})\) d’un faisceau localement constant \(\overline{\mathscr{F}}\) sur \(\mathrm{X}_n\), que nous appellerons le prolongement localement constant de \(\mathscr{F}\) à \(\mathrm{X}_n\). Nous noterons \(\mathscr{F}[\mathrm{E}]\) la restriction de \(\overline{\mathscr{F}}\) à \(\mathrm{E}\). Elle est munie d’une action de \(\mu_n^{\mathrm{I}}\), déduite de l’action de \(\mu_n^{\mathrm{I}}\) sur \(\mathrm{X}_n\). Posant \(\mathbf{Z}_{\mathbf{L}}(1)=\prod\limits_{\ell\in\mathbf{L}}\mathbf{Z}_\ell(1)_{\mathrm{X}}\) (un pro-faisceau sur \(\mathrm{E}\)), on obtient ainsi un foncteur \(\mathscr{F}\to\mathscr{F}[\mathrm{E}]\) des faisceaux d’ensembles sur \(\mathrm{X}-\mathrm{D}\), localement constants et \(\mathbf{L}\)-ramifiés le long de \(\mathrm{D}\), dans les faisceaux d’ensembles, localement constants sur \(\mathrm{E}\), munis d’une action de \(\mathbf{Z}_{\mathbf{L}}(1)^{\mathrm{I}}\).
Ce foncteur passe aux \(\overline{\mathbf{Q}}_\ell\)-faisceaux constructibles par passages à la limite (et on a une variante pour les faisceaux de Weil). Dans bien des cas (voir (1.7.12.1)), et notamment si \(\mathrm{X}\) est de type fini sur \(\mathbf{Z}\), l’action de \(\widehat{\mathbf{Z}}_{\mathbf{L}}(1)^{\mathrm{I}}\) est quasi-unipotente : il existe des endomorphismes nilpotents commutant entre eux \(\mathrm{N}_i\) dans \((\operatorname{End}(\mathscr{F}[\mathrm{E}]))(-1)\) tels que, si \(\sigma\) est dans un sous-faisceau « ouvert » \(n\mathbf{Z}_{\mathbf{L}}(1)^{\mathrm{I}}\) convenable de \(\mathbf{Z}_{\mathbf{L}}(1)^{\mathrm{I}}\), il agisse par \(\exp(\sum\mathrm{N}_i\sigma_i)\). On appellera \(\mathrm{N}_i\) le « logarithme de la partie unipotente de la monodromie,
\Page{173}
autour de \(\mathrm{D}_i\) ». Dans le cas d’un diviseur lisse \((|\mathrm{I}|=1)\), on utilise cet endomorphisme comme précédemment pour définir la filtration de monodromie locale de \(\mathscr{F}[\mathrm{D}]\) (resp. rel. \(\mathrm{W}\)).

\textbf{(1.7.9)} Au voisinage d’une intersection \(p\) à \(p\) des \(\mathrm{D}_j\), \(\mathrm{D}\) apparaît comme la réunion de \(p\) diviseurs lisses, et on peut refaire les constructions qui précèdent. Quelques notations : pour \(\mathrm{J}\subset\mathrm{I}\), \(\mathrm{D}_{\mathrm{J}}\) est l’intersection des \(\mathrm{D}_j\) \((j\in\mathrm{J})\), \(\mathrm{X}_{\mathrm{J}}\) est le complément de la réunion des \(\mathrm{D}_i\) \((i\notin\mathrm{J})\) et \(\mathrm{D}_{\mathrm{J}}^*=\mathrm{D}_{\mathrm{J}}\cap\mathrm{X}_{\mathrm{J}}\). On a : (a) les \(\mathrm{D}_i\) \((i\notin\mathrm{J})\) découpent sur le schéma régulier \(\mathrm{D}_{\mathrm{J}}\) un diviseur à croisements normaux, de complément \(\mathrm{D}_{\mathrm{J}}^*\); (b) les \(\mathrm{D}_j\) \((j\in\mathrm{J})\) découpent sur \(\mathrm{X}_{\mathrm{J}}\) un diviseur à croisements normaux, de complément \(\mathrm{X}-\mathrm{D}\); l’intersection des traces des \(\mathrm{D}_j\) est \(\mathrm{D}_{\mathrm{J}}^*\).

Soit donc \(\mathscr{F}\) un faisceau d’ensembles sur \(\mathrm{X}-\mathrm{D}\), localement constant et \(\mathbf{L}\)-ramifié le long de \(\mathrm{D}\). Appliquant (1.7.8) à (b) ci-dessus, on obtient pour chaque \(\mathrm{J}\) un faisceau \(\mathscr{F}[\mathrm{D}_{\mathrm{J}}^*]\) sur \(\mathrm{D}_{\mathrm{J}}^*\), muni d’une action de \(\mathbf{Z}_{\mathbf{L}}^{\mathrm{J}}(1)\). Pour \(n\) convenable (à facteurs premiers dans \(\mathbf{L}\)), c’est la restriction à \(\mathrm{D}_{\mathrm{J}}^*\) du prolongement localement constant de \(\mathscr{F}\) à \(\mathrm{X}_{\mathrm{J}}[t_j^{1/n}]\) \((j\in\mathrm{J})\). Choisissons \(n\) de sorte que \(\mathscr{F}\) admette un prolongement localement constant à \(\mathrm{X}_n=\mathrm{X}[t_j^{1/n}(j\in\mathrm{J})][t_i^{1/n}(i\in\mathrm{I}-\mathrm{J})]\). Restreignant ce prolongement à l’image inverse réduite \(\mathrm{D}_{\mathrm{J}}[t_i^{1/n}(i\in\mathrm{I}-\mathrm{J})]\) de \(\mathrm{D}_{\mathrm{J}}\), on voit que \(\mathscr{F}[\mathrm{D}_{\mathrm{J}}^*]\) admet un prolongement localement constant sur \(\mathrm{D}_{\mathrm{J}}[t_i^{1/n}(i\in\mathrm{I}-\mathrm{J})]\) : le faisceau \(\mathscr{F}[\mathrm{D}_{\mathrm{J}}^*]\) est \(\mathbf{L}\)-ramifié le long de la trace des \(\mathrm{D}_i\) \((i\notin\mathrm{J})\). Ceci permet d’itérer la construction et, pour \(\mathrm{K}\subset\mathrm{I}-\mathrm{J}\), de définir \(\mathscr{F}[\mathrm{D}_{\mathrm{J}}^*][\mathrm{D}_{\mathrm{J}\cup\mathrm{K}}^*]\) : un faisceau localement constant sur \(\mathrm{D}_{\mathrm{J}\cup\mathrm{K}}^*\), muni d’une action de \(\mathbf{Z}_{\mathbf{L}}^{\mathrm{J}}(1)\) (par fonctorialité), et d’une action de \(\mathbf{Z}_{\mathbf{L}}^{\mathrm{K}}(1)\). La transitivité des images inverses fournit un isomorphisme
\[
\mathscr{F}[\mathrm{D}_{\mathrm{J}}^*][\mathrm{D}_{\mathrm{J}\cup\mathrm{K}}^*]
\simeq
\mathscr{F}[\mathrm{D}_{\mathrm{J}\cup\mathrm{K}}^*],
\tag*{\textbf{(1.7.9.1)}}
\]
compatible aux actions de \(\mathbf{Z}_{\mathbf{L}}(1)^{\mathrm{J}\cup\mathrm{K}}\). Pour \(\mathrm{J}\subset\mathrm{K}\subset\mathrm{L}\), le carré :
\[
\begin{tikzcd}[ampersand replacement=\&, row sep=large, column sep=large]
\mathscr{F}[\mathrm{D}_{\mathrm{J}}^*][\mathrm{D}_{\mathrm{K}}^*][\mathrm{D}_{\mathrm{L}}^*]
  \arrow[r, equal] \arrow[d, equal]
\&
\mathscr{F}[\mathrm{D}_{\mathrm{K}}^*][\mathrm{D}_{\mathrm{L}}^*]
  \arrow[d, equal]
\\
\mathscr{F}[\mathrm{D}_{\mathrm{J}}^*][\mathrm{D}_{\mathrm{L}}^*]
  \arrow[r, equal]
\&
\mathscr{F}[\mathrm{D}_{\mathrm{L}}^*]
\end{tikzcd}
\]
est commutatif.

Tout ceci s’étend aux \(\overline{\mathbf{Q}}_\ell\)-faisceaux constructibles (et aux faisceaux de Weil) par passages à la limite.

\textbf{(1.7.10)} Les constructions (1.7.8) dépendent du choix des coordonnées \(t_i\). Indiquons brièvement comment on peut les canonifier : on introduit le fibré normal \(\mathrm{T}\) de \(\mathrm{E}\) dans \(\mathrm{X}\) (un fibré en espaces affines sur \(\mathrm{E}\)) et les sous-fibrés \(\mathrm{T}_i\) tangents aux \(\mathrm{D}_i\). On remplace la construction de \(\mathscr{F}[\mathrm{E}]\) par celle d’un faisceau \(\mathscr{F}[\mathrm{E}]'\) sur \(\mathrm{T}-\cup\mathrm{T}_i\).
\Page{174}
Un système d’équations \(t_i\) définit une section de \(\mathrm{T}-\cup\mathrm{T}_i\), et le faisceau \(\mathscr{F}[\mathrm{E}]\) défini en terme des \(t_i\) s’identifie à l’image inverse de \(\mathscr{F}[\mathrm{E}]'\) par cette section. La construction est par ailleurs compatible au changement de base par un morphisme lisse \(\mathrm{X}_1\to\mathrm{X}\). Les \(\mathrm{N}_i\), et \(\mathbf{Z}_{\mathbf{L}}(1)^{\mathrm{I}}\), agissent sur \(\mathscr{F}[\mathrm{E}]'\), qui joue le rôle de la restriction de \(\mathscr{F}\) à un voisinage tubulaire de \(\mathrm{E}\) dans \(\mathrm{X}\). On peut aussi le voir comme un analogue des « nilpotent orbits » en théorie des variations de structures de Hodge.

\textbf{(1.7.11)} Supposons \(\mathrm{X}\) hensélien local, \(\mathrm{X}=\operatorname{Spec}(\mathrm{R})\), et prenons pour \(\mathbf{L}\) l’ensemble de tous les nombres premiers inversibles sur \(\mathrm{X}\). Nous allons interpréter en termes galoisiens la construction (1.7.8). Soient donc \(k\) le corps résiduel de \(\mathrm{R}\), \(\mathrm{K}\) son corps des fractions, \(\mathrm{R}_1\) un hensélisé strict de \(\mathrm{R}\), de corps résiduel \(k_1\) (une clôture séparable de \(k\)) et de corps des fractions \(\mathrm{K}_1\), \(\overline{\mathrm{K}}\) une clôture algébrique de \(\mathrm{K}_1\) et \(\mathrm{K}_2\subset\overline{\mathrm{K}}\) l’extension de \(\mathrm{K}_1\) obtenue en adjoignant à \(\mathrm{K}_1\) les racines \(n\)-ièmes des \(t_i\), pour \(n\) inversible dans \(\mathrm{R}\). Pour qu’un faisceau d’ensembles sur \(\mathrm{X}-\mathrm{D}\), localement constant, soit modérément ramifié le long de \(\mathrm{D}\), il faut et il suffit qu’il devienne trivial sur \(\operatorname{Spec}(\mathrm{K}_2)\), et ceci identifie le groupe fondamental modéré de \(\mathrm{X}-\mathrm{D}\) (rel. au point base \(\operatorname{Spec}(\overline{\mathrm{K}})\)) à \(\operatorname{Gal}(\mathrm{K}_2/\mathrm{K})\). C’est une extension de \(\operatorname{Gal}(\mathrm{K}_1/\mathrm{K})=\operatorname{Gal}(k_1/k)\) par \(\operatorname{Gal}(\mathrm{K}_2/\mathrm{K}_1)=\widehat{\mathbf{Z}}_{\mathbf{L}}(1)^{\mathrm{I}}\) :
\[
0\to\mathbf{Z}_{\mathbf{L}}(1)^{\mathrm{I}}
\to\pi_1^{\mathrm{mod}}(\mathrm{X}-\mathrm{D},\operatorname{Spec}(\overline{\mathrm{K}}))
\to\operatorname{Gal}(k_1/k)\to0.
\]

Le choix des \(t_i\) et des \(t_i^{1/n}\) scinde cette extension : prendre pour relèvement de \(\operatorname{Gal}(k_1/k)\) le fixateur des \(t_i^{1/n}\). D’un faisceau d’ensemble \(\mathscr{F}\) sur \(\mathrm{X}-\mathrm{D}\), localement constant et modérément ramifié le long de \(\mathrm{D}\), identifié à un ensemble \(\mathrm{F}\) sur lequel \(\pi_1^{\mathrm{mod}}(\mathrm{X}-\mathrm{D},\operatorname{Spec}(\overline{\mathrm{K}}))\) agit, on déduit donc une action de \(\operatorname{Gal}(k_1/k)\) sur \(\mathrm{F}\), et une action \(\operatorname{Gal}(k_1/k)\)-équivariante de \(\widehat{\mathbf{Z}}_{\mathbf{L}}(1)^{\mathrm{I}}\). Ceci définit un faisceau sur \(\operatorname{Spec}(k)\) (ou un faisceau localement constant sur \(\mathrm{E}\), cela revient au même), muni d’une action de \(\mathbf{Z}_{\mathbf{L}}(1)^{\mathrm{I}}\). C’est le faisceau \(\mathscr{F}[\mathrm{E}]\) de (1.7.8).

\textbf{(1.7.12)} Si \(k\) est fini, on définit \(\mathrm{W}^{\mathrm{mod}}(\mathrm{X}-\mathrm{D},\operatorname{Spec}(\overline{\mathrm{K}}))\) comme l’image inverse de \(\mathrm{W}(k_1/k)\) dans \(\pi_1^{\mathrm{mod}}(\mathrm{X}-\mathrm{D},\operatorname{Spec}(\overline{\mathrm{K}}))\). On a, généralisant (1.7.3) à (1.7.7) (et avec la même démonstration) :

\textbf{(1.7.12.1)}La restriction à \(\mathbf{Z}_{\mathbf{L}}(1)^{\mathrm{I}}\) d’une représentation \(\ell\)-adique \(\mathrm{V}\) de \(\mathrm{W}^{\mathrm{mod}}(\mathrm{X}-\mathrm{D},\operatorname{Spec}(\overline{\mathrm{K}}))\) est quasi-unipotente.

\textbf{(1.7.12.2)} Si \(\mathrm{F}'\) et \(\mathrm{F}''\) sont deux relèvements dans \(\mathrm{W}^{\mathrm{mod}}\) de \(\mathrm{F}\in\mathrm{W}(k_1/k)\), les valeurs propres de \(\mathrm{F}'\) et \(\mathrm{F}''\) coïncident, à multiplication près par des racines de l’unité. Ceci permet de définir les \(\iota\)-poids..., comme en (1.7.4).

\textbf{(1.7.12.3)} Si les \(\iota\)-poids de \(\mathrm{V}\) sont entiers, il existe sur \(\mathrm{V}\) une et une seule filtration par le poids \(\mathrm{M}\), comme en (1.7.5). Les analogues de (1.7.6) et (1.7.7) sont vrais.

\Page{175}
\section*{(1.8) Monodromie locale des faisceaux purs}

\textit{Les conventions \textup{(0.7)} sont en vigueur.}

Soient \(\mathrm{X}_0\) une courbe lisse absolument irréductible sur \(\mathbf{F}_q\) (cf. (0.5)) et \(j:\mathrm{U}_0\to\mathrm{X}_0\) un ouvert de \(\mathrm{X}_0\), complément de l’ensemble fini de points \(\mathrm{S}_0\).

\textit{Lemme} \textbf{(1.8.1)}. — \textit{Si \(\mathscr{F}_0\) est un faisceau lisse ponctuellement \(\iota\)-pur de poids \(\beta\) sur \(\mathrm{U}_0\), alors, pour tout \(x\in|\mathrm{S}_0|\) et toute valeur propre \(\alpha\) de \(\mathrm{F}_x\) sur \(j_*\mathscr{F}_0\), on a \(w_{\mathrm{N}(x)}(\alpha)\leq\beta\).}

Quitte à ôter de \(\mathrm{X}_0\) un point dans \(\mathrm{U}_0\), on se ramène à supposer \(\mathrm{X}_0\) affine. On a alors une égalité (1.4.7.2) :
\begin{align*}
&\prod_{x\in|\mathrm{U}_0|}
  \iota\det(1-\mathrm{F}_x t^{\deg x},\mathscr{F}_0)^{-1}
  \,.\,
  \prod_{x\in|\mathrm{S}_0|}
  \iota\det(1-\mathrm{F}_x t^{\deg x},j_*\mathscr{F}_0)^{-1}
  \tag*{\textbf{(1.8.1.1)}}\\
&\hspace{8em}=
  \frac{\iota\det(1-\mathrm{F}t,\mathrm{H}_c^1(\mathrm{X},j_*\mathscr{F}))}
       {\iota\det(1-\mathrm{F}t,\mathrm{H}_c^2(\mathrm{X},j_*\mathscr{F}))}.
\end{align*}

Dans cette formule :

\textit{a)} D’après (1.4.6), le premier facteur au premier membre converge pour \(|t|<q^{-(\beta+2)/2}\), et n’a ni zéro ni pôle dans cette région.

\textit{b)} D’après (1.4.3), le second membre n’a pas de pôle pour \(|t|<q^{-(\beta+2)/2}\).

\textit{c)} Comparant \textit{a)} et \textit{b)}, on trouve que les facteurs restants \(\iota\det(1-\mathrm{F}_x t^{\deg x},j_*\mathscr{F}_0)^{-1}\) \((x\in|\mathrm{S}_0|)\) sont sans pôle pour \(|t|<q^{-(\beta+2)/2}\). Si \(\alpha\) est une valeur propre de \(\mathrm{F}_x\) sur \(j_*\mathscr{F}_0\), on a donc \(|\iota\alpha|\leq\mathrm{N}(x)^{(\beta+2)/2}\), soit :
\[
w_{\mathrm{N}(x)}(\alpha)\leq\beta+2.
\tag*{\textbf{(1.8.1.2)}}
\]

Pour conclure, on applique (1.8.1.2) aux puissances tensorielles de \(\mathscr{F}_0\). Puisque \(\mathop{\otimes}\limits^k j_*\mathscr{F}_0\) est un sous-faisceau de \(j_*\mathop{\otimes}\limits^k\mathscr{F}_0\), on trouve que les valeurs propres \(\alpha\) de \(\mathrm{F}_x\) sur \(j_*\mathscr{F}_0\) satisfont à \(w_{\mathrm{N}(x)}(\alpha^k)\leq k\beta+2\), soit \(w_{\mathrm{N}(x)}(\alpha)\leq\beta+2/k\), et on fait tendre \(k\) vers l’infini.

\textit{Remarque} \textbf{(1.8.2)}. — Comparons les zéros des deux membres de (1.8.1.1). A gauche, il n’y en a pas pour \(|t|<q^{-(\beta+2)/2}\). A droite, le dénominateur ne s’annule pas dans cette même région. Le numérateur ne peut donc s’y annuler, et, pour toute valeur propre \(\alpha\) de \(\mathrm{F}\) sur \(\mathrm{H}_c^1(\mathrm{X},j_*\mathscr{F})\), on a \(w_q(\alpha)\leq\beta+2\).

\textbf{(1.8.3)} Soient \(\bar{s}\in\mathrm{S}\), d’image \(s\) dans \(\mathrm{S}_0\), \(\eta\) le point générique du trait hensélien à corps résiduel fini \(\mathrm{X}_{0(s)}\), et \(\bar{\eta}\) un point générique géométrique de son hensélisé strict \(\mathrm{X}_{(\bar{s})}\) (cf. (1.7.1)). La fibre \(\mathscr{F}_{\bar{\eta}}\) de \(\mathscr{F}_0\) en \(\bar{\eta}\) est une représentation du groupe de Weil \(\mathrm{W}(\bar{\eta}/\eta)\). Nous noterons \(\mathrm{M}\) sa filtration de monodromie locale (1.7.2).

\textit{Théorème} \textbf{(1.8.4)}. — \textit{Avec les notations de \textup{(1.8.1)} et \textup{(1.8.3)}, si \(\mathscr{F}_0\) est ponctuellement \(\iota\)-pur de poids \(\beta\), la représentation \(\operatorname{Gr}_i^{\mathrm{M}}(\mathscr{F}_{0\eta})\) de \(\mathrm{W}(\bar{\eta}/\eta)\) est \(\iota\)-pure \textup{(1.7.4)} de poids \(\beta+i\).}

\Page{176}
Par torsion (1.2.7), on se ramène à supposer que \(\beta=0\). Remplaçons \(\mathrm{X}_0\) par un revêtement fini, \(\mathrm{U}_0\) et \(\mathscr{F}_0\) par leurs images réciproques, et \(s\) par un point au-dessus de \(s\). Ceci ne modifie ni les poids, ni \(\mathscr{F}_{0\bar{\eta}}\), ni la filtration \(\mathrm{M}\), et nous ramène au cas où la formule (1.7.2.1), pour la représentation \(\rho\) de \(\mathrm{W}(\bar{\eta}/\eta)\) sur \(\mathscr{F}_{0\bar{\eta}}\), vaut dans \(\mathrm{I}\) tout entier :
\(\rho(\sigma)=\exp(t_\ell(\sigma).\mathrm{N})\quad\text{pour}\quad\sigma\in\mathrm{I}\).

Sous ces hypothèses, \(\operatorname{Ker}(\mathrm{N})=\mathscr{F}_{0\bar{\eta}}^{\mathrm{I}}\); cet espace vectoriel est la fibre de \(j_*\mathscr{F}_0\) en \(\bar{s}\). D’après (1.8.1), si \(\alpha\) est une valeur propre de \(\mathrm{F}\) sur \(\operatorname{Ker}(\mathrm{N})\), on a donc \(|\iota\alpha|\leq1\).

Soit \(\alpha\) une valeur propre de \(\mathrm{F}\) sur \(\mathrm{P}_{-j}(\mathscr{F}_{0\bar{\eta}})\). D’après (1.6.14.2), on a \(|\iota\alpha|\leq1\). Appliquons ce résultat à \(\mathscr{F}_0\otimes\mathscr{F}_0\). D’après (1.6.14.4), \(\mathrm{P}_{-j}(\mathscr{F}_{0\bar{\eta}})\otimes\mathrm{P}_{-j}(\mathscr{F}_{0\bar{\eta}})(-j)\) est facteur direct de \(\mathrm{P}_0(\mathscr{F}_{0\bar{\eta}}\otimes\mathscr{F}_{0\bar{\eta}})\), de sorte que \(\alpha^2q^j\) est valeur propre de \(\mathrm{F}\) sur \(\mathrm{P}_0(\mathscr{F}_{0\bar{\eta}}\otimes\mathscr{F}_{0\bar{\eta}})\), et \(|\iota(\alpha^2q^j)|\leq1\), i.e. \(|\iota\alpha|\leq q^{-j/2}\). Appliquons ce dernier résultat au faisceau dual. Appliquant (1.6.14.5), on trouve que \(|\iota(\alpha^{-1}q^{-j})|\leq q^{-j/2}\), i.e. \(|\iota\alpha|\geq q^{-j/2}\).

Les valeurs propres \(\alpha\) de \(\mathrm{F}\) sur \(\mathrm{P}_{-j}(\mathscr{F}_{0\bar{\eta}})\) vérifient donc \(|\iota\alpha|=q^{-j/2}\), et on conclut par (1.6.14.3).

\textit{Corollaire} \textbf{(1.8.5)}. — \textit{Supposons que \(\mathscr{F}_0\) soit un faisceau lisse sur \(\mathrm{U}_0\), et qu’il admette une filtration finie croissante \(\mathrm{W}\) par des sous-faisceaux lisses, telle que, pour chaque \(i\), \(\operatorname{Gr}_i^{\mathrm{W}}(\mathscr{F}_0)\) soit ponctuellement \(\iota\)-pur de poids \(i\). Notons encore \(\mathrm{W}\) la filtration qui s’en déduit sur \(\mathscr{F}_{\bar{\eta}}\). Alors, \(\mathscr{F}_{\bar{\eta}}\) est \(\iota\)-mixte à poids entiers, la filtration de monodromie locale de \(\mathscr{F}_{\bar{\eta}}\), rel. \(\mathrm{W}\) \textup{(1.7.2)} existe, et elle coïncide avec la filtration par le poids de \(\mathscr{F}_{\bar{\eta}}\) \textup{(1.7.5)}.}

D’après (1.7.5), et (1.8.4) appliqué aux \(\operatorname{Gr}_i^{\mathrm{W}}(\mathscr{F}_0)\), les propriétés caractéristiques (1.6.13) sont en effet satisfaites par la filtration par le poids.

\textbf{(1.8.6)} Soient \(\mathrm{D}_0\) un diviseur lisse dans un schéma \(\mathrm{X}_0\) lisse sur \(\mathbf{F}_q\), \(\mathrm{U}_0\) l’ouvert complémentaire, et \(\mathscr{F}_0\) un faisceau lisse sur \(\mathrm{U}_0\) modérément ramifié le long de \(\mathrm{D}_0\). On suppose que \(\mathscr{F}_0\) admet une filtration finie croissante \(\mathrm{W}\) par des sous-faisceaux lisses, telle que \(\operatorname{Gr}_i^{\mathrm{W}}\mathscr{F}_0\) soit ponctuellement \(\iota\)-pur de poids \(i\) (nous verrons en (3.4) que tout faisceau lisse \(\iota\)-mixte de poids ponctuels entiers admet une telle filtration). Supposons \(\mathrm{D}_0\) défini par une équation \(t\), et soit \(\mathscr{F}_0[\mathrm{D}_0]\) le faisceau sur \(\mathrm{D}_0\) muni d’une action de \(\mathbf{Z}_\ell(1)\), qui s’en déduit par la construction (1.7.8).

En restreignant \(\mathscr{F}_0\) à des courbes transverses à \(\mathrm{D}_0\), on déduit de (1.8.5) le résultat suivant :

\textit{Corollaire} \textbf{(1.8.7)}. — \textit{Avec les notations de \textup{(1.8.6)}, la filtration de monodromie locale autour de \(\mathrm{D}\) de \(\mathscr{F}_0[\mathrm{D}_0]\), rel. \(\mathrm{W}\) existe. Pour cette filtration, notée \(\mathrm{M}\), \(\operatorname{Gr}_i^{\mathrm{M}}\mathscr{F}_0[\mathrm{D}_0]\) est ponctuellement \(\iota\)-pur de poids \(i\).}

\textit{Remarque} \textbf{(1.8.8)}. — 1) Avec les notations de (1.8.1), si \(\mathscr{F}_0\) est un faisceau lisse sur \(\mathrm{U}_0\), on a \((j_*\mathscr{F})_{\bar{s}}=(\mathscr{F}_{\bar{\eta}})^{\mathrm{I}}\); ce groupe se déduit de \(\operatorname{Ker}\mathrm{N}\) en prenant les invariants par un groupe fini \(\mathrm{I}/\mathrm{I}_1\). Sous les hypothèses de (1.8.4), on sait que \(\operatorname{Ker}\mathrm{N}\subset\mathrm{M}_0\) (1.6). La filtration \(\mathrm{M}\) induit donc sur \((j_*\mathscr{F})_{\bar{s}}\) une filtration pour laquelle \(\operatorname{Gr}_i^{\mathrm{M}}(j_*\mathscr{F})_{\bar{s}}\) est \(\iota\)-pur de poids \(\beta+i\), et nul pour \(i>0\).

% ---- 177-184 ----
\Page{177}
2) De même, avec les notations de (1.8.6), si \(\mathscr{F}_0\) est un faisceau lisse sur \(\mathrm{U}_0\), modérément ramifié le long de \(\mathrm{D}_0\), la restriction de \(j_*\mathscr{F}_0\) à \(\mathrm{D}_0\) est le sous-faisceau des invariants d’un schéma en groupes finis \(\mu_n\) dans \(\operatorname{Ker}\mathrm{N}\subset\mathscr{F}'_0[\mathrm{D}_0]\), et si \(\mathscr{F}_0\) est ponctuellement \(\iota\)-pur de poids \(\beta\), on déduit par torsion de (1.8.7) que, pour la filtration de monodromie \(\mathrm{M}\), \(\operatorname{Gr}_i^{\mathrm{M}}(j_*\mathscr{F}_0\vert\mathrm{D}_0)\) est ponctuellement \(\iota\)-pur de poids \(\beta+i\), et nul pour \(i>0\).

3) Sous les hypothèses de (1.8.6), on voit aussi que \(\mathrm{W}\) et la monodromie définissent sur \(j_*\mathscr{F}_0\vert\mathrm{D}_0\) une filtration \(\mathrm{M}\) pour laquelle \(\operatorname{Gr}_i^{\mathrm{M}}(j_*\mathscr{F}_0\vert\mathrm{D}_0)\) est lisse et \(\iota\)-pur de poids \(i\).

\textit{Corollaire} \textbf{(1.8.9)}. — \textit{Soient \(\mathrm{X}_0\) un schéma de type fini sur \(\mathbf{F}_q\), \(j:\mathrm{U}_0\hookrightarrow\mathrm{X}_0\) un ouvert de \(\mathrm{X}_0\), et \(\mathscr{F}_0\) un faisceau \(\iota\)-mixte de poids ponctuels \(\leq\beta\) sur \(\mathrm{U}_0\). Alors \(j_*\mathscr{F}_0\) est encore \(\iota\)-mixte de poids ponctuels \(\leq\beta\). Si les poids ponctuels de \(\mathscr{F}_0\) sont entiers, ceux de \(j_*\mathscr{F}_0\) le sont aussi.}

Nous utiliserons les dévissages suivants (justifiés par (1.2.5)).

\textit{a)} Dévissage en \(\mathscr{F}_0\) : si \(\mathscr{F}_0\) est une extension de \(\mathscr{F}''_0\) par \(\mathscr{F}'_0\), la suite exacte \(0\to j_*\mathscr{F}'_0\to j_*\mathscr{F}_0\to j_*\mathscr{F}''_0\) montre qu’il suffit de vérifier (1.8.9) pour \(\mathscr{F}'_0\) et \(\mathscr{F}''_0\).

\textit{b)} Support : si \(i:\mathrm{V}_0\hookrightarrow\mathrm{U}_0\) est un sous-schéma localement fermé de \(\mathrm{U}_0\), d’adhérence \(\overline{\mathrm{V}}_0\) dans \(\mathrm{X}_0\), et que \(\mathscr{F}_0\hookrightarrow i_*i^*\mathscr{F}_0\), on peut remplacer \((\mathrm{X}_0,\mathrm{U}_0,\mathscr{F}_0)\) par \((\overline{\mathrm{V}}_0,\mathrm{V}_0,i^*\mathscr{F}_0)\).

\textit{c)} Normalisation : si \(\varepsilon:\mathrm{X}'_0\to\mathrm{X}_0\) est fini et surjectif :
\[
\begin{tikzcd}[column sep=2.5em,row sep=3.5em]
\varepsilon^{-1}(\mathrm{U}_0)
  \arrow[r,hook,"j"] \arrow[d,"\varepsilon"]
  & \mathrm{X}'_0 \arrow[d,"\varepsilon"] \\
\mathrm{U}_0 \arrow[r,hook,"j"]
  & \mathrm{X}_0
\end{tikzcd}
\]
on a \(j_*\mathscr{F}_0\hookrightarrow\varepsilon_*j_*\varepsilon^*\mathscr{F}_0\). Ceci permet de remplacer \((\mathrm{X}_0,\mathrm{U}_0,\mathscr{F}_0)\) par \((\mathrm{X}'_0,\varepsilon^{-1}(\mathrm{U}_0),\varepsilon^*\mathscr{F}_0)\).

\textit{d)} Faire de l’infini un diviseur : si \(\mathscr{F}_0\) est lisse, et que \(k:\mathrm{V}_0\hookrightarrow\mathrm{U}_0\) est un ouvert dense, on a \(j_*\mathscr{F}_0\hookrightarrow(jk)_*k^*\mathscr{F}_0\). Ceci permet de remplacer \(\mathrm{U}_0\) et \(\mathscr{F}_0\) par \(\mathrm{V}_0\) et \(k^*\mathscr{F}_0\).

\textit{e)} Localisation au point générique à l’infini : soit \(i:\mathrm{F}_0\hookrightarrow\mathrm{X}_0\) le complément de \(\mathrm{U}_0\), et \(k:\mathrm{V}_0\hookrightarrow\mathrm{F}_0\) un ouvert dense de \(\mathrm{F}_0\). Si \(\mathrm{X}_0\) est normal, et \(\mathscr{F}_0\) lisse, on a \(i^*j_*\mathscr{F}_0\hookrightarrow k_*k^*j_*\mathscr{F}_0\).

Soient en effet \(x\) un point géométrique de \(\mathrm{F}_0\), \(y\) une générisation de \(x\) dans \(\mathrm{V}_0\) et \(z\) une générisation de \(y\) dans \(\mathrm{U}_0\). L’hypothèse assure que l’image inverse de \(\mathrm{U}_0\) dans l’hensélisé strict de \(\mathrm{X}_0\) en \(x\) est connexe, donc que \((i^*j_*\mathscr{F}_0)_x\hookrightarrow(\mathscr{F}_0)_z\). La factorisation de cette flèche en :
\[
(i^*j_*\mathscr{F}_0)_x
\to(k_*k^*j_*\mathscr{F}_0)_x
\to(j_*\mathscr{F}_0)_y
\to(\mathscr{F}_0)_z
\]
fournit l’injectivité voulue.

Prouvons (1.8.9) par récurrence sur \(\dim\mathrm{U}_0\). Les dévissages \textit{a)} à \textit{d)} nous ramènent à supposer \(\mathrm{X}_0\) normal, \(\mathrm{U}_0\) lisse et le complément d’un diviseur de Weil \(\mathrm{F}_0\) et \(\mathscr{F}_0\) lisse ponctuellement \(\iota\)-pur et modérément ramifié aux points génériques de \(\mathrm{F}_0\). Dans \textit{e)},
\Page{178}
le théorème est acquis par récurrence pour \(k_*\), et il suffit donc de prouver (1.8.9) après avoir remplacé \(\mathrm{X}_0\) par des ouverts dont la réunion des traces dans \(\mathrm{F}_0\) soit dense dans \(\mathrm{F}_0\). Ceci nous ramène à (1.8.8) 2).

\textit{Corollaire} \textbf{(1.8.10)}. — \textit{Soient \(\mathscr{F}_0\) un faisceau lisse sur \(\mathrm{X}_0\) de type fini sur \(\mathbf{F}_q\), et \(j:\mathrm{U}_0\hookrightarrow\mathrm{X}_0\) un ouvert dense de \(\mathrm{X}_0\). Si la restriction de \(\mathscr{F}_0\) à \(\mathrm{U}_0\) est ponctuellement \(\iota\)-pure, alors \(\mathscr{F}_0\) lui-même est ponctuellement \(\iota\)-pur.}

Soit \(\beta\) le poids de \(\mathscr{F}_0\vert\mathrm{U}_0\). On a \(\mathscr{F}_0\hookrightarrow j_*j^*\mathscr{F}_0\) et \(\mathscr{F}_0^\vee\hookrightarrow j_*j^*\mathscr{F}_0^\vee\), puisque \(\mathscr{F}_0\) et son dual \(\mathscr{F}_0^\vee\) sont lisses. D’après (1.8.9), \(\mathscr{F}_0\) est donc \(\iota\)-mixte de poids ponctuels \(\leq\beta\). Puisque son dual est de même \(\iota\)-mixte de poids ponctuels \(\leq-\beta\), ses poids ponctuels sont aussi \(\geq\beta\), et (1.8.10) en résulte.

Il n’est en fait pas nécessaire d’invoquer (1.8.9) : par restriction à des courbes et normalisation, on se ramène facilement au cas où \(\mathrm{X}_0\) est une courbe lisse, et on invoque directement (1.8.1).

\textit{Corollaire} \textbf{(1.8.11)}. — \textit{Tout faisceau lisse \(\iota\)-mixte \(\mathscr{F}_0\) sur \(\mathrm{X}_0\) normal de type fini sur \(\mathbf{F}_q\) est extension successive de faisceaux lisses ponctuellement \(\iota\)-purs.}

On peut supposer \(\mathrm{X}_0\) connexe, et on se ramène à supposer le faisceau irréductible. Puisque \(\mathrm{X}_0\) est normal, sa restriction à un quelconque ouvert non vide \(\mathrm{U}_0\) est encore irréductible — donc ponctuellement \(\iota\)-pure pour \(\mathrm{U}_0\) convenable, par définition des faisceaux \(\iota\)-mixtes — et on applique (1.8.10).

Nous verrons plus tard (3.4.1) que (1.8.11) vaut sans hypothèse de normalité.

\textit{Corollaire} \textbf{(1.8.12)}. — \textit{Soit \(\mathscr{F}_0\) un faisceau lisse \(\iota\)-mixte sur \(\mathrm{X}_0\) connexe et de type fini sur \(\mathbf{F}_q\). Si \(\mathscr{F}_0\) est \(\iota\)-pur en un point de \(\mathrm{X}_0\), il est ponctuellement \(\iota\)-pur.}

Supposons d’abord \(\mathrm{X}_0\) normal. Il résulte alors de (1.8.11) que les poids de \(\mathscr{F}_0\) se lisent sur sa fibre en un point arbitraire de \(\mathrm{X}_0\). Dans le cas général, en appliquant (1.8.11) à l’image inverse de \(\mathscr{F}_0\) sur le normalisé de \(\mathrm{X}_0\), on trouve que l’ensemble des points de \(\mathrm{X}_0\) où \(\mathscr{F}_0\) est pur d’un poids \(\beta\), contient avec un point fermé \(x_0\) toutes les composantes irréductibles de \(\mathrm{X}_0\) passant par \(x_0\). Il est donc ouvert et fermé.

\textit{Variante} \textbf{(1.8.13)}. — Dans tout ce paragraphe, on peut remplacer \(\iota\)-pur et \(\iota\)-mixte par « pur » et « mixte », ainsi que par « \(\iota\)-pur de poids entiers » et « \(\iota\)-mixte de poids entiers ».

\textbf{(1.8.14)} Dans le dictionnaire heuristique de [2], I, le théorème (1.8.4) correspond à la théorie de W. Schmid du comportement asymptotique des variations de structures de Hodge polarisables sur le disque épointé \(\mathrm{D}^*\). Le corollaire (1.8.5) suggère le problème suivant :

\Page{179}
Appelons « variation de structures de Hodge mixtes » sur un espace analytique lisse \(\mathrm{S}\) la donnée de :

\noindent\hangindent=1.5em
\textit{a)} un système local \(\mathrm{V}_{\mathbf{Z}}\) de \(\mathbf{Z}\)-modules libres de type fini;

\noindent\hangindent=1.5em
\textit{b)} une filtration finie croissante \(\mathrm{W}\) de \(\mathrm{V}_{\mathbf{Q}}=\mathrm{V}_{\mathbf{Z}}\otimes\mathbf{Q}\) par des sous-systèmes locaux d’espaces vectoriels sur \(\mathbf{Q}\);

\noindent\hangindent=1.5em
\textit{c)} une filtration finie décroissante \(\mathrm{F}\) de \(\mathrm{V}=\mathrm{V}_{\mathbf{Z}}\otimes\mathscr{O}\) par des sous-fibrés vectoriels (i.e. la filtration \(\mathrm{F}\) varie de façon holomorphe); on exige que \(\mathrm{F}\) vérifie l’axiome de transversalité : \(\nabla\mathrm{F}^i\subset\Omega^1\otimes\mathrm{F}^{i-1}\) et que \(\mathrm{W}\) et \(\mathrm{F}\) définissent en chaque point de \(\mathrm{S}\) une structure de Hodge mixte.

\noindent
Les variations qui apparaissent en géométrie algébrique ont des propriétés additionnelles; par exemple, les variations de structures de Hodge \(\operatorname{Gr}_i^{\mathrm{W}}(\mathrm{V})\) sont polarisables.

Soient \((\mathrm{V},\mathrm{W},\mathrm{F})\) une variation de structures de Hodge mixtes sur \(\mathrm{D}^*\), \(t\in\mathrm{D}^*\), \(\mathrm{V}\) la fibre de \(\mathrm{V}\) en \(t\) et \(\mathrm{T}\in\operatorname{End}(\mathrm{V})\) la transformation de monodromie. On suppose pour simplifier que \(\mathrm{T}\) est unipotente, et on pose \(\mathrm{N}=\log\mathrm{T}\).

\textit{Problème} \textbf{(1.8.15)}. — Dégager une classe de « bonnes » variations de structures de Hodge mixtes sur \(\mathrm{D}^*\) telle que, pour \((\mathrm{V},\mathrm{W},\mathrm{F})\) bonne (et unipotente), il existe sur \(\mathrm{V}\) une filtration finie croissante \(\mathrm{M}\), telle que \(\mathrm{N}\mathrm{M}_a\subset\mathrm{M}_{a-2}\), et que \(\mathrm{N}\) induise des isomorphismes :
\[
\mathrm{N}^b:
\mathrm{G}_{a+b}^{\mathrm{M}}\mathrm{G}_a^{\mathrm{W}}\mathrm{V}
\xrightarrow{\sim}
\mathrm{G}_{a-b}^{\mathrm{M}}\mathrm{G}_a^{\mathrm{W}}\mathrm{V}.
\]

On aimerait aussi que :

\noindent\hangindent=1.5em
\textit{a)} la variation \(\mathrm{V}\) soit asymptotique, en un sens convenable, à une « orbite nilpotente » : une bonne variation pour laquelle \(\mathrm{F}\) en \(t\exp(u)\) est le transformé par \(\exp(-u\mathrm{N})\) de la filtration déduite de \(\mathrm{F}_t\) par transport horizontal de \(t\) à \(t\exp(u)\);

\noindent\hangindent=1.5em
\textit{b)} pour une orbite nilpotente, \((\mathrm{V},\mathrm{M},\mathrm{F})\) soit encore une variation de structures de Hodge mixtes.

\section*{(1.9) Monodromie locale des faisceaux mixtes}

Dans ce numéro, qui ne sera pas utilisé par la suite, nous généralisons (1.8.7) au cas des diviseurs à croisements normaux.

\textbf{(1.9.1)} Soient \(\mathrm{X}\) un schéma lisse sur \(\mathbf{F}_q\), \(\mathrm{D}\) un diviseur à croisements normaux de \(\mathrm{X}\), réunion transverse de diviseurs lisses \(\mathrm{D}_i\) (\(i\in\mathrm{I}\)) et \(\mathrm{E}\) l’intersection des \(\mathrm{D}_i\). On suppose chaque \(\mathrm{D}_i\) défini par une équation globale \(t_i=0\), pour pouvoir appliquer la construction (1.7.8) (cf. toutefois (1.7.10)).

Soit par ailleurs \(\mathscr{F}\) un faisceau lisse sur \(\mathrm{X}-\mathrm{D}\), modérément ramifié le long de \(\mathrm{D}\). On suppose qu’il admet une filtration finie croissante \(\mathrm{W}\) par des sous-faisceaux lisses, telle que \(\operatorname{Gr}_i^{\mathrm{W}}\mathscr{F}\) soit ponctuellement \(\iota\)-pur de poids \(i\) (nous verrons en (3.4.1) que tout faisceau lisse \(\iota\)-mixte de poids ponctuels entiers admet une telle filtration). Soit \(\mathbf{L}\)
\Page{180}
l’ensemble de tous les nombres premiers \(\neq p\), et définissons \(\mathscr{F}[\mathrm{E}]\) comme en (1.7.8). C’est un faisceau lisse sur \(\mathrm{E}\), muni d’une action de \(\mathbf{Z}_{\mathbf{L}}(1)^{\mathrm{I}}\); cette action est quasi-unipotente, d’où des « endomorphismes » \(\mathrm{N}_i:\mathscr{F}[\mathrm{E}](1)\to\mathscr{F}[\mathrm{E}]\) (1.7.8).

Etant donnés un faisceau \(\mathscr{G}\), deux filtrations finies croissantes \(\mathrm{W}'\) et \(\mathrm{W}''\) de \(\mathscr{G}\) et une famille finie d’ « endomorphismes » \(\mathrm{N}_\alpha:\mathscr{G}(1)\to\mathscr{G}\) (\(\alpha\in\mathrm{A}\)), nous aurons à considérer la condition :

\(\mathrm{L}(\mathrm{W}',\mathrm{W}'',(\mathrm{N}_\alpha)_{\alpha\in\mathrm{A}})\) : Les \(\mathrm{N}_\alpha\) stabilisent \(\mathrm{W}'\), vérifient \(\mathrm{N}_\alpha\mathrm{W}''_a(1)\subset\mathrm{W}''_{a-2}\) et quels que soient les nombres rationnels \(c_\alpha>0\), l’endomorphisme \(\mathrm{N}=\sum c_\alpha\mathrm{N}_\alpha\) induit des isomorphismes :
\[
\mathrm{N}^b:
(\operatorname{Gr}_{a+b}^{\mathrm{W}''}\mathrm{G}_a^{\mathrm{W}'}\mathscr{G})(b)
\xrightarrow{\sim}
\operatorname{Gr}_{a-b}^{\mathrm{W}''}\mathrm{G}_a^{\mathrm{W}'}\mathscr{G}.
\]

Il résulte de (1.6.13) que, \(\mathrm{W}'\) et la famille des \(\mathrm{N}_\alpha\) étant donnés, il existe au plus une filtration \(\mathrm{W}''\) vérifiant cette condition. Dans son énoncé, il est loisible de remplacer « nombre rationnel \(>0\) » par « nombre entier \(>0\) » : remplacer les \(c_\alpha\) par \(nc_\alpha\) et \(\mathrm{N}\) par \(n\mathrm{N}\).

\textit{Théorème} \textbf{(1.9.2)}. — \textit{Avec les hypothèses et notations de \textup{(1.9.1)}, il existe sur \(\mathscr{F}[\mathrm{E}]\) un et un seul système de filtrations finies croissantes \(\mathrm{W}(\mathrm{J})\) (\(\mathrm{J}\subset\mathrm{I}\)) par des sous-faisceaux lisses stables par \(\mathbf{Z}_{\mathbf{L}}(1)^{\mathrm{I}}\) (donc par les \(\mathrm{N}_i\)), tel que :}

\noindent a) \textit{\(\mathrm{W}(\varnothing)\) est la filtration déduite de \(\mathrm{W}\) par fonctorialité;}

\noindent b) \textit{pour \(\mathrm{J}\subset\mathrm{K}\), la condition \(\mathrm{L}(\mathrm{W}(\mathrm{J}),\mathrm{W}(\mathrm{K}),(\mathrm{N}_k)_{k\in\mathrm{K}-\mathrm{J}})\) est vérifiée.}

\noindent\textit{En outre, \(\operatorname{Gr}_a^{\mathrm{W}(\mathrm{I})}\mathscr{F}[\mathrm{E}]\) est ponctuellement \(\iota\)-pur de poids \(a\).}

L’unicité des \(\mathrm{W}(\mathrm{J})\) résulte aussitôt de (1.6.13) (et il suffit de disposer des conditions a), et b) pour \(\varnothing\subset\mathrm{K}\)). Pour prouver l’existence, nous traiterons d’abord un cas particulier.

\textit{Lemme} \textbf{(1.9.3)}. — \textit{Soit \((c_k)_{k\in\mathrm{I}}\) un système de nombres entiers \(>0\). Le faisceau \(\mathscr{F}[\mathrm{E}]\) admet une (unique) filtration finie croissante \(\mathrm{W}(\mathrm{I})\) par des sous-faisceaux lisses, telle que \(\mathrm{N}=\sum c_k\mathrm{N}_k\) vérifie la condition \(\mathrm{L}(\mathrm{W},\mathrm{W}(\mathrm{I}),\mathrm{N})\).}

\textit{En outre \(\operatorname{Gr}_a^{\mathrm{W}(\mathrm{I})}\mathscr{F}[\mathrm{E}]\) est ponctuellement \(\iota\)-pur de poids \(a\).}

Vu l’unicité (1.6.13), si pour tout point géométrique \(\bar{x}\) de \(\mathrm{E}\), localisé en un point fermé \(x\), il existe une filtration finie croissante \(\mathrm{M}\) sur \(\mathscr{F}[\mathrm{E}]_{\bar{x}}\) satisfaisant à \(\mathrm{L}(\mathrm{W},\mathrm{M},\mathrm{N})\) alors ces filtrations se recollent en une filtration \(\mathrm{M}\) de \(\mathscr{F}[\mathrm{E}]\) par des sous-faisceaux lisses, satisfaisant à \(\mathrm{L}(\mathrm{W},\mathrm{W}(\mathrm{I}),\mathrm{N})\). Il suffit même de prouver l’existence de \(\mathrm{M}\) en un point de chaque composante connexe de \(\mathrm{E}\).

Ceci nous ramène, pour prouver (1.9.3), à montrer que pour \(\bar{x}\) et \(x\) comme ci-dessus, \(\mathrm{M}\) existe et que \(\operatorname{Gr}_a^{\mathrm{M}}\mathscr{F}[\mathrm{E}]_{\bar{x}}\) est pur de poids \(a\). Nous le vérifierons par réduction à (1.8.5). Une extension préliminaire du corps fini de base nous ramène à supposer que \(x\) est un point rationnel. On complète alors les \(t_i\) en un système de coordonnées locales \((t_i,u_j)\), centré en \(x\) (une application étale d’un voisinage de \(x\) dans l’espace
\Page{181}
affine type \(\mathbf{A}^n\) (\(n\) convenable) envoyant \(x\) sur l’origine). Restreignons \(\mathscr{F}\) à la courbe d’équation paramétrique \(u_j=0\), \(t_i=t^{c_i}\). Plus précisément : on envoie la droite affine (coordonnée \(t\)) dans l’espace affine \(\mathbf{A}^n\) de coordonnées \(t_i,u_j\) par \(t\mapsto(t^{c_i},0)\), et on prend l’image inverse de \(\mathscr{F}\) sur la courbe \(\mathrm{C}\) produit fibré d’un voisinage de \(x\), et de la droite affine, au-dessus de \(\mathbf{A}^n\). Notons encore \(x\) (resp. \(\bar{x}\)) le point (resp. le point géométrique) de \(\mathrm{C}\) au-dessus de \(x\) (resp. \(\bar{x}\)) dans \(\mathrm{X}\), et au-dessus de \(0\) dans la droite affine, et passons au localisé strict \(\mathrm{C}_{(\bar{x})}\). La fibre générique géométrique sur \(\mathrm{C}_{(\bar{x})}\) de \(\mathscr{F}\) s’identifie à \(\mathscr{F}[\mathrm{E}]_{\bar{x}}\). Via cette identification, le logarithme de la monodromie devient \(\mathrm{N}\), et (1.9.3) résulte de (1.8.5).

\textbf{(1.9.4)} Déduisons le théorème de son cas particulier (1.9.3). On remarque d’abord que la filtration \(\mathrm{W}(\mathrm{I})\) de (1.9.3) ne dépend pas du choix des \(c_k\), puisque caractérisée indépendamment de ceux-ci comme étant la filtration par le poids.

Pour \(\mathrm{J}\subset\mathrm{I}\), soient \(\mathrm{X}_{\mathrm{J}}\), \(\mathrm{D}_{\mathrm{J}}\) et \(\mathrm{D}_{\mathrm{J}}^*\) comme en (1.7.9), et appliquons (1.9.3) à \(\mathrm{X}_{\mathrm{J}}\), aux traces des \(\mathrm{D}_j\) (\(j\in\mathrm{J}\)), et à leur intersection \(\mathrm{D}_{\mathrm{J}}^*\). On obtient sur \(\mathscr{F}[\mathrm{D}_{\mathrm{J}}^*]\) une filtration par le poids \(\mathrm{W}(\mathrm{J})\), et elle satisfait \(\mathrm{L}(\mathrm{W},\mathrm{W}(\mathrm{J}),(\mathrm{N}_j)_{j\in\mathrm{J}})\).

Réappliquons cette construction à \(\mathrm{D}_{\mathrm{J}}\), aux traces des diviseurs \(\mathrm{D}_i\) (\(i\notin\mathrm{J}\)) et à \(\mathscr{F}[\mathrm{D}_{\mathrm{J}}^*]\), muni de la filtration \(\mathrm{W}(\mathrm{J})\). Pour \(\mathrm{J}\subset\mathrm{K}\), on trouve sur \(\mathscr{F}[\mathrm{D}_{\mathrm{J}}^*][\mathrm{D}_{\mathrm{K}}^*]\) une filtration par le poids \(\mathrm{W}(\mathrm{K})\), satisfaisant \(\mathrm{L}(\mathrm{W}(\mathrm{J}),\mathrm{W}(\mathrm{K}),(\mathrm{N}_i)_{i\in\mathrm{K}-\mathrm{J}})\). Via l’isomorphisme (1.7.9.1) de \(\mathscr{F}[\mathrm{D}_{\mathrm{J}}^*][\mathrm{D}_{\mathrm{K}}^*]\) avec \(\mathscr{F}[\mathrm{D}_{\mathrm{K}}^*]\), cette filtration par le poids s’identifie à celle construite par une application directe de (1.9.3), satisfaisant \(\mathrm{L}(\mathrm{W},\mathrm{W}(\mathrm{K}),(\mathrm{N}_i)_{i\in\mathrm{K}})\) (unicité de la filtration par le poids).

On a \(\mathrm{E}=\mathrm{D}_{\mathrm{I}}\), et, sur \(\mathscr{F}[\mathrm{E}]\), on définit enfin \(\mathrm{W}(\mathrm{J})\) comme déduit par fonctorialité de la filtration \(\mathrm{W}(\mathrm{J})\) de \(\mathscr{F}[\mathrm{D}_{\mathrm{J}}^*]\), via l’isomorphisme \(\mathscr{F}[\mathrm{E}]=\mathscr{F}[\mathrm{D}_{\mathrm{J}}^*][\mathrm{E}]\). On a a) par définition, et la condition b) se déduit de \(\mathrm{L}(\mathrm{W}(\mathrm{J}),\mathrm{W}(\mathrm{K}),(\mathrm{N}_k)_{k\in\mathrm{K}-\mathrm{J}})\) pour \(\mathscr{F}[\mathrm{D}_{\mathrm{K}}^*]\).

\textit{Variante} \textbf{(1.9.5)}. — Soient \(\mathrm{S}\) un schéma de type fini sur \(\mathbf{Z}\), \(f:\mathrm{X}\to\mathrm{S}\) un morphisme lisse, \(\mathrm{D}\subset\mathrm{X}\) un diviseur à croisements normaux relatif, et \(\mathscr{F}\) un faisceau lisse sur \(\mathrm{X}-\mathrm{D}\), modérément ramifié le long de \(\mathrm{D}\). Une variante de (1.9.2) vaut dans ce cadre. Elle se vérifie fibre par fibre, par réduction à (1.9.2).

Le cas d’un diviseur à croisements normaux dans un schéma \(\mathrm{S}\) régulier de type fini sur \(\mathbf{Z}\) semble par contre inaccessible (déjà pour \(\mathrm{S}=\operatorname{Spec}(\mathbf{Z}[1/\ell])\)).

\textit{Remarque} \textbf{(1.9.6)} \textit{(ajoutée sur épreuves)}. — Comme expliqué dans [2] I, on s’attend à ce que la théorie des variations de structures de Hodge (toujours supposées polarisables) soit souvent parallèle à celle des faisceaux lisses ponctuellement purs. Notons \(\mathrm{D}^*\) le disque épointé. Cattani et Kaplan viennent de prouver un analogue de (1.9.2) (restreint au cas pur) pour les variations de structures de Hodge sur \((\mathrm{D}^*)^{\mathrm{I}}\). Dans leur contexte, les \(c_\alpha\) de la condition \(\mathrm{L}(\mathrm{W}',\mathrm{W}'',(\mathrm{N}_\alpha)_{\alpha\in\mathrm{A}})\) peuvent même être pris réels \(>0\), plutôt que rationnels \(>0\).

\Page{182}
\section*{(1.10) Valeurs absolues non archimédiennes}

\textit{Les conventions \textup{(0.7)} sont en vigueur.}

Soient \(\mathrm{X}_0\) une courbe lisse absolument irréductible sur \(\mathbf{F}_q\), \(j:\mathrm{U}_0\hookrightarrow\mathrm{X}_0\) un ouvert de \(\mathrm{X}_0\), complément de l’ensemble fini de points \(\mathrm{S}_0\), et \(\mathscr{F}_0\) un faisceau lisse sur \(\mathrm{U}_0\). Au n\textsuperscript{o} (1.8), nous avons exploité la formule des traces, sous la forme (1.4.7.2), et la structure de la monodromie locale, pour étudier la valeur absolue archimédienne des valeurs propres de Frobenius sur \(\mathscr{F}_0\), en les points de \(\mathrm{S}_0\) (au sens (1.10.2) ci-dessous). Dans ce numéro, nous étudierons des valeurs absolues non nécessairement archimédiennes.

Choisissons donc soit \textit{a)} un isomorphisme \(\iota\) de \(\overline{\mathbf{Q}}_\ell\) avec \(\mathbf{C}\) (cas archimédien), soit \textit{b)} un nombre premier \(\ell'\), une clôture algébrique \(\overline{\mathbf{Q}}_{\ell'}\) de \(\mathbf{Q}_{\ell'}\) et un isomorphisme \(\iota\) de \(\overline{\mathbf{Q}}_\ell\) avec \(\overline{\mathbf{Q}}_{\ell'}\) (cas non archimédien). On normalisera la valeur absolue naturelle de \(\overline{\mathbf{Q}}_{\ell'}\) par \(\lVert\ell'\rVert=1/\ell'\), par exemple.

Le lemme (1.8.1) se généralise en le

\textit{Lemme} \textbf{(1.10.1)}. — \textit{Soit \(b>0\). Si, pour tout point \(x\) de \(\mathrm{U}_0\) et toute valeur propre \(\alpha\) de \(\mathrm{F}_x\) sur \(\mathscr{F}_0\) on a \(\lVert\iota\alpha\rVert\leq b^{\deg(x)}\), alors, pour tout point \(x\) de \(\lvert\mathrm{S}_0\rvert\), et toute valeur propre \(\alpha\) de \(\mathrm{F}_x\) sur \(j_*\mathscr{F}_0\), on a encore \(\lVert\iota\alpha\rVert\leq b^{\deg(x)}\).}

Reprenons la preuve de (1.8.1). Nous nous placerons dans le cas non archimédien, et indiquerons sous forme de resp. les modifications à apporter pour traiter le cas archimédien (qui est une légère généralisation de (1.8.1)).

Dans (1.8.1.1), le premier facteur au premier membre converge pour \(\lvert t\rvert<1/b\) (resp. \(\lvert t\rvert<1/qb\) (1.4.6)). Le second membre n’a pas de pôle pour \(\lvert t\rvert<1/\lVert q\rVert b\), donc pour \(\lvert t\rvert<1/b\), par (1.4.4) (resp. pour \(\lvert t\rvert<1/qb\)). Les facteurs restant au premier membre :
\[
\iota\det(1-\mathrm{F}_xt^{\deg x},j_*\mathscr{F}_0)^{-1}
\qquad (x\in\lvert\mathrm{S}_0\rvert)
\]
sont donc sans pôle pour \(\lvert t\rvert<1/b\) (resp. \(\lvert t\rvert<1/qb\)), ce qui prouve (1.10.1) (resp. prouve \(\lVert\iota\alpha\rVert\leq(qb)^{\deg x}\), résultat dont on déduit (1.10.1) en remplaçant \(\mathscr{F}_0\) par ses puissances tensorielles, comme en (1.8.1)).

\textbf{(1.10.2)} Soient \(s\in\lvert\mathrm{S}_0\rvert\), \(\bar{s}\), \(\eta\), \(\bar{\eta}\) comme en (1.8.3), et la représentation \(\mathscr{F}_{\bar{\eta}}\) du groupe de Weil \(\mathrm{W}(\bar{\eta}/\eta)\). Nous appellerons « valeurs propres de \(\mathrm{F}_s\) sur \(\mathscr{F}_0\) » les valeurs propres d’un relèvement de \(\mathrm{F}_s\in\mathrm{W}(\bar{s}/s)\) dans \(\mathrm{W}(\bar{\eta}/\eta)\), agissant sur \(\mathscr{F}_{\bar{\eta}}\). Cette terminologie ne sera utilisée que lorsque le choix du relèvement n’importe pas (cf. (1.7.4)).

\textit{Théorème} \textbf{(1.10.3)}. — \textit{On se place dans le cas non archimédien. Soit \(0\leq b\leq c\leq\infty\). Si, pour tout point \(x\in\lvert\mathrm{U}_0\rvert\) et pour toute valeur propre \(\alpha\) de \(\mathrm{F}_x\) sur \(\mathscr{F}_0\), on a \(b^{\deg x}\leq\lVert\iota\alpha\rVert\leq c^{\deg x}\), alors pour tout point \(s\in\mathrm{S}_0\) et pour toute valeur propre \(\alpha\) de \(\mathrm{F}_s\) sur \(\mathscr{F}_0\), on a encore :}
\[
b^{\deg s}\leq\lVert\iota\alpha\rVert\leq c^{\deg s}.
\]

\Page{183}
Soit \(s\in\lvert\mathrm{S}_0\rvert\). Comme en (1.8.4), on passe à un revêtement fini pour supposer la monodromie locale unipotente : l’action de \(\mathrm{I}\subset\mathrm{W}(\bar{\eta}/\eta)\) sur \(\mathscr{F}_{\bar{\eta}}\) s’écrit :
\[
\sigma\mapsto\exp(t_\ell(\sigma).\mathrm{N}).
\]
Le lemme (1.10.1) assure que les valeurs propres de \(\mathrm{F}_s\) sur \(\operatorname{Ker}(\mathrm{N})\) satisfont à \(\lVert\iota\alpha\rVert\leq c^{\deg s}\). Puisque \(\lVert q\rVert\leq1\), il en va de même des valeurs propres de \(\mathrm{F}_s\) sur \(\mathscr{F}_{\bar{\eta}}\) tout entier (1.6.14.3). L’inégalité \(b^{\deg s}\leq\lVert\iota\alpha\rVert\) s’obtient par passage au faisceau dual.

\textbf{(1.10.4)} Si l’on applique le théorème (1.10.3) aux puissances extérieures de \(\mathscr{F}_0\), on obtient une relation entre les polygones de Newton attachés aux \(\mathrm{F}_x\) (\(x\in\lvert\mathrm{U}_0\rvert\)), et aux \(\mathrm{F}_s\) (\(s\in\lvert\mathrm{S}_0\rvert\)). Nous n’énoncerons le résultat que pour \(\ell'=p\) (1.10.7) : pour \(\ell'\neq p\), les valeurs propres de Frobenius sont en pratique toujours des unités \(\ell'\)-adiques, et (1.10.3) n’est jusqu’ici utile que sous la forme suivante :

\textit{Corollaire} \textbf{(1.10.5)}. — \textit{Si, pour tout point \(x\in\lvert\mathrm{U}_0\rvert\), les valeurs propres \(\alpha\) de \(\mathrm{F}_x\) sur \(\mathscr{F}_0\) sont des unités \(\ell'\)-adiques (\textup{i.e.} si \(\lVert\iota\alpha\rVert=1\)), alors, pour tout point \(s\in\lvert\mathrm{S}_0\rvert\), les valeurs propres de \(\mathrm{F}_s\) sur \(\mathscr{F}_0\) sont encore des unités \(\ell'\)-adiques.}

En rapetissant \(\mathrm{U}_0\), on peut remplacer l’hypothèse « pour tout \(x\in\lvert\mathrm{U}_0\rvert\) », par « pour tous sauf un nombre fini », et récupérer gratuitement les \(x\) manquants.

\textbf{(1.10.6)} Pour \(\ell'=p\), il est utile de passer du langage des valeurs absolues à celui des valuations. Pour \(\mathrm{N}\) une puissance de \(p\), soit donc \(v_{\mathrm{N}}\) la valuation de \(\overline{\mathbf{Q}}^{p}\), normalisée par \(v_{\mathrm{N}}(\mathrm{N})=1\).

Soit \((\alpha_i)_{i\in\mathrm{I}}\) une famille finie d’éléments de \(\overline{\mathbf{Q}}_p\). Son \textit{polygone de Newton} rel. \(\mathrm{N}\) est le graphe de la fonction continue linéaire par morceau \(n\) de \([0,\lvert\mathrm{I}\rvert]\) dans \(\mathbf{R}\) caractérisée comme suit : \(n(0)=0\), et si l’on range les \(v_{\mathrm{N}}(\alpha_i)\) en une suite croissante, indexée par les entiers entre \(1\) et \(\lvert\mathrm{I}\rvert\), la pente de \(n\) entre \(j-1\) et \(j\) est la \(j\)-ième de ces valuations. La valeur de \(n\) en l’entier \(j\) (\(0\leq j\leq\lvert\mathrm{I}\rvert\)) est donc la plus petite valuation d’un produit de \(j\) \(\alpha_i\) d’indices distincts.

Pour \(x\in\lvert\mathrm{X}_0\rvert\), le \textit{polygone de Newton} de \(\mathscr{F}_0\) \textit{en \(x\), rel. \(\iota\)}, est le polygone de Newton, rel. \(\mathrm{N}(x)\), de l’image par \(\iota\) du système des valeurs propres de \(\mathrm{F}_x\) sur \(\mathscr{F}_0\). La valeur en l’entier \(j\) (\(0\leq j\leq\operatorname{rg}\mathscr{F}_0\)) de la fonction correspondante est le plus petit des nombres \(v_{\mathrm{N}(x)}(\iota\alpha)\), pour \(\alpha\) valeur propre de \(\mathrm{F}_x\) sur \(\bigwedge\limits^{j}\mathscr{F}_0\).

\textit{Corollaire} \textbf{(1.10.7)}. — \textit{Supposons que, pour presque tout \(x\in\lvert\mathrm{U}_0\rvert\), \(\mathscr{F}_0\) ait en \(x\) le même polygone de Newton \(n\), rel. \(\iota\). Alors, pour \(s\in\lvert\mathrm{S}_0\rvert\), le polygone de Newton \(n_s\) de \(\mathscr{F}_0\) en \(s\) est au-dessus de \(n\) : \(n_s\geq n\), et a les mêmes extrémités.}

La minoration \(n_s\geq n\) reformule la majoration (1.10.3) des valeurs absolues des valeurs propres de \(\mathrm{F}_s\), pour les faisceaux \(\bigwedge\limits^{j}\mathscr{F}_0\). Le complément « même extrémité » : \(n_s(\operatorname{rg}\mathscr{F}_0)=n(\operatorname{rg}\mathscr{F}_0)\), résulte de ce que pour un faisceau de rang un (en l’occurrence,
\Page{184}
la puissance extérieure maximale de \(\mathscr{F}_0\)), la valuation \(v_{\mathrm{N}(x)}\) de \(\iota\alpha\), pour \(\alpha\) la valeur propre de \(\mathrm{F}_x\), est indépendante de \(x\) (cf. (1.3.3)).

\textbf{(1.10.8)} On obtient des résultats plus précis en tenant compte de la filtration de monodromie \(\mathrm{M}\). Par exemple :

\textit{Corollaire} \textbf{(1.10.9)}. — \textit{Soient \(\beta\leq\gamma\), et supposons que, pour presque tout \(x\in\lvert\mathrm{U}_0\rvert\), les valeurs propres \(\alpha\) de \(\mathrm{F}_x\) sur \(\mathscr{F}_0\) satisfont à \(\beta\leq v_{\mathrm{N}(x)}(\iota\alpha)\leq\gamma\). Alors, si \(n\) est la partie entière de \(\gamma-\beta\), le logarithme \(\mathrm{N}\) de la partie unipotente de la monodromie locale, en \(s\in\lvert\mathrm{S}_0\rvert\), satisfait à \(\mathrm{N}^{n+1}=0\).}

En effet, si \(m\) est le plus petit entier tel que \(\mathrm{N}^{m+1}=0\), \(\mathrm{F}_s\) admet des valeurs propres \(\alpha\) et \(\alpha.\mathrm{N}(s)^m\), et on applique (1.10.3).

\section*{(1.11) Spécialisation de la monodromie}

\textit{Proposition} \textbf{(1.11.1)}. — \textit{Soient \(\mathrm{S}\) un schéma irréductible, \(f:\mathrm{X}\to\mathrm{S}\) un morphisme lisse purement de dimension relative \(1\), à fibres géométriquement connexes, \(g\) une section de \(f\) et \(\mathscr{F}\) un \(\mathbf{Z}_\ell\)-faisceau lisse sur \(\mathrm{X}\). Pour chaque point géométrique \(\bar{s}\) de \(\mathrm{S}\), soit \(\mathrm{M}_{\bar{s}}^n\) l’image de \(\pi_1(\mathrm{X}_{\bar{s}},g(\bar{s}))\) dans \(\operatorname{Aut}((\mathscr{F}/\ell^n\mathscr{F})_{g(\bar{s})})\). Alors, il existe un ouvert \(\mathrm{U}\) de \(\mathrm{S}\) tel que, pour tout \(n\), et tout point géométrique \(\bar{s}\) de \(\mathrm{U}\), \(\mathrm{M}_{\bar{s}}^n\) soit la fibre en \(\bar{s}\) d’un sous-faisceau lisse \(\mathrm{M}^n\) de \(\operatorname{Aut}(g^*\mathscr{F}/\ell^n\mathscr{F})\) sur \(\mathrm{U}\).}

Soit \(\mathrm{I}_n\) le revêtement fini étale \(\mathscr{I}\!som(\mathscr{F}/\ell^n\mathscr{F},f^*g^*\mathscr{F}/\ell^n\mathscr{F})\) de \(\mathrm{X}\). On a :
\[
g^*\mathrm{I}_n=\mathscr{A}\!ut(g^*\mathscr{F}/\ell^n\mathscr{F}).
\]
L’automorphisme identique nous fournit une section \(g_n\) de \(\mathrm{I}_n\) au-dessus de \(g\). Après avoir rapetissé \(\mathrm{S}\), nous construisons simultanément les \(\mathrm{M}^n\), et une tour de revêtements finis étales \(\mathrm{X}_n\subset\mathrm{I}_n\), contenant la section \(g_n\), stables par \(\mathrm{M}^n\), et tels que, fibre géométrique à fibre géométrique, \(\mathrm{X}_n\) soit le revêtement galoisien de \(\mathrm{X}\) qui trivialise \(\mathscr{F}/\ell^n\mathscr{F}\). On aura \(\mathrm{M}^n=g^*(\mathrm{X}_n/\mathrm{X})\).

Soit \(\mathrm{F}\) la fibre de \(\mathscr{F}\) en un point géométrique. La preuve repose sur les faits que \(\operatorname{Aut}(\mathrm{F}/\ell\mathrm{F})\) est fini, tandis que le groupe des automorphismes de \(\mathrm{F}\) triviaux sur \(\mathrm{F}/\ell\mathrm{F}\) est un pro-\(\ell\)-groupe.

Prenons la composante neutre de la section \(g_1\) de \(\mathrm{I}_1\). Au-dessus d’un ouvert convenable \(\mathrm{U}_1\) de \(\mathrm{S}\), cette construction est compatible au changement de base (cf. EGA IV (9.7.7)). Elle nous fournit, sur \(\mathrm{U}_1\), le revêtement \(\mathrm{X}_1\) cherché. On prend \(\mathrm{M}^1=g^*\mathrm{X}_1\).

Soient \(\mathbf{L}\) un ensemble de nombres premiers inversibles sur \(\mathrm{U}_1\), et, pour \(\pi_1\) un groupe fondamental, soit \(\pi_1^{(\mathbf{L})}\) son plus grand quotient qui soit un pro-\(\mathbf{L}\)-groupe. On sait que, quitte à remplacer \(\mathrm{U}_1\) par un ouvert plus petit \(\mathrm{U}\), les \(\pi_1^{(\mathbf{L})}(\mathrm{X}'_{\bar{s}},g'(\bar{s}))\) pour \(\bar{s}\) point géométrique de \(\mathrm{U}\) s’organisent en un « système local » sur \(\mathrm{U}\) (SGA 1 XIII (3.1), (3.3) et (4.5), où le lecteur trouvera un énoncé précis). Pour \(\mathbf{L}=\{\ell\}\), il agit sur \(g^*\mathscr{F}\). A la tour des quotients qui agissent trivialement sur les \(g^*\mathscr{F}/\ell^n\mathscr{F}\) correspond la tour cherchée de revêtements finis étales \(\mathrm{X}_n\subset\mathrm{I}_n\). On prend \(\mathrm{M}^n=g^*\mathrm{X}_n\).

% ---- 185-192 ----
\Page{185}
\textit{Remarque} \textbf{(1.11.2)}. — Pour que les \(\pi_1^{(\mathbf{L})}(\mathrm{X}_{\bar{s}},g(\bar{s}))\) forment un « système local » sur \(\mathrm{S}\), il suffit qu’il existe un plongement \(\mathrm{X}\hookrightarrow\overline{\mathrm{X}}\), où \(\overline{\mathrm{X}}\) est propre et lisse sur \(\mathrm{S}\), de dimension relative \(1\), et où \(\mathrm{X}\) est le complément d’un sous-schéma \(\mathrm{T}\) fini étale sur \(\mathrm{S}\) (SGA 1 XIII (4.4), (4.5)). Il suffit même qu’un tel plongement existe après un changement de base \(\mathrm{S}'\to\mathrm{S}\) fini, radiciel et surjectif (que la topologie étale ne voit pas). Ceci est toujours possible sur un ouvert de \(\mathrm{S}\) (se ramener par passages à la limite au cas où \(\mathrm{S}\) est spectre d’un corps parfait).

\textit{Remarque} \textbf{(1.11.3)}. — Avec les notations de (1.11.1), remplaçons \(\mathrm{S}\) par \(\mathrm{U}\), et supposons qu’il existe des plongements (1.11.2) :
\[
\begin{tikzcd}[column sep=2em,row sep=2em]
\mathrm{X}_1 \arrow[r,hook] \arrow[d]
  & \overline{\mathrm{X}}_1 \arrow[d]
  &[-2em] {,\ \text{complément de }\mathrm{T}_1} \\
\mathrm{X} \arrow[r,hook]
  & \overline{\mathrm{X}}
  &[-2em] {,\ \text{complément de }\mathrm{T}.}
\end{tikzcd}
\]
Les \(\mathrm{X}_n/\mathrm{X}_1\) sont des revêtements modérés. On dispose donc encore d’une tour de compactifications (1.11.2) : \(\mathrm{X}_n\hookrightarrow\overline{\mathrm{X}}_n\), complément de \(\mathrm{T}_n\) (SGA 1 XIII 5).

Soient \(\bar{s}\) un point géométrique de \(\mathrm{S}\), \(t\in\mathrm{T}_{\bar{s}}\), \(\eta_t\) le point générique de \(\overline{\mathrm{X}}_{\bar{s}}\), \(k(\bar{\eta}_t)\) une clôture algébrique de \(k(\eta_t)\) et \(\mathrm{I}_t=\operatorname{Gal}(\bar{\eta}_t/\eta_t)\). On dispose d’une classe de conjugaison d’applications de \(\mathrm{I}_t\) dans \(\pi_1(\mathrm{X}_{\bar{s}},g(\bar{s}))\). Soit \(\mathrm{M}_t^n\) l’image de \(\mathrm{I}_t\) dans \(\mathrm{M}_{\bar{s}}^n\) (un sous-groupe bien défini à conjugaison près). Nous nous proposons de vérifier que si \(\bar{s}_1\) se spécialise en \(\bar{s}_2\), et \(\bar{t}_1\) en \(\bar{t}_2\), alors \(\mathrm{M}_{t_1}^n\) se spécialise en (un conjugué de) \(\mathrm{M}_{t_2}^n\).

On peut supposer \(\mathrm{S}\) strictement local. Dans ce cas, \(t_1\) et \(t_2\) sont sur une même section \(\tau\) de \(\mathrm{T}\), et celle-ci se relève en une section \(\tau^n\) de \(\mathrm{T}_n\). Le faisceau \(\mathrm{M}^n\) agit sur \(\mathrm{T}_n\), et il suffit d’observer que les \(\mathrm{M}_{t_i}^n\) sont les fibres en \(s_1\) et \(s_2\) du stabilisateur de \(\tau^n\).

Tout ceci vaut encore si les plongements voulus n’existent que sur \(\mathrm{S}'/\mathrm{S}\) comme en (1.11.2), donc en tout cas sur un ouvert non vide de \(\mathrm{S}\).

Si \(\bar{s}\) est de caractéristique \(p\), et que \(\mathscr{F}_{\bar{s}}\) sur \(\mathrm{X}_{\bar{s}}\) est modérément ramifié en \(t\in\mathrm{T}_{\bar{s}}\), l’application de \(\mathrm{I}_t\) dans \(\mathrm{M}_{\bar{s}}^n\) se factorise par le groupe d’inertie modéré \(\mathrm{I}_t^{(p')}\sim\mathbf{Z}_{(p')}(1)\). Si \(\mathscr{F}\) est modérément ramifié le long de la section \(\tau\) de \(\mathrm{T}/\mathrm{S}\), avec \(t=\tau(\bar{s})\), elle se factorise même par \(\mathbf{Z}_{(\mathbf{L})}(1)=\displaystyle\prod_{\ell'\in\mathbf{L}}\mathbf{Z}_{\ell'}(1)\), pour \(\mathbf{L}\) l’ensemble des nombres premiers inversibles sur \(\mathrm{S}\). On dispose alors d’une classe de conjugaison d’applications \(\varphi_t^n\) de \(\mathbf{Z}_{(\mathbf{L})}(1)\) dans \(\mathrm{M}_t^n\). Ici encore, si \(\bar{s}_1\) se spécialise en \(\bar{s}_2\) et \(\bar{t}_1\) en \(\bar{t}_2\), \(\varphi_{t_1}^n\) se spécialise en \(\varphi_{t_2}^n\). Pour la preuve, sinon pour l’énoncé, voir SGA 1 XIII (2.11).

\textbf{(1.11.4)} \textit{Preuve de} (1.3.1). — Les notations sont celles de (1.3.1). On peut supposer, et on suppose, que \(\dim\mathrm{X}_0\geq1\). Observons qu’il suffit de prouver l’assertion après une extension finie préliminaire du corps des scalaires \(\mathbf{F}_q\).

Notons \(\mathrm{H}_1(\mathrm{X},\widehat{\mathbf{Z}})\) le plus grand quotient abélien de \(\pi_1(\mathrm{X},\bar{x})\), et \(\mathrm{H}_1(\mathrm{X},\widehat{\mathbf{Z}}_{(p')})\) le plus grand quotient premier à \(p\). L’image de \(\pi_1(\mathrm{X},\bar{x})\) dans \(\mathrm{W}(\mathrm{X}_0,\bar{x})^{ab}\) est le plus grand
\Page{186}
quotient de \(\mathrm{H}_1(\mathrm{X},\widehat{\mathbf{Z}})\) sur lequel Frobenius agit trivialement, et il nous faut montrer que le plus grand quotient de \(\mathrm{H}_1(\mathrm{X},\widehat{\mathbf{Z}}_{(p')})\) sur lequel Frobenius agit trivialement est fini.

Soit \(\mathrm{U}_0\) un ouvert dense quasi-projectif de \(\mathrm{X}_0\) — on prend par exemple \(\mathrm{U}_0\) affine — et plongeons \(\mathrm{U}_0\) dans un espace projectif \(\mathbf{P}_0\). Soit \(\mathrm{V}\) une section linéaire générique de dimension \(1\) de \(\mathrm{U}\) : on prend un point générique géométrique \(\bar{\eta}\) de la grassmannienne des sous-espaces linéaires de codimension \((\dim\mathrm{U}-1)\) de \(\mathbf{P}\), le sous-espace \(\mathrm{D}_{\bar{\eta}}\) (sur \(k(\bar{\eta})\)) correspondant, et, sur \(k(\bar{\eta})\), l’intersection \(\mathrm{V}=(\mathrm{U}\otimes_{\mathbf{F}}k(\bar{\eta}))\cap\mathrm{D}_{\bar{\eta}}\). Le théorème de Bertini assure que, pour \(v\) un point base dans \(\mathrm{V}\), le groupe fondamental \(\pi_1(\mathrm{V},v)\) s’envoie \textit{sur} \(\pi_1(\mathrm{U},v)\). On sait par ailleurs que \(\pi_1(\mathrm{U},v)\) s’envoie sur \(\pi_1(\mathrm{X},v)\). \textit{A fortiori}, \(\mathrm{H}_1(\mathrm{V},\widehat{\mathbf{Z}}_{(p')})\) s’envoie sur \(\mathrm{H}_1(\mathrm{X},\widehat{\mathbf{Z}}_{(p')})\). Le groupe \(\mathrm{H}_1(\mathrm{V},\widehat{\mathbf{Z}}_{(p')})\) ne change pas si on remplace la section générique \(\mathrm{V}\) par une section assez générale \(\mathrm{W}\) (SGA 1 XIII déjà cité), et la remarque liminaire permet de supposer \(\mathrm{W}\) défini sur \(\mathbf{F}_q\) (i.e. provenant de \(\mathrm{W}_0\) sur \(\mathbf{F}_q\)). Puisque \(\mathrm{H}_1(\mathrm{W},\widehat{\mathbf{Z}}_{(p')})\) s’envoie \textit{sur} \(\mathrm{H}_1(\mathrm{X},\widehat{\mathbf{Z}}_{(p')})\), il suffit de prouver (1.3.1) pour \(\mathrm{W}_0\); ceci a été fait en (1.3.2).

\textit{Remarque} \textbf{(1.11.5)}. — La proposition (1.11.1) vaut pour tout morphisme de type fini \(f\) : localisation sur \(\mathrm{S}\) et un argument à la Mayer-Vietoris ramènent à supposer \(f\) plat à fibres normales; les arguments de restriction à un ouvert (1.11.4) permettent ensuite de supposer \(f\) lisse et quasi-projectif; l’argument de section hyperplane nous ramène alors au cas de dimension relative \(1\) de (1.11.1).

\Page{187}
\section*{II. — LA MÉTHODE DE HADAMARD-DE LA VALLÉE-POUSSIN}

\section*{(2.1) La méthode}

\textbf{(2.1.1)} Soient \(\Gamma\) un groupe isomorphe à \(\mathbf{Z}\) ou à \(\mathbf{R}\), \(\omega_1\) un quasi-caractère non trivial \(\omega_1:\Gamma\to\mathbf{R}_+^*\), \(\mathrm{G}\) un groupe localement compact extension de \(\Gamma\) par un groupe compact \(\mathrm{G}^0\), \(\Sigma\) un ensemble infini dénombrable, et \((x_v)_{v\in\Sigma}\) une famille indexée par \(\Sigma\) de classes de conjugaison dans \(\mathrm{G}\). On fait les hypothèses (A), (B) suivantes. Le lecteur trouvera des exemples en (2.1.2), (2.1.3), (2.1.9) et au numéro suivant.

\((\mathrm{A}')\) \textit{Si \(\Gamma\) est isomorphe à \(\mathbf{R}\), l’extension \(\mathrm{G}\) est triviale.}

Cette hypothèse est en fait automatiquement vérifiée.

\((\mathrm{A}'')\) \textit{Si \(\Gamma\) est isomorphe à \(\mathbf{Z}\), le centre de \(\mathrm{G}\) s’envoie sur un sous-groupe d’indice fini de \(\Gamma\).}

Pour démontrer (2.2.10), seul utilisé dans la preuve du théorème principal, on pourrait se contenter du cas où \(\mathrm{G}\) est le produit de \(\Gamma=\mathbf{Z}\) par un groupe de Lie compact \(\mathrm{G}^0\), de composante connexe \(\mathrm{G}^{00}\) semi-simple.

La condition \((\mathrm{A}'')\) n’est pas nécessaire pour prouver le théorème (2.1.4) ci-dessous : le groupe n’y apparaît que via ses représentations linéaires (continues complexes, de dimension finie), et si \(\mathrm{G}\) admet une représentation linéaire dont la restriction à \(\mathrm{G}^0\) a un noyau fini, \((\mathrm{A}'')\) est automatiquement satisfaite (1.3.10).

Pour \(s\in\mathbf{C}\), posons \(\omega_s=\omega_1^s\). Notant encore \(\omega_s\) le morphisme composé \(\mathrm{G}\to\Gamma\xrightarrow{\omega_s}\mathbf{C}^*\), on pose \(\mathrm{N}v=\omega_{-1}(x_v)\). Si \(\Gamma\) est isomorphe à \(\mathbf{Z}\), on note \(q\) et \(\deg\) le nombre \(>1\) et l’isomorphisme de \(\Gamma\) avec \(\mathbf{Z}\) tels que \(\omega_1(\gamma)=q^{-\deg(\gamma)}\). On a \(\omega_s=\omega_{s+2\pi i/\log q}\). Notant encore \(\deg\) le morphisme composé \(\mathrm{G}\to\Gamma\xrightarrow{\deg}\mathbf{Z}\), on pose \(\deg(v)=\deg(x_v)\).

Soit \(g\) un élément du centre de \(\mathrm{G}\) d’image non triviale dans \(\Gamma\). Il en existe par hypothèse. Une représentation linéaire complexe \(\tau:\mathrm{G}\to\operatorname{GL}(\mathrm{V})\) est unitarisable si et seulement si \(\tau(g)\) l’est : d’une structure hermitienne invariante par \(g\) on en déduit une invariante par \(\mathrm{G}\) par intégration sur le groupe compact \(\mathrm{G}/g^{\mathbf{Z}}\). Si \(\tau\) est irréductible, donc \(\tau(g)\) scalaire, il existe un unique nombre réel \(\sigma\) tel que \(|\tau(g)|=\omega_\sigma(g)\), et \(\tau.\omega_{-\sigma}\) est unitarisable. On appelle \(\sigma\) la \textit{partie réelle} \(\mathscr{R}(\tau)\) de \(\tau\). On a \(\mathscr{R}(\tau.\omega_s)=\mathscr{R}(\tau)+\mathscr{R}(s)\).

Soient \(\widetilde{\mathrm{G}}\) l’ensemble des classes d’isomorphie de représentations irréductibles de \(\mathrm{G}\), et \(\widehat{\mathrm{G}}\) l’ensemble de celles qui sont unitaires. Les ensembles \(\{\tau.\omega_s\mid s\in\mathbf{C}\}\) forment une partition de \(\widetilde{\mathrm{G}}\), et l’application \(s\mapsto\tau.\omega_s\) identifie \(\{\tau.\omega_s\mid s\in\mathbf{C}\}\) au quotient de \(\mathbf{C}\) par
\Page{188}
un sous-groupe discret de \(i\mathbf{R}\) si \(\Gamma\simeq\mathbf{Z}\), et à \(\mathbf{C}\) si \(\Gamma\simeq\mathbf{R}\). Nous munirons \(\widetilde{\mathrm{G}}\) de la structure de surface de Riemann pour laquelle il est la somme disjointe de ces quotients.

\((\mathrm{B}')\) \textit{Pour tout \(v\in\Sigma\), on a \(\mathrm{N}v>1\).}

\((\mathrm{B}'')\) \textit{Le produit infini \(\displaystyle\prod_{v\in\Sigma}(1-\mathrm{N}v^{-s})^{-1}\) converge absolument pour \(\mathscr{R}(s)>1\).}

Pour \(\Gamma\) isomorphe à \(\mathbf{Z}\), ces conditions peuvent se récrire : \(\deg(v)>0\), et, pour tout \(\varepsilon>0\) :
\[
\#\{v\mid\deg(v)=n\}\leq\mathrm{O}(q^{(1+\varepsilon)n}).
\]
L’hypothèse \((\mathrm{B}'')\) assure que, pour \(\tau\in\widetilde{\mathrm{G}}\), le produit infini :
\[
\mathrm{L}(\tau)=\prod_{v\in\Sigma}\det(1-\tau(x_v))^{-1}
\]
converge absolument pour \(\mathscr{R}(\tau)>1\). Chaque facteur est holomorphe en \(\tau\) pour \(\mathscr{R}(\tau)>0\), et \(\mathrm{L}(\tau)\) est fonction holomorphe de \(\tau\) pour \(\mathscr{R}(\tau)>1\). On pose \(\mathrm{L}(\tau,s)=\mathrm{L}(\tau.\omega_s)\).

\textit{Exemple} \textbf{(2.1.2)}. — Supposons que \(\Gamma\simeq\mathbf{R}\), et donnons-nous un isomorphisme de l’extension \(\mathrm{G}\) avec \(\mathrm{G}^0\times\Gamma\). Soit \(x_v^0\) la composante de \(x_v\) dans \(\mathrm{G}^0\). La situation est entièrement décrite par le groupe compact \(\mathrm{G}^0\), la famille des classes de conjugaison \(x_v^0\) dans \(\mathrm{G}^0\) et la famille des nombres réels \(\mathrm{N}v>1\). Les axiomes (A) (B) deviennent : le produit infini \(\displaystyle\prod_v(1-\mathrm{N}v^{-s})^{-1}\) converge absolument pour \(\mathscr{R}(s)>1\). Cette situation est considérée par Serre dans [11], IA2, et nous renvoyons le lecteur aux exemples que donne Serre.

\textit{Exemple} \textbf{(2.1.3)}. — Soient \(\mathrm{X}_0\) une courbe lisse absolument irréductible sur \(\mathbf{F}_q\), \(\mathrm{K}\) son corps des fractions et \(\bar{x}\) un point géométrique de \(\mathrm{X}_0\). Soit \(\mathrm{G}\) l’extension :
\[
\begin{tikzcd}[column sep=1.2em]
0 \arrow[r]
  & \pi_1(\mathrm{X},\bar{x}) \arrow[r]
  & \mathrm{W}(\mathrm{X}_0,\bar{x}) \arrow[r,"\deg"]
  & \mathbf{Z} \arrow[r]
  & 0,
\end{tikzcd}
\]
\(\omega_1(n)=q^{-n}\), \(\Sigma=|\mathrm{X}_0|\) et \(x_v\) le Frobenius géométrique \(\mathrm{F}_v\). On a \(\omega_{-1}(x_v)=\mathrm{N}v\). La condition \((\mathrm{A}'')\) n’est en général pas satisfaite, mais cf. les remarques après \((\mathrm{A}'')\). Les conditions (B) sont satisfaites. Les fonctions \(\mathrm{L}(\tau,s)\) sont les fonctions \(\mathrm{L}\) d’Artin associées aux représentations irréductibles de \(\operatorname{Gal}(\overline{\mathrm{K}}/\mathrm{K})\), non ramifiées sur \(\mathrm{X}_0\), ou plutôt ces fonctions \(\mathrm{L}\) privées des facteurs locaux hors de \(|\mathrm{X}_0|\) et translatées en \(s\).

\textit{Théorème} \textbf{(2.1.4)}. — \textit{Avec les hypothèses et notations de \textup{(2.1.1)}, supposons que la fonction \(\mathrm{L}(\tau)\) se prolonge en une fonction méromorphe de \(\tau\) pour \(\mathscr{R}(\tau)\geq1\) et que, dans cette région \(\mathscr{R}(\tau)\geq1\), elle est holomorphe sauf pour un pôle simple en \(\omega_1\). Alors, la fonction \(\mathrm{L}(\tau)\) ne s’annule pas pour \(\mathscr{R}(\tau)=1\), sauf pour au plus une représentation \(\tau_0\), de dimension \(1\) et définie par un caractère \(\omega_1\varepsilon\) avec \(\varepsilon\) d’ordre \(2\).}

Pour une représentation \(\tau\) non nécessairement irréductible, on définira encore \(\mathrm{L}(\tau)\) et \(\mathrm{L}(\tau,s)\) comme en (2.1.1). On pose \(\mathrm{L}'(\tau)=\left.\dfrac{d}{ds}\mathrm{L}(\tau,s)\right|_{s=0}\). On a dans le domaine de convergence :
\[
-\frac{\mathrm{L}'}{\mathrm{L}}(\tau)
 =\sum_{\substack{v\in\Sigma\\ n>0}}
   \log\mathrm{N}v.\operatorname{Tr}\tau(x_v^n).
\tag*{\textbf{(2.1.4.1)}}
\]

\Page{189}
On peut plus généralement prendre pour \(\tau\) une représentation virtuelle, i.e. un élément du groupe de Grothendieck de la catégorie des représentations de \(\mathrm{G}\). Pour \(\tau\) unitaire et \(\sigma\) réel \(>1\), \(\omega_\sigma\tau\) est dans le domaine de convergence et (2.1.4.1), appliqué à \(\omega_\sigma.\tau\), se récrit :
\[
-\frac{\mathrm{L}'}{\mathrm{L}}(\omega_\sigma\tau)
 =\sum_{\substack{v\in\Sigma\\ n>0}}
   \bigl(\log\mathrm{N}v.(\mathrm{N}v)^{-n\sigma}\bigr)
   \operatorname{Tr}\tau(x_v^n).
\tag*{\textbf{(2.1.4.2)}}
\]

Soit \(\mu\) une mesure sur \(\mathrm{G}\), voire sur l’ensemble des classes de conjugaison de \(\mathrm{G}\). Pour toute représentation unitaire virtuelle \(\tau\) de \(\mathrm{G}\), nous poserons :
\[
\widehat{\mu}(\tau)=\int\operatorname{Tr}\tau(g)\,d\mu.
\tag*{\textbf{(2.1.4.3)}}
\]
L’intégrale converge si la masse totale de \(|\mu|\) est finie. Nous appellerons alors la fonction \(\tau\mapsto\widehat{\mu}(\tau)\) la \textit{transformée de Fourier} de \(\mu\). Si l’on n’impose pas à \(\tau\) d’être unitaire, on parle de même de \textit{transformée de Fourier-Laplace}. Si \(\mu\) est positive, de masse totale finie, on a pour toute représentation unitaire virtuelle \(\rho\) :
\[
\widehat{\mu}(\rho\otimes\bar{\rho})\geq0
\qquad(\text{pour }\mu\geq0).
\tag*{\textbf{(2.1.4.4)}}
\]

La formule (2.1.4.1) exprime que, pour \(\sigma>1\), \(\Lambda_\sigma(\tau)=-\dfrac{\mathrm{L}'}{\mathrm{L}}(\omega_\sigma\tau)\) est la transformée de Fourier de la mesure \textit{positive} de masse totale finie :
\[
\mu_\sigma
 =\sum_{\substack{v\in\Sigma\\ n>0}}
   \log\mathrm{N}v.\mathrm{N}v^{-n\sigma}.\delta[x_v^n]
\tag*{\textbf{(2.1.4.5)}}
\]
sur \(\mathrm{G}\) (on écrit \(\delta[a]\) pour la mesure de Dirac en \(a\)). Pour toute représentation unitaire virtuelle \(\rho\) de \(\mathrm{G}\), de caractère réel et \(\geq0\), on a donc \(\Lambda_\sigma(\rho)\geq0\) pour \(\sigma>1\); en particulier, pour toute représentation unitaire virtuelle \(\rho\), on a :
\[
\Lambda_\sigma(\rho\otimes\bar{\rho})\geq0
\qquad(\text{pour }\sigma>1).
\tag*{\textbf{(2.1.4.6)}}
\]

Pour \(\tau\in\widehat{\mathrm{G}}\), soit \(\nu(\tau)\) l’ordre du pôle (l’opposé de l’ordre du zéro) de \(\mathrm{L}\) en \(\tau\omega_1\). On prolonge \(\nu\) par additivité au groupe de Grothendieck de la catégorie des représentations unitaires de \(\mathrm{G}\). Pour \(\tau\) dans ce groupe, la fonction \(-\dfrac{\mathrm{L}'}{\mathrm{L}}(\tau.\omega_s)\) n’a que des pôles simples, et son résidu en \(\tau\omega_1\) \((\tau\in\widehat{\mathrm{G}})\) est \(\nu(\tau)\). Faisant tendre \(\sigma\) vers \(1\) dans (2.1.4.6), on trouve donc que pour toute représentation unitaire virtuelle de \(\mathrm{G}\) :
\[
\nu(\rho\otimes\bar{\rho})\geq0.
\]
Le théorème résulte du lemme suivant

\textit{Lemme} \textbf{(2.1.5)}. — \textit{Soit \(\nu\) une fonction à valeurs entières sur \(\widehat{\mathrm{G}}\). On suppose que :}

\noindent a) \textit{pour la représentation triviale \(1\), \(\nu(1)=1\);}

\noindent b) \(\nu(\tau)=\nu(\bar{\tau})\);

\noindent c) \(\nu(\tau)\leq0\) \textit{pour} \(\tau\neq1\);

\Page{190}
\noindent d) \textit{prolongeant \(\nu\) par additivité au groupe de Grothendieck des représentations unitaires de \(\mathrm{G}\), on a, pour toute représentation unitaire \(\rho\) de \(\mathrm{G}\) :}
\[
\nu(\rho\otimes\bar{\rho})\geq0.
\]

\textit{Alors, \(\nu(\tau)=0\) pour \(\tau\neq1\), sauf pour au plus une représentation \(\tau_0\) de dimension un et définie par un caractère d’ordre deux.}

Dans ce lemme, \(\mathrm{G}\) n’apparaît que via ses représentations unitaires. Ceci nous permet de le remplacer par un groupe compact (son adhérence dans le produit de ses représentations unitaires) et de ramener (2.1.5) à la variante suivante.

\textit{Variante} \textbf{(2.1.6)}. — \textit{Même énoncé que \textup{(2.1.5)}, mais cette fois \(\mathrm{G}\) est un groupe compact (et \(\widehat{\mathrm{G}}\) l’ensemble de ses classes d’isomorphie de représentations irréductibles).}

Dans la preuve de (2.1.6), si \(\rho\) est une représentation de \(\mathrm{G}\), nous écrirons \(\rho(g)\) pour la valeur en \(g\) du \textit{caractère} de \(\rho\), et nous identifierons représentations et caractères. Pour \(\tau\in\widehat{\mathrm{G}}\), nous notons \([\rho:\tau]\) la multiplicité avec laquelle \(\tau\) apparaît dans \(\rho\). On a \([\rho:\tau]=\int\overline{\rho(g)}\tau(g)\,dg\), la mesure de Haar \(dg\) étant normalisée pour être de masse totale \(1\).

\textit{Lemme} \textbf{(2.1.7)}. — \textit{Soit \(\mathrm{T}\) un ensemble fini de représentations irréductibles du groupe compact \(\mathrm{G}\). Pour tout \(\varepsilon>0\), il existe une représentation \(\rho\) de \(\mathrm{G}\) telle que :}
\[
[\rho\otimes\bar{\rho}:\tau]
 \geq(1-\varepsilon)\dim\tau.[\rho\otimes\bar{\rho}:1]
 \qquad\text{\itshape pour tout }\tau\in\mathrm{T}.
\]

Soit \(\mathrm{X}\) l’algèbre des combinaisons linéaires à coefficients complexes de caractères. On sait que \(\mathrm{X}\) est stable par \(\varphi(g)\mapsto\bar{\varphi}(g)\) et par \(\varphi(g)\mapsto\varphi(g^{-1})\), et est dense dans l’algèbre des fonctions continues centrales sur \(\mathrm{G}\). Quels que soient \(\varepsilon_1,\varepsilon_2>0\), il existe donc \(\rho_0\in\mathrm{X}\) qui est :

\noindent\textit{a)} réel, strictement positif, tel que \(\rho_0(g)=\rho_0(g^{-1})\), et tel que sa masse soit concentrée autour de \(e\) au sens suivant :

\noindent\textit{b)} il existe un voisinage \(\mathrm{V}\) de \(e\) dans lequel \(|\tau(g)-\tau(e)|<\varepsilon_1\tau(e)\) pour tout \(\tau\in\mathrm{T}\), et tel que :
\[
\int_{\mathrm{G}-\mathrm{V}}|\rho_0(g)|^2\,dg
 \leq\varepsilon_2\int_{\mathrm{G}}|\rho_0(g)|^2.
\]

D’après \textit{a)}, \(\rho_0\) est combinaison linéaire à coefficients \textit{réels} de caractères irréductibles. On peut même supposer les coefficients rationnels, voire entiers (multiplier \(\rho_0\) par un entier).

Posons :
\begin{align*}
\rho_0&=\sum n(\chi)\chi
  \qquad\text{(somme sur les caractères irréductibles)}\\
\rho^\pm&=\sum_{\pm n(\chi)>0}\pm n(\chi)\chi
  \qquad\text{(on a }\rho_0=\rho^+-\rho^-\text{)}\\
\rho&=\rho^++\rho^-.
\tag*{et :}
\end{align*}

\Page{191}
Prenons \(\varepsilon_1+\varepsilon_2<\varepsilon\). Pour \(\tau\in\mathrm{T}\), on a alors :
\[
\begin{aligned}
&[\rho\otimes\bar{\rho}:\tau]
 =[\rho_0\otimes\bar{\rho}_0:\tau]
   +2[\rho^+\otimes\bar{\rho}^-+\rho^-\otimes\bar{\rho}^+:\tau]\\
&\quad\geq[\rho_0\otimes\bar{\rho}_0:\tau]
 =\int|\rho_0(g)|^2\tau(g)\,dg\\
&\quad\geq(1-\varepsilon)\tau(e)\int|\rho_0(g)|^2\,dg
 =(1-\varepsilon)\tau(e)[\rho_0\otimes\bar{\rho}_0:1]
 =(1-\varepsilon)\dim(\tau)[\rho\otimes\bar{\rho}:1],
\end{aligned}
\]
d’où le lemme.

\textbf{(2.1.8)} Prouvons (2.1.6). Soit \(\mathrm{T}\) une partie finie de \(\widehat{\mathrm{G}}\), avec \(1\notin\mathrm{T}\), et soit \(\rho\) comme dans (2.1.7). Puisque \(\nu(\sigma)\leq0\) pour \(\sigma\neq1\), on a :
\[
\begin{aligned}
0\leq\nu(\rho\otimes\bar{\rho})
 =\sum_\sigma[\rho\otimes\bar{\rho}:\sigma]\nu(\sigma)
 &\leq[\rho\otimes\bar{\rho}:1]\nu(1)
   +\sum_{\tau\in\mathrm{T}}[\rho\otimes\bar{\rho}:\tau]\nu(\tau)\\
 &\leq[\rho\otimes\bar{\rho}:1].
 \biggl(\nu(1)+(1-\varepsilon)
   \sum_{\tau\in\mathrm{T}}\dim(\tau).\nu(\tau)\biggr).
\end{aligned}
\]
Faisons tendre \(\varepsilon\) vers \(0\) et prenons \(\mathrm{T}\) de plus en plus grand.

On trouve :
\[
\sum_{\tau\in\widehat{\mathrm{G}}}\dim\tau.\nu(\tau)\geq0
\]
et (2.1.6) en résulte.

\textbf{(2.1.9)} Expliquons le titre du chapitre. Avec les notations de (2.1.2), le cas de la fonction \(\zeta\) de Riemann correspond à \(\mathrm{G}^0=\{e\}\), \(\Sigma=\) l’ensemble des nombres premiers, et \(\mathrm{N}p=p\). On a :
\[
-\frac{\zeta'(s)}{\zeta(s)}
 =\sum_{p,n}\log p.p^{-ns}
\]
de sorte que, pour \(\sigma>1\), la fonction det \(\Lambda_\sigma(t)=-\dfrac{\zeta'(\sigma+it)}{\zeta(\sigma+it)}\) est la transformée de Fourier de la mesure positive \(\displaystyle\sum_{p,n}\log p.p^{-n\sigma}.\delta[n\log p]\). C’est donc une fonction positive au sens de Bochner; en particulier, elle satisfait la sempiternelle inégalité :
\[
3\mathscr{R}\Lambda_\sigma(0)
 +4\mathscr{R}\Lambda_\sigma(t)
 +2\mathscr{R}\Lambda_\sigma(2t)\geq0
\]
dont Mertens déduit le théorème de Hadamard et de La Vallée-Poussin selon lequel \(\zeta\) ne s’annule pas pour \(\mathscr{R}s=1\).

\textbf{(2.1.10)} Il est classique que la non-annulation de \(\mathrm{L}(\tau)\) pour \(\mathscr{R}(\tau)=1\) implique une équidistribution des \(x_v\). Le cas où \(\Gamma\) est isomorphe à \(\mathbf{R}\) est traité en détail dans Serre [11], IA2. Je traiterai ici du cas où \(\Gamma\) est isomorphe à \(\mathbf{Z}\). Faisons donc les hypothèses :

\noindent (C) la fonction \(\mathrm{L}(\tau)\) admet un prolongement méromorphe pour \(\mathscr{R}(\tau)\geq1\); dans ce domaine, elle a un pôle simple en \(\omega_1\), et aucun autre pôle, et elle ne s’annule pas pour \(\mathscr{R}(\tau)=1\);

\noindent (D) \(\Gamma\) est isomorphe à \(\mathbf{Z}\).

\Page{192}
Soit \(z\) un élément central de \(\mathrm{G}\), de degré \(d>0\). L’espace \(\mathrm{G}^{\natural}\) des classes de conjugaison de \(\mathrm{G}\) est le quotient de \(\mathrm{G}\) par le groupe compact \(\mathrm{G}/z^{\mathbf{Z}}\) agissant par automorphismes intérieurs. Il est somme disjointe des \(\mathrm{G}_n^{\natural}\), pour \(\mathrm{G}_n^{\natural}\) l’ensemble des classes de conjugaison de degré \(n\), et la multiplication par \(z^k\) induit un isomorphisme de \(\mathrm{G}_n^{\natural}\) avec \(\mathrm{G}_{n+kd}^{\natural}\).

Considérons les mesures suivantes sur \(\mathrm{G}\) et \(\mathrm{G}^{\natural}\) :
\[
\mu^{\natural}
 =\sum_{\substack{v\in\Sigma\\ n>0}}
   \deg(v)q^{-n\deg(v)}\delta[x_v^n],
\]
\noindent \(dg=\) la mesure de Haar sur \(\mathrm{G}\) pour laquelle \(\mathrm{G}^0\) est de volume \(1\),

\noindent \(\mu_0=\) le produit de \(dg\) par la fonction caractéristique de l’ouvert où \(\omega_{-1}>1\),

\noindent \(\mu_0^{\natural}=\) la projection de \(\mu_0\) sur \(\mathrm{G}^{\natural}\).

\textit{Proposition} \textbf{(2.1.11)}. — \textit{Si \textup{(C)} et \textup{(D)} sont satisfaites, la transformée de Fourier-Laplace de \(\mu^{\natural}-\mu_0^{\natural}\) (qui converge pour \(\tau\in\widetilde{\mathrm{G}}\), \(\mathscr{R}(\tau)>0\)) se prolonge en une fonction holomorphe de \(\tau\) pour \(\mathscr{R}(\tau)\geq0\).}

Par (2.1.4.1), la transformée de Fourier-Laplace de \(\mu^{\natural}\) est \(-\dfrac{1}{\log q}\dfrac{\mathrm{L}'}{\mathrm{L}}(\omega_1.\tau)\). Celle de \(\mu_0^{\natural}\) s’annule en dehors des \(\omega_s\), et sa valeur en \(\omega_s\) est \(\dfrac{q^{-s}}{1-q^{-s}}\). Dans le domaine \(\mathscr{R}(\tau)\geq0\), \(\widehat{\mu}^{\natural}\) et \(\widehat{\mu}_0^{\natural}\) sont donc méromorphes, avec pour seul pôle un pôle simple de résidu \(1/\log q\) en \(\omega_0\). La proposition en résulte.

Pour comparer les mesures \(\mu^{\natural}\) et \(\mu_0^{\natural}\) sur \(\mathrm{G}_n^{\natural}\), nous les ramènerons dans un \(\mathrm{G}_i^{\natural}\) avec \(0\leq i<d\), par translation par une puissance de \(z\). On a

\textit{Théorème} \textbf{(2.1.12)}. — \textit{Sous les hypothèses \textup{(C)} et \textup{(D)}, quel que soit \(i\), la mesure \(z^{-n}(\mu^{\natural}\mid\mathrm{G}_{i+nd}^{\natural})\) sur \(\mathrm{G}_i^{\natural}\) converge vaguement vers \(z^{-n}(\mu_0^{\natural}\mid\mathrm{G}_{i+nd}^{\natural})\) pour \(n\to\infty\).}

Des arguments bien connus ([11], IA1) montrent que (2.1.12) équivaut à l’assertion suivante, où \(z\) n’apparaît plus :

\textit{Corollaire} \textbf{(2.1.13)}. — \textit{Pour toute représentation unitaire \(\tau\) de \(\mathrm{G}\), de caractère encore noté \(\tau\), on a \(\displaystyle\int_{\mathrm{G}_n^{\natural}}\tau(g)(\mu^{\natural}-\mu_0^{\natural})\to0\) pour \(n\to\infty\).}

Prouvons (2.1.13). Pour \(n>0\), posons \(a_n=\displaystyle\int_{\mathrm{G}_n^{\natural}}\tau(g)(\mu^{\natural}-\mu_0^{\natural})\). En \(\tau.\omega_s\), la transformée de Fourier-Laplace de \(\mu^{\natural}-\mu_0^{\natural}\) vaut \(\sum a_nq^{-ns}\). D’après (2.1.11), la série de puissance \(\sum a_nt^n\) est une fonction holomorphe de \(t\) pour \(|t|\leq1\), donc aussi pour \(t\) dans un disque un peu plus grand, de sorte que \(a_n\) décroît comme une série géométrique.

\section*{(2.2) — Pureté et compacité}

\textit{Les conventions \textup{(0.7)} sont en vigueur.}

\textbf{(2.2.1)} Soit \(\mathrm{G}_{\mathbf{C}}\) un schéma en groupes sur \(\mathbf{C}\), extension de \(\mathbf{Z}\) par un groupe algébrique \(\mathrm{G}_{\mathbf{C}}^0\), dont la composante neutre \(\mathrm{G}_{\mathbf{C}}^{00}\) est semi-simple :
\[
\begin{tikzcd}[column sep=1.2em]
0 \arrow[r]
  & \mathrm{G}_{\mathbf{C}}^0 \arrow[r]
  & \mathrm{G}_{\mathbf{C}} \arrow[r]
  & \mathbf{Z} \arrow[r]
  & 0.
\end{tikzcd}
\tag*{\textbf{(2.2.1.1)}}
\]

% ---- 193-200 ----
\Page{193}

D’après (1.3.10), le centre \(\mathrm Z_{\mathbf C}\) de \(\mathrm G_{\mathbf C}\) s’envoie sur un sous-groupe d’indice fini de \(\mathbf Z\). Le schéma en groupe \(\mathrm G_{\mathbf C}/\mathrm Z_{\mathbf C}\) est donc un groupe algébrique (i.e. est de type fini). Choisissons un sous-groupe compact maximal \(\mathrm U\) dans \(\mathrm G_{\mathbf C}/\mathrm Z_{\mathbf C}\), et notons \(\mathrm G_{\mathbf R}\) son image inverse dans \(\mathrm G_{\mathbf C}\). L’intersection \(\mathrm G_{\mathbf R}^0\) de \(\mathrm G_{\mathbf R}\) avec \(\mathrm G_{\mathbf C}^0\) est un sous-groupe compact maximal de \(\mathrm G_{\mathbf C}^0\), et l’on dispose d’une suite exacte :
\[
0\to\mathrm G_{\mathbf R}^0\to\mathrm G_{\mathbf R}\to\mathbf Z\to0.
\tag*{\textbf{(2.2.1.2)}}
\]

\textit{Lemme} \textbf{(2.2.2)}. — \textit{Deux éléments de \(\mathrm G_{\mathbf R}\) conjugués dans \(\mathrm G_{\mathbf C}\) sont conjugués dans \(\mathrm G_{\mathbf R}\).}

Si \(\mathrm K\) est un groupe compact, ses classes de conjugaison sont séparées par les caractères de ses représentations complexes irréductibles, et on sait que celles-ci se prolongent en des représentations algébriques de son complexifié \(\mathrm K_{\mathbf C}\). Deux éléments de \(\mathrm K\) conjugués dans \(\mathrm K_{\mathbf C}\) sont donc conjugués dans \(\mathrm K\). Si \(\mathrm Z'\) est le sous-groupe de \(\mathrm G_{\mathbf R}\) engendré par un élément central de degré non nul, on prouve (2.2.2) en appliquant ce qui précède au groupe compact \(\mathrm G_{\mathbf R}/\mathrm Z'\), et en observant que deux éléments de \(\mathrm G_{\mathbf R}\) de même degré sont conjugués si leurs images dans \(\mathrm G_{\mathbf R}/\mathrm Z'\) le sont.

\textbf{(2.2.3)} Un élément de \(\mathrm G_{\mathbf C}\) est dit \textit{semi-simple} si son image dans le groupe algébrique \(\mathrm G_{\mathbf C}/\mathrm Z_{\mathbf C}\) l’est, \textit{unipotent} s’il est contenu dans \(\mathrm G_{\mathbf C}^0\) et unipotent dans ce groupe algébrique. Tout élément \(g\) de \(\mathrm G_{\mathbf C}\) s’écrit de façon unique comme produit \(g=g_sg_u=g_ug_s\), avec \(g_s\) semi-simple et \(g_u\) unipotent. Cela résulte du même fait pour \(\mathrm G_{\mathbf C}/\mathrm Z_{\mathbf C}\), et de ce que les éléments unipotents de \(\mathrm G_{\mathbf C}\) s’envoient bijectivement sur ceux de \(\mathrm G_{\mathbf C}/\mathrm Z_{\mathbf C}\).

\textbf{(2.2.4)} Soit \(\mathrm X_0\) un schéma de type fini sur \(\mathbf F_q\), normal et géométriquement connexe, muni d’un point base \(\bar{x}\in\mathrm X\). Le groupe de Weil \(\mathrm W(\mathrm X_0,\bar{x})\) est une extension (1.1.13.1) de \(\mathbf Z\) par le groupe fondamental géométrique \(\pi_1(\mathrm X,\bar{x})\). Soit un morphisme d’extensions :
\[
\begin{tikzcd}[ampersand replacement=\&, column sep=1.5em, row sep=2em]
0 \arrow[r] \&
\pi_1(\mathrm X,\bar{x}) \arrow[r] \arrow[d] \&
\mathrm W(\mathrm X_0,\bar{x}) \arrow[r] \arrow[d] \&
\mathbf Z \arrow[r] \arrow[d,equal] \& 0 \\
0 \arrow[r] \&
\mathrm G^0 \arrow[r] \&
\mathrm G \arrow[r] \&
\mathbf Z \arrow[r] \& 0,
\end{tikzcd}
\]
satisfaisant aux conditions \textit{a)}, \textit{b)}, \textit{c)} suivantes :

\textit{a)} \(\mathrm G^0\) est un groupe algébrique sur \(\overline{\mathbf Q}_\ell\), extension d’un groupe fini par un groupe semi-simple.

Cette hypothèse implique que \(\mathrm G\) admet des représentations linéaires fidèles.

\textit{b)} Le groupe \(\mathrm G^0\) peut être défini sur une extension finie suffisamment grande \(\mathrm E\) de \(\mathbf Q_\ell\) dans \(\overline{\mathbf Q}_\ell\), et \(\mathrm G^0(\overline{\mathbf Q}_\ell)\) est la réunion des groupes topologiques \(\mathrm G^0(\mathrm F)\), pour \(\mathrm F\) une extension finie de \(\mathbf Q_\ell\) dans \(\overline{\mathbf Q}_\ell\) contenant \(\mathrm E\). On suppose que la première flèche verticale : \(\pi_1(\mathrm X,\bar{x})\to\mathrm G^0(\overline{\mathbf Q}_\ell)\) se factorise par un homomorphisme continu de \(\pi_1(\mathrm X,\bar{x})\) dans un \(\mathrm G(\mathrm F)\), et que l’image de \(\pi_1(\mathrm X,\bar{x})\) est Zariski-dense dans \(\mathrm G^0\).

\Page{194}

En conséquence, pour toute représentation linéaire \((\mathrm V,\tau)\) de \(\mathrm G^0\) (resp. de \(\mathrm G\)), on a \(\mathrm V^{\mathrm G^0}=\mathrm V^{\pi_1(\mathrm X,\bar{x})}\). Chaque représentation linéaire \((\mathrm V,\tau)\) de \(\mathrm G\) définit une représentation \(\ell\)-adique de \(\mathrm W(\mathrm X_0,\bar{x})\), d’où un faisceau lisse \(\mathscr F_0(\tau)\), muni d’un isomorphisme de \(\mathrm W(\mathrm X_0,\bar{x})\)-représentations \(\mathscr F_0(\tau)_{\bar{x}}\sim\mathrm V\). Via cet isomorphisme, on a :
\[
\mathrm V^{\mathrm G^0}
=\mathrm V^{\pi_1(\mathrm X,\bar{x})}
=\mathrm H^0(\mathrm X,\mathscr F(\tau)).
\]

\textit{c)} Il existe une représentation linéaire \(\mathrm V\) de \(\mathrm G\), dont la restriction à \(\mathrm G^0\) soit de noyau fini, telle que le faisceau lisse correspondant soit \(\iota\)-mixte.

D’après (2.2.8) (i), pour toute représentation \((\mathrm V,\tau)\), \(\mathscr F_0(\tau)\) est alors \(\iota\)-mixte. La conjecture (1.2.9) implique que la condition \textit{c)} est toujours remplie.

\textit{Exemple} \textbf{(2.2.5)}. — Soient \(\mathscr F_0\) un faisceau lisse \(\iota\)-mixte sur \(\mathrm X_0\), \(\mathrm V=\mathscr F_{0\bar{x}}\) et supposons que la restriction à \(\pi_1(\mathrm X,\bar{x})\) de la représentation \(\mathrm V\) de \(\mathrm W(\mathrm X_0,\bar{x})\) soit semi-simple (par exemple : \(\mathscr F_0\) semi-simple). On définit \(\mathrm G\) comme en (1.3.7), et la preuve de (1.3.9) montre que les hypothèses de (2.2.4) sont satisfaites.

Étendons les scalaires de \(\overline{\mathbf Q}_\ell\) à \(\mathbf C\), par l’isomorphisme \(\iota\). On obtient une suite exacte (2.2.1.1). Choisissons \(\mathrm G_{\mathbf R}\) comme en (2.2.1).

\textit{Proposition} \textbf{(2.2.6)}. — \textit{Soit \(v\in|\mathrm X_0|\), et soit \(\iota\mathrm F_v\) l’image d’un Frobenius géométrique \(\mathrm F_v\in\mathrm W(\mathrm X_0,\bar{x})\) par l’application composée \(\mathrm W(\mathrm X_0,\bar{x})\to\mathrm G\xrightarrow{\iota}\mathrm G_{\mathbf C}\). La composante semi-simple \((\iota\mathrm F_v)_s\) de \(\iota\mathrm F_v\) est alors conjuguée à un élément de \(\mathrm G_{\mathbf R}\), et ce dernier est unique à \(\mathrm G_{\mathbf R}\)-conjugaison près.}

L’unicité résulte de (2.2.2). Prouvons l’existence. Soient \(\mathrm V\) comme en (2.2.4) \textit{c)} et \(\mathrm V_\alpha\) les constituants de la représentation \(\mathrm V\). Si \(\mathrm V'_\alpha\) se déduit de \(\mathrm V_\alpha\) par torsion (1.2.7) et que \(\mathrm V'=\bigoplus\mathrm V'_\alpha\), la représentation \(\mathrm V'\) satisfait encore les hypothèses de \textit{c)}. Remplaçant \(\mathrm V\) par \(\mathrm V'\) convenable, on peut donc supposer que \(\mathrm G\) agit sur \(\mathrm V\) via un quotient \(\mathrm G/\mathrm Z'\), avec \(\deg\mathrm Z'\) d’indice fini dans \(\mathbf Z\). Les faisceaux déduits des constituants de \(\mathrm V'\) sont alors ponctuellement \(\iota\)-purs, et de poids déterminantiel \(0\) (cf. (1.3.12)), donc ponctuellement \(\iota\)-purs, de poids \(0\) : pour \(v\in|\mathrm X_0|\) et \(\alpha\) une valeur propre de \(\mathrm F_v\) sur un quelconque constituant de \(\mathrm V\), on a \(|\iota\alpha|=1\). Notons \(\tau\) la représentation de \(\mathrm G\) sur \(\mathrm V\), et soit \(\tau(\iota\mathrm F_v)_s\) la composante semi-simple de \(\tau(\iota\mathrm F_v)\). Puisque \(\tau(\iota\mathrm F_v)_s\) a pour valeurs propres des nombres de valeur absolue \(1\), le groupe \((\tau(\iota\mathrm F_v)_s)^{\mathbf Z}\) est relativement compact dans \(\mathrm{GL}(\mathrm V_{\mathbf C})\). Puisque \(\mathrm V_{\mathbf C}\) est une représentation linéaire de noyau fini du groupe algébrique de type fini \(\mathrm G_{\mathbf C}/\mathrm Z'_{\mathbf C}\), on a \(\tau(\iota\mathrm F_v)_s=\tau((\iota\mathrm F_v)_s)\), et l’image dans \(\mathrm G_{\mathbf C}/\mathrm Z'_{\mathbf C}\) du groupe \((\iota\mathrm F_v)_s^{\mathbf Z}\) est relativement compacte donc contenue dans un sous-groupe compact maximal; ceux-ci étant tous conjugués, \((\iota\mathrm F_v)_s\) est contenu dans un conjugué de \(\mathrm G_{\mathbf R}\).

\textbf{(2.2.7)} Soient maintenant \(\mathrm G_{\mathbf R}\) comme ci-dessus, extension de \(\mathbf Z\) par le groupe compact \(\mathrm G_{\mathbf R}^0\) et :
\[
\begin{aligned}
&\Sigma=|\mathrm X_0|,\\
&x_v=\text{la classe de conjugaison, dans }\mathrm G_{\mathbf R},
  \text{ d’un conjugué de }(\iota\mathrm F_v)_s,\\
&\omega_1:g\mapsto q^{-\deg(g)}.
\end{aligned}
\]

\Page{195}

Ces données sont du type (2.1.1), et satisfont aux axiomes (A) et \((\mathrm B')\). Le produit infini en \((\mathrm B'')\) est celui qui définit la fonction \(\zeta\) de \(\mathrm X_0\). Il converge pour \(\mathscr R s>\dim\mathrm X_0\) (1.4.6). Dès lors :

— si \(\dim\mathrm X_0=1\), les axiomes (A), (B) sont vérifiés;

— si \(\dim\mathrm X_0=\mathrm N\), ils le deviennent si on prend pour caractère \(\omega_1\) non pas \(q^{-\deg g}\) comme ci-dessus, mais \(q^{-\mathrm N\deg g}\). Dans la suite, nous garderons le choix (2.2.7) de \(\omega_1\), et modifierons en conséquence les énoncés du n\textsuperscript{o} (2.1) lorsque \(\mathrm N>1\). Par exemple, pour \(\tau\) unitaire, les produits \(\mathrm L(\tau,s)\) convergent pour \(\mathscr R s>\mathrm N\).

La relation entre les fonctions \(\mathrm L\) de (2.1) et les séries formelles \(\ell\)-adiques \(\mathrm Z\) de I.\,(1.14) est donnée par la proposition suivante, de vérification laissée au lecteur.

\textit{Proposition} \textbf{(2.2.8)}. — (i) \textit{Les foncteurs « extension des scalaires par \(\iota\) » et « restriction à \(\mathrm G_{\mathbf R}\) » sont des équivalences de catégories : (représentations linéaires de \(\mathrm G\))\(\to\)(représentations linéaires de \(\mathrm G_{\mathbf C}\))\(\to\)(représentations complexes continues de \(\mathrm G_{\mathbf R}\)). Pour \(\tau\) une représentation complexe continue de \(\mathrm G_{\mathbf R}\), soit \(\mathscr F_0(\tau_\ell)\) le faisceau déduit de la représentation \(\ell\)-adique correspondante \(\tau_\ell\) de \(\mathrm W(\mathrm X_0,\bar{x})\). Si \(\tau\) est irréductible, ce faisceau est ponctuellement \(\iota\)-pur de poids \(2\mathscr R(\tau)\).}

(ii) \textit{Posons \(\mathrm Z(\mathscr F_0(\tau_\ell),t)=\prod\limits_{x\in|\mathrm X_0|}\det(1-\mathrm F_xt^{\deg x},\mathscr F_0(\tau_\ell))^{-1}\). La fonction \(\mathrm L(\tau,s)\), définie pour \(\mathscr R s\) assez grand, admet le développement en puissances de \(q^{-s}\) suivant :}
\[
\mathrm L(\tau,s)
=\left.\iota\mathrm Z(\mathscr F_0(\tau_\ell),t)\right|_{t=q^{-s}}
\tag*{\textbf{(2.2.8.1)}}
\]

\textit{Corollaire} \textbf{(2.2.9)}. — \textit{Si \(\dim(\mathrm X_0)=1\), la fonction \(\mathrm L(\tau)\) remplit la condition \textup{(C)} de \textup{(2.1.10)}.}

Vérifions les hypothèses de (2.1.4), i.e. que :

\textit{a)} si la représentation irréductible \(\tau\) est unitaire, et n’est pas de forme \(\omega_{it}\), la fonction \(\mathrm L(\tau,s)\) est holomorphe pour \(\mathscr R s\geq1\);

\textit{b)} la fonction \(\mathrm L(\omega_s)\) est méromorphe pour \(\mathscr R s\geq1\), avec pour seul pôle dans cette région un pôle simple en \(\omega_s=\omega_1\).

D’après Grothendieck, on a :
\[
\mathrm Z(\mathscr F_0(\tau_\ell),t)
=
\frac{
 \det(1-\mathrm F t,\mathrm H_c^1(\mathrm X,\mathscr F(\tau_\ell))
}{
 \det(1-\mathrm F t,\mathrm H_c^0(\mathrm X,\mathscr F(\tau_\ell))
 \det(1-\mathrm F t,\mathrm H_c^2(\mathrm X,\mathscr F(\tau_\ell)))
}.
\tag*{\textbf{(2.2.9.1)}}
\]

Dans le cas \textit{a)}, la restriction de \(\tau\) à \(\mathrm G_{\mathbf R}^0\) est non triviale. La restriction de \(\tau_\ell\) à \(\pi_1(\mathrm X,\bar{x})\) est donc non triviale, et le dénominateur vaut \(1\). On conclut par (2.2.8) (ii).

Dans le cas \textit{b)}, \(\mathrm L(\omega_s)=\zeta_{\mathrm X_0}(s)\), et l’assertion est classique. On peut aussi la prouver par voie cohomologique : la formule (2.2.9.1) assure que \(\mathrm Z(\mathrm X_0,t)\) est une fonction rationnelle avec au plus des pôles simples en \(t=1\) et \(t=q^{-1}\), et s’il n’y avait pas de pôle en \(q^{-1}\), les arguments de (2.1.12) montreraient que le nombre \(\#\mathrm X_0(\mathbf F_{q^n})\) est borné indépendamment de \(n\), ce qui est absurde.

\Page{196}

Ceci nous assure, par (2.1.4), que \(\mathrm L(\tau)\) vérifie (C), sauf peut-être pour un zéro en \(\omega_1.\varepsilon\), avec \(\varepsilon\) d’ordre \(2\). Le caractère \(\varepsilon\) de \(\mathrm W(\mathrm X_0,\bar{x})\) correspond à un revêtement double \(\mathrm X'_0\) de \(\mathrm X_0\). La fonction :
\[
\zeta_{\mathrm X'_0}(s)=\zeta_{\mathrm X_0}(s).\mathrm L(\omega_s.\varepsilon)
\]
ayant un pôle simple en \(s=1\), la fonction \(\mathrm L\) ne peut s’annuler en ce point.

On peut déduire de (2.1.12) et (2.2.9) des résultats d’équidistribution à la Sato-Tate. Nous en repoussons la discussion à (3.5), pour pouvoir traiter le cas de dimension \(>1\).

La variante suivante de (2.2.9) nous sera d’un usage plus commode :

\textit{Corollaire} \textbf{(2.2.10)}. — \textit{Soient \(\mathrm X_0\) une courbe lisse sur \(\mathbf F_q\) et \(\mathscr F_0\) un faisceau lisse \(\iota\)-pur de poids \(\beta\) sur \(\mathrm X_0\). Alors, si \(\alpha\) est une valeur propre de \(\mathrm F\) sur \(\mathrm H_c^1(\mathrm X,\mathscr F)\), on a \(w_q(\alpha)<\beta+2\).}

Ce corollaire améliore la majoration triviale \(w_q(\alpha)\leq\beta+2\) de (1.8.2). Il sera amélioré en (3.2).

La validité de (2.2.10) est invariante par une extension préliminaire des scalaires de \(\mathbf F_q\) à \(\mathbf F_{q^n}\) : ceci remplace \(\mathrm F\) par \(\mathrm F^n\), et \(w_q\) par \(w_{q^n}\). Puisque \(\mathrm H_c^1(\mathrm X,\mathscr F)\) est somme sur les composantes connexes \(\mathrm Y_0\) de \(\mathrm X_0\) des \(\mathrm H_c^1(\mathrm Y,\mathscr F)\), nous sommes ramenés au cas où \(\mathrm X_0\) est absolument irréductible.

Si \(\mathscr F_0\) figure dans une suite exacte courte \(0\to\mathscr F'_0\to\mathscr F_0\to\mathscr F''_0\to0\), la suite exacte longue de cohomologie :
\[
\to\mathrm H_c^1(\mathrm X,\mathscr F')
\to\mathrm H_c^1(\mathrm X,\mathscr F)
\to\mathrm H_c^1(\mathrm X,\mathscr F'')\to
\]
nous ramène à prouver (2.2.10) pour \(\mathscr F'_0\) et \(\mathscr F''_0\) : on peut supposer \(\mathscr F_0\) irréductible.

Les hypothèses de (2.2.5) sont maintenant vérifiées par \((\mathrm X_0,\mathscr F_0)\). Si on y tient, et que, non seulement on procède à une extension du corps de base, mais encore qu’on remplace \(\mathrm X_0\) par un revêtement (ce qui est loisible, car la cohomologie de \(\mathrm X\) s’injecte dans celle d’un revêtement), on remplace \(\mathrm G\) par un sous-groupe d’indice fini, et on peut supposer \(\mathrm G\) de la forme \(\mathrm G^0\times\mathbf Z\).

Soient \(\bar{x}\in|\mathrm X|\), \(\mathrm G\) comme en (1.3.7) et \(\tau_\ell\) la représentation naturelle de \(\mathrm G\) sur \(\mathscr F_{0\bar{x}}\). On a \(\mathscr F_0=\mathscr F_0(\tau_\ell)\). On se ramène par torsion (1.2.7) à supposer que \(\beta=0\). Dans ce cas, avec les notations de (2.2), \(\tau_\ell\) définit par extension des scalaires par \(\iota\) une représentation \(\tau\) de \(\mathrm G\) dont la restriction à \(\mathrm G_{\mathbf R}\) est unitaire. D’après (2.2.9), \(\mathrm L(\tau.\omega_s)\) est donc holomorphe inversible pour \(\mathscr R s\geq1\), sauf qu’elle a un pôle simple pour \(\tau\omega_s=\omega_1\). Il reste à appliquer (2.2.8.1), (2.2.9.1).

\Page{197}

\section*{III. — LE THÉORÈME FONDAMENTAL}

\section*{(3.1) — Un calcul de cycles évanescents}

Les résultats de ce paragraphe serviront à calculer, modulo entiers, les poids de certains groupes de cycles évanescents.

\textbf{(3.1.1)} Soient \(\mathrm S\) une surface projective et lisse sur un corps algébriquement clos \(k\), \(\mathrm D\) un diviseur à croisements normaux sur \(\mathrm S\), \(\mathrm V=\mathrm S-\mathrm D\), \(j\) l’inclusion de \(\mathrm V\) dans \(\mathrm S\) et \(\mathscr F\) un faisceau sur \(\mathrm V\), modérément ramifié le long de \(\mathrm D\). Nous nous proposons d’étudier par la méthode des pinceaux de Lefschetz les groupes de cohomologie \(\mathrm H_c^*(\mathrm V,\mathscr F)\).

Comme dans I.5, on plonge \(\mathrm S\) dans un espace projectif \(\mathbf P\), et on la balaie par un pinceau d’hyperplans \((\mathrm H_t)_{t\in\mathrm A^*}\). Notations : \(\mathrm A^*\) est une droite dans l’espace projectif dual \(\check{\mathbf P}\), elle paramétrise les hyperplans contenant un sous-espace de codimension \(2\) de \(\mathbf P\), l’\textit{axe} \(\mathrm A\) du pinceau; pour \(t\in\mathrm A^*\), \(\mathrm S_t\) est la section hyperplane \(\mathrm S\cap\mathrm H_t\) de \(\mathrm S\); \(\widetilde{\mathrm S}\subset\mathrm S\times\mathrm A^*\) est l’espace des couples \((x,t)\) tels que \(x\in\mathrm H_t\); \(\widetilde{\mathrm V}\) est l’image inverse de \(\mathrm V\) dans \(\widetilde{\mathrm S}\), et les applications premières et secondes coordonnées définissent un diagramme :
\[
\begin{tikzcd}[ampersand replacement=\&, column sep=small, row sep=large]
\& \mathrm V \arrow[dl,hook,"j"'] \& \&
\widetilde{\mathrm V}
  \arrow[ll,"\pi"']
  \arrow[dl,hook,"j"']
  \arrow[ddl,"f"] \\
\mathrm S \& \&
\widetilde{\mathrm S}\arrow[ll,"\pi"']\arrow[d,"f"'] \& \\
\& \& \mathrm A^* \&
\end{tikzcd}
\tag*{\textbf{(3.1.1.1)}}
\]

Que des lettres y désignent deux applications ne devrait pas créer de confusion. Les fibres de \(f:\widetilde{\mathrm S}\to\mathrm A^*\) sont les \(\mathrm S_t\).

On fait les hypothèses de position générale (A) à (D) suivantes.

(A) L’axe \(\mathrm A\) est transverse à \(\mathrm S\) et disjoint de \(\mathrm D\); l’espace \(\widetilde{\mathrm S}\) est donc lisse, déduit de \(\mathrm S\) par éclatement des points de l’ensemble fini \(\mathrm S\cap\mathrm A\). Aucun de ces points n’est sur \(\mathrm D\), ce qui nous permet d’identifier \(\mathrm D\) à un diviseur de \(\widetilde{\mathrm S}\).

(B) Les seules singularités de \(f\) sont des points quadratiques ordinaires; aucun de ces points critiques n’est sur \(\mathrm D\).

(C) Sur le normalisé \(\mathrm D'\) de \(\mathrm D\), les seules singularités de \(f\) sont des points quadratiques ordinaires. Aucun de ces points n’est au-dessus d’un point singulier de \(\mathrm D\).

Un point de \(\mathrm S\) sera dit \textit{exceptionnel} s’il est de l’un des trois types \textit{(a)} un point critique de \(f\) sur \(\mathrm S\), \textit{(b)} un point critique de \(f\) sur \(\mathrm D'\), \textit{(c)} un point singulier de \(\mathrm D\). Une fibre \(\mathrm S_t\)
\Page{198}
\((t\in\mathrm A^*)\) de \(f\) sera dite exceptionnelle si elle contient un point exceptionnel et on dira alors que \(t\) est exceptionnel.

(D) Chaque fibre exceptionnelle ne contient qu’un point exceptionnel.

Les fibres exceptionnelles sont donc de l’un des trois types suivants :
\[
\begin{tikzcd}[
  ampersand replacement=\&,
  column sep=1.4em,
  row sep=1em,
  baseline=(current bounding box.north),
  cells={nodes={inner sep=0pt}}
]
{
\begin{tikzpicture}[x=0.8cm,y=0.8cm,baseline=0pt]
\draw[line width=0.5pt]
 (1.05,1.05)
 .. controls (0.45,0.05) and (-0.20,-0.60) .. (-0.80,-0.60)
 .. controls (-1.50,-0.60) and (-1.50,0.65) .. (-0.85,0.65)
 .. controls (0.05,0.62) and (0.50,-0.20) .. (1.10,-1.20);
\node[anchor=west] at (1.15,-1.30) {\(\mathrm S_t\)};
\end{tikzpicture}
}
\&
{
\begin{tikzpicture}[x=0.8cm,y=0.8cm,baseline=0pt]
\draw[line width=0.5pt] (0.60,-1.50)--(0.60,1.50);
\draw[line width=0.5pt]
 (-1.00,0.90)
 .. controls (0.15,0.72) and (0.60,0.40) .. (0.60,0)
 .. controls (0.60,-0.50) and (0.15,-0.66) .. (-1.25,-0.80);
\node[anchor=north west] at (-1.20,-0.85) {\(\mathrm D\)};
\node[anchor=west] at (0.65,-1.30) {\(\mathrm S_t\)};
\end{tikzpicture}
}
\&
{
\begin{tikzpicture}[x=0.8cm,y=0.8cm,baseline=0pt]
\draw[line width=0.5pt] (0,-1.30)--(0,1.45);
\draw[line width=0.5pt] (-1.55,-1.05)--(1.55,1.05);
\draw[line width=0.5pt] (-1.55,0.85)--(1.60,-0.88);
\node[anchor=north east] at (1.60,0.90) {\(\mathrm D\)};
\node[anchor=north east] at (1.60,-0.90) {\(\mathrm D\)};
\node[anchor=west] at (0.05,-1.25) {\(\mathrm S_t\)};
\end{tikzpicture}
}
\\
{
\begin{array}{@{}ll@{}}
\text{\small\textit{a)}} & \text{\small courbe ayant un point double}\\
 & \text{\small à tangentes distinctes;}
\end{array}
}
\&
{
\begin{array}{@{}ll@{}}
\text{\small\textit{b)}} & \text{\small tangence de \(\mathrm D\)}\\
 & \text{\small avec \(\mathrm S_t\);}
\end{array}
}
\&
{
\begin{array}{@{}ll@{}}
\text{\small\textit{c)}} & \text{\small deux branches de \(\mathrm D\)}\\
 & \text{\small se coupent sur \(\mathrm S_t\).}
\end{array}
}
\end{tikzcd}
\tag*{\textbf{(3.1.1.2)}}
\]

Les cohomologies de \(\mathrm V\) et \(\widetilde{\mathrm V}\) sont liées par un isomorphisme canonique :
\[
\mathrm H_c^*(\widetilde{\mathrm V},\pi^*\mathscr F)
\cong
\mathrm H_c^*(\mathrm V,\mathscr F)
\oplus
(\mathrm H^0(\mathrm V\cap\mathrm A,\mathscr F)(-1)
 \text{ placé en degré }2).
\tag*{\textbf{(3.1.1.3)}}
\]

Pour la suite, il nous suffirait de savoir l’injectivité :
\[
\pi^*:\mathrm H_c^*(\mathrm V,\mathscr F)
\hookrightarrow
\mathrm H_c^*(\widetilde{\mathrm V},\pi^*\mathscr F);
\tag*{\textbf{(3.1.1.4)}}
\]
le transposé par dualité de Poincaré de \(\pi^*:\mathrm H^*(\mathrm V,\mathscr F)\to\mathrm H^*(\widetilde{\mathrm V},\pi^*\mathscr F)\) est une rétraction de (3.1.1.4).

Pour étudier la cohomologie de \(\widetilde{\mathrm V}\), nous utiliserons la suite spectrale de Leray :
\[
\mathrm E_2^{pq}
=\mathrm H^p(\mathrm A^*,\mathrm R^qf_!\pi^*\mathscr F)
\Rightarrow
\mathrm H_c^{p+q}(\widetilde{\mathrm V},\pi^*\mathscr F).
\tag*{\textbf{(3.1.1.5)}}
\]

Les faisceaux \(\mathrm R^qf_!\pi^*\mathscr F=\mathrm R^qf_*(j_!\pi^*\mathscr F)\) sont lisses, sauf en les valeurs exceptionnelles de \(f\).

Soient \(t\) une valeur exceptionnelle de \(f\), \(\mathrm A^*_{(t)}\) l’hensélisé de \(\mathrm A^*\) en \(t\) (un trait strictement hensélien) et \(\bar{\eta}\) un point générique géométrique de \(\mathrm A^*_{(t)}\). Appliquons la théorie des cycles évanescents à l’image inverse de \(\widetilde{\mathrm S}\), \(j_!\pi^*\mathscr F\), sur \(\mathrm A^*_{(t)}\). Les faisceaux de cycles évanescents \(\Phi^q=\Phi^q(j_!\pi^*\mathscr F)\) sont concentrés au point exceptionnel \(x\) de \(\mathrm S_t\) et l’on trouve une suite exacte longue :
\[
\ldots\xrightarrow{\partial}
(\mathrm R^qf_!\pi^*\mathscr F)_t
\to
(\mathrm R^qf_!\pi^*\mathscr F)_{\bar{\eta}}
\to\Phi_x^q\xrightarrow{\partial}\ldots
\tag*{\textbf{(3.1.1.6)}}
\]

\textbf{(3.1.2)} Nous nous proposons de calculer les \(\Phi_x^q\), ou plutôt leur gradué pour une filtration convenable, sous l’hypothèse additionnelle suivante.

(E) La monodromie locale de \(\mathscr F\) autour de \(\mathrm D\) est unipotente.

Soient \(d\in\mathrm D\), \(\mathrm S_{(d)}\) le localisé strict de \(\mathrm S\) en \(d\) et \(\mathrm V_{(d)}\) l’image inverse de \(\mathrm V\) dans \(\mathrm S_{(d)}\). L’hypothèse (E) assure que l’image inverse de \(\mathscr F\) sur \(\mathrm V_{(d)}\) admet une filtration finie \(\mathrm F\) (par des sous-faisceaux lisses) telle que les faisceaux \(\operatorname{Gr}_{\mathrm F}^i(\mathscr F)\) soient constants sur \(\mathrm V_{(d)}\). On note \(\operatorname{Gr}_{\mathrm F}^i(\mathscr F)_d\) la fibre en \(d\) du prolongement constant de \(\operatorname{Gr}_{\mathrm F}^i(\mathscr F)\) sur \(\mathrm S_{(d)}\) : c’est \(\mathrm H^0(\mathrm V_{(d)},\operatorname{Gr}_{\mathrm F}^i(\mathscr F))\).

\Page{199}

Nous utiliserons la notation suivante : si \(\mathrm B\) est un ensemble à deux éléments, \(\varepsilon(\mathrm B)\) est un groupe muni de deux isomorphismes opposés avec \(\mathbf Z\), indexés par \(\mathrm B\). Par exemple, \(\bigwedge\limits^{2}\mathbf Z^{\mathrm B}\) ou \(\mathbf Z^{\mathrm B}/(\mathbf Z\text{ diagonal})\).

Distinguons trois cas, selon la nature du point exceptionnel \(x\) de \(\mathrm S_t\) (3.1.1.2).

\textbf{(3.1.3)} \textit{Cas} (a). — Puisque \(j_!\pi^*\mathscr F\) est lisse en \(x\), on a \(\Phi_x^*(j_!\pi^*\mathscr F)=\Phi_x^*(\overline{\mathbf Q}_\ell)\otimes\mathscr F_x\), et \(\Phi_x^*(\overline{\mathbf Q}_\ell)\) est donné par la théorie de Picard-Lefschetz (en dimension relative \(1\)) : \(\Phi_x^q(\overline{\mathbf Q}_\ell)=0\) si \(q\ne1\), et si \(\mathrm B\) est l’ensemble à deux éléments des branches de \(\widetilde{\mathrm S}_t\) en \(x\), on a \(\Phi_x^1(\overline{\mathbf Q}_\ell)=\overline{\mathbf Q}_\ell(-1)\otimes\varepsilon(\mathrm B)\). Au total :
\[
\Phi_x^q(j_!\pi^*\mathscr F)=0
\qquad\text{pour}\quad q\ne1
\quad\text{et}
\tag*{\textbf{(3.1.3.1)}}
\]
\[
\Phi_x^1(j_!\pi^*\mathscr F)
\cong\mathscr F_x(-1)\otimes\varepsilon(\mathrm B).
\tag*{\textbf{(3.1.3.2)}}
\]

On peut donner de l’isomorphisme (3.1.3.2) la description suivante (que nous n’utiliserons pas). Si \(\mathrm W\) est la variété des cycles évanescents en \(x\), il provient d’isomorphismes :
\[
\mathbf Z/\ell^n
\xrightarrow{\sim}
\mathrm H^1(\mathrm W,\mu_{\ell^n})
=\Phi_x^1(\mathbf Z/\ell^n(1));
\tag*{\textbf{(3.1.3.3)}}
\]
si \(u\) et \(v\) sont des équations locales pour les deux branches de \(\mathrm S_t\) en \(x\), le générateur canonique de \(\mathrm H^1(\mathrm W,\mu_{\ell^n})\) est défini par le revêtement kummérien \(\mathrm T^{\ell^n}=uv^{-1}\).

\textbf{(3.1.4)} \textit{Cas} (b). — Supposons d’abord que \(\mathscr F\) soit le faisceau constant \(\overline{\mathbf Q}_\ell\). Soit la suite exacte courte :
\[
0\to j_!\overline{\mathbf Q}_\ell
\to\overline{\mathbf Q}_\ell
\to\overline{\mathbf Q}_{\ell\mathrm D}\to0.
\tag*{\textbf{(3.1.4.1)}}
\]

Les groupes de cycles évanescents \(\Phi_x^*\) sont nuls pour le faisceau constant \(\overline{\mathbf Q}_\ell\), puisque \(f\) est lisse en \(x\). Pour \(\overline{\mathbf Q}_{\ell\mathrm D}\), ils coïncident avec le groupe analogue calculé sur \(\mathrm D\), et la suite exacte longue déduite de (3.1.4.1) par application du foncteur cohomologique \(\Phi_x^*\) fournit des isomorphismes :
\[
\Phi_x^q(j_!\overline{\mathbf Q}_\ell)
\cong\Phi_x^{q-1}(\mathrm D,\overline{\mathbf Q}_\ell).
\tag*{\textbf{(3.1.4.2)}}
\]

Si \(\mathrm B\) est l’ensemble des deux points de l’hensélisé \(\mathrm D_{(x)}\) au-dessus de \(\bar{\eta}\), on a donc :
\[
\Phi_x^q(j_!\overline{\mathbf Q}_\ell)=0
\qquad\text{pour}\quad q\ne1
\quad\text{et}
\tag*{\textbf{(3.1.4.3)}}
\]
\[
\Phi_x^1(j_!\overline{\mathbf Q}_\ell)
=\overline{\mathbf Q}_\ell\otimes\varepsilon(\mathrm B).
\tag*{\textbf{(3.1.4.4)}}
\]

Dans le cas général, soit \(\mathrm F\) une filtration comme en (3.1.2). On trouve par dévissage que :
\[
\Phi_x^q(j_!\pi^*\mathscr F)=0
\qquad\text{pour}\quad q\ne1
\tag*{\textbf{(3.1.4.5)}}
\]
et que \(\Phi_x^1(j_!\pi^*\mathscr F)\) admet une filtration \(\mathrm F\), pour laquelle :
\[
\operatorname{Gr}_{\mathrm F}^i\Phi_x^1(j_!\pi^*\mathscr F)
=\operatorname{Gr}_{\mathrm F}^i(\mathscr F)_x\otimes\varepsilon(\mathrm B).
\tag*{\textbf{(3.1.4.6)}}
\]

\Page{200}

\textbf{(3.1.5)} \textit{Cas} (c). — Soit \(\mathrm B\) l’ensemble à deux éléments des branches de \(\mathrm D\) par \(x\). Notant \(i\) la projection dans \(\mathrm S\) ou \(\widetilde{\mathrm S}\) de la normalisée \(\mathrm D'\) de \(\mathrm D\), on a au voisinage de \(x\) une suite exacte :
\[
0\to j_!\overline{\mathbf Q}_\ell
\to\overline{\mathbf Q}_\ell
\to i_*\overline{\mathbf Q}_\ell
\to\overline{\mathbf Q}_{\ell x}\otimes\varepsilon(\mathrm B)
\to0.
\tag*{\textbf{(3.1.5.1)}}
\]

Puisque \(f\) est lisse en \(x\), et que \(f\circ i\) est lisse en les deux points de \(\mathrm D'\) au-dessus de \(x\), les faisceaux \(\Phi_x^*\) sont nuls pour \(\overline{\mathbf Q}_\ell\) et \(i_*\overline{\mathbf Q}_\ell\). On obtient donc des isomorphismes :
\[
\Phi_x^q(j_!\overline{\mathbf Q}_\ell)
=\Phi_x^{q-2}(x,\overline{\mathbf Q}_{\ell x}\otimes\varepsilon(\mathrm B)).
\tag*{\textbf{(3.1.5.2)}}
\]

Les \(\Phi_x^q(x,\overline{\mathbf Q}_{\ell x}\otimes\varepsilon(\mathrm B))\) sont nuls pour \(q\ne-1\), et \(\Phi_x^{-1}(x,\overline{\mathbf Q}_{\ell x}\otimes\varepsilon(\mathrm B))=\overline{\mathbf Q}_\ell\otimes\varepsilon(\mathrm B)\). Ceci fournit la valeur des \(\Phi_x^q(j_!\overline{\mathbf Q}_\ell)\) et, par dévissage, celle de \(\Phi_x^q(j_!\pi^*\mathscr F)\) :
\[
\Phi_x^q(j_!\pi^*\mathscr F)=0
\qquad\text{pour}\quad q\ne1,
\tag*{\textbf{(3.1.5.3)}}
\]
et, pour \(\mathrm F\) comme en (3.1.2) :
\[
\operatorname{Gr}_{\mathrm F}^i\Phi_x^1(j_!\pi^*\mathscr F)
=\operatorname{Gr}_{\mathrm F}^i(\mathscr F)_x\otimes\varepsilon(\mathrm B).
\tag*{\textbf{(3.1.5.4)}}
\]

\section*{(3.2) — Dimension 1}

\textit{Les conventions \textup{(0.7)} sont en vigueur.}

\textit{Proposition} \textbf{(3.2.1)}. — \textit{Soient \(\mathrm X_0\) une courbe lisse absolument irréductible sur \(\mathbf F_q\), et \(\mathscr F_0\) un faisceau lisse ponctuellement \(\iota\)-pur et \(\iota\)-réel sur \(\mathrm X_0\). Alors, les polynômes \(\iota\det(1-\mathrm F t,\mathrm H_c^i(\mathrm X,\mathscr F))\) sont à coefficients réels.}

Dans la formule de Grothendieck, (1.4.7.1), les facteurs au membre de gauche sont par hypothèse des polynômes réels et la fonction rationnelle au membre de droite est réelle. Soit \(\beta\) le poids de \(\mathscr F_0\). D’après (1.4.4) et (2.2.10), les zéros de \(\iota\det(1-\mathrm F t,\mathrm H_c^2(\mathrm X,\mathscr F))\) (pris avec leur multiplicité) sont ceux des pôles de cette fonction rationnelle qui sont de valeur absolue \(q^{-(\beta+2)/2}\). Leur ensemble est stable par conjugaison complexe, et le polynôme \(\iota\det(1-\mathrm F t,\mathrm H_c^2(\mathrm X,\mathscr F))\) est réel. Si \(\mathrm X_0\) est affine, le \(\mathrm H_c^0\) est nul. Si \(\mathrm X_0\) est projective, c’est le dual de \(\mathrm H_c^2(\mathrm X,\mathscr F^\vee(-1))\) (dualité de Poincaré), et le polynôme \(\iota\det(1-\mathrm F t,\mathrm H^2(\mathrm X,\mathscr F^\vee(-1)))\) étant réel, on trouve dans tous les cas que le polynôme \(\iota\det(1-\mathrm F t,\mathrm H_c^0(\mathrm X,\mathscr F))\) est réel. Le dénominateur au membre de droite de (1.4.7.1) est donc réel, et le numérateur \(\iota\det(1-\mathrm F t,\mathrm H_c^1(\mathrm X,\mathscr F))\) doit l’être également.

\textit{Remarque} \textbf{(3.2.2)}. — Les arguments de (0.5) montrent que (3.2.1) vaut sans hypothèse d’irréductibilité absolue.

\textit{Théorème} \textbf{(3.2.3)}. — \textit{Soient \(\mathrm X_0\) une courbe projective et lisse sur \(\mathbf F_q\), \(j:\mathrm U_0\hookrightarrow\mathrm X_0\) un ouvert dense de \(\mathrm X_0\), et \(\mathscr F_0\) un faisceau lisse ponctuellement \(\iota\)-pur de poids \(\beta\) sur \(\mathrm U_0\). Alors les valeurs propres \(\alpha\) de \(\mathrm F\) sur \(\mathrm H^i(\mathrm X,j_*\mathscr F)\) satisfont à \(w_q(\alpha)=\beta+i\).}

% ---- 201-208 ----
\Page{201}

Pour une description des grandes lignes de la démonstration, je renvoie à l’introduction. Pour \(i=1\), le groupe \(\mathrm{H}^1(\mathrm{X},j_*\mathscr{F})\) est aussi l’image de \(\mathrm{H}_c^1(\mathrm{U},\mathscr{F})\) dans \(\mathrm{H}^1(\mathrm{U},\mathscr{F})\) (cf. les groupes de cohomologie « parabolique » rencontrés dans l’étude des formes modulaires).

Nous commencerons par prouver par récurrence sur \(k\) la proposition suivante qui, pour \(k=0\), résulte de (1.8.2).

\textit{Lemme} \textbf{(3.2.4)} \((k)\). — \textit{Soient \(\mathrm{U}_0\) une courbe lisse sur \(\mathbf{F}_q\) et \(\mathscr{F}_0\) un faisceau lisse ponctuellement \(\iota\)-pur de poids \(0\) sur \(\mathrm{U}_0\). Alors, les valeurs propres \(\alpha\) de \(\mathrm{F}\) sur \(\mathrm{H}_c^1(\mathrm{U},\mathscr{F})\) satisfont à \(w_q(\alpha)\leq 1+2^{-k}\).}

\textbf{(3.2.5)} \((k)\) Si \(\mathrm{U}_0\) est un ouvert dense d’une courbe \(\mathrm{X}_0\) : \(j:\mathrm{U}_0\hookrightarrow\mathrm{X}_0\) et que \(\mathscr{G}_0\) est un faisceau sur \(\mathrm{X}_0\), la suite exacte longue de cohomologie définie par la suite exacte courte \(0\to j_!j^*\mathscr{G}_0\to\mathscr{G}_0\to\mathscr{G}_0/j_!j^*\mathscr{G}_0\to0\) fournit une surjection :
\[
\mathrm{H}_c^1(\mathrm{U},j^*\mathscr{G})
=\mathrm{H}_c^1(\mathrm{X},j_!j^*\mathscr{G})
\twoheadrightarrow\mathrm{H}_c^1(\mathrm{X},\mathscr{G}).
\]
Si \(\mathrm{U}_0\) est lisse, et que \(j^*\mathscr{G}_0\) est lisse et ponctuellement \(\iota\)-pur de poids \(\beta\), il résulte donc de (3.2.4) \((k)\) et d’un argument de torsion (1.2.7) que les valeurs propres de \(\mathrm{F}\) sur \(\mathrm{H}_c^1(\mathrm{X},\mathscr{G})\) satisfont à \(w_q(\alpha)\leq\beta+1+2^{-k}\).

\textbf{(3.2.6)} Soit \(k\geq0\). Admettons (3.2.4) \((k)\) (et son corollaire (3.2.5) \((k)\)) et prouvons (3.2.4) \((k+1)\). Nous le ferons d’abord sous les hypothèses (A) à (C) ci-dessous. Soit \(\mathrm{X}_0\) la courbe projective et lisse dont \(\mathrm{U}_0\) est un ouvert dense.

(A) Le faisceau \(\mathscr{F}\) sur \(\mathrm{U}\) est modérément ramifié, à monodromie locale unipotente, en les points de \(\mathrm{X}-\mathrm{U}\).

Soient \(\mathrm{S}_0=\mathrm{X}_0\times\mathrm{X}_0\), \(\mathrm{V}_0=\mathrm{U}_0\times\mathrm{U}_0\), \(\mathrm{D}_0\) le diviseur à croisements normaux de \(\mathrm{S}_0\) dont \(\mathrm{V}_0\) est le complément, plongeons \(\mathrm{S}_0\) dans un espace projectif \(\mathbf{P}_0\) et balayons \(\mathrm{S}_0\) par un pinceau d’hyperplans d’axe \(\mathrm{A}_0\subset\mathbf{P}_0\). On suppose que le plongement et le pinceau peuvent être pris de sorte que

(B) \(\mathrm{S}\), \(\mathrm{D}\) et le pinceau de sections hyperplanes d’axe \(\mathrm{A}\) remplissent les conditions de position générale (A) à (D) de (3.1.1).

On fera enfin l’hypothèse de commodité :

(C) Les points exceptionnels (3.1.1) sont définis sur \(\mathbf{F}_q\).

Les flèches du diagramme (3.1.1.1) proviennent par extension des scalaires de morphismes de \(\mathbf{F}_q\)-schémas, que nous désignerons par les mêmes lettres \(f\), \(\pi\), \(j\). Nous noterons \(\mathrm{T}_0\) l’ensemble fini des valeurs exceptionnelles de \(f\), \(\mathrm{W}_0\) le complément de \(\mathrm{T}_0\) dans \(\mathrm{A}_0^*\) et \(w\) l’inclusion de \(\mathrm{W}_0\) dans \(\mathrm{A}_0^*\).

Soit \(\mathscr{G}_0=\mathscr{F}_0\boxtimes\mathscr{F}_0=\operatorname{pr}_1^*\mathscr{F}_0\otimes\operatorname{pr}_2^*\mathscr{F}_0\) le produit tensoriel externe de \(\mathscr{F}_0\) avec lui-même. Nous appliquerons la théorie (3.1) au calcul de la cohomologie de \(\mathscr{G}\). Les hypothèses (A) à (E) de (3.1) sont satisfaites par hypothèse par \((\mathrm{S},\mathrm{D},\mathscr{G})\) et le pinceau d’axe \(\mathrm{A}\).

\Page{202}

\textit{Lemme} \textbf{(3.2.7)}. — \textit{Le faisceau lisse \(w^*\mathrm{R}^1f_!\pi^*\mathscr{G}_0\) sur \(\mathrm{W}_0\) est facteur direct d’un faisceau \(\iota\)-réel.}

\textit{Preuve.} — Puisque le faisceau \(\mathscr{G}\) est lisse et ponctuellement \(\iota\)-pur de poids \(0\), le faisceau \(\mathscr{H}_0=\mathscr{G}_0\oplus\mathscr{G}_0^\vee\) est ponctuellement \(\iota\)-pur et \(\iota\)-réel, et il nous suffit de vérifier que sur \(\mathrm{W}_0\) le faisceau \(\mathrm{R}^1f_!\pi^*\mathscr{H}_0\) est \(\iota\)-réel. C’est une conséquence de (3.2.1) : si \(x\in|\mathrm{W}_0|\), si \(\mathrm{Y}_1\) est la fibre de \(f:\widetilde{\mathrm{V}}_0\to\mathrm{A}_0^*\) en \(x\) (une courbe sur \(k(x)\)) et si \(\mathrm{Y}\) se déduit de \(\mathrm{Y}_1\) par l’extension des scalaires de \(k(x)\) à \(\mathbf{F}\) définie par un point géométrique \(\bar{x}\) localisé en \(x\), on a :
\[
(\mathrm{R}^1f_!\pi^*\mathscr{H}_0)_{\bar{x}}
=\mathrm{H}_c^1(\mathrm{Y},\pi^*\mathscr{H})
\tag*{\textbf{(3.2.7.1)}}
\]
et l’on applique (3.2.2) pour \(i=1\) et pour la restriction de \(\pi^*\mathscr{H}\) à \(\mathrm{Y}_1\).

Nous prouverons ensemble les lemmes (3.2.8) et (3.2.9) ci-dessous.

\textit{Lemme} \textbf{(3.2.8)}. — \textit{Le faisceau \(\mathrm{R}^1f_!\pi^*\mathscr{G}\) n’a pas de section à support dans \(\mathrm{T}\).}

Ce lemme permet d’identifier \(\mathrm{R}^1f_!\pi^*\mathscr{G}_0\) à un sous-faisceau de \(w_*w^*\mathrm{R}^1f_!\pi^*\mathscr{G}_0\). D’après (3.2.7) et (1.5.1), le faisceau \(w^*\mathrm{R}^1f_!\pi^*\mathscr{G}_0\) admet une filtration finie \(\mathrm{G}^\bullet\) à quotients successifs lisses et ponctuellement \(\iota\)-purs. On la prolonge en une filtration finie de \(\mathrm{R}^1f_!\pi^*\mathscr{G}_0\) par la formule :
\[
\mathrm{G}^i\mathrm{R}^1f_!\pi^*\mathscr{G}_0
=w_*\mathrm{G}^iw^*\mathrm{R}^1f_!\pi^*\mathscr{G}_0
\cap\mathrm{R}^1f_!\pi^*\mathscr{G}_0
\]
(intersection dans \(w_*w^*\mathrm{R}^1f_!\pi^*\mathscr{G}_0\)).

\textit{Lemme} \textbf{(3.2.9)}. — \textit{Si \(\operatorname{Gr}_{\mathrm{G}}^i\mathrm{R}^1f_!\pi^*\mathscr{G}_0\) ne devient pas constant sur \(\mathrm{A}^*\), sa restriction à \(\mathrm{W}_0\) est ponctuellement \(\iota\)-pure de poids entier.}

Soient \(\bar{x}\in\mathrm{X}-\mathrm{U}\), \(x\) son image dans \(\mathrm{X}_0-\mathrm{U}_0\), \(\eta\) le point générique de \(\mathrm{X}_{0(x)}\) et \(\bar{\eta}\) un point géométrique localisé au point générique de \(\mathrm{X}_{(\bar{x})}\). D’après (1.8.4), la représentation \(\mathscr{F}_{0\bar{\eta}}\) du groupe de Weil \(\mathrm{W}(\bar{\eta}/\eta)\) admet une filtration de quotients successifs \(\iota\)-purs de poids entiers. L’hypothèse (A) donne de plus que sur ces quotients l’inertie n’agit pas. Pour \(x\in\mathrm{U}_0\), la même conclusion vaut par hypothèse pour la filtration triviale, à un seul cran.

Soient \(x\in\mathrm{S}_0-\mathrm{V}_0\), \(\bar{x}\) un point géométrique localisé en \(x\), \(\mathrm{S}_{0(x)}\) l’hensélisé de \(\mathrm{S}_0\) en \(x\) et \(\mathrm{V}_{0(x)}\) l’image inverse de \(\mathrm{V}_0\) dans \(\mathrm{S}_{0(x)}\). Par produit tensoriel, on trouve que \(\mathscr{G}_0\) admet sur \(\mathrm{V}_{0(x)}\) une filtration finie \(\mathrm{F}\), telle que les \(\operatorname{Gr}_{\mathrm{F}}^i\mathscr{G}_0\) soient constants sur \(\mathrm{V}_{(x)}\) et que les \(\operatorname{Gr}_{\mathrm{F}}^i(\mathscr{G}_0)_{\bar{x}}\) (notation (3.1.2)) soient \(\iota\)-purs de poids entiers. Pour \(x\in\mathrm{V}_0\), \(\mathscr{G}_{\bar{x}}\) lui-même est \(\iota\)-pur de poids entier (savoir, \(0\)).

Soient \(\bar{t}\in\mathrm{T}\), \(x\in\widetilde{\mathrm{S}}\) le point exceptionnel au-dessus (3.1.1), \(t\) et \(x\) leurs images dans \(\mathrm{T}_0\) et \(\mathrm{S}_0\), \(\bar{\eta}\) un point géométrique localisé au point générique de \(\mathrm{A}_{(t)}^*\), et \(\eta\) le point générique de \(\mathrm{A}_{0(t)}^*\). Les faisceaux de cycles évanescents, calculés au n\textsuperscript{o} (3.1), étant concentrés en \(x\) et nuls pour \(i\neq1\), on trouve une suite exacte à cinq termes :
\[
0\to(\mathrm{R}^1f_!\pi^*\mathscr{G}_0)_{\bar{t}}
\to(\mathrm{R}^1f_!\pi^*\mathscr{G}_0)_{\bar{\eta}}
\to\Phi_{\bar{x}}^1
\to(\mathrm{R}^2f_!\pi^*\mathscr{G}_0)_{\bar{t}}
\to(\mathrm{R}^2f_!\pi^*\mathscr{G}_0)_{\bar{\eta}}\to0.
\tag*{\textbf{(3.2.9.1)}}
\]

\Page{203}

L’injectivité à gauche assure que \(\mathrm{R}^1f_!\bar{u}_*\mathscr{G}\) n’a pas de section à support dans \(t\). Ceci prouve (3.2.8). On trouve aussi des suites exactes courtes :
\[
0\to\operatorname{Gr}_{\mathrm{G}}^i
(\mathrm{R}^1f_!\pi^*\mathscr{G}_0)_{\bar{t}}
\to\operatorname{Gr}_{\mathrm{G}}^i
(\mathrm{R}^1f_!\pi^*\mathscr{G}_0)_{\bar{\eta}}
\to\mathrm{A}_t^i\to0,
\]
où \(\mathrm{A}_t^i\) est un sous-quotient de \(\Phi_{\bar{x}}^1\). Le calcul (3.1) assure que \(\Phi_{\bar{x}}^1\), donc \(\mathrm{A}_t^i\), admet une filtration de quotients successifs \(\iota\)-purs de poids entier. D’après (1.8.4), si \(\mathrm{A}_t^i\neq0\), le faisceau \(\iota\)-pur \(w^*\operatorname{Gr}_{\mathrm{G}}^i(\mathrm{R}^1f_!\pi^*\mathscr{G}_0)\) est donc de poids entier.

Pour conclure, on note que si, pour \(i\) fixé, tous les \(\mathrm{A}_t^i\) sont nuls, le faisceau \(\operatorname{Gr}_{\mathrm{G}}^i\mathrm{R}^1f_!\pi^*\mathscr{G}_0\) est lisse, donc constant sur \(\mathrm{A}^*\) puisque la droite projective \(\mathrm{A}^*\) est simplement connexe.

\textit{Lemme} \textbf{(3.2.10)}. — \textit{Les constituants de \(w^*\mathrm{R}^1f_!\pi^*\mathscr{G}_0\) sont ponctuellement \(\iota\)-purs de poids \(<2\).}

On sait que les constituants sont ponctuellement \(\iota\)-purs. Il suffit de tester les poids en un point. Le lemme résulte de (2.2.10) (cf. (3.2.7.1)).

\textit{Lemme} \textbf{(3.2.11)}. — \textit{Les valeurs propres \(\alpha\) de \(\mathrm{F}\) sur \(\mathrm{E}_2^{11}=\mathrm{H}^1(\mathrm{A}^*,\mathrm{R}^1f_!\pi^*\mathscr{G})\) satisfont à \(w_q(\alpha)\leq2+2^{-k}\).}

Il suffit de le vérifier pour les valeurs propres de \(\mathrm{F}\) sur chaque \(\mathrm{H}^1(\mathrm{A}^*,\operatorname{Gr}_{\mathrm{G}}^i\mathrm{R}^1f_!\pi^*\mathscr{G})\). Si \(\operatorname{Gr}_{\mathrm{G}}^i\mathrm{R}^1f_!\pi^*\mathscr{G}\) est constant, ce \(\mathrm{H}^1\) est nul, puisque \(\mathrm{A}^*\) est une droite projective. Sinon, \(w^*\operatorname{Gr}_{\mathrm{G}}^i\mathrm{R}^1f_!\pi^*\mathscr{G}\) est de poids entier \(<2\) ((3.2.9) et (3.2.10)), donc de poids \(\leq1\), et l’on applique (3.2.5) \((k)\).

\textit{Lemme} \textbf{(3.2.12)}. — \textit{Les valeurs propres \(\alpha\) de \(\mathrm{F}\) sur \(\mathrm{E}_2^{02}=\mathrm{H}^0(\mathrm{A}^*,\mathrm{R}^2f_!\pi^*\mathscr{G})\) et sur \(\mathrm{E}_2^{20}=\mathrm{H}^2(\mathrm{A}^*,\mathrm{R}^0f_!\pi^*\mathscr{G})\) satisfont à \(w_q(\alpha)\leq2\).}

La formule (3.2.7.1), et une double application de (1.4.4), montrent que, pour \(x\in|\mathrm{A}_0^*|\), les valeurs propres \(\alpha\) de \(\mathrm{F}_x\) sur \(\mathrm{R}^if_!\pi^*\mathscr{G}_0\) \((i=0,2)\) satisfont à \(w_{\mathrm{N}(x)}(\alpha)\leq i\), puis l’énoncé.

\textbf{(3.2.13)} \textit{Preuve de} (3.2.4) \((k+1)\) \textit{sous les hypothèses} (3.2.6).

Il résulte de (3.2.11), (3.2.12) et de la suite spectrale (3.1.1.5) que les valeurs propres \(\alpha\) de \(\mathrm{F}\) sur \(\mathrm{H}_c^2(\widetilde{\mathrm{V}},\pi^*\mathscr{G})\) satisfont à \(w_q(\alpha)\leq2+2^{-k}\).

La même estimation vaut \textit{a fortiori} pour les valeurs propres de \(\mathrm{F}\) sur :
\[
\mathrm{H}_c^1(\mathrm{U},\mathscr{F})\otimes
\mathrm{H}_c^1(\mathrm{U},\mathscr{F})
\subset\mathrm{H}_c^2(\mathrm{V},\mathscr{G})
\subset\mathrm{H}_c^2(\widetilde{\mathrm{V}},\pi^*\mathscr{G})
\]
(Künneth et (3.1.1.4)). En particulier, si \(\alpha\) est une valeur propre de \(\mathrm{F}\) sur \(\mathrm{H}_c^1(\mathrm{U},\mathscr{F})\), on a \(w_q(\alpha^2)\leq2+2^{-k}\), donc \(w_q(\alpha)\leq1+2^{-(k+1)}\).

\textbf{(3.2.14)} \textit{Fin de la preuve de} (3.2.4) \((k+1)\). — La validité de l’assertion (3.2.4) \((k+1)\) est invariante par une extension préliminaire du corps fini de base \(\mathbf{F}_q\) (cf. preuve de (2.2.10)). Puisque, sur \(\mathbf{F}\), il existe des plongements et des pinceaux satis-\Page{204}faisant à (B), il en existe aussi sur une extension finie de \(\mathbf{F}_q\). Sur une extension finie plus grande, (C) sera satisfait. Ceci prouve (3.2.4) \((k+1)\) sous la seule hypothèse (A).

Le théorème de monodromie assure qu’il existe un revêtement fini surjectif \(g:\mathrm{U}'_0\to\mathrm{U}_0\), avec \(\mathrm{U}'_0\) lisse, tel que \(g^*\mathscr{F}_0\) sur \(\mathrm{U}'_0\) vérifie (A); on conclut en notant que le morphisme \(g^*\) injecte \(\mathrm{H}_c^*(\mathrm{U},\mathscr{F})\) dans \(\mathrm{H}_c^*(\mathrm{U}',g^*\mathscr{F})\).

\textbf{(3.2.15)} \textit{Preuve de} (3.2.3). — L’assertion est triviale pour \(i=0\) ou \(2\) (1.4). Pour \(i=1\), la validité de (3.2.5) \((k)\) pour tout \(k\) assure que \(w_q(\alpha)\leq\beta+1\). Appliquons ce résultat au faisceau \(\mathscr{F}_0^\vee(1)\). Puisque \(\mathrm{H}^1(\mathrm{X},j_*\mathscr{F})\) et \(\mathrm{H}^1(\mathrm{X},j_*\mathscr{F}^\vee(1))\) sont en dualité (SGA \(4\frac{1}{2}\) [dualité] (1.3) et (2.1)), on trouve que \(w_q(\alpha^{-1})\leq(-\beta-2)+1\), i.e. que \(w_q(\alpha)\geq\beta+1\). Le théorème en résulte.

\section*{(3.3) — Le cas général}

\textit{Théorème} \textbf{(3.3.1)}. — \textit{Soient \(f:\mathrm{X}\to\mathrm{Y}\) un morphisme de schémas de type fini sur \(\mathbf{Z}\), et \(\mathscr{F}\) un faisceau sur \(\mathrm{X}\). Si \(\mathscr{F}\) est mixte de poids \(\leq n\), alors, pour chaque \(i\), \(\mathrm{R}^if_!\mathscr{F}\) est mixte de poids \(\leq n+i\).}

La preuve utilise les dévissages suivants :

a) \textit{Dévissages en \(\mathscr{F}\)} : Soit une suite exacte \(0\to\mathscr{F}'\to\mathscr{F}\to\mathscr{F}''\to0\). Si les \(\mathrm{R}^if_!\mathscr{F}'\) et \(\mathrm{R}^if_!\mathscr{F}''\) sont mixtes de poids \(\leq n+i\), alors les \(\mathrm{R}^if_!\mathscr{F}\) le sont également. Si les \(\mathrm{R}^if_!\mathscr{F}\) et \(\mathrm{R}^if_!\mathscr{F}''\) sont mixtes de poids \(\leq n+i\), alors les \(\mathrm{R}^if_!\mathscr{F}'\) le sont également (appliquer la suite exacte longue de cohomologie). Si la suite est scindée et que \(\mathrm{R}^if_!\mathscr{F}\) est mixte de poids \(\leq n+i\), alors \(\mathrm{R}^if_!\mathscr{F}'\) l’est également.

b) \textit{Dévissage en \(\mathrm{X}\)} : Soit \(j:\mathrm{U}\hookrightarrow\mathrm{X}\) un ouvert de \(\mathrm{X}\), de complément \(i:\mathrm{S}\to\mathrm{X}\). Si les \(\mathrm{R}^i(fj)_!j^*\mathscr{F}\) et les \(\mathrm{R}^i(fi)_!i^*\mathscr{F}\) sont mixtes de poids \(\leq n+i\), les \(\mathrm{R}^if_!\mathscr{F}\) le sont également (appliquer (a) à la suite exacte courte \(0\to j_!j^*\mathscr{F}\to\mathscr{F}\to i_*i^*\mathscr{F}\to0\)).

c) \textit{Dévissage en \(\mathrm{Y}\)} : Soit \(j:\mathrm{V}\hookrightarrow\mathrm{Y}\) un ouvert de \(\mathrm{Y}\), de complément \(i:\mathrm{T}\to\mathrm{Y}\). Si la conclusion de (3.3.1) devient vraie après les changements de base par \(i\) et \(j\), elle l’est au départ (appliquer le théorème de changement de base).

d) \textit{Transitivité} : Si \(f=gh\), et que les \(\mathrm{R}^ig_!\mathrm{R}^jh_!\mathscr{F}\) sont mixtes de poids \(\leq n+i+j\), alors les \(\mathrm{R}^kf_!\mathscr{F}\) sont mixtes de poids \(\leq n+k\) (appliquer la suite spectrale de Leray \(\mathrm{R}^pg_!\mathrm{R}^qh_!\Rightarrow\mathrm{R}^{p+q}f_!\)).

e) Si \(\mathrm{Y}'\xrightarrow{g}\mathrm{Y}\) est un homéomorphisme universel, il suffit de vérifier (3.3.1) après le changement de base par \(g\) (la cohomologie étale ne « voit » pas un tel changement de base). Exemples : \(\mathrm{Y}'=\mathrm{Y}_{\mathrm{red}}\), ou \(\mathrm{Y}'=\) le normalisé de \(\mathrm{Y}\), supposé normal et intègre, dans une extension inséparable de son corps de fonctions.

f) Si \(f\) est de dimension relative \(0\), et que \(\mathscr{F}\) est ponctuellement pur, alors \(f_!\mathscr{F}=\mathrm{R}^0f_!\mathscr{F}\) est ponctuellement pur du même poids, et les \(\mathrm{R}^if_!\mathscr{F}\) sont nuls pour \(i\neq0\). Pour un tel \(f\), on conclut par (a).

\Page{205}

Les dévissages (b), (c), (d) nous ramènent à supposer que \(f\) est de dimension relative \(1\), puis (a), (b) nous ramènent à supposer que \(\mathscr{F}\) est lisse et ponctuellement pur de poids \(n\).

Notons que si \(\mathrm{C}\) est une courbe sur un corps \(\mathrm{K}\), et que \(\mathscr{F}\) est un faisceau lisse sur \(\mathrm{C}\), les assertions suivantes sont vraies pour \(\mathrm{K}\) parfait, donc dans le cas général après une extension finie purement inséparable de \(\mathrm{K}\) : \((\alpha)\) Il existe dans \(\mathrm{C}_{\mathrm{red}}\) une partie finie dont le complément \(\mathrm{C}'\) est lisse sur \(\mathrm{K}\). \((\beta)\) Il existe un revêtement fini surjectif \(u:\mathrm{D}\to\mathrm{C}'\), où \(\mathrm{D}\) est le complément dans une courbe projective et lisse \(\overline{\mathrm{D}}\) d’un diviseur \(\mathrm{E}\) étale sur \(\mathrm{K}\), tel que \(u^*\mathscr{F}\) soit modérément ramifié le long de \(\mathrm{E}\); \(\mathscr{F}\) est facteur direct de \(u_*u^*\mathscr{F}=\mathrm{R}u_*u^*\mathscr{F}\). Nous appliquerons ceci à la fibre de \(f\) en un point générique \(\eta\) de \(\mathrm{Y}\). Les arguments habituels de passage à la limite montrent que \(\mathrm{C}'\), \(\mathrm{D}\)\dots{} existent, avec les mêmes propriétés, au-dessus d’un voisinage de \(\eta\) dans \(\mathrm{Y}_{\mathrm{red}}\). On se ramène par (c), (e), à supposer que ce voisinage est \(\mathrm{Y}_{\mathrm{red}}\) tout entier et que \(\mathrm{Y}_{\mathrm{red}}\) est lisse, puis par (b), (f), (a), à supposer que \(\mathrm{X}\) soit le complément dans \(\overline{\mathrm{X}}\), projectif et lisse purement de dimension relative \(1\) sur \(\mathrm{Y}\), d’un diviseur \(\mathrm{D}\) fini étale sur \(\mathrm{Y}\), et que \(\mathscr{F}\) soit modérément ramifié le long de \(\mathrm{D}\).
\[
\begin{tikzcd}[ampersand replacement=\&,column sep=1.2em,row sep=1.5em]
\mathrm{X} \arrow[rr,hook,"j"] \arrow[dr,"f"']
\& \&
\overline{\mathrm{X}} \arrow[dl,"\bar f"]
\& \&
\mathrm{D} \arrow[ll,hook',"i"']
\\
\& \mathrm{Y} \& \& \&
\end{tikzcd}
\]
On a \(\mathrm{R}f_!\mathscr{F}=\mathrm{R}\bar{f}_*(j_!\mathscr{F})\). Considérons la suite exacte :
\[
0\to j_!\mathscr{F}\to j_*\mathscr{F}
\to i_!i^*j_*\mathscr{F}\to0
\]
et la filtration induite par la filtration de monodromie locale \(\mathrm{M}\) sur \(i_!i^*j_*\mathscr{F}\) ((1.8.6) à (1.8.8)). La formation de ces faisceaux et de \(\mathrm{M}\) est compatible au passage aux fibres. Le théorème (3.2.3), appliqué à tous les isomorphismes \(\iota\), assure donc que \(\mathrm{R}^i\bar{f}_*j_*\mathscr{F}\) est ponctuellement pur de poids \(n+i\), et le théorème (1.8.4) (lui aussi pour tout \(\iota\)) assure que \(\operatorname{Gr}_{\mathrm{M}}^k(i^*j_*\mathscr{F})\), nul pour \(k>0\), est pur de poids \(n+k\) sur \(\mathrm{D}\) (1.8.8). Le morphisme \(\bar{f}i\) est fini, et il ne reste qu’à appliquer (f), (a).

\textbf{(3.3.2)} Pour \textit{minorer} les poids, on peut utiliser les résultats de SGA 7 XXI 5 sur les valuations \(p\)-adiques des valeurs propres de Frobenius. Disons qu’un faisceau est \textit{entier} si les valeurs propres des Frobenius sont des entiers algébriques. Les poids d’un faisceau mixte entier sont nécessairement positifs : si un entier algébrique \(\alpha\) est pur de poids \(n\), rel. \(q\), et de degré \(d\), sa norme est un entier ordinaire, de valeur absolue le produit \(q^{nd/2}\) des valeurs absolues de ses conjugués complexes, de sorte que \(n\geq0\).

\textit{Corollaire} \textbf{(3.3.3)}. — \textit{Soient \(f:\mathrm{X}\to\mathrm{Y}\) un morphisme de schémas de type fini sur \(\mathbf{Z}\) et \(\mathscr{F}\) un faisceau entier sur \(\mathrm{X}\). Si \(\mathscr{F}\) est mixte de poids \(\leq n\), alors, pour chaque \(i\), \(\mathrm{R}^if_!\mathscr{F}\) est mixte de poids entre \(0\) et \(n+i\). Si \(i\) dépasse la dimension maximum \(d\) des fibres, les poids sont compris entre \(2(i-d)\) et \(n+i\).}

D’après SGA 7 XXI (5.2.2), appliqué aux fibres de \(f\), les faisceaux \(\mathrm{R}^if_!\mathscr{F}\) et \(\mathrm{R}^if_!\mathscr{F}(i-d)\) sont en effet entiers.

\Page{206}

Lorsque \(\mathrm{X}\) et \(\mathrm{Y}\) sont de caractéristique \(p>0\), on peut remplacer les faisceaux étales par des faisceaux de Weil. Le cas le plus important est celui où \(\mathrm{Y}=\operatorname{Spec}(\mathbf{F}_q)\). Dans ce cas, (3.3.1) et (3.3.3) signifient ceci — avec les notations (0.7) :

\textit{Corollaire} \textbf{(3.3.4)}. — \textit{Soient \(\mathrm{X}_0\) un schéma de type fini sur \(\mathbf{F}_q\) et \(\mathscr{F}_0\) un faisceau mixte de poids \(\leq n\) sur \(\mathrm{X}_0\). Alors, \(\mathrm{H}_c^i(\mathrm{X},\mathscr{F})\) est mixte de poids \(\leq n+i\) : pour chaque \(i\) et chaque valeur propre \(\alpha\) de \(\mathrm{F}\) sur \(\mathrm{H}_c^i(\mathrm{X},\mathscr{F})\), \(\alpha\) est algébrique et il existe un entier \(w\leq n+i\) tel que tous les conjugués complexes de \(\alpha\) soient de valeur absolue \(q^{w/2}\). Si \(\mathscr{F}_0\) est entier, on a \(w\geq0\). Si de plus \(i\) dépasse la dimension \(d\) de \(\mathrm{X}\), on a \(w\geq2(i-d)\).}

\textit{Corollaire} \textbf{(3.3.5)}. — \textit{Soient \(\mathrm{X}_0\) lisse de type fini sur \(\mathbf{F}_q\) et \(\mathscr{F}_0\) un faisceau lisse mixte de poids \(\geq n\) sur \(\mathrm{X}_0\). Alors, \(\mathrm{H}^i(\mathrm{X},\mathscr{F})\) est mixte de poids \(\geq n+i\).}

On se ramène à supposer \(\mathrm{X}_0\) purement d’une dimension \(\mathrm{N}\). Soit \(\mathscr{F}_0^\vee\) le faisceau dual de \(\mathscr{F}_0\). Il est mixte de poids \(\leq-n\). D’après (3.3.3), le groupe \(\mathrm{H}_c^{2\mathrm{N}-i}(\mathrm{X},\mathscr{F}^\vee)(\mathrm{N})\) est mixte de poids \(\leq-n+(2\mathrm{N}-i)-2\mathrm{N}=-n-i\). La dualité de Poincaré identifie \(\mathrm{H}^i(\mathrm{X},\mathscr{F})\) à son dual, mixte de poids \(\geq n+i\).

Combinant les inégalités opposées de (3.3.4) et (3.3.5), on obtient le

\textit{Corollaire} \textbf{(3.3.6)}. — \textit{Soient \(\mathrm{X}_0\) lisse de type fini sur \(\mathbf{F}_q\) et \(\mathscr{F}_0\) un faisceau lisse pur de poids \(n\) sur \(\mathrm{X}_0\). Alors, l’image de \(\mathrm{H}_c^i(\mathrm{X},\mathscr{F})\) dans \(\mathrm{H}^i(\mathrm{X},\mathscr{F})\) est pure de poids \(n+i\).}

Dans SGA \(4\frac{1}{2}\) [Sommes trigonométriques] j’explique comment ces résultats permettent de majorer certaines sommes trigonométriques.

\textbf{(3.3.7)} Soit \(\alpha\) un nombre algébrique, pur d’un poids \(n\), rel. \(q\), et normalisons les valuations \(p\)-adiques de \(\mathbf{Q}(\alpha)\) en exigeant que \(v(q)=1\). A \(\alpha\), et à chaque valuation \(v\), attachons le couple de nombres rationnels \((v(\alpha),v(q^n\alpha^{-1}))\), où \(q^n\alpha^{-1}=\bar{\alpha}\). Si \(\alpha\) est entier, ces nombres sont \(\geq0\). Leur somme vaut \(n\). Pour le faisceau constant \(\overline{\mathbf{Q}}_\ell\), (3.3.4) et la méthode de dualité (3.3.5) permettent de localiser comme suit les couples \((r,s)\) attachés aux valeurs propres de Frobenius sur un \(\mathrm{H}^i\).

\textit{Corollaire} \textbf{(3.3.8)}. — \textit{On suppose \(\mathrm{X}\) de dimension \(\leq d\). En cohomologie à support propre (ou en cohomologie ordinaire pour \(\mathrm{X}\) propre), \((r,s)\) est dans un triangle inférieur du diagramme ci-dessous. En cohomologie ordinaire, et pour \(\mathrm{X}\) lisse, \((r,s)\) est dans un triangle supérieur du même diagramme.}
\[
\begin{tikzcd}[
  ampersand replacement=\&,
  baseline=(current bounding box.south),
  row sep=0pt,
  column sep=0pt,
  cells={nodes={inner sep=0pt,outer sep=0pt,minimum size=0pt,text height=0pt,text depth=0pt}},
  arrows={shorten <=0pt,shorten >=0pt}
]
{} \&[1.8cm] {} \&[1.6cm] {}
\\[0.55cm]
{} \arrow[r,no head] \arrow[dr,no head]
\& {} \arrow[d,no head]
\& {}
\\[1.8cm]
{} \arrow[uu,"\textstyle s"{pos=.9,left}]
   \arrow[rr,"\textstyle r"{pos=.9,below}]
   \arrow[r,phantom,"\textstyle i"{pos=1,below}]
   \arrow[r,phantom,"\textstyle\text{si }i\leq d"{pos=1,below,yshift=-.55cm}]
   \arrow[ur,phantom,"\textstyle\mathrm{H}_c^i"{description,pos=.18}]
   \arrow[ur,phantom,"{\begin{gathered}\textstyle\mathrm{H}^i,\\[-.3ex]\textstyle\mathrm{X}\ \text{lisse}\end{gathered}}"{description,pos=.6,xshift=.25cm}]
\& {}
\& {}
\end{tikzcd}
\qquad
\begin{tikzcd}[
  ampersand replacement=\&,
  baseline=(current bounding box.south),
  row sep=0pt,
  column sep=0pt,
  cells={nodes={inner sep=0pt,outer sep=0pt,minimum size=0pt,text height=0pt,text depth=0pt}},
  arrows={shorten <=0pt,shorten >=0pt}
]
{} \&[.6cm] {} \&[1.8cm] {} \&[.6cm] {} \&[1.6cm] {}
\\[.6cm]
{} \arrow[r,no head]
\& {} \arrow[r,no head]
      \arrow[dr,no head,shorten <=-4mm,shorten >=-18mm]
\& {} \arrow[d,no head]
\& {}
\& {}
\\[1.8cm]
{}
\& {} \arrow[u,no head] \arrow[r,no head]
      \arrow[ur,phantom,"\textstyle\mathrm{H}_c^i"{description,pos=.18}]
      \arrow[ur,phantom,"{\begin{gathered}\textstyle\mathrm{H}^i,\\[-.3ex]\textstyle\mathrm{X}\ \text{lisse}\end{gathered}}"{description,pos=.65,xshift=.15cm}]
\& {} \arrow[d,no head]
\& {}
\& {}
\\[.6cm]
{} \arrow[uuu,"\textstyle s"{pos=.9,left}]
   \arrow[rrrr,"\textstyle r"{pos=.9,below}]
   \arrow[rr,phantom,"\textstyle d"{pos=1,below}]
   \arrow[rrr,phantom,"\textstyle i"{pos=1,below}]
   \arrow[rrr,phantom,"\textstyle\text{si }i\geq d"{pos=1,below,yshift=-.7cm}]
\& {}
\& {}
\& {}
\& {}
\end{tikzcd}
\]

\Page{207}

La vérification est facile. On peut montrer que, en cohomologie ordinaire, \((r,s)\) est toujours dans le carré réunion des deux triangles mis en évidence.

Prenons \(\mathrm{X}_0\) propre et lisse. On trouve que \(\mathrm{H}^i(\mathrm{X},\mathbf{Q}_\ell)\) est pur de poids \(i\). Raisonnant comme dans I (preuve de \((1.7)\Rightarrow(1.6)\)), on retrouve le théorème principal (1.6) de I, avec « projectif » remplacé par « propre » :

\textit{Corollaire} \textbf{(3.3.9)}. — \textit{Soit \(\mathrm{X}_0\) une variété propre et lisse sur \(\mathbf{F}_q\). Pour chaque \(i\), le polynôme caractéristique \(\det(t\cdot1-\mathrm{F}^*,\mathrm{H}^i(\mathrm{X},\mathbf{Q}_\ell))\) est à coefficients entiers, indépendants de \(\ell\) \((\ell\neq p)\). Les racines complexes \(\alpha\) de ce polynôme (les conjugués complexes des valeurs propres de \(\mathrm{F}^*\)) sont de valeur absolue \(|\alpha|=q^{i/2}\).}

\textbf{(3.3.10)} Le théorème (3.3.1) vaut, avec la même démonstration, avec mixte remplacé par \(\iota\)-mixte, et \(n\) remplacé par \(\beta\in\mathbf{R}\). Les seuls poids qui peuvent apparaître dans \(\mathrm{R}^if_!\mathscr{F}\) sont les nombres \(\leq\beta+i\) congrus mod \(\mathbf{Z}\) à l’un des poids qui apparaît dans \(\mathscr{F}\).

L’argument (3.3.2) pour minorer les poids ne s’applique plus dans ce contexte. Mais (3.3.4) — sauf ce qui concerne les faisceaux entiers — et (3.3.5), (3.3.6) valent tels quels.

\textbf{(3.3.11)} Dans tout ceci, l’hypothèse « lisse » n’apparaît que pour justifier la dualité de Poincaré. Ceci permet de remplacer « lisse » par « rational homology manifold » — dans le contexte \(\ell\)-adique : lisse purement de dimension \(n\) peut être remplacé par la condition :
\[
\text{si }a\text{ est la projection de }\mathrm{X}_0
\text{ sur }\mathbf{F}_q,\text{ on a }\quad
\mathrm{R}a^!\mathbf{Q}_\ell=\mathbf{Q}_\ell(n)[2n].
\tag*{(*)}
\]
Par exemple, un schéma \(\mathrm{X}_0\) qui localement pour la topologie étale est quotient d’un schéma lisse de dimension \(n\) par un groupe fini satisfait à (*).

\section*{(3.4) — Application : la structure des faisceaux mixtes}

\textit{Les notations} (0.7) \textit{sont en vigueur.}

\textit{Théorème} \textbf{(3.4.1)}. — \textit{Soit \(\mathscr{F}_0\) un faisceau \(\iota\)-mixte sur un schéma \(\mathrm{X}_0\) de type fini sur \(\mathbf{F}_q\).}

(i) \textit{\(\mathscr{F}_0\) admet une unique décomposition \(\mathscr{F}_0=\displaystyle\bigoplus_{b\in\mathbf{R}/\mathbf{Z}}\mathscr{F}_0(b)\), la \textup{décomposition selon le poids mod \(\mathbf{Z}\)}, dans laquelle les \(\iota\)-poids ponctuels de \(\mathscr{F}_0(b)\) sont tous dans \(b\). Cette décomposition, dans laquelle les \(\mathscr{F}_0(b)\) sont bien sûr presque tous nuls, est fonctorielle en \(\mathscr{F}_0\). Remarquons que chaque \(\mathscr{F}_0(b)\) se déduit par torsion d’un faisceau \(\iota\)-mixte de poids ponctuels entiers.}

(ii) \textit{Si les poids ponctuels de \(\mathscr{F}_0\) sont entiers et que \(\mathscr{F}_0\) est lisse, \(\mathscr{F}_0\) admet une unique filtration finie croissante \(\mathrm{W}\) par des sous-faisceaux lisses, la \textup{filtration par le poids ponctuel}, telle que \(\operatorname{Gr}_i^{\mathrm{W}}(\mathscr{F}_0)\) soit ponctuellement \(\iota\)-pur de poids \(i\). Cette filtration est fonctorielle en \(\mathscr{F}_0\). Plus précisément, tout morphisme entre faisceaux lisses \(\iota\)-mixtes de poids ponctuels entiers est strictement compatible à leurs filtrations par le poids ponctuel.}

(iii) \textit{Si \(\mathscr{F}_0\) est lisse ponctuellement \(\iota\)-pur, et que \(\mathrm{X}_0\) est normal, alors le faisceau \(\mathscr{F}\) sur \(\mathrm{X}\) est semi-simple.}

\Page{208}

\textit{Lemme} \textbf{(3.4.2)}. — \textit{Si \(\mathscr{F}_0\) et \(\mathscr{G}_0\) sont deux faisceaux lisses sur \(\mathrm{X}_0\), on a une suite exacte :}
\[
0\to\mathrm{H}^0(\mathrm{X},\mathscr{H}\!om(\mathscr{F},\mathscr{G}))_{\mathrm{F}}
\to\operatorname{Ext}^1(\mathscr{F}_0,\mathscr{G}_0)
\to\mathrm{H}^1(\mathrm{X},\mathscr{H}\!om(\mathscr{F},\mathscr{G}))^{\mathrm{F}}.
\]

Dans ce lemme, \(\mathrm{F}\) en exposant (resp. en indice) indique le groupe des invariants (resp. co-invariants), pour l’action de \(\mathrm{W}(\bar{k}/\mathbf{F}_q)\). Le \(\operatorname{Ext}^1\) est le groupe des classes d’extensions, dans la catégorie abélienne des faisceaux sur \(\mathrm{X}_0\). La flèche de droite est l’image réciproque sur \(\mathrm{X}\) : \(\operatorname{Ext}^1(\mathscr{F}_0,\mathscr{G}_0)\to\operatorname{Ext}^1(\mathscr{F},\mathscr{G})=\mathrm{H}^1(\mathrm{X},\mathscr{H}\!om(\mathscr{F},\mathscr{G}))\). On tombe dans la partie invariante par \(\mathrm{F}\). Si une extension \(\mathscr{E}_0\) de \(\mathscr{F}_0\) par \(\mathscr{G}_0\) est géométriquement triviale, elle admet sur \(\mathrm{X}\) un scindage \(\varphi:\mathscr{F}\to\mathscr{E}\). Les autres scindages sont de la forme \(\varphi-f\), avec \(f\in\operatorname{Hom}(\mathscr{F},\mathscr{G})\). L’extension est triviale sur \(\mathrm{X}_0\) si et seulement si \(\varphi-f\) peut être choisi invariant par \(\mathrm{F}\), i.e. si \(\mathrm{F}\varphi-\varphi\in\operatorname{Hom}(\mathscr{F},\mathscr{G})\) est de la forme \(\mathrm{F}f-f\), i.e. est d’image nulle dans \(\operatorname{Hom}(\mathscr{F},\mathscr{G})_{\mathrm{F}}\), et la construction \(\mathscr{E}_0\mapsto(\mathrm{F}\varphi-\varphi)\) met en bijection classes d’extensions géométriquement triviales et \(\operatorname{Hom}(\mathscr{F},\mathscr{G})_{\mathrm{F}}=\mathrm{H}^0(\mathrm{X},\mathscr{H}\!om(\mathscr{F},\mathscr{G}))_{\mathrm{F}}\). Le lemme en résulte.

\textit{Lemme} \textbf{(3.4.3)}. — \textit{Si \(\mathscr{F}_0\) et \(\mathscr{G}_0\) sont lisses sur \(\mathrm{X}_0\) lisse, et ponctuellement \(\iota\)-purs de poids \(\beta\) et \(\gamma\), alors, il ne peut exister d’extension \(\mathscr{E}_0\) de \(\mathscr{F}_0\) par \(\mathscr{G}_0\), géométriquement non triviale, que si \(\beta\equiv\gamma(\mathrm{mod}\,\mathbf{Z})\) et que \(\beta>\gamma\).}

Il suffit de vérifier que \(\mathrm{F}\), agissant sur \(\mathrm{H}^1(\mathrm{X},\mathscr{H}\!om(\mathscr{F},\mathscr{G}))\), n’admet pas la valeur propre \(1\). Le faisceau \(\mathscr{H}\!om(\mathscr{F},\mathscr{G})\) est en effet lisse, de poids \(\gamma-\beta\), et d’après (3.3.5), amplifié par (3.3.10), les seuls poids qui peuvent apparaître dans son \(\mathrm{H}^1\) sont de la forme \(\gamma-\beta+1+n\), avec \(n\) entier \(\geq0\) : le poids \(0\) ne peut apparaître que pour \(\beta\equiv\gamma(\mathrm{mod}\,\mathbf{Z})\) et \(\beta>\gamma\).

La même démonstration donne :

\textit{Lemme} \textbf{(3.4.4)}. — \textit{Si \(\mathscr{F}_0\) et \(\mathscr{G}_0\) sont lisses sur \(\mathrm{X}_0\) lisse, et ponctuellement \(\iota\)-purs de poids \(\beta\) et \(\gamma\), alors, on ne peut avoir \(\operatorname{Ext}^1(\mathscr{F}_0,\mathscr{G}_0)\neq0\) que si \(\beta\equiv\gamma(\mathrm{mod}\,\mathbf{Z})\) et \(\beta\geq\gamma\).}

\textbf{(3.4.5)} \textit{Preuve de} (iii).

Si \(\mathrm{U}_0\) est un ouvert de \(\mathrm{X}_0\), et \(\bar{u}\) un point géométrique de \(\mathrm{U}_0\), \(\pi_1(\mathrm{U}_0,\bar{u})\) s’envoie sur \(\pi_1(\mathrm{X}_0,\bar{u})\). Remplaçant \(\mathrm{X}_0\) par \(\mathrm{U}_0\), ceci nous ramène à supposer que \(\mathrm{X}_0\) est lisse.

Soit \(\mathscr{F}'\) le plus grand sous-faisceau lisse semi-simple de \(\mathscr{F}\) (la somme des sous-faisceaux irréductibles de \(\mathscr{F}\)). Par transport de structures, il est stable par Frobenius, donc provient d’un sous-faisceau \(\mathscr{F}'_0\) de \(\mathscr{F}_0\). Posons \(\mathscr{F}''_0=\mathscr{F}_0/\mathscr{F}'_0\).

D’après (3.4.3), puisque \(\mathscr{F}'_0\) et \(\mathscr{F}''_0\) sont \(\iota\)-purs de même poids, l’extension \(\mathscr{F}_0\) de \(\mathscr{F}''_0\) par \(\mathscr{F}'_0\) est géométriquement triviale. Si \(\mathscr{F}''\neq0\), ceci contredit la maximalité de \(\mathscr{F}'\) (relever dans \(\mathscr{F}\) un sous-faisceau simple de \(\mathscr{F}''\)). On a donc \(\mathscr{F}''=0\), et ceci prouve (3.4.1) (iii).

% ---- 209-216 ----
\Page{209}
\textbf{(3.4.6)} \textit{Preuve de l’unicité dans} (i), (ii).

Soit \(x\in|\mathrm X|\), au-dessus de \(x_0\in|\mathrm X_0|\). Pour chaque \(\beta\in\mathbf R\), soit \(\mathscr F_x(\beta)\) la somme des sous-espaces propres généralisés de \(\mathrm F_{x_0}\) relatifs aux valeurs propres \(\alpha\) de \(\iota\)-poids \(\beta\), rel. \(q^{\deg(x_0)}\). On a nécessairement :
\[
\mathscr F_0(b)_x=\bigoplus_{\beta\in b}\mathscr F_x(\beta)
\]
et si les poids ponctuels sont entiers :
\[
(\mathrm W_i\mathscr F_0)_x=\bigoplus_{\beta\leq i}\mathscr F_x(\beta).
\]
Ceci prouve l’unicité, et la fonctorialité.

\textbf{(3.4.7)} Prouvons (i), (ii) sous les hypothèses additionnelles que \(\mathrm X_0\) est lisse et que \(\mathscr F_0\) est extension successive de faisceaux lisses ponctuellement \(\iota\)-purs.

C’est une conséquence formelle de (3.4.4) et de la semi-exactitude du foncteur \(\operatorname{Ext}^1\). On procède par récurrence sur le nombre d’extensions requises pour construire \(\mathscr F_0\). Ceci nous permet de supposer que \(\mathscr F_0\) est extension d’un faisceau lisse \(\mathscr F''_0\), vérifiant (i), (ii), par un faisceau lisse ponctuellement \(\iota\)-pur \(\mathscr F'_0\) de poids \(\beta\).

\textit{Preuve de} (i). — Soit \(b\) la classe de \(\beta\) dans \(\mathbf R/\mathbf Z\). D’après (3.4.4), si \(b'\ne b\), l’image inverse de \(\mathscr F''_0(b')\) dans \(\mathscr F_0\) est une extension triviale de \(\mathscr F''_0(b')\) par \(\mathscr F'_0\). On prend pour \(\mathscr F_0(b)\) l’image inverse de \(\mathscr F''_0(b)\), et pour \(\mathscr F_0(b')\) (\(b'\ne b\)) un relèvement dans \(\mathscr F''_0(b')\).

\textit{Preuve de} (ii). — Sous les hypothèses (ii), \(\beta\) est un entier. D’après (3.4.4), l’image inverse de \(\mathrm W_{\beta-1}\mathscr F''_0\) dans \(\mathscr F_0\) est une extension triviale de \(\mathrm W_{\beta-1}\mathscr F''_0\) par \(\mathscr F'_0\). Pour \(i<\beta\), on prend pour \(\mathrm W_i\mathscr F_0\) un relèvement de \(\mathrm W_i\mathscr F''_0\subset\mathrm W_{\beta-1}\mathscr F''_0\) dans \(\mathscr F_0\). Pour \(i\geq\beta\), on prend pour \(\mathrm W_i\mathscr F_0\) l’image inverse de \(\mathrm W_i\mathscr F''_0\).

\textbf{(3.4.8)} \textit{Preuve de} (i), (ii).

On procède par récurrence sur \(\dim\mathrm X_0\). Quitte à remplacer \(\mathrm X_0\) par \(\mathrm X_{0\mathrm{red}}\), il existe un ouvert dense \(j:\mathrm U_0\hookrightarrow\mathrm X_0\) sur lequel \(\mathscr F_0\) satisfait aux hypothèses additionnelles de (3.4.7). L’hypothèse de récurrence s’applique au complément \(\mathrm F_0\). Rappelons enfin que le foncteur \(\mathscr F_0\mapsto(j^*\mathscr F_0,i^*\mathscr F_0,\text{flèche de spécialisation})\) est une équivalence de la catégorie des faisceaux sur \(\mathrm X_0\) avec la catégorie des triples \((\mathscr F'_0,\mathscr F''_0,s)\) formés d’un faisceau sur \(\mathrm U_0\), d’un faisceau sur \(\mathrm F_0\), et d’un morphisme \(s:\mathscr F''_0\to i^*j_*\mathscr F'_0\).

\textit{Preuve de} (i). — Il s’agit de voir que le morphisme de spécialisation :
\[
s:i^*\mathscr F_0\to i^*j_*j^*\mathscr F_0
\]
envoie \((i^*\mathscr F_0)(b)\) dans \(i^*j_*((j^*\mathscr F_0)(b))\). On déduit en effet de (1.8.9) par torsion que les poids ponctuels de \(i^*j_*((j^*\mathscr F_0)(b))\) sont dans \(b\), de sorte que ce faisceau coïncide avec \((i^*j_*j^*\mathscr F_0)(b)\).

\Page{210}
\textit{Preuve de} (ii). — L’unicité étant assurée, on peut supposer par descente que \(\mathrm X_0\) est normal. On a alors \(\mathscr F_0=j_*j^*\mathscr F_0\), et l’on prend \(\mathrm W_i\mathscr F_0=j_*\mathrm W_i j^*\mathscr F_0\); ceci donne \(\operatorname{Gr}_i^{\mathrm W}\mathscr F_0=j_*\operatorname{Gr}_i^{\mathrm W}j^*\mathscr F_0\), et l’on applique (1.8.10).

\textit{Variante} \textbf{(3.4.9)}. — On déduit aussitôt de (3.4.1) que tout faisceau mixte lisse admet une \textit{filtration par le poids ponctuel} par des sous-faisceaux lisses, et que tout faisceau pur est géométriquement semi-simple.

\textbf{(3.4.10)} Soit \(\mathrm P\) une propriété de faisceaux sur les schémas de type fini sur un corps fini (p.e. être mixte, être ponctuellement \(\iota\)-pur...). Un faisceau \(\mathscr F\) sur un schéma \(\mathrm X\) de type fini sur \(k\) algébriquement clos est dit avoir \textit{potentiellement} la propriété \(\mathrm P\) si \((\mathrm X,\mathscr F)/k\) se déduit par un changement de base \(\bar x:\operatorname{Spec}(k)\to\mathrm S\) d’un système \((\mathrm X_{\mathrm S},\mathscr F_{\mathrm S})\) sur \(\mathrm S\) — avec \(\mathrm X_{\mathrm S}\) de type fini sur \(\mathrm S\) de type fini sur \(\mathbf Z\) — tel que les restrictions de \(\mathscr F_{\mathrm S}\) aux fibres de \(\mathrm X_{\mathrm S}/\mathrm S\) en les points fermés de \(\mathrm S\) aient la propriété \(\mathrm P\).

\textit{Exemple} \textbf{(3.4.11)}. — Si \(f:\mathrm Y\to\mathrm X\) est propre et lisse, les faisceaux \(\mathrm R^if_*\mathbf Q_\ell\) sont potentiellement ponctuellement purs. Ceci résulte d’un argument standard de passage à la limite, pour construire \(f_0:\mathrm Y_0\to\mathrm X_0/\mathrm S_0\) propre et lisse, et de (3.3.9) appliqué aux fibres de \(f_0\).

Le théorème de spécialisation (1.11.1), (1.11.5) permet de généraliser (3.4.1) (iii) comme suit :

\textit{Corollaire} \textbf{(3.4.12)}. — \textit{Soient \(\mathrm X\) un schéma normal de type fini sur \(k\) algébriquement clos, et \(\mathscr F\) un faisceau lisse sur \(\mathrm X\). Si \(\mathscr F\) est potentiellement ponctuellement \(\iota\)-pur, il est semi-simple.}

En particulier, l’exemple (3.4.11) donne le

\textit{Corollaire} \textbf{(3.4.13)}. — \textit{Soient \(\mathrm S\) un schéma normal connexe sur \(k\) algébriquement clos, et \(f:\mathrm X\to\mathrm S\) un morphisme propre et lisse. Alors, les faisceaux \(\mathrm R^if_*\mathbf Q_\ell\) sont semi-simples.}

\textit{Remarque} \textbf{(3.4.14)}. — Le point clef dans la démonstration du théorème de Lefschetz difficile donnée en (4.1.1) est un cas particulier de (3.4.13) dans lequel \(f\) est projectif et \(\mathrm S\) une courbe (un ouvert de \(\mathbf P^1\)). Pour traiter ce cas, il suffit de (1.1.6) (pour assurer la pureté de \(\mathrm R^if_{0*}\mathbf Q_\ell\)), (1.11.1), pour l’argument de spécialisation, et (2.2.10) (Hadamard-de la Vallée-Poussin), pour interdire la valeur propre \(1\) dans la preuve de (3.4.1).

\section*{(3.5) — Application : théorèmes d’équidistribution}

\textbf{(3.5.1)} Reprenons les notations et hypothèses de (2.2) : on a :
\[
\begin{tikzcd}[column sep=large,row sep=large]
0 \arrow[r] &
\pi_1(\mathrm X,\bar x) \arrow[r] \arrow[d,hook] &
\mathrm W(\mathrm X_0,\bar x) \arrow[r] \arrow[d,hook] &
\mathbf Z \arrow[r] \arrow[d,equal] & 0 \\
0 \arrow[r] &
\mathrm G^0 \arrow[r] &
\mathrm G \arrow[r] &
\mathbf Z \arrow[r] & 0
\end{tikzcd}
\]
\Page{211}
avec \(\mathrm X_0\) normal géométriquement connexe de dimension \(\mathrm N\), et, d’après (2.2.8), pour toute représentation irréductible \(\tau\) de \(\mathrm G_{\mathbf R}\), le faisceau \(\mathscr F(\tau_\ell)\) correspondant est ponctuellement \(\iota\)-pur de poids \(2\mathscr R(\tau)\). On note \(\mathrm G_{\mathbf R i}^{\natural}\) l’espace des classes de conjugaison de degré \(i\) de \(\mathrm G_{\mathbf R}^{\natural}\), et \(\mu_0\) la mesure sur \(\mathrm G_{\mathbf R}^{\natural}\) définie en (2.1.10).

Si \(\tau\in\widetilde{\mathrm G}\) est une représentation complexe irréductible de \(\mathrm G\), le produit qui définit \(\mathrm L(\tau)\) converge pour \(\mathscr R(\tau)>\mathrm N\). La formule (2.2.8.1), l’interprétation cohomologique de Grothendieck des fonctions \(\mathrm L\) (1.4.5.1), et (3.3.4) modulé par la deuxième variation (3.3.10), montrent que \(\mathrm L(\tau)\) se prolonge en une fonction méromorphe sur \(\widetilde{\mathrm G}\), qui n’a de zéros ou de pôles que pour \(\mathscr R(\tau)\) entier ou demi-entier, et que pour \(\mathscr R(\tau)>\mathrm N-1/2\), \(\mathrm L(\tau)\) est holomorphe inversible sauf pour un pôle simple en \(\tau=\omega_{\mathrm N}\). Ceci permet d’appliquer les résultats d’équidistribution (2.1.10) à (2.1.12) (avec \(\omega_1\) remplacé par \(\omega_{\mathrm N}\), cf. (2.2.7)).

\textbf{(3.5.2)} Pour tout entier \(n>0\), soit \(\mathbf F_{q^n}\) l’extension de degré \(n\) de \(\mathbf F_q\) dans \(\bar k\), et \(\mathrm X_{(n)}\) le schéma sur \(\mathbf F_{q^n}\) déduit de \(\mathrm X_0\) par extension des scalaires. Le schéma \(\mathrm X_{(n)}\) est un revêtement de \(\mathrm X_0\); il correspond au sous-groupe de \(\mathrm W(\mathrm X_0,\bar x)\) formé des éléments de degré divisible par \(n\). Chaque point \(x\) de \(\mathrm X_0\) à coordonnées dans \(\mathbf F_{q^n}\) définit un point rationnel de \(\mathrm X_{(n)}\), donc un point fermé de \(\mathrm X_{(n)}\). On note \(\mathrm F_x\) un élément de Frobenius géométrique correspondant, et son image dans \(\mathrm W(\mathrm X_0,\bar x)\). Cette image est de degré \(n\). On note encore \(\mathrm F_x\) son image dans \(\mathrm G\), et \((\iota\mathrm F_x)_s\) la classe de conjugaison dans \(\mathrm G_{\mathbf R}\) des éléments conjugués à la composante semi-simple de \(\iota\mathrm F_x\in\mathrm G_{\mathbf C}\) (2.2.4).

Si un point fermé \(x_0\) de \(\mathrm X_0\) est image d’un point rationnel de \(\mathrm X_{(n)}\), son degré \(d\) divise \(n\). Réciproquement, si \(d\mid n\), il y a \(d\) points rationnels \(x\) de \(\mathrm X_{(n)}\) au-dessus de \(x_0\), et pour chacun d’eux \(\mathrm F_x=\mathrm F_{x_0}^{n/d}\) (à conjugaison près). Sur l’espace des classes de conjugaison dans \(\mathrm W(\mathrm X_0,x)\), on a donc une identité entre mesures :
\[
\sum_{\substack{n>0\\x_0\in|\mathrm X_0|}}
\deg(x_0)\delta[\mathrm F_{x_0}^{n}]
=
\sum_{\substack{n>0\\x\in\mathrm X_0(\mathbf F_{q^n})}}
\delta[\mathrm F_x].
\tag*{\textbf{(3.5.2.1)}}
\]

Choisissons un élément central \(z\), de degré \(d>0\), pour identifier entre eux les \(\mathrm G_{\mathbf R i}^{\natural}\) lorsque \(i\) parcourt une progression arithmétique de raison \(d\). Compte tenu de (3.5.2.1), le théorème (2.1.12) se spécialise en le suivant :

\textit{Théorème} \textbf{(3.5.3)}. — \textit{Avec les notations de (3.5.1), quel que soit \(i\), la mesure \(\displaystyle z^{-n}\bigl(q^{-\mathrm N(nd+i)}\sum_{x\in\mathrm X_0(\mathbf F_{q^{nd+i}})}\delta[(\iota\mathrm F_x)_s]\bigr)\) sur \(\mathrm G_{\mathbf R i}^{\natural}\) converge vaguement vers \(\mu_0^{\natural}\vert\mathrm G_{\mathbf R i}^{\natural}=z^{-n}(\mu_0^{\natural}\vert\mathrm G_{\mathbf R(nd+i)}^{\natural})\) pour \(n\to\infty\).}

\textit{Exemple} \textbf{(3.5.4)}. — Soit \(f:\mathrm E_0\to\mathrm X_0\) une famille de courbes elliptiques paramétrée par \(\mathrm X_0\), et prenons \(\mathscr F_0=\mathrm R^1f_*\mathbf Q_\ell\). Posons \(\mathrm V=\mathscr F_{0\bar x}\). On a un isomorphisme naturel \(\bigwedge\limits^{2}\mathrm V\sim\mathbf Q_\ell(-1)\), compatible à l’action de \(\mathrm G\). Puisque \(\mathrm G\) agit sur \(\mathbf Q_\ell(-1)\) par \(\omega_{-1}\), les éléments de degré \(n\) de \(\mathrm G\) agissent sur \(\mathrm V\) avec pour déterminant \(q^n\).

\textit{Lemme} \textbf{(3.5.5)}. — \textit{Si l’invariant modulaire \(j\) de \(\mathrm E/\mathrm X\) est non constant, le groupe \(\mathrm G^0\) est le groupe \(\operatorname{SL}(\mathrm V)\) tout entier.}

\Page{212}
Soient \(n\) un entier supérieur ou égal à \(3\) premier à \(p\), et \(\mathrm X'\) l’une des composantes connexes du revêtement \(\mathscr I\mathit{som}(\mathrm E_n,(\mathbf Z/n)^2)\) de \(\mathrm X\) : c’est un revêtement galoisien étale de \(\mathrm X\), de groupe de Galois un sous-groupe de \(\operatorname{SL}(2,\mathbf Z/n)\), sur lequel le faisceau \(\mathrm E_n\) est trivial : on dispose de \(\alpha:\mathrm E_n\to(\mathbf Z/n)^2\). Soit \(\bar x'\) un point base au-dessus de \(\bar x\). Le couple \((\mathrm E,\alpha)\) définit un morphisme \(u\) de \(\mathrm X'\) dans le schéma de module \(\mathrm M_n\) des courbes elliptiques munies d’une structure de niveau \(n\). L’hypothèse sur \(j\) assure que le morphisme \(u\) est dominant, donc que l’image de \(u:\pi_1(\mathrm X',\bar x')\to\pi_1(\mathrm M_n,u(\bar x'))\) est d’indice fini. On sait par voie transcendante que \(\pi_1(\mathrm M_n,\bar x')\) s’envoie sur le sous-groupe de \(\operatorname{SL}(\mathrm H^1(\mathrm E_{\bar x},\mathbf Z_\ell))\) formés des éléments \(\equiv1\bmod n\). Le groupe de monodromie géométrique de \(\mathscr F/\mathrm X\) contient donc un sous-groupe d’indice fini de \(\operatorname{SL}(\mathrm H^1(\mathrm E_{\bar x},\mathbf Z_\ell))\), et (3.5.5) en résulte.

\textbf{(3.5.6)} On suppose \(j\) non constant. Le lemme (3.5.5) montre alors que \(\mathrm G_{\mathbf R}\) n’est autre que \(\operatorname{SU}(2)\times\mathbf Z\), agissant sur \(\iota\mathrm V\sim\mathbf C^2\) par \((g,n)\to q^{n/2}\cdot g\). Identifions l’ensemble des classes de conjugaisons dans \(\operatorname{SU}(2)\) à \([0,\pi]\), par :
\[
\text{classe de }
\begin{pmatrix}
e^{i\theta}&0\\
0&e^{-i\theta}
\end{pmatrix}
\longmapsto\theta.
\]
On sait que l’image directe de la mesure de Haar de \(\operatorname{SU}(2)\) est la mesure \(\displaystyle\frac{2}{\pi}\sin^2\theta\,d\theta\) sur \([0,\pi]\).

Pour \(x\in\mathrm X_0(\mathbf F_{q^n})\), on définit l’angle de Frobenius \(\theta(x)\in[0,\pi]\) en demandant que les valeurs propres de \(\iota\mathrm F_x\) soient les nombres \(e^{\pm i\theta(x)}\cdot q^{n/2}\), i.e. en demandant que la courbe \(\mathrm E_x\) sur \(\mathbf F_{q^n}\) ait \(1-2\cos\theta(x)\cdot q^{n/2}+q^n\) points rationnels. Dans \(\mathrm G_{\mathbf R}\sim\operatorname{SU}(2)\times\mathbf Z\), on a \((\iota\mathrm F_x)_s=(\theta(x),n)\). Le théorème fournit donc, \textit{en égale caractéristique}, le théorème d’équidistribution conjecturé par Sato-Tate :

\textit{Proposition} \textbf{(3.5.7)}. — \textit{Avec les notations précédentes (et pour \(j\) non constant), lorsque \(n\to\infty\), la mesure \(\displaystyle\frac{1}{q^{n\mathrm N}}\sum_{x\in\mathrm X_0(\mathbf F_{q^n})}\delta[\theta(x)]\) tend vaguement vers \(\displaystyle\frac{2}{\pi}\sin^2\theta\,d\theta\).}

Ce résultat avait été obtenu par H. Yoshida (On an analogue of the Sato conjecture, \textit{Inv. Math., 19}, 4 (1973), p. 261-277) pour certaines familles modulaires de courbes elliptiques. C’est originellement pour étendre son résultat au cas d’une famille à un paramètre quelconque que j’ai démontré le théorème à la Hadamard-de la Vallée-Poussin du § 2.

\section*{(3.6) — Application : le théorème local des cycles invariants}

Soient \(k\) un corps algébriquement clos, \(\mathrm S\) le spectre de l’hensélisé de \(k[\mathrm T]\) en l’idéal \((\mathrm T)\), \(s\) le point fermé de \(\mathrm S\), \(\eta\) le point générique et \(\bar\eta\) un point générique géométrique. Soit \(f:\mathrm X\to\mathrm S\) un morphisme propre. On dispose alors d’un morphisme de spécialisation \(\operatorname{sp}^*:\mathrm H^*(\mathrm X_s)\to\mathrm H^*(\mathrm X_{\bar\eta})\); son image est contenue dans les invariants sous l’action du groupe d’inertie \(\operatorname{Gal}(\bar\eta/\eta)\).

\Page{213}
\textit{Théorème} \textbf{(3.6.1)}. — \textit{Supposons \(\mathrm X\) essentiellement lisse sur \(k\), et \(\mathrm X_{\bar\eta}\) lisse sur \(\bar\eta\). Alors, en \(\mathbf Q_\ell\)-cohomologie, le morphisme :}
\[
\operatorname{sp}^*:\mathrm H^*(\mathrm X_s)\to
\mathrm H^*(\mathrm X_{\bar\eta})^{\operatorname{Gal}(\bar\eta/\eta)}
\]
\textit{est surjectif.}

Pour pouvoir parler de poids, on commence par se ramener à une situation arithmétique. Un argument de passage à la limite permet tout d’abord de supposer que \(\mathrm S\) est l’hensélisé en un point \(s\) d’une courbe \(\mathrm S'\) lisse sur \(k\), et que \(f\) se déduit par le changement de base \(g:(\mathrm S,s)\to(\mathrm S',s)\) d’un morphisme propre \(f':\mathrm X'\to\mathrm S'\), avec \(\mathrm X'\) lisse sur \(k\), et lisse sur \(\mathrm S'\) sauf au-dessus de \(s\). Utilisant (1.11.3), on se ramène ensuite à supposer que \(k\) est la clôture algébrique de \(\mathbf F_q\). Pour \(\mathbf F_q\) assez grand, on peut enfin supposer que \((f',\mathrm X',\mathrm S',s)\) sur \(k\) provient par extension des scalaires de \((f'_0,\mathrm X'_0,\mathrm S'_0,s_0)\) sur \(\mathbf F_q\). Soient \((\mathrm S_0,\eta_0,s_0)\) l’hensélisé de \(\mathrm S'_0\) en \(s_0\), et \(f_0:\mathrm X_0\to\mathrm S_0\) déduit de \(f'_0\) par le changement de base \(g_0:(\mathrm S_0,s_0)\to(\mathrm S'_0,s_0)\). On a :

\noindent — sur \(\mathbf F_q\) : \((f'_0,\mathrm X'_0,\mathrm S'_0,s_0)\) global, qui se localise en \((f_0,\mathrm X_0,\mathrm S_0,s_0)\);

\noindent — sur \(k=\mathbf F\) : \((f',\mathrm X',\mathrm S',s)\) global, qui se localise en \((f,\mathrm X,\mathrm S,s)\).

Le groupe de Galois \(\operatorname{Gal}(\bar\eta/\eta)\) est le groupe d’inertie de \(\operatorname{Gal}(\bar\eta/\eta_0)\); ce dernier agit sur la cohomologie de \(\mathrm X_{\bar\eta}\), et ceci permet de parler de poids (cf. (1.7.4)).

Calculons la \(\mathbf Z/\ell^n\) cohomologie de \(\mathrm X_\eta\) à l’aide de la suite spectrale de Leray de \(\mathrm X_\eta\to\eta\), appelée aussi suite spectrale de Hochschild-Serre :
\[
\mathrm E_2^{pq}=\mathrm H^p(\mathrm I,\mathrm H^q(\mathrm X_{\bar\eta}))
\Rightarrow\mathrm H^{p+q}(\mathrm X_\eta).
\tag*{\textbf{(1)}}
\]
On sait que \(\mathrm I\) est une extension de \(\mathbf Z_\ell(1)\) (\(\ell\)-quotient modéré de \(\mathrm I\)) par un groupe profini \(\mathrm I'\) d’ordre premier à \(\ell\). Pour tout \(\mathbf Z/\ell^n[\mathrm I]\)-module \(\mathrm V\), on a donc une suite spectrale :
\[
\mathrm E_2^{pq}=\mathrm H^p(\mathbf Z_\ell(1),\mathrm H^q(\mathrm I',\mathrm V))
\Rightarrow\mathrm H^{p+q}(\mathrm I,\mathrm V).
\tag*{\textbf{(2)}}
\]
L’espace \(\mathrm V^{\mathrm I'}\) des invariants de \(\mathrm I'\) dans \(\mathrm V\) a un unique supplémentaire \(\mathrm I'\)-stable. Sur ce dernier, \(\mathrm I'\) n’a ni invariant, ni co-invariant. Ceci fournit un isomorphisme \(\mathrm V^{\mathrm I'}\sim\mathrm V_{\mathrm I'}\) entre invariants et co-invariants. De plus, \(\mathrm H^q(\mathrm I',\mathrm V)=0\) pour \(q>0\), de sorte que (2) se réduit à :
\[
\mathrm H^p(\mathrm I,\mathrm V)
=\mathrm H^p(\mathbf Z_\ell(1),\mathrm V^{\mathrm I'})
=\mathrm H^p(\mathbf Z_\ell(1),\mathrm V_{\mathrm I'}).
\tag*{\textbf{(3)}}
\]
Pour toute représentation \(\ell\)-adique \(\mathrm V\) de \(\mathbf Z_\ell(1)\), on a :
\begin{align*}
\mathrm H^0(\mathbf Z_\ell(1),\mathrm V)&=\mathrm V^{\mathbf Z_\ell(1)},
\tag*{\textbf{(4)}}\\
\mathrm H^1(\mathbf Z_\ell(1),\mathrm V)&=\mathrm V_{\mathbf Z_\ell(1)}(-1),\\
\mathrm H^p(\mathbf Z_\ell(1),\mathrm V)&=0\quad\text{pour}\quad p\geq2.
\end{align*}
Combinant (1), (3), (4), on trouve des suites exactes (analogues aux suites exactes de Wang) :
\[
0\to\mathrm H^{i-1}(\mathrm X_{\bar\eta})_{\mathrm I}(-1)
\to\mathrm H^i(\mathrm X_\eta)
\to\mathrm H^i(\mathrm X_{\bar\eta})^{\mathrm I}\to0.
\tag*{\textbf{(5)}}
\]
On dispose d’une suite exacte longue de \(\mathbf Z/\ell^n\)-cohomologie à support :
\[
\to\mathrm H^i_{\mathrm X_s}(\mathrm X)
\to\mathrm H^i(\mathrm X)
\to\mathrm H^i(\mathrm X_\eta)\to
\tag*{\textbf{(6)}}
\]

\Page{214}
Le théorème de changement de base pour le morphisme propre \(f\) donne un isomorphisme \(\mathrm H^*(\mathrm X)\xrightarrow{\sim}\mathrm H^*(\mathrm X_s)\), et \(\operatorname{sp}^*\) est le composé :
\[
\operatorname{sp}^*:\mathrm H^*(\mathrm X_s)
\xleftarrow{\sim}\mathrm H^*(\mathrm X)
\to\mathrm H^*(\mathrm X_\eta)
\to\mathrm H^*(\mathrm X_{\bar\eta}).
\tag*{\textbf{(7)}}
\]
Puisque \(\mathrm X\) est essentiellement étale sur \(\mathrm X'\), et que \(\mathrm X_s\xrightarrow{\sim}\mathrm X'_s\), on a \(\mathrm H^i_{\mathrm X'_s}(\mathrm X')\xrightarrow{\sim}\mathrm H^i_{\mathrm X_s}(\mathrm X)\). Remplaçant \(\mathrm X'_0\) par une de ses composantes connexes, on se ramène à supposer que \(\mathrm X'\) est purement d’une dimension \(\mathrm N\). Le groupe de cohomologie \(\mathrm H^i_{\mathrm X'_s}(\mathrm X')\) peut alors s’interpréter comme un groupe d’homologie de \(\mathrm X'_s\) : si \(a\) est la projection de \(\mathrm X\) sur \(\operatorname{Spec}(k)\), l’adjonction entre \(\mathrm R a_!\) et \(\mathrm R a^!\) appliquée à \((\mathbf Z/\ell^n)_{\mathrm X'_s}\) et aux \(\mathbf Z/\ell^n[i]\) sur \(\operatorname{Spec}(k)\) fournit une dualité parfaite à valeurs dans \(\mathbf Z/\ell^n(-\mathrm N)\) entre \(\mathrm H^i_{\mathrm X'_s}(\mathrm X')\) et \(\mathrm H^{2\mathrm N-i}(\mathrm X'_s)\), i.e. entre \(\mathrm H^i_{\mathrm X_s}(\mathrm X)\) et \(\mathrm H^{2\mathrm N-i}(\mathrm X_s)\). Les suites (5) et (6) nous fournissent donc une croix de suites exactes :
\[
\begin{tikzcd}[column sep=small,row sep=large]
& & \mathrm H^{2\mathrm N-i-1}(\mathrm X_s)^{\vee}(-\mathrm N) & & \\
0 \arrow[r] &
\mathrm H^{i-1}(\mathrm X_{\bar\eta})_{\mathrm I}(-1) \arrow[r] &
\mathrm H^i(\mathrm X_\eta) \arrow[r] \arrow[u] &
\mathrm H^i(\mathrm X_{\bar\eta})^{\mathrm I} \arrow[r] & 0 \\
& & \mathrm H^i(\mathrm X_s) \arrow[u] \arrow[ur,"\operatorname{sp}^*"'] & &
\end{tikzcd}
\tag*{\textbf{(8)}}
\]

Par passage à la limite, cette croix subsiste en \(\mathbf Z_\ell\)-cohomologie, avec en haut le groupe \(\varprojlim\mathrm H^{i+1}_{\mathrm X'_s}(\mathrm X',\mathbf Z/\ell^n)\). La présence possible de torsion complique l’interprétation de ce groupe par dualité, mais cette complication disparaît par passage à la \(\mathbf Q_\ell\)-cohomologie, où la croix (8) est à nouveau valable. Le groupe de Galois arithmétique \(\operatorname{Gal}(\bar\eta/\eta_0)\) agit sur cette croix via son quotient \(\operatorname{Gal}(\mathbf F/\mathbf F_q)\). Calculons les poids.

\textit{Lemme} \textbf{(3.6.2)}. — \textit{\(\mathrm H^i(\mathrm X_{\bar\eta})^{\mathrm I}\) est mixte de poids \(\leq i\).}

Le groupe \(\mathrm H^i(\mathrm X_{\bar\eta})\) est la fibre en \(\bar\eta\) du faisceau \(\mathrm R^if_{0*}\mathbf Q_\ell\) sur \(\mathrm S_0\). D’après (3.3.9) appliqué aux fibres de \(f_0\), ce faisceau est pur de poids \(i\), et on applique (1.8.8) (i).

\textit{Lemme} \textbf{(3.6.3)}. — \textit{\(\mathrm H^{2\mathrm N-i-1}(\mathrm X_s)^{\vee}(-\mathrm N)\) est mixte de poids \(\geq i+1\).}

Résulte aussitôt de (3.3.4).

Bien que ce ne nous soit pas indispensable, notons encore que \(\mathrm H^i(\mathrm X_s)\) est mixte de poids \(\leq i\).

Le foncteur \(\mathrm W_i\) (partie de poids \(\leq i\)) est exact. En l’appliquant à (8), on obtient un diagramme exact :
\[
\begin{tikzcd}[column sep=large,row sep=large]
0 & & \\
\mathrm W_i\mathrm H^i(\mathrm X_\eta) \arrow[u] \arrow[r] &
\mathrm H^i(\mathrm X_{\bar\eta})^{\mathrm I} \arrow[r] & 0 \\
\mathrm H^i(\mathrm X_s) \arrow[u] & &
\end{tikzcd}
\]
dont le théorème résulte.

\Page{215}
\textbf{(3.6.4)} Soient \(\mathrm D\subset\mathbf C\) le disque unité, \(f:\mathrm X\to\mathrm D\) un morphisme propre, et supposons \(\mathrm X\) lisse, et \(f\) lisse au-dessus de \(\mathrm D^*:=\mathrm D-\{0\}\). Si \(f\) se factorise par :
\[
\operatorname{pr}_2:\mathbf P^n(\mathbf C)\times\mathrm D\to\mathrm D,
\]
on dispose d’un analogue de (3.6.1) : notant \(\mathrm X_0\) la fibre spéciale \(f^{-1}(0)\), et \(\mathrm X_t:=f^{-1}(t)\) une autre fibre (\(t\in\mathrm D^*\)), on a un homomorphisme surjectif :
\[
\mathrm H^*(\mathrm X_0,\mathbf Q)\twoheadrightarrow
\mathrm H^*(\mathrm X_t,\mathbf Q)^{\pi_1(\mathrm D^*,t)}.
\]
Ceci se démontre par un argument parallèle à (3.6.1), la filtration par le poids de la théorie de Hodge mixte se substituant à celle déduite de Frobenius. Voir J. Steenbrink, Mixed Hodge structure on the vanishing cohomology, \textit{Symp. in Math.}, Oslo, 1976.

\section*{(3.7) — Retour à (I.8)}

Les résultats des n\textsuperscript{os} (3.3) et (1.8) permettent de simplifier les preuves de (I.8.2) et (I.8.4).

\textbf{(3.7.1)} Soient \(\mathrm S_0\) un schéma lisse sur \(\mathbf F_q\), \(f:\mathrm E_0\to\mathrm S_0\) une famille de courbes elliptiques paramétrée par \(\mathrm S_0\), et :
\[
{}^k\mathrm W:=\operatorname{Im}\bigl(
\mathrm H_c^1(\mathrm S,\operatorname{Sym}^k(\mathrm R^1f_*\mathbf Q_\ell))
\to\mathrm H^1(\mathrm S,\operatorname{Sym}^k(\mathrm R^1f_*\mathbf Q_\ell))
\bigr).
\]
Pour \(\mathrm S_0\) une courbe, de complétion projective \(j:\mathrm S_0\hookrightarrow\overline{\mathrm S}_0\), \({}^k\mathrm W\) peut encore se décrire comme \(\mathrm H^1(\overline{\mathrm S},j_*\operatorname{Sym}^k(\mathrm R^1f_*\mathbf Q_\ell))\). Le faisceau \(\mathrm R^1f_*\mathbf Q_\ell\) est ponctuellement pur de poids \(1\). Sa puissance symétrique \(k\)-ième est donc pure de poids \(k\), et d’après (3.3.6), \({}^k\mathrm W\) est pur de poids \(k+1\). L’assertion (5.1) de [4] est un cas particulier de ce résultat. Ceci rend inutiles les arguments de \textit{loc. cit.} pour ramener [4] (5.1) à la conjecture de Weil sous sa forme usuelle (I.1.6). Rappelons que [4] (5.1) est un des points essentiels de la preuve de la conjecture de Ramanujan-Petersson citée en (I.8.2).

\textbf{(3.7.2)} Voici une version simplifiée de la preuve de (I.8.4). Soit \(\Psi:\mathbf F_p\to\mathbf C^*\) un caractère additif non trivial de \(\mathbf F_p\). Si \(\mathrm Q\in\mathbf F_q[\mathrm X_1,\ldots,\mathrm X_n]\) est un polynôme de degré \(d\) premier à \(p\), dont la partie homogène de degré \(d\), \(\mathrm Q_d\), définit une hypersurface projective et lisse, il s’agit de prouver que :
\[
\left|
\sum_{x_1,\ldots,x_n\in\mathbf F_q}
\Psi\operatorname{Tr}_{\mathbf F_q/\mathbf F_p}\mathrm Q(x_1,\ldots,x_n)
\right|
\leq(d-1)^nq^{n/2}.
\tag*{\textbf{(3.7.2.1)}}
\]
On regarde \(\Psi\) comme étant à valeurs \(\ell\)-adiques plutôt que complexes, à l’aide de \(\iota\), et, comme dans SGA \(4\frac{1}{2}\) [Sommes trigonométriques], on définit un faisceau \(\mathscr F_0(\Psi\mathrm Q)\) sur l’espace affine \(\mathrm A_0^n\) de dimension \(n\) sur \(\mathbf F_q\), à l’aide de \(\Psi\) et du revêtement d’Artin-Schreier \(\mathrm T^p-\mathrm T=\mathrm Q\) de \(\mathrm A_0^n\). On a :
\[
\sum_{x_1,\ldots,x_n\in\mathbf F_q}
\Psi\operatorname{Tr}_{\mathbf F_q/\mathbf F_p}\mathrm Q(x_1,\ldots,x_n)
=
\sum(-1)^i\operatorname{Tr}(\mathrm F,\mathrm H_c^i(\mathrm A^n,\mathscr F(\Psi\mathrm Q))),
\tag*{\textbf{(3.7.2.2)}}
\]
\Page{216}
et on va prouver le résultat plus précis que (3.7.2.1) :
\begin{align*}
&\mathrm H_c^i(\mathrm A^n,\mathscr F(\Psi\mathrm Q))=0
\quad\text{pour}\quad i\ne n,
\tag*{\textbf{(3.7.2.3)}}\\
&\dim\mathrm H_c^n(\mathrm A^n,\mathscr F(\Psi\mathrm Q))=(d-1)^n,
\quad\text{et}\\
&\mathrm H_c^n(\mathrm A^n,\mathscr F(\Psi\mathrm Q))
\quad\text{est pur de poids }n.
\end{align*}

Soient \(\mathrm S'_0\) l’espace affine sur \(\mathbf F_q\) qui paramétrise les polynômes à \(n\) variables de degré \(\leq d\), et \(\mathrm S_0\) l’ouvert de \(\mathrm S'_0\) qui paramétrise ceux dont la partie homogène de degré \(d\) définit une hypersurface non singulière (i.e. est de discriminant \(\ne0\)). Soit :
\[
\mathrm Q_{\mathrm S}\in
\mathrm H^0(\mathrm S_0,\mathscr O_{\mathrm S_0}[\mathrm X_1,\ldots,\mathrm X_n])
\]
le polynôme universel sur \(\mathrm S_0\); c’est une section du faisceau structural de \(\mathrm S_0\times\mathrm A_0^n\); il définit un faisceau \(\mathscr F_0(\Psi\mathrm Q_{\mathrm S})\) sur \(\mathrm S_0\times\mathrm A_0^n\).

Soit \(f\) la projection de \(\mathrm S_0\times\mathrm A_0^n\) sur \(\mathrm S_0\). Le polynôme \(\mathrm Q\) correspond à un point rationnel de \(\mathrm S_0\), et \((\mathrm A_0^n,\mathscr F_0(\Psi\mathrm Q))\) est la fibre de \((\mathrm S_0\times\mathrm A_0^n,\mathscr F_0(\Psi\mathrm Q_{\mathrm S}))\) en ce point.

\textit{Lemme} \textbf{(3.7.3)}. — \textit{Les faisceaux \(\mathrm R^if_!\mathscr F_0(\Psi\mathrm Q_{\mathrm S})\) sur \(\mathrm S_0\) sont lisses.}

Soient \(\mathbf P_0^n\) l’espace projectif sur \(\mathbf F_q\) déduit de l’espace affine \(\mathrm A_0^n\) par adjonction d’un hyperplan à l’infini \(\mathrm H_0\), \(\bar f\) la projection de \(\mathrm S_0\times\mathbf P_0^n\) sur \(\mathrm S_0\) et \(j\) l’inclusion de \(\mathrm S_0\times\mathrm A_0^n\) dans \(\mathrm S_0\times\mathbf P_0^n\). Localement pour la topologie étale, le \(\mathrm S\)-schéma \(\mathrm S\times\mathbf P^n\) muni du faisceau \(j_!\mathscr F(\Psi\mathrm Q_{\mathrm S})\) est constant, i.e. isomorphe au produit avec \(\mathrm S\) d’un espace muni d’un faisceau. C’est clair sur \(\mathrm S\times\mathrm A^n\), où \(j_!\mathscr F(\Psi\mathrm Q_{\mathrm S})\) est localement constant, et près d’un point \(h\) de \(\mathrm H\) où \(\mathrm Q_{\mathrm S d}\) ne s’annule pas (resp. s’annule), on peut trouver des coordonnées locales \(t\) (un \(\mathrm S\)-morphisme étale d’un voisinage étale de \(h\) dans \(\mathrm S\times\mathrm A^n\)) dans lesquelles \(\mathrm Q_{\mathrm S}=t_1^{-d}\) (resp. \(\mathrm Q_{\mathrm S}=t_2\cdot t_1^{-d}\)). Le couple \((\mathrm S\times\mathbf P^n,j_!\mathscr F(\Psi\mathrm Q_{\mathrm S}))\) est donc localement acyclique sur \(\mathrm S_0\) (SGA \(4\frac{1}{2}\) [Th. finitude], 2.16) et les :
\[
\mathrm R^if_!\mathscr F_0(\Psi\mathrm Q_{\mathrm S})
=
\mathrm R^i\bar f_!j_!\mathscr F_0(\Psi\mathrm Q_{\mathrm S})
=
\mathrm R^i\bar f_*j_!\mathscr F_0(\Psi\mathrm Q_{\mathrm S})
\]
sont donc lisses (\textit{loc. cit.}, A2).

\textbf{(3.7.4)} Le schéma \(\mathrm S_0\) est connexe, et les faisceaux \(\mathrm R^if_!\mathscr F_0(\Psi\mathrm Q_{\mathrm S})\) sont mixtes. Pour prouver (3.7.2.3), il suffit donc, d’après (1.8.12), de le faire pour un polynôme \(\mathrm Q\) particulier. On prend \(\mathrm Q=\sum\mathrm X_i^d\). Pour ce choix, le faisceau \(\mathscr F_0(\Psi\mathrm Q)\) est le produit tensoriel des images réciproques de faisceaux analogues sur les \(n\) facteurs \(\mathrm A_0^1\) de \(\mathrm A_0^n\), et la formule de Künneth ramène (3.7.2.3) au cas où \(n=1\) et \(\mathrm Q=\mathrm X^d\). Pour ce cas, je renvoie à (I.8.11).

% ---- 217-224 ----
\Page{217}
\section*{IV. — PINCEAUX DE LEFSCHETZ}

Dans ce paragraphe, sauf au n\textsuperscript{o}~5, nous travaillerons au-dessus d’un corps algébriquement clos \(k\), de caractéristique \(p\). Soit \(\ell\) un nombre premier distinct de \(p\). Les groupes de cohomologie considérés seront à coefficients dans \(\mathbf{Z}_\ell\), \(\mathbf{Q}_\ell\) ou \(\mathbf{Z}/m\), pour \(m\) une puissance de \(\ell\).

\section*{(4.1) Le théorème de Lefschetz difficile}

Dans ce numéro, nous supposerons choisi un isomorphisme, sur \(k\), entre \(\mathbf{Z}_\ell\) et \(\mathbf{Z}_\ell(1)\). Ceci nous dispensera d’écrire les twists à la Tate.

\textit{Théorème} \textbf{(4.1.1)}. — \textit{Soient \(\mathrm{X}\) une variété projective et lisse sur \(k\), \(\mathscr{L}\) un faisceau inversible ample sur \(\mathrm{X}\) et \(\eta=c_1(\mathscr{L})\in\mathrm{H}^2(\mathrm{X},\mathbf{Q}_\ell)\). On suppose \(\mathrm{X}\) purement de dimension \(n\). Alors, pour chaque \(i\geq0\), le cup produit par \(\eta^i\) : \(\eta^i\cup:\mathrm{H}^{n-i}(\mathrm{X},\mathbf{Q}_\ell)\to\mathrm{H}^{n+i}(\mathrm{X},\mathbf{Q}_\ell)\) est un isomorphisme.}

Prouvons le théorème par récurrence sur \(n\). Le cas \(n=0\) est trivial, et on suppose que \(n\geq1\). Puisque \(c_1(\mathscr{L}^{\otimes m})=mc_1(\mathscr{L})\), l’assertion du théorème pour \(\mathscr{L}\) équivaut à la même assertion pour ses puissances tensorielles. Ceci nous permet de supposer \(\mathscr{L}\) très ample. Soit \(\mathrm{Y}\) une section hyperplane lisse de \(\mathrm{X}\) (dans le plongement projectif défini par \(\mathscr{L}\)). Nous allons tout d’abord déduire de l’hypothèse de récurrence que le théorème, pour \(\mathrm{X}\) et \(\mathscr{L}\), équivaut au

\textit{Lemme} \textbf{(4.1.2)} (cf. SGA7 XVIII (5.2.4)). — \textit{La forme d’intersection sur \(\mathrm{H}^{n-1}(\mathrm{Y})\) est non dégénérée sur l’image de \(\mathrm{H}^{n-1}(\mathrm{X})\).}

Pour \(i=0\), le théorème est trivial. Pour \(i\geq1\), on interprète le cup-produit par \(\eta=c\ell(\mathrm{Y})\) comme le composé du morphisme de restriction \(\mathrm{H}^*(\mathrm{X})\to\mathrm{H}^*(\mathrm{Y})\), et de son transposé par dualité de Poincaré \(\mathrm{H}^*(\mathrm{Y})\to\mathrm{H}^{*+2}(\mathrm{X})\) (SGA \(4\,\frac12\)), pour factoriser \(\eta^i\cup:\mathrm{H}^{n-i}(\mathrm{X})\to\mathrm{H}^{n+i}(\mathrm{X})\) en :
\[
\eta^i\cup:\mathrm{H}^{n-i}(\mathrm{X})
\to\mathrm{H}^{(n-1)-(i-1)}(\mathrm{Y})
\xrightarrow{\eta^{i-1}\cup}
\mathrm{H}^{(n-1)+(i-1)}(\mathrm{Y})
\to\mathrm{H}^{n+i}(\mathrm{X}).
\]
Pour \(i\geq2\), les flèches extrêmes sont des isomorphismes, par le théorème de Lefschetz faible rappelé ci-dessous, et la flèche médiane est un isomorphisme d’après l’hypothèse de récurrence, d’où le théorème. Pour \(i=1\), le cas crucial, Lefschetz faible dit seulement que \(\mathrm{H}^{n-1}(\mathrm{X})\) s’injecte dans \(\mathrm{H}^{n-1}(\mathrm{Y})\). Son transposé \(\mathrm{H}^{n-1}(\mathrm{Y})\to\mathrm{H}^{n+1}(\mathrm{X})\) est donc surjectif, et pour que le composé soit bijectif, il faut et il suffit que l’énoncé du lemme soit
\Page{218}vrai. Remplaçant \(\mathscr{L}\) par une de ses puissances tensorielles, on peut supposer que, dans le plongement projectif défini par \(\mathscr{L}\), \(\mathrm{X}\) admet un pinceau de Lefschetz de sections hyperplanes (SGA7 XVII (2.5)). Prenons \(\mathrm{Y}\) dans un tel pinceau \((\mathrm{Y}_t)_{t\in\mathbf{P}^1}:\mathrm{Y}=\mathrm{Y}_{t_0}\). Nous allons prouver (4.1.2), pour \(\mathrm{Y}\), en interprétant l’image de \(\mathrm{H}^{n-1}(\mathrm{X})\) dans \(\mathrm{H}^{n-1}(\mathrm{Y})\) en terme de monodromie.

Soit \(\mathrm{S}\) l’ensemble des valeurs de \(t\) pour lesquelles \(\mathrm{Y}_t\) est singulière. Les \(\mathrm{H}^i(\mathrm{Y}_t)\) (\(t\in\mathbf{P}^1-\mathrm{S}\)) sont alors les fibres d’un \(\mathbf{Q}_\ell\)-faisceau lisse sur \(\mathbf{P}^1-\mathrm{S}\), et \(\pi_1(\mathbf{P}^1-\mathrm{S},t_0)\) agit sur \(\mathrm{H}^i(\mathrm{Y})\). On a

\textbf{(4.1.3)} L’image de \(\mathrm{H}^{n-1}(\mathrm{X})\) dans \(\mathrm{H}^{n-1}(\mathrm{Y})\) est le sous-espace \(\mathrm{H}^{n-1}(\mathrm{Y})^{\pi_1}\) des invariants par la monodromie.

Ce résultat est prouvé dans SGA7 XVIII (5.6.10) (sous l’hypothèse que \(\mathrm{Y}\) vérifie le théorème — notre hypothèse de récurrence). Sur \(\mathbf{C}\), et en cohomologie entière, une démonstration directe est donnée dans SGA7 XIX (4) (y conjuguer (4.3.1) et la formule de Picard-Lefschetz). Au n\textsuperscript{o}~(4.3), nous transposerons cette dernière au cas de la \(\mathbf{Z}_\ell\)-cohomologie, sur \(k\) quelconque.

Nous savons par ailleurs (3.4.13) que \(\mathrm{H}^{n-1}(\mathrm{Y})\) est une représentation complètement réductible de \(\pi_1(\mathbf{P}^1-\mathrm{S},t_0)\) et il ne reste qu’à appliquer le lemme algébrique suivant.

\textit{Lemme} \textbf{(4.1.4)}. — \textit{Soit \(\mathrm{V}\) une représentation linéaire complètement réductible d’un groupe \(\pi\), munie d’une forme bilinéaire \(\Phi\), invariante par \(\pi\) et non dégénérée. Alors, la restriction de \(\Phi\) à \(\mathrm{V}^{\pi}\) est non dégénérée.}

Par complète réductibilité, on a \(\mathrm{V}=\mathrm{V}^{\pi}\oplus\mathrm{W}\), pour une sous-représentation \(\mathrm{W}\) ne contenant pas la représentation triviale en quotient. Le sous-espace \(\mathrm{W}\) est donc \(\Phi\)-orthogonal à \(\mathrm{V}^{\pi}\) : \(\Phi=(\Phi\mid\mathrm{V}^{\pi})\oplus(\Phi\mid\mathrm{W})\), et \(\Phi\mid\mathrm{V}^{\pi}\) est non dégénérée.

\textit{Corollaire} \textbf{(4.1.5)}. — \textit{Les nombres de Betti d’indice impair d’une variété projective non singulière \(\mathrm{X}\) sont pairs.}

On se ramène à supposer la variété connexe, purement de dimension \(n\). Si \(\eta=c_1(\mathscr{L})\), avec \(\mathscr{L}\) ample, la forme \(\operatorname{Tr}(\eta^ixy)\) sur \(\mathrm{H}^{n-i}(\mathrm{X})\) est non dégénérée, d’après (4.1.1) et la dualité de Poincaré. Si \(n-i\) est impair, elle est alternée, de sorte que \(b_{n-i}=b_{n+i}\) est pair.

Ce résultat devrait valoir sans hypothèse de projectivité, mais je ne sais pas le prouver en caractéristique \(\neq0\).

\textbf{(4.1.6)} \textit{Lefschetz faible.} — Soient \(\mathrm{X}\) et \(\mathrm{Y}\) comme ci-dessus, à cela près qu’on suppose \(\mathrm{X}\) purement de dimension \(n+1\) (donc \(\mathrm{Y}\) purement de dimension \(n\)). Le résultat clef est celui qui donne la dimension cohomologique des variétés affines (SGA4 XIV (3.2)). Il donne, en \(\mathbf{Z}/m\)-cohomologie, que :
\[
\mathrm{H}^i(\mathrm{X}-\mathrm{Y})=0
\qquad\text{pour}\qquad i>n+1.
\]

\Page{219}
Puisque \(\mathrm{X}-\mathrm{Y}\) est lisse, on a aussi, par dualité de Poincaré :
\[
\mathrm{H}_c^i(\mathrm{X}-\mathrm{Y})=0
\qquad\text{pour}\qquad i\leq n.
\]
La suite exacte longue de cohomologie relative :
\[
\xrightarrow{\partial}\mathrm{H}_c^i(\mathrm{X}-\mathrm{Y})
\to\mathrm{H}^i(\mathrm{X})
\to\mathrm{H}^i(\mathrm{Y})\xrightarrow{\partial}
\]
permet d’en déduire que :
\[
\mathrm{H}^i(\mathrm{X})\xrightarrow{\sim}\mathrm{H}^i(\mathrm{Y})
\qquad\text{pour}\qquad i<n=\dim(\mathrm{Y})
\qquad\text{et que}\qquad
\mathrm{H}^n(\mathrm{X})\hookrightarrow\mathrm{H}^n(\mathrm{Y}).
\]
Les morphismes de Gysin, transposés par dualité de Poincaré des morphismes de restriction, sont donc tels que :
\[
\mathrm{H}^i(\mathrm{Y})\xrightarrow{\sim}\mathrm{H}^{i+2}(\mathrm{X})
\qquad\text{pour}\qquad i>n
\qquad\text{et}\qquad
\mathrm{H}^n(\mathrm{Y})\twoheadrightarrow\mathrm{H}^{n+2}(\mathrm{X}).
\]

Par passage à la limite, ces résultats restent vrais en \(\mathbf{Z}_\ell\)- et \(\mathbf{Q}_\ell\)-cohomologie. En \(\mathbf{Z}_\ell\)-cohomologie, la formule des coefficients universels :
\[
\mathrm{H}_c^i(\mathrm{X}-\mathrm{Y},\mathbf{Z}/\ell)
=
\mathrm{H}_c^i(\mathrm{X}-\mathrm{Y},\mathbf{Z}_\ell)
\otimes\mathbf{Z}/\ell
\oplus
\operatorname{Tor}_1^{\mathbf{Z}_\ell}
\bigl(\mathrm{H}_c^{i+1}(\mathrm{X}-\mathrm{Y},\mathbf{Z}_\ell),\mathbf{Z}/\ell\bigr)
\]
fournit le renseignement supplémentaire que \(\mathrm{H}_c^{n+1}(\mathrm{X}-\mathrm{Y},\mathbf{Z}_\ell)\) est sans torsion. Ceci implique que \(\mathrm{H}^n(\mathrm{Y},\mathbf{Z}_\ell)/\mathrm{H}^n(\mathrm{X},\mathbf{Z}_\ell)\) est sans torsion.

\section*{(4.2) Pinceaux de Lefschetz}

Dans (1.3), (1.4), je résumais les résultats essentiels de la théorie des pinceaux de Lefschetz, en excluant le cas où \(p=2\) et où les sections hyperplanes sont de dimension paire. Le but de ce numéro est de réparer en partie cette omission et, ce faisant, de compléter SGA7 XVIII : dans le cas sauvage, la théorie vaut pour des pinceaux de Lefschetz généraux, et non seulement pour ceux qui sont génériques.

\textbf{(4.2.1)} \textit{Théorie locale.} — Soit \(f:\mathrm{X}\to\mathrm{S}\) comme en (1.4.2), de dimension relative \(n\). On pose \(n=2n'\) ou \(2n'+1\). Soit \(\delta\in\mathrm{H}^n(\mathrm{X}_{\bar\eta})(n')\) le cycle évanescent. Il est bien défini au signe près. Dans le cas où \(n\) est pair et où \(p=2\), exclu dans (1.4.3), le groupe d’inertie \(\mathrm{I}\) admet plusieurs caractères d’ordre 2. Pour l’un d’eux, dépendant de la singularité quadratique ordinaire de \(f\), la monodromie locale est encore donnée par la formule de \textit{loc. cit.} Tous ces résultats valent aussi bien en \(\mathbf{Z}/m\)-, en \(\mathbf{Z}_\ell\)- ou en \(\mathbf{Q}_\ell\)-cohomologie. La preuve, dans le cas de la \(\mathbf{Z}/m\)-cohomologie, dont découlent les autres, est dans SGA7. Notons le

\textit{Corollaire} \textbf{(4.2.2)}. — \textit{L’image de \(\mathrm{H}^n(\mathrm{X}_s)\) dans \(\mathrm{H}^n(\mathrm{X}_{\bar\eta})\) est l’orthogonal de \(\delta\). En \(\mathbf{Z}_\ell\)- ou \(\mathbf{Q}_\ell\)-cohomologie, c’est aussi le sous-espace des invariants par le groupe d’inertie \(\mathrm{I}\).}

La seconde assertion résulte de ce que l’expression \((x,\delta)\delta\) qui apparaît dans la formule de Picard-Lefschetz donnant l’action de \(\mathrm{I}\) est nulle si et seulement si \((x,\delta)=0\) : c’est clair si \(\delta\) n’est pas de torsion, et si \(\delta\) est de torsion, \((x,\delta)\) est identiquement nul.

\Page{220}
\textbf{(4.2.3)} \textit{Théorie globale.} — On reprend les notations de (1.5.6); en particulier \(\mathrm{X}\subset\mathbf{P}\) est une variété projective non singulière connexe et de dimension \(n+1\), \((\mathrm{X}_t)_{t\in\mathrm{D}}\) est un pinceau de Lefschetz de sections hyperplanes de \(\mathrm{X}\), \(\mathrm{S}\) est l’ensemble des valeurs exceptionnelles \(t\in\mathrm{D}\) pour lesquelles la section hyperplane \(\mathrm{X}_t\) est singulière, et \(j:\mathrm{U}\to\mathrm{D}\) l’inclusion de l’ouvert complémentaire. On pose \(n=2n'\) ou \(2n'+1\). Identifions \(\mathrm{D}\) à une droite de l’espace projectif dual \(\mathbf{P}^{\vee}\), et notons \(\mathrm{X}^{\vee}\subset\mathbf{P}^{\vee}\) l’image de \(\mathrm{X}\) par la correspondance « l’hyperplan \(\mathrm{H}\) est tangent à \(\mathrm{X}\) en \(x\) ». Un pinceau de Lefschetz est dit \textit{transverse} si \(\mathrm{D}\) est transverse à \(\mathrm{X}^{\vee}\) (en particulier, \(\mathrm{D}\cap\mathrm{X}^{\vee}=\varnothing\) si \(\mathrm{X}^{\vee}\) n’est pas une hypersurface). Dans le cas d’exception, un pinceau de Lefschetz n’est pas automatiquement transverse.

Rappelons (cf. SGA1 V (5.7)) qu’un \textit{chemin} entre deux points géométriques \(\bar a\) et \(\bar b\) d’un schéma \(\mathrm{S}\) est un isomorphisme entre les foncteurs « fibre en \(\bar a\) » et « fibre en \(\bar b\) », de la catégorie des revêtements étales de \(\mathrm{S}\) dans (Ens). Pour tout chemin \(c\), tracé sur \(\mathrm{U}\), de \(u\) à un point générique géométrique \(\bar\eta_t\) de \(\mathrm{D}_{(t)}\), on a \(\mathrm{H}^n(\mathrm{X}_u,\mathbf{Z}_\ell)\xrightarrow{\sim}\mathrm{H}^n(\mathrm{X}_{\bar\eta_t},\mathbf{Z}_\ell)\), d’où un cycle évanescent \(\pm\delta_c\) dans \(\mathrm{H}^n(\mathrm{X}_u,\mathbf{Z}_\ell)(n')\).

\textit{Définition} \textbf{(4.2.4)} (i) \textit{L’ensemble des \textup{cycles évanescents} est l’ensemble de ces cycles \(\pm\delta_c\) pour \(c\) et \(t\in\mathrm{S}\) variables.}

(ii) \textit{La \textup{partie évanescente} de \(\mathrm{H}^n(\mathrm{X}_u)\) est le sous-espace engendré par les \(\lambda\delta\) (\(\lambda\in\mathbf{Z}_\ell(-n')\), \(\delta\) évanescent).}

\textbf{(4.2.5)} Puisque la droite projective \(\mathrm{D}\) est simplement connexe, le groupe fondamental \(\pi_1(\mathrm{U},u)\) est topologiquement engendré par les conjugués des groupes d’inertie \(\mathrm{I}_t\) pour \(t\in\mathrm{S}\). Le \textit{groupe de monodromie}, image de \(\pi_1(\mathrm{U},u)\) dans \(\operatorname{GL}(\mathrm{H}^n(\mathrm{X}_u,\mathbf{Z}_\ell))\), est donc topologiquement engendré par les images des conjugués des \(\mathrm{I}_t\) :
\begin{quote}
\(n\) impair : par les transvections \(x\mapsto x+\lambda(x\delta)\delta\), \(\lambda\in\mathbf{Z}_\ell(1)\), le quotient \(\ell\)-primaire du groupe d’inertie modéré;

\(n\) pair : par les réflexions \(x\mapsto x-(-1)^{n'}(x\delta)\delta\) (on a \((\delta,\delta)=(-1)^{n'}\cdot2\)).
\end{quote}

Il devrait résulter d’une théorie — non écrite — des cycles évanescents pour une base de dimension \(>1\) que, pour un pinceau de Lefschetz transverse, les cycles évanescents pris au signe près \(\pm\delta\) sont tous conjugués sous le groupe de monodromie et que la géométrie de la situation (en particulier, le groupe de monodromie) ne dépend pas du pinceau de Lefschetz transverse choisi. Nous nous contenterons de résultats plus faibles (cf. SGA7 XVIII (6)) :

\textit{Proposition} \textbf{(4.2.6)}. — \textit{Pour un pinceau de Lefschetz générique, les cycles évanescents, pris au signe près, sont tous conjugués sous le groupe de monodromie.}

Si \(\mathrm{X}^{\vee}\) n’est pas une hypersurface de \(\mathbf{P}^{\vee}\), alors, pour un pinceau général, on a \(\mathrm{S}=\varnothing\). L’assertion est triviale dans ce cas. Supposons donc que \(\mathrm{X}^{\vee}\) soit une hypersurface. Puisque \(\mathrm{X}\) est supposé irréductible, \(\mathrm{X}^{\vee}\) est irréductible.

\Page{221}
Soit \(\eta\) le point générique de la grassmannienne des droites dans \(\mathbf{P}^{\vee}\), \(\bar\eta\) un point générique géométrique, \(\mathrm{D}_\eta\) la droite projective sur \(k(\eta)\) correspondante, et \(\mathrm{D}_{\bar\eta}\) celle qui s’en déduit par extension des scalaires à \(k(\bar\eta)\). Sur \(\mathrm{D}_\eta\), le lieu critique \(\mathrm{S}_\eta\) est l’intersection de \(\mathrm{D}_\eta\) avec \(\mathrm{X}^{\vee}\). Le schéma \(\mathrm{S}_\eta\) s’envoie sur le point générique de \(\mathrm{X}^{\vee}\), avec une fibre géométrique connexe. Le schéma \(\mathrm{S}_\eta\) est donc \textit{connexe}.

Posons \(\mathrm{U}_\eta=\mathrm{D}_\eta-\mathrm{S}_\eta\) et \(\mathrm{U}_{\bar\eta}=\mathrm{D}_{\bar\eta}-\mathrm{S}_{\bar\eta}\). Pour \(u\in\mathrm{U}_{\bar\eta}\), considérons les groupes fondamentaux :
\[
\pi_1(\mathrm{U}_{\bar\eta},u)
\to\pi_1(\mathrm{U}_\eta,u)
\to\pi_1(\mathbf{P}^{\vee}-\mathrm{X}^{\vee},u).
\]
Le théorème de Bertini assure que l’application composée est surjective. Par ailleurs, l’action de \(\pi_1(\mathrm{U}_{\bar\eta},u)\) sur la cohomologie de la section hyperplane \(\mathrm{X}_u\) correspondant à \(u\) se factorise par \(\pi_1(\mathbf{P}^{\vee}-\mathrm{X}^{\vee},u)\). Pour prouver la conjugaison des cycles évanescents sous \(\pi_1(\mathrm{U}_{\bar\eta},u)\) il suffit donc de prouver une conjugaison sous \(\pi_1(\mathrm{U}_\eta,u)\). Ce groupe s’envoie sur \(\operatorname{Gal}(\bar\eta/\eta)\), et la conjugaison résulte par transport de structure de ce que \(\operatorname{Gal}(\bar\eta/\eta)\) agit transitivement sur \(\mathrm{S}_{\bar\eta}\) (connexité de \(\mathrm{S}_\eta\)).

\textit{Corollaire} \textbf{(4.2.7)}. — \textit{Pour \(p=2\), \(n\) pair, on suppose que le pinceau de Lefschetz \((\mathrm{X}_t)_{t\in\mathrm{D}}\) est assez général (\(=\) d’axe dans un ouvert convenable de la grassmannienne). Alors :}

(i) \textit{\(n\) impair : les homomorphismes de \(\mathbf{Z}_\ell(1)\) dans le groupe de monodromie, donnés par les formules de Picard-Lefschetz \(x\mapsto x+\lambda(x\delta)\delta\), pour \(\delta\) un cycle évanescent, sont conjugués entre eux;}

(ii) \textit{\(n\) pair : les réflexions \(x\mapsto x-(-1)^{n'}(x\delta)\delta\), pour \(\delta\) un cycle évanescent, sont conjuguées entre elles.}

Pour un pinceau générique, ce corollaire résulte de (4.2.6) et des formules de Picard-Lefschetz donnant l’action des groupes d’inertie locaux. On passe de là au cas général grâce à (1.11).

\textit{Corollaire} \textbf{(4.2.8)}. — \textit{Sous les hypothèses de \((4.2.7)\) les cycles évanescents pris au signe près sont tous conjugués dans \(\mathrm{H}^n(\mathrm{X}_u,\mathbf{Q}_\ell)\).}

Si l’on néglige la torsion, \(\pm\delta\) est en effet déterminé, au signe près, par la transformation de Picard-Lefschetz correspondante.

\section*{(4.3) Complément à SGA7 XIX (4)}

\textit{Nous supposons choisi un isomorphisme, sur \(k\), entre \(\mathbf{Z}_\ell\) et \(\mathbf{Z}_\ell(1)\). Ceci nous dispensera d’écrire les twists à la Tate.}

Dans ce numéro, nous donnons une démonstration directe de (4.1.3), généralisant SGA7 XIX (4) au cas d’un corps de base quelconque, pour la \(\mathbf{Z}_\ell\)-cohomologie. La démonstration est une transposition de celle de \textit{loc. cit.} que, bien que ce ne soit pas logiquement nécessaire, le lecteur est invité à lire au préalable.

\Page{222}
\textbf{(4.3.1)} Soient \(\mathbf{P}\) un espace projectif sur \(k\), \(\mathrm{X}\subset\mathbf{P}\) projectif et lisse, purement de dimension \(n+1\), \(\mathrm{A}\) un sous-espace de \(\mathbf{P}\) de codimension 2, qui coupe \(\mathrm{X}\) transversalement, et supposons que les hyperplans par \(\mathrm{A}\) balayent sur \(\mathrm{X}\) un pinceau de Lefschetz de sections hyperplanes (SGA7 XVII (2.2)). Soient \(\Delta=\mathrm{X}\cap\mathrm{A}\), \(\widetilde{\mathrm{X}}\) déduit de \(\mathrm{X}\) en éclatant \(\Delta\), et \(\widetilde{\Delta}\sim\Delta\times\mathbf{P}^1\) le diviseur exceptionnel. On sait que les sections hyperplanes du pinceau sont les fibres d’un morphisme \(f:\widetilde{\mathrm{X}}\to\mathbf{P}^1\) (SGA7 XVIII (4.1.1)). Soient \(\mathrm{S}\) l’ensemble des valeurs critiques de \(f\), et \(0\notin\mathrm{S}\) un point rationnel de \(\mathbf{P}^1\). On note \(\mathrm{Y}_t\) la fibre de \(f\) en \(t\), et l’on pose \(\mathrm{Y}=\mathrm{Y}_0\) :
\[
\begin{tikzcd}[column sep=large,row sep=large]
\Delta \arrow[r,hook] \arrow[d,hook]
  & \widetilde{\Delta} \arrow[d,hook] & {} \\
\mathrm{Y} \arrow[r,hook] \arrow[d]
  & \widetilde{\mathrm{X}} \arrow[r] \arrow[d,"f"]
  & \mathrm{X} \\
0 \arrow[r,hook] & \mathbf{P}^1 & {}
\end{tikzcd}
\tag*{\textbf{(4.3.1.1)}}
\]

La partie évanescente \(\operatorname{Ev}(\mathrm{Y})\) de \(\mathrm{H}^n(\mathrm{Y})\) a été définie en (4.2.4). Nous noterons \(\operatorname{Ev}(\mathrm{Y})^\perp\) son orthogonal pour la forme \((x,y)=\operatorname{Tr}(x\cup y)\).

\textit{Lemme} \textbf{(4.3.2)}. — \textit{En \(\mathbf{Z}/m\)-cohomologie, on a \(\mathrm{H}^0(\mathbf{P}^1,\mathrm{R}^n f_!\mathbf{Z}/m)\xrightarrow{\sim}\operatorname{Ev}(\mathrm{Y})^\perp\subset\mathrm{H}^n(\mathrm{Y})\). En \(\mathbf{Z}_\ell\)- (ou \(\mathbf{Q}_\ell\)-) cohomologie, ces groupes coïncident encore avec \(\mathrm{H}^n(\mathrm{Y})^{\pi_1(\mathbf{P}^1-\mathrm{S},0)}\).}

Le faisceau \(\mathrm{R}^n f_!\mathbf{Z}/m\) est sans section à support ponctuel (injectivité des morphismes de spécialisation). Son \(\mathrm{H}^0\) s’injecte donc dans \(\mathrm{H}^n(\mathrm{Y})\), et l’image est la fibre en \(0\) du plus grand sous-faisceau constant de \(\mathrm{R}^n f_!\mathbf{Z}/m\). La théorie locale montre par ailleurs que \(\operatorname{Ev}(\mathrm{Y})^\perp\) est la fibre en \(0\) du plus grand sous-faisceau localement constant de \(\mathrm{R}^n f_!\mathbf{Z}/m\). Puisque \(\mathbf{P}^1\) est simplement connexe, il est constant, d’où l’assertion. Enfin, en \(\mathbf{Z}_\ell\)-cohomologie, la théorie locale montre que \(\mathrm{R}^n f_!\mathbf{Z}_\ell\) est l’image directe de sa restriction à \(\mathbf{P}^1-\mathrm{S}\), d’où \(\mathrm{H}^0(\mathbf{P}^1,\mathrm{R}^n f_!\mathbf{Z}_\ell)=\mathrm{H}^0(\mathbf{P}^1-\mathrm{S},\mathrm{R}^n f_!\mathbf{Z}_\ell)=\mathrm{H}^n(\mathrm{Y})^{\pi_1(\mathbf{P}^1-\mathrm{S},0)}\).

Rappelons que, par Lefschetz faible, \(\mathrm{H}^n(\mathrm{X})\) s’injecte dans \(\mathrm{H}^n(\mathrm{Y})\). On a la

\textit{Proposition} \textbf{(4.3.3)}. — \textit{\(\operatorname{Ev}(\mathrm{Y})^\perp\) est encore l’image de \(\mathrm{H}^n(\mathrm{X})\) dans \(\mathrm{H}^n(\mathrm{Y})\).}

La démonstration, pour la \(\mathbf{Z}/m\)-cohomologie, sera donnée en (4.3.7). Auparavant, trois préliminaires : une réduction au cas de la \(\mathbf{Z}/m\)-cohomologie, ci-dessous; le formalisme de la cohomologie relative, en (4.3.4); un raisonnement qui se substituera à la théorie de Morse, pour une boule qui grossit dans \(\mathbf{P}^1(\mathbf{C})\), en (4.3.5).

On peut écrire (4.3.3) comme une suite exacte \(\mathrm{H}^n(\mathrm{X})\to\mathrm{H}^n(\mathrm{Y})\to\operatorname{Ev}(\mathrm{Y})^*\). Pour les \(\mathbf{Z}/m\)-modules, le passage au dual est exact; le transposé par dualité de Poincaré du
\Page{223}morphisme de restriction est un morphisme de Gysin et, en \(\mathbf{Z}/m\)-cohomologie, (4.3.3) fournit une croix de suites exactes :
\[
\begin{tikzcd}[column sep=large,row sep=large]
{} & {} & \operatorname{Ev}(\mathrm{Y})^* & {} & {} \\
0 \arrow[r]
  & \operatorname{Ev}(\mathrm{Y}) \arrow[r]
  & \mathrm{H}^n(\mathrm{Y}) \arrow[r] \arrow[u,two heads]
  & \mathrm{H}^{n+2}(\mathrm{X}) \arrow[r]
  & 0 \\
{} & {}
  & \mathrm{H}^n(\mathrm{X}) \arrow[u,hook] \arrow[ur,"\eta\cup"']
  & {} & {}
\end{tikzcd}
\tag*{\textbf{(4.3.3.1)}}
\]
où le composé \(\mathrm{H}^n(\mathrm{X})\to\mathrm{H}^n(\mathrm{Y})\to\mathrm{H}^{n+2}(\mathrm{X})\) est le cup-produit avec \(\eta=c\ell(\mathrm{Y})\). On ne peut en déduire la même croix en \(\mathbf{Z}_\ell\)-cohomologie par passage à la limite projective, les \(\operatorname{Ev}(\mathrm{Y})^*\), pris en \(\mathbf{Z}/\ell^j\)-cohomologie, ne formant pas toujours un système projectif. Toutefois, les \(\operatorname{Ev}(\mathrm{Y})\) s’envoient les uns \textit{sur} les autres, et si les \(\delta_i\) (\(i\in\mathrm{B}\)) engendrent \(\operatorname{Ev}(\mathrm{Y})\) en cohomologie \(\ell\)-adique, on peut passer à la limite dans la suite exacte :
\[
0\to\mathrm{H}^n(\mathrm{X})
\to\mathrm{H}^n(\mathrm{Y})
\to(\mathbf{Z}/\ell^j)^{\mathrm{B}}
\]
d’où en \(\mathbf{Z}_\ell\)-cohomologie une croix :
\[
\begin{tikzcd}[column sep=large,row sep=large]
{} & {} & \operatorname{Ev}(\mathrm{Y})^* & {} & {} \\
0 \arrow[r]
  & \operatorname{Ev}(\mathrm{Y}) \arrow[r]
  & \mathrm{H}^n(\mathrm{Y}) \arrow[r] \arrow[u]
  & \mathrm{H}^{n+2}(\mathrm{X}) \arrow[r]
  & 0. \\
{} & {}
  & \mathrm{H}^n(\mathrm{X}) \arrow[u,hook] \arrow[ur,"\eta\cup"']
  & {} & {}
\end{tikzcd}
\tag*{\textbf{(4.3.3.2)}}
\]
En \(\mathbf{Q}_\ell\)-cohomologie, on retrouve la surjectivité du morphisme \(\mathrm{H}^n(\mathrm{Y})\to\operatorname{Ev}(\mathrm{Y})^*\), transposé de l’inclusion de \(\operatorname{Ev}(\mathrm{Y})\) dans \(\mathrm{H}^n(\mathrm{Y})\).

\textbf{(4.3.4)} \textit{Cohomologie relative.} — Soit \(u:\mathrm{T}'\to\mathrm{T}''\) un morphisme de topos. Les triples (\(\mathscr{F}'\) sur \(\mathrm{T}'\), \(\mathscr{F}''\) sur \(\mathrm{T}''\), \(\varphi:u^*\mathscr{F}''\to\mathscr{F}'\)) forment alors un nouveau topos \(\mathrm{T}\).

Si \(\mathscr{F}=(\mathscr{F}',\mathscr{F}'',\varphi)\) est un faisceau abélien de \(\mathrm{T}\), on pose :
\[
\mathrm{H}^0(\mathrm{T}''\mathbin{\mathrm{mod}}\mathrm{T}',\mathscr{F})
=
\operatorname{Ker}
\bigl(\mathrm{H}^0(\varphi):
\mathrm{H}^0(\mathrm{T}'',\mathscr{F}'')
\to\mathrm{H}^0(\mathrm{T}',\mathscr{F}')\bigr).
\]

Par dérivation de ce foncteur, on obtient la \textit{cohomologie relative} de \(\mathrm{T}''\mathbin{\mathrm{mod}}\mathrm{T}'\). On peut encore le construire par dérivation du foncteur \(\mathscr{F}\mapsto(\mathrm{H}^0(\mathrm{T}'',\mathscr{F}'')\to\mathrm{H}^0(\mathrm{T}',\mathscr{F}'))\) des faisceaux abéliens sur \(\mathrm{T}\) dans les complexes de faisceaux réduits aux degrés 0 et 1. Cela revient au même, parce que tout faisceau \(\mathscr{F}\) s’injecte dans un faisceau \(\mathscr{F}_1\) pour lequel \(\mathrm{H}^0(\varphi)\) est surjectif : une somme de faisceaux \((0,\mathscr{F}'')\) et \((\mathscr{F}',u_*\mathscr{F}')\).

La seconde description montre que, pour \(\mathrm{K}\) dans \(\mathrm{D}^+(\mathrm{T})\), on a un triangle :
\[
\to\mathrm{R}\Gamma(\mathrm{T}''\mathbin{\mathrm{mod}}\mathrm{T}',\mathrm{K})
\to\mathrm{R}\Gamma(\mathrm{T}'',\mathrm{K}'')
\to\mathrm{R}\Gamma(\mathrm{T}',\mathrm{K}')
\to
\]
donnant lieu à une suite exacte longue de cohomologie.

\Page{224}
\textbf{(4.3.5)} \textit{Faisceaux sur \(\mathbf{P}^1\).} — Soient \(\mathrm{X}\) une variété irréductible complète sur \(k\), \(x,y\) deux points fermés de \(\mathrm{X}\), et \(\mathrm{K}\) dans \(\mathrm{D}^+(\mathrm{X})\). Soient \(\bar\eta_x\) et \(\bar\eta_y\) des points géométriques localisés en des points génériques \(\eta_x\) et \(\eta_y\) des hensélisés \(\mathrm{X}_{(x)}\) et \(\mathrm{X}_{(y)}\). On se propose d’associer, à un \(\mathrm{X}\)-isomorphisme \(\sigma\) de \(\bar\eta_x\) avec \(\bar\eta_y\), un morphisme
\[
\mathbf{H}_c^*(\mathrm{X}-\{x\},\mathrm{K})
=
\mathbf{H}^*(\mathrm{X}\mathbin{\mathrm{mod}}\{x\},\mathrm{K})
\to
\mathbf{H}^*(\mathrm{X}_{(y)}\mathbin{\mathrm{mod}}\bar\eta_y,\mathrm{K}).
\tag*{\textbf{(4.3.5.1)}}
\]
On le fait en suivant le diagramme d’espaces :
\[
\begin{tikzcd}[column sep=large,row sep=large]
\{x\} \arrow[r] \arrow[d]
  & \mathrm{X}_{(x)} \arrow[d]
  & \bar\eta_x \arrow[l] \arrow[r,leftrightarrow] \arrow[d]
  & \bar\eta_y \arrow[r,equal] \arrow[d]
  & \bar\eta_y \arrow[d] \\
\mathrm{X} \arrow[r,equal]
  & \mathrm{X} \arrow[r,equal]
  & \mathrm{X} \arrow[r,equal]
  & \mathrm{X} \arrow[r,leftarrow]
  & \mathrm{X}_{(y)}
\end{tikzcd}
\]
dont on peut remonter la première flèche parce que \(\mathbf{H}^*(\mathrm{X}_{(x)},\mathrm{K})\xrightarrow{\sim}\mathbf{H}^*(\{x\},\mathrm{K})\).

\textit{Proposition} \textbf{(4.3.6)}. — \textit{Soit \(t\in\mathbf{P}^1\), et soit \(\mathrm{K}\) dans \(\mathrm{D}^b(\mathbf{P}^1)\) tel que les \(\mathscr{H}^i\mathrm{K}\) soient constructibles et que le support de \(\mathscr{H}^i(\mathrm{K})\) soit fini si \(i<0\). Soit \(\mathrm{S}\) une partie finie de \(\mathbf{P}^1\) telle que les \(\mathscr{H}^i\mathrm{K}\) soient localement constants en dehors de \(\mathrm{S}\). Alors, le produit des morphismes \((4.3.5.1)\), pour \(\mathrm{X}=\mathbf{P}^1\), \(x=t\) et \(y\in\mathrm{S}\), est injectif :}
\[
\mathbf{H}^1(\mathbf{P}^1\mathbin{\mathrm{mod}}t,\mathrm{K})
\hookrightarrow
\prod\mathbf{H}^1(\mathbf{P}^1_{(s)}\mathbin{\mathrm{mod}}\bar\eta_s,\mathrm{K}).
\]

Dans l’application que nous ferons de (4.3.6), on aura \(t\notin\mathrm{S}\).

Soit \(\mathrm{K}'\) le sous-complexe de \(\mathrm{K}\) formé des \(\mathrm{K}^i\) pour \(i<0\), et de l’image inverse dans \(\operatorname{Ker}(d)\subset\mathrm{K}^0\) du plus grand sous-faisceau de \(\mathscr{H}^0(\mathrm{K})\) dont le support est fini. Les groupes d’hypercohomologie considérés satisfont à \(\mathrm{H}^i(\mathrm{K}')=0\) pour \(i>0\), donc \(\mathrm{H}^1(\mathrm{K})\xrightarrow{\sim}\mathrm{H}^1(\mathrm{K}/\mathrm{K}')\). Remplaçant \(\mathrm{K}\) par \(\mathrm{K}/\mathrm{K}'\), ceci nous ramène à supposer que \(\mathrm{K}^i=0\) pour \(i<0\), et que \(\mathscr{H}^0\mathrm{K}\) n’a pas de section à support ponctuel. On le suppose.

Les groupes de cohomologie considérés donnent lieu à une suite exacte :
\[
0\to\mathrm{H}^1(\mathscr{H}^0\mathrm{K})
\to\mathbf{H}^1(\mathrm{K})
\to\mathrm{H}^0(\mathscr{H}^1\mathrm{K});
\]
il suffit donc de vérifier que :
\[
\mathrm{H}^0(\mathbf{P}^1\mathbin{\mathrm{mod}}t,\mathscr{H}^1\mathrm{K})
\hookrightarrow
\prod\mathrm{H}^0(\mathbf{P}^1_{(s)}\mathbin{\mathrm{mod}}\bar\eta_{(s)},\mathscr{H}^1\mathrm{K})
\]
et
\[
\mathrm{H}^1(\mathbf{P}^1\mathbin{\mathrm{mod}}t,\mathscr{H}^0\mathrm{K})
\hookrightarrow
\prod\mathrm{H}^1(\mathbf{P}^1_{(s)}\mathbin{\mathrm{mod}}\bar\eta_{(s)},\mathscr{H}^0\mathrm{K}).
\]
Le cas du \(\mathrm{H}^0\) est laissé au lecteur (il suffirait, ici, à droite, d’un produit étendu à \(\mathrm{S}\)). Posons \(\mathscr{H}=\mathscr{H}^0\mathrm{K}\). Il nous reste à vérifier que, pour \(\mathscr{H}\) sans section à support ponctuel, et localement constant en dehors de \(\mathrm{S}\), on a :
\[
\mathrm{H}^1(\mathbf{P}^1\mathbin{\mathrm{mod}}t,\mathscr{H})
\hookrightarrow
\prod\mathrm{H}^1(\mathbf{P}^1_{(s)}\mathbin{\mathrm{mod}}\bar\eta_s,\mathscr{H}).
\]
Le groupe à gauche classifie les \(\mathscr{H}\)-torseurs \(\mathrm{P}\) sur \(\mathbf{P}^1\), trivialisés en \(t\), donc sur \(\mathbf{P}^1_{(t)}\). Ces torseurs n’ont pas d’automorphismes. Ils peuvent se construire comme suit, par étapes :

\noindent\textit{a)} sur \(\mathbf{P}^1_{(t)}\), c’est le torseur trivial, avec sa trivialisation naturelle;

\noindent\textit{b)} le prolonger à \(\mathbf{P}^1-(\mathrm{S}-\{t\})\) revient à se donner une action continue \(\ast\) de \(\pi_1(\mathbf{P}^1-\mathrm{S},\bar\eta_t)\) sur sa fibre \(\mathrm{P}_{\bar\eta_t}\sim\mathscr{H}_{\bar\eta_t}\) en \(\bar\eta_t\). Cette action doit être telle que \(\sigma\ast(h+p)=\sigma(h)+\sigma\ast p\);

% ---- 225-232 ----
\Page{225}
\noindent\textit{c)} le prolonger en outre au-dessus de \(s\in\mathrm{S}-t\) revient à prolonger à \(\mathbf{P}^1_{(s)}\) le torseur déjà donné sur \(\eta_s\).

Pour tout \(\sigma:\bar{\eta}_t\to\bar{\eta}_s\), le groupe de monodromie locale en \(s\), \(\operatorname{Gal}(\bar{\eta}_s/\eta_s)\) s’envoie dans \(\pi_1(\mathbf{P}^1-\mathrm{S},\bar{\eta}_t)\). L’image (4.3.5.1) de \(\mathrm{P}\) dans \(\mathbf{H}^1(\mathbf{P}^1_{(s)}\ \mathrm{mod}\ \bar{\eta},\mathscr{H})\) détermine la \(*\)-action du groupe image sur \(\mathrm{P}_{\bar{\eta}_t}\), ainsi que le prolongement du torseur de \(\eta_s\) à \(\mathbf{P}^1_{(s)}\). On conclut en notant que, parce que \(\mathbf{P}^1\) est simplement connexe, les images des groupes de monodromie locale engendrent topologiquement \(\pi_1(\mathbf{P}^1-\mathrm{S},\bar{\eta}_t)\).

\textbf{(4.3.7)} \textit{Plan de la preuve de} (4.3.3) (les notations sont celles de (4.3.1)).

L’injection de \(\mathrm{H}^n(\mathrm{X})\) dans \(\mathrm{H}^n(\mathrm{Y})\) se factorise par \(\mathrm{H}^0(\mathbf{P}^1,\mathrm{R}^nf_!\mathbf{Z}/m)\subset\mathrm{H}^n(\mathrm{Y})\), et il nous faut vérifier que si \(\alpha\) est dans cet \(\mathrm{H}^0\), l’obstruction \(d\alpha\in\mathrm{H}_c^{n+1}(\mathrm{X}-\mathrm{Y})\) à prolonger \(\alpha\) de \(\mathrm{Y}\) à \(\mathrm{X}\) est nulle. On a \(\mathrm{X}-\mathrm{Y}=\widetilde{\mathrm{X}}-\mathrm{Y}-\widetilde{\Delta}\), et si \(\varphi\) est la projection de \(\widetilde{\mathrm{X}}-\widetilde{\Delta}\) sur \(\mathbf{P}^1\), on a :
\[
\mathrm{H}_c^{n+1}(\mathrm{X}-\mathrm{Y})
=\mathbf{H}^1(\mathbf{P}^1\ \mathrm{mod}\ o,\mathrm{R}\varphi_!\mathbf{Z}/m[n]).
\]
Si \(t\notin\mathrm{S}\), \(\varphi^{-1}(t)\) est une variété affine lisse purement de dimension \(n\). Pour \(i<n\), son \(\mathrm{H}_c^i\) est donc nul (théorème de Lefschetz faible). Le support du faisceau \(\mathrm{R}^i\varphi_!\mathbf{Z}/m\) est donc fini (\(\subset\mathrm{S}\)) pour \(i<n\) (en fait, le faisceau est même nul, mais peu importe). Ceci permet d’appliquer (4.3.6), et nous ramène à calculer des obstructions locales, qui ne sont autres que les \((\alpha,\delta_\sigma)\).

\textbf{(4.3.8)} \textit{Preuve de} (4.3.3). — Soit \(\alpha\in\mathrm{H}^n(\mathrm{Y})\), et \(\widetilde{\alpha}\) son image réciproque sur l’image réciproque \(\mathrm{Y}\cup\widetilde{\Delta}\) de \(\mathrm{Y}\subset\mathrm{X}\) dans \(\widetilde{\mathrm{X}}\). L’obstruction \(d\alpha\in\mathrm{H}_c^{n+1}(\mathrm{X}-\mathrm{Y})\) à étendre \(\alpha\) à \(\mathrm{X}\) coïncide avec l’obstruction à étendre \(\widetilde{\alpha}\) à \(\widetilde{\mathrm{X}}\). Pour l’exprimer, nous utiliserons la version « catégorie dérivée » de la suite spectrale de Leray de \(f\). Soit donc \((\mathbf{Z}/m)^*\) une résolution flasque de \(\mathbf{Z}/m\) sur \(\widetilde{\mathrm{X}}\), \(\mathrm{K}=f_*(\mathbf{Z}/m)^*\) et \(\mathrm{L}=f_*((\mathbf{Z}/m)^*|\widetilde{\Delta})\). Les groupes de cohomologie qui nous intéressent s’expriment tous en terme du morphisme de complexes \(u:\mathrm{K}\to\mathrm{L}\) sur \(\mathbf{P}^1\). Ainsi, la suite exacte longue de cohomologie relative pour \(\widetilde{\mathrm{X}}\) et \(\mathrm{Y}\cup\widetilde{\Delta}\) provient de la suite exacte courte des complexes simples associés à des complexes multiples :
\[
\Gamma(\mathbf{P}^1,\mathrm{K})
\longleftarrow
\mathbf{s}\left(
\begin{tikzcd}[ampersand replacement=\&, column sep=1.5em, row sep=2.2em]
\Gamma(\mathbf{P}^1,\mathrm{K}) \arrow[r] \arrow[d]
  \& \Gamma(\mathbf{P}^1,\mathrm{L}) \arrow[d] \\
\mathrm{K}_o \arrow[r] \& \mathrm{L}_o
\end{tikzcd}
\right)
\longleftarrow
\mathbf{s}\left(
\begin{tikzcd}[ampersand replacement=\&, column sep=1.5em, row sep=2.2em]
0 \arrow[r] \arrow[d]
  \& \Gamma(\mathbf{P}^1,\mathrm{L}) \arrow[d] \\
\mathrm{K}_o \arrow[r] \& \mathrm{L}_o
\end{tikzcd}
\right)
\]
donnant respectivement la cohomologie de \(\widetilde{\mathrm{X}}\), celle, à support propre, de \(\mathrm{X}-\mathrm{Y}\), et celle de \(\mathrm{Y}\cup\widetilde{\Delta}\).

En particulier, l’obstruction à étendre \(\widetilde{\alpha}\) à \(\widetilde{\mathrm{X}}\) apparaît bien comme étant dans \(\mathbf{H}^1(\mathbf{P}^1\ \mathrm{mod}\ t,\mathbf{s}(\mathrm{K}\to\mathrm{L})[n])\), où \(\mathbf{s}(\mathrm{K}\to\mathrm{L})\) représente \(\mathrm{R}\varphi_!\mathbf{Z}/m\). Comme expliqué en (4.3.7), les hypothèses de la proposition (4.3.6) sont satisfaites. Appliquons-la. Si on suit la
\Page{226}
définition (4.3.5), on trouve pour l’image de \(\widetilde{\alpha}\) par (4.3.5.1) une nouvelle description comme obstruction : on a une suite exacte courte :
\[
\Gamma(\mathbf{P}^1_{(s)},\mathrm{K})
\longleftarrow
\mathbf{s}\left(
\begin{tikzcd}[ampersand replacement=\&, column sep=1.5em, row sep=2.2em]
\Gamma(\mathbf{P}^1_{(s)},\mathrm{K}) \arrow[r] \arrow[d]
  \& \Gamma(\mathbf{P}^1_{(s)},\mathrm{L}) \arrow[d] \\
\mathrm{K}_{\bar{\eta}_s} \arrow[r]
  \& \mathrm{L}_{\bar{\eta}_s}
\end{tikzcd}
\right)
\longleftarrow
\mathbf{s}\left(
\begin{tikzcd}[ampersand replacement=\&, column sep=1.5em, row sep=2.2em]
0 \arrow[r] \arrow[d]
  \& \Gamma(\mathbf{P}^1_{(s)},\mathrm{L}) \arrow[d] \\
\mathrm{K}_{\bar{\eta}_s} \arrow[r]
  \& \mathrm{L}_{\bar{\eta}_s}
\end{tikzcd}
\right)
\]
une classe \(\widetilde{\alpha}'\) à droite, qui induit dans \(\mathrm{H}^n(\mathrm{K}_{\bar{\eta}_s})=\mathrm{H}^n(\mathrm{Y}_{\bar{\eta}_s})\) une classe transportée de \(\alpha\in\mathrm{H}^n(\mathrm{Y})\), et on cherche l’obstruction à relever \(\widetilde{\alpha}'\) dans \(\mathrm{R}^n\Gamma(\mathbf{P}^1_{(s)},\mathrm{K})=\mathrm{H}^n(\widetilde{\mathrm{X}}_{\mathbf{P}^1_{(s)}},\mathbf{Z}/m)\). Puisque \(\widetilde{\Delta}\) est lisse sur \(\mathbf{P}^1\), les faisceaux \(\mathscr{H}^i(\mathrm{L})\) sont constants, et \(\mathrm{R}\Gamma(\mathbf{P}^1_{(s)},\mathrm{L})\to\mathrm{L}_{\bar{\eta}_s}\) est un quasi-isomorphisme. Ceci nous permet de remplacer la suite exacte courte ci-dessus par son image quasi-isomorphe :
\[
\Gamma(\mathbf{P}^1_{(s)},\mathrm{K})
\longleftarrow
\mathbf{s}\left(
\begin{tikzcd}[ampersand replacement=\&, row sep=1.8em]
\Gamma(\mathbf{P}^1_{(s)},\mathrm{K}) \arrow[d] \\
\mathrm{K}_{\bar{\eta}_s}
\end{tikzcd}
\right)
\longleftarrow
\mathbf{s}\left(
\begin{tikzcd}[ampersand replacement=\&, row sep=1.8em]
0 \arrow[d] \\
\mathrm{K}_{\bar{\eta}_s}
\end{tikzcd}
\right).
\]
L’obstruction cherchée devient celle à prolonger \(\alpha_\sigma\in\mathrm{H}^n(\mathrm{Y}_{\bar{\eta}_s})\) au-dessus de l’image inverse de \(\mathbf{P}^1_{(s)}\) dans \(\widetilde{\mathrm{X}}\). Ces obstructions sont nulles si \(\alpha\) vient de \(\mathrm{H}^0(\mathbf{P}^1,\mathrm{R}^nf_!\mathbf{Z}/m)\).

\textit{Corollaire} \textbf{(4.3.9)}. — \textit{En \(\mathbf{Q}_\ell\)-cohomologie, on a une décomposition en somme directe orthogonale (pour la forme d’intersection \(\operatorname{Tr}(x\cup y)\)) :}
\[
\mathrm{H}^n(\mathrm{Y})=\mathrm{H}^n(\mathrm{X})\oplus\operatorname{Ev}(\mathrm{Y}).
\]

Cette reformulation du théorème de Lefschetz difficile résulte de (4.3.3) et (4.1.2).

\textbf{(4.3.10)} En \(\mathbf{Z}_\ell\)-cohomologie, la croix (4.3.3.2) montre que l’intersection de la partie évanescente \(\operatorname{Ev}(\mathrm{Y})\) de \(\mathrm{H}^n(\mathrm{Y})\) et de la partie fixe \(\mathrm{H}^n(\mathrm{X})=\mathrm{H}^n(\mathrm{Y})^{\pi_1(\mathbf{P}^1-\mathrm{S},o)}\) est le sous-groupe :
\[
\operatorname{Ker}(\eta\cup:\mathrm{H}^n(\mathrm{X})\to\mathrm{H}^{n+2}(\mathrm{X})).
\]
En \(\mathbf{Q}_\ell\)-cohomologie, ce noyau est nul (c’est le théorème de Lefschetz difficile pour \(\mathrm{X}\)) et on retrouve ce que Lefschetz appelle son lemme fondamental : « un cycle à la fois invariant et évanescent est nul ». Toutefois, ainsi que J. Morgan me l’a fait remarquer, ce noyau peut être non nul en \(\mathbf{Z}_\ell\)-cohomologie (donc, sur \(\mathbf{C}\), en cohomologie entière) : si on remplace \(\eta\) par un multiple (en changeant de plongement projectif), on peut faire en sorte qu’il coïncide avec le sous-groupe de torsion de \(\mathrm{H}^n(\mathrm{X})\) tout entier. Le « lemme fondamental » de Lefschetz est donc faux en \(\mathbf{Z}_\ell\)-cohomologie (et, sur \(\mathbf{C}\), en cohomologie entière). Ceci, joint au fait que personne ne l’ait récemment comprise, permet de mettre en doute la démonstration que Lefschetz croyait en avoir.

\Page{227}
\section*{(4.4) La monodromie des pinceaux de Lefschetz}

Les notations sont celles de (4.3.1).

Travaillons en \(\mathbf{Q}_\ell\)-cohomologie, et soit \(\mathrm{G}\) le groupe des automorphismes de \(\operatorname{Ev}(\mathrm{Y})\), muni de la forme d’intersection \((,)\), non dégénérée d’après le théorème de Lefschetz difficile (4.3.9). Le groupe \(\pi_1(\mathbf{P}^1-\mathrm{S},o)\) agit sur \(\mathrm{Y}\) en respectant la forme \((,)\). Notons \(\mathrm{M}\) son image dans \(\mathrm{G}\).

\textit{Théorème} \textbf{(4.4.1)}. — \textit{On suppose \(\mathrm{X}\) irréductible. Si \(p=2\) et que \(n\) est pair, on suppose de plus le pinceau assez général (\(=\)d’axe dans un ouvert convenable d’une grassmannienne). Alors, de deux choses l’une :}
\par\noindent a) \textit{\(\mathrm{M}\) est ouvert dans \(\mathrm{G}\).}
\par\noindent b) \textit{\(n\) est pair, et \(\mathrm{M}\) est fini.}

Pour \(p=2\), et \(n\) pair, « assez général » signifie : tel que (4.2.8) soit applicable.

Le groupe \(\mathrm{G}\) est un groupe algébrique semi-simple sur \(\mathbf{Q}_\ell\). Pour que le sous-groupe compact \(\mathrm{M}\) de \(\mathrm{G}\) soit ouvert, il suffit donc qu’il soit Zariski-dense. La densité de Zariski se vérifiant après extension des scalaires de \(\mathbf{Q}_\ell\) à \(\overline{\mathbf{Q}}_\ell\), nous sommes ramenés à appliquer les lemmes algébriques (4.4.2\textsuperscript{a}) (pour \(n\) impair) et (4.4.2\textsuperscript{s}) (pour \(n\) pair) à la clôture de Zariski de \(\mathrm{M}\) dans \(\mathrm{G}(\overline{\mathbf{Q}}_\ell)\) et à la classe de conjugaison des transformations de Picard-Lefschetz. Les hypothèses en sont vérifiées grâce à (4.2.5) et (4.2.8). Dans ces lemmes, on a subrepticement remplacé \(\overline{\mathbf{Q}}_\ell\) par le corps isomorphe \(\mathbf{C}\), pour se sentir plus à l’aise.

\textit{Lemme} \textbf{(4.4.2\textsuperscript{a})}. — \textit{Soit \(\mathrm{V}\) un \(\mathbf{C}\)-espace vectoriel de dimension finie muni d’une forme alternée non dégénérée \((,)\), \(\mathrm{M}\) un sous-groupe algébrique de \(\operatorname{Sp}(\mathrm{V})\) et \(\mathrm{R}\) une orbite de \(\mathrm{M}\) dans \(\mathrm{V}\), qui engendre \(\mathrm{V}\). On suppose que \(\mathrm{M}\) est le plus petit sous-groupe algébrique de \(\operatorname{Sp}(\mathrm{V})\) contenant les transvections \(x\mapsto x+(x\delta)\delta\) \((\delta\in\mathrm{R})\). En d’autres termes, on suppose que \(\mathrm{M}\) est engendré par les sous-groupes à un paramètre de transvections \(\mathrm{U}_\delta=\{x\mapsto x+\lambda(x\delta)\delta\mid\lambda\in\mathbf{C}\}\) pour \(\delta\in\mathrm{R}\). Alors, \(\mathrm{M}=\operatorname{Sp}(\mathrm{V})\).}

\textit{Lemme} \textbf{(4.4.3\textsuperscript{a})}. — \textit{Soient \(\delta',\delta''\in\mathrm{R}\), non orthogonaux. Alors, les combinaisons linéaires non nulles de \(\delta'\) et \(\delta''\) sont dans \(\mathrm{R}\).}

Le sous-groupe de \(\mathrm{G}\) engendré par \(\mathrm{U}_{\delta'}\) et \(\mathrm{U}_{\delta''}\) est le groupe \(\operatorname{SL}(2)\) du plan tendu par \(\delta'\) et \(\delta''\). Il est dans \(\mathrm{M}\), et les vecteurs non nuls de ce plan forment une seule orbite pour \(\operatorname{SL}(2)\).

\textbf{(4.4.4\textsuperscript{a})} Prouvons (4.4.2\textsuperscript{a}). Le cas \(\mathrm{R}=\{0\}\) est trivial. Excluons-le. Si \(\delta\in\mathrm{R}\), on a alors \(\delta\neq0\), et il existe \(\delta'\in\mathrm{R}\) tel que \((\delta,\delta')\neq0\) (puisque \(\mathrm{R}\) engendre \(\mathrm{V}\) et que \((,)\) est non dégénérée). D’après (4.4.3\textsuperscript{a}), les multiples non nuls de \(\delta\) sont donc dans \(\mathrm{R}\). Soit \(\mathrm{W}\) un sous-espace maximal de \(\mathrm{V}\) tel que \(\mathrm{R}\cap\mathrm{W}\) soit dense dans \(\mathrm{W}\). On a \(\mathrm{W}\neq0\). S’il existe \(\delta\in\mathrm{R}-\mathrm{W}\) qui n’est pas orthogonal à \(\mathrm{W}\), il résulte de (4.4.3\textsuperscript{a}) que \(\mathrm{R}\) est dense
\Page{228}
dans \(\mathrm{W}\oplus\mathbf{C}.\delta\). Ceci contredit la maximalité de \(\mathrm{W}\). On a donc \(\mathrm{R}=(\mathrm{R}\cap\mathrm{W})\cup(\mathrm{R}\cap\mathrm{W}^{\perp})\). Puisque \(\mathrm{R}\) engendre \(\mathrm{V}\), on a \(\mathrm{V}=\mathrm{W}+\mathrm{W}^{\perp}\), et \(\mathrm{W}\) est non isotrope. Le groupe \(\mathrm{M}\) est engendré par les \(\mathrm{U}_\delta\) \((\delta\in\mathrm{R})\), donc contenu dans \(\operatorname{Sp}(\mathrm{W})\times\operatorname{Sp}(\mathrm{W}^{\perp})\). Puisque \(\mathrm{R}\) est une seule orbite de \(\mathrm{M}\), on a \(\mathrm{R}\subset\mathrm{W}\) et \(\mathrm{W}=\mathrm{V}\) : \(\mathrm{R}\) est dense dans \(\mathrm{V}\), les \(\mathrm{U}_\delta\) \((\delta\in\mathrm{V})\) sont tous dans \(\mathrm{M}\), et \(\mathrm{M}=\operatorname{Sp}(\mathrm{V})\).

\textit{Lemme} \textbf{(4.4.2\textsuperscript{s})}. — \textit{Soit \(\mathrm{V}\) un \(\mathbf{C}\)-espace vectoriel de dimension finie muni d’une forme bilinéaire symétrique non dégénérée \((,)\), \(\mathrm{M}\) un sous-groupe algébrique de \(\mathrm{O}(\mathrm{V})\) et \(\mathrm{R}\) une orbite de \(\mathrm{M}\) dans \(\mathrm{V}\), qui engendre \(\mathrm{V}\). On suppose que les éléments de \(\mathrm{R}\) satisfont à \((\delta,\delta)=2\), et que \(\mathrm{M}\) est le plus petit sous-groupe algébrique de \(\mathrm{O}(\mathrm{V})\) contenant les réflexions \(s_\delta:x\mapsto x-(x\delta)\delta\) pour \(\delta\in\mathrm{R}\). Alors, \(\mathrm{M}\) est fini ou \(\mathrm{M}=\mathrm{O}(\mathrm{V})\).}

\textit{Lemme} \textbf{(4.4.3\textsuperscript{s})}. — \textit{Soient \(\delta',\delta''\in\mathrm{R}\). Si \((\delta',\delta'')\) n’est pas de la forme \(2\cos\theta\), avec \(\theta\) multiple rationnel de \(2\pi\), les combinaisons linéaires \(\delta\) de \(\delta'\) et \(\delta''\) telles que \((\delta,\delta)=2\) sont dans \(\mathrm{R}\). La même conclusion vaut si \((\delta',\delta'')=\pm2\) (mais nous n’en aurons pas besoin).}

Dans le premier cas, la forme \((,)\) est non dégénérée sur le plan tendu par \(\delta'\) et \(\delta''\), et \(s_{\delta'}s_{\delta''}\) est une rotation d’ordre infini dans ce plan. Le plus petit sous-groupe algébrique contenant \(s_{\delta'}s_{\delta''}\) est donc le groupe \(\operatorname{SO}\) de ce plan (de dimension 1) et on conclut comme en (4.4.3\textsuperscript{a}). Dans le second cas, on se ramène à supposer que \((\delta',\delta'')=2\), en remplaçant si besoin est \(\delta'\) par \(s_{\delta'}(\delta')=-\delta'\). On peut aussi supposer que \(\delta'\neq\delta''\). La restriction de \((,)\) au plan tendu par \(\delta'\) et \(\delta''\) a alors pour noyau \(\mathbf{C}(\delta''-\delta')\), et dans ce plan \(s_{\delta'}s_{\delta''}\) est une transvection non nulle parallèle à ce noyau. L’enveloppe algébrique est ici le groupe de toutes ces transvections, et l’on conclut comme auparavant.

\textbf{(4.4.4\textsuperscript{s})} Prouvons (4.4.2\textsuperscript{s}). Comme en (4.4.2\textsuperscript{a}), soit \(\mathrm{W}\) un sous-espace non totalement isotrope de \(\mathrm{V}\), maximal parmi ceux tels que \(\mathrm{R}\cap\mathrm{W}\) soit dense dans la sphère \((\delta,\delta)=2\) de \(\mathrm{W}\). Si \(\delta\in\mathrm{R}-\mathrm{W}\), et que la forme linéaire \((\delta,x)\) n’est pas localement constante sur la sphère \((x,x)=2\) de \(\mathrm{W}\), elle prendra des valeurs non exclues dans (4.4.3\textsuperscript{s}) sur une partie topologiquement dense de \(\mathrm{R}\cap\mathrm{W}\) (car \(\mathrm{R}\cap\mathrm{W}\) est localement fermé et Zariski-dense dans ladite sphère). Comme en (4.4.4\textsuperscript{a}), ceci contredit la maximalité de \(\mathrm{W}\). La forme \((\delta,x)\) est donc localement constante sur la sphère \((x,x)=2\) de \(\mathrm{W}\). Ceci ne peut arriver que dans les cas suivants :
\par\noindent\textit{a)} \(\mathrm{W}\) est de dimension 1;
\par\noindent\textit{b)} \(\delta\perp\mathrm{W}\);
\par\noindent\textit{c)} \(\mathrm{W}\) est isotrope, \(\mathrm{W}\cap\mathrm{W}^{\perp}\) est de codimension 1 dans \(\mathrm{W}\), et \(\delta\perp(\mathrm{W}\cap\mathrm{W}^{\perp})\).

Si \(\mathrm{W}\cap\mathrm{W}^{\perp}\neq\{0\}\), on doit être dans les cas \textit{b)} ou \textit{c)}, et \(\mathrm{R}\subset(\mathrm{W}\cap\mathrm{W}^{\perp})^{\perp}\); ceci contredit l’hypothèse que \(\mathrm{R}\) engendre \(\mathrm{V}\). Si \(\dim\mathrm{W}\neq1\), on doit donc être dans le cas \textit{b)}, et \(\mathrm{R}\subset\mathrm{W}\cup\mathrm{W}^{\perp}\). Comme en (4.4.4\textsuperscript{a}), on en déduit que \(\mathrm{W}=\mathrm{V}\) et que \(\mathrm{M}=\mathrm{O}(\mathrm{V})\). Reste le cas où \(\mathrm{W}\) est toujours de dimension 1. D’après (4.4.3\textsuperscript{s}), les produits scalaires \((\delta',\delta'')\) \((\delta',\delta''\in\mathrm{R})\) sont alors tous de la forme \(2\cos\theta\) avec \(\theta\) multiple rationnel
\Page{229}
de \(2\pi\). Puisque \(\mathrm{R}\) est algébrique, ceci n’est possible que si \((\delta',\delta'')\) ne prend qu’un nombre fini de valeurs. Si les \(\delta_i\in\mathrm{R}\) forment une base de \(\mathrm{V}\), on n’a qu’un nombre fini de possibilités pour les \((\delta_i,\delta)\), \((\delta\in\mathrm{R})\), et \(\mathrm{R}\) est donc fini. Le groupe \(\mathrm{M}\), qui s’injecte dans le groupe des permutations de \(\mathrm{R}\), l’est également.

Nous allons étudier plus en détail le cas où le groupe \(\mathrm{M}\) est fini.

\textit{Proposition} \textbf{(4.4.5)}. — \textit{Si \(\mathrm{M}\) est fini, le groupe analogue en cohomologie \(\ell'\)-adique est également fini, pour tout \(\ell'\neq p\). En tant que quotient de \(\pi_1(\mathbf{P}^1-\mathrm{S},o)\), il est indépendant de \(\ell'\), et le caractère de sa représentation dans \(\operatorname{Ev}(\mathrm{Y},\mathbf{Q}_{\ell'})\) est à valeurs entières indépendantes de \(\ell'\).}

Pour tout \(\ell'\), le groupe fondamental \(\pi_1=\pi_1(\mathbf{P}^1-\mathrm{S},o)\) agit trivialement sur \(\mathrm{H}^i(\mathrm{Y})\), pour \(i\neq n\), et \(\mathrm{H}^n(\mathrm{Y})\) est somme directe d’un sous-espace sur lequel il agit trivialement, et de la partie évanescente. La représentation de \(\pi_1\) sur la partie évanescente est semi-simple et disjointe de la représentation triviale. On retrouve donc la représentation de \(\pi_1\) sur \(\operatorname{Ev}(\mathrm{Y})\) en débarrassant la représentation virtuelle de \(\pi_1\) sur \(\Sigma(-1)^i\mathrm{H}^i(\mathrm{Y})\) de ses facteurs triviaux : il suffit de vérifier que le caractère de cette représentation virtuelle est à valeurs rationnelles indépendantes de \(\ell'\).

Un argument facile de spécialisation (utilisant (1.11)) nous ramène à supposer que \(k\) est la clôture algébrique d’un corps fini. La variété \(\mathrm{X}\subset\mathbf{P}\), \(\mathrm{Y}\), le pinceau de Lefschetz, et donc le morphisme \(f:\widetilde{\mathrm{X}}\to\mathbf{P}^1\) sont alors définis sur un corps fini \(\mathbf{F}_q\), et, avec les notations (0.7), la représentation virtuelle de \(\pi_1(\mathbf{P}^1-\mathrm{S},o)\) sur \(\Sigma(-1)^i\mathrm{H}^i(\mathrm{Y})\) est induite par une représentation virtuelle de \(\pi_1(\mathbf{P}^1_0-\mathrm{S}_0,o)\) (groupe fondamental arithmétique). Pour tout point \(x\) de \(\mathbf{P}^1_0-\mathrm{S}_0\), rationnel sur \(\mathbf{F}_{q^m}\), la valeur de son caractère en \(\mathrm{F}_x\) est entière et indépendante de \(\ell'\) : d’après la formule de traces de Lefschetz, c’est le nombre de points \(\mathbf{F}_{q^m}\)-rationnels de la fibre de \(\widetilde{\mathrm{X}}\) en \(x\). La proposition résulte dès lors du lemme suivant.

\textit{Lemme} \textbf{(4.4.6)}. — \textit{Soient \(\mathrm{U}_0\) une variété algébrique normale et géométriquement connexe sur \(\mathbf{F}_q\), munie d’un point base \(u\), \(\mathrm{G}\) le quotient de \(\mathrm{W}(\mathrm{U}_0,u)\) par un sous-groupe d’indice fini de \(\pi_1(\mathrm{U},u)\) et \(\mathrm{V}\) une représentation \(\ell\)-adique virtuelle de \(\mathrm{G}\). Soit aussi \(\mathrm{V}'\) une représentation \(\ell'\)-adique virtuelle de \(\mathrm{W}(\mathrm{U}_0,u)\). On suppose que, sur les puissances des Frobenius \(\mathrm{F}_x\in\mathrm{W}(\mathrm{U}_0,u)\) \((x\in|\mathrm{U}_0|)\), les caractères de \(\mathrm{V}\) et \(\mathrm{V}'\) prennent des valeurs rationnelles, et coïncident. Alors \(\mathrm{V}'\) provient par inflation d’une représentation virtuelle de \(\mathrm{G}\), et les caractères de \(\mathrm{V}\) et \(\mathrm{V}'\) sont à valeurs rationnelles et égaux.}

\textit{Preuve} : D’après \v{C}ebotarev, tout élément de \(\mathrm{G}\) de degré assez grand est de la forme \(\mathrm{F}_x\) (cf. (3.5.3)). Un caractère de \(\mathrm{G}\) est uniquement déterminé par sa restriction aux éléments de degré assez grand : pour \(g\) central de degré \(>0\), et \(h\in\mathrm{G}\), la fonction de \(n:n\mapsto\chi(g^nh)\) est en effet combinaison linéaire de fonctions \(n\mapsto\lambda^n\), et la non-nullité d’un déterminant de Vandermonde montre que, quel que soit \(\mathrm{N}\), les restrictions de ces fonctions à \(n\geqq\mathrm{N}\) sont linéairement indépendantes. Il résulte donc des hypothèses que le caractère de \(\mathrm{V}\) est rationnel (le comparer à un de ses conjugués sous \(\operatorname{Aut}(\mathbf{C}/\mathbf{Q})\)).
\Page{230}
La représentation virtuelle \(\mathrm{V}\) est donc réalisable sur un corps de nombres (neutralisant quelques algèbres simples), et il existe une représentation virtuelle \(\ell'\)-adique \(\mathrm{V}_1\) de \(\mathrm{G}\) de même caractère. Remplaçant \(\mathrm{V}\) par \(\mathrm{V}_1\), on peut supposer que \(\ell=\ell'\). Dans ce cas, on a plus généralement :

\textit{Lemme} \textbf{(4.4.7)}. — \textit{Soient \(\mathrm{V}'\) et \(\mathrm{V}''\) deux représentations \(\ell\)-adiques virtuelles de \(\mathrm{W}(\mathrm{U}_0,u)\). Si leurs caractères coïncident sur les puissances des Frobenius, ils coïncident.}

Soit \(\mathrm{V}\) la somme des constituants de \(\mathrm{V}'\) et \(\mathrm{V}''\), \(\mathrm{G}^0\) l’adhérence de Zariski de l’image de \(\pi_1(\mathrm{U},u)\) dans \(\operatorname{GL}(\mathrm{V})\), et \(\mathrm{G}\) défini comme en (1.3.7). Les représentations virtuelles \(\mathrm{V}'\) et \(\mathrm{V}''\) sont définies par des représentations virtuelles de \(\mathrm{G}\). Soit \(g\) un élément central de \(\mathrm{G}\) de degré \(d>0\), et découpons \(\mathrm{V}'\) et \(\mathrm{V}''\) selon les valuations \(\ell\)-adiques \(v\) des valeurs propres de \(g\). Le polynôme caractéristique de \(\mathrm{F}_x\) agissant sur \(\mathrm{V}'\) détermine le polynôme caractéristique de \(\mathrm{F}_x\) agissant sur chaque morceau \(\mathrm{V}'_v\) : des valeurs propres de \(\mathrm{F}_x\), prises avec leur multiplicité, on ne garde que celles de valuation \(\dfrac{\deg(x)}{\deg(g)}.v\) (cf. (1.3.14)). Le polynôme caractéristique ne dépendant que des traces des puissances :
\[
\log\det(1-\mathrm{F}t,\mathrm{V}')
=-\Sigma\operatorname{Tr}(\mathrm{F}^n,\mathrm{V}')\frac{t^n}{n},
\]
\(\mathrm{V}'_v\) et \(\mathrm{V}''_v\) ont mêmes caractères sur les puissances des Frobenius. Par torsion (1.2.7), ceci nous ramène au cas où \(\mathrm{V}'_v\) et \(\mathrm{V}''_v\) proviennent de représentations de \(\pi_1(\mathrm{U}_0,u)\). On applique enfin le théorème de \v{C}ebotarev, selon lequel les Frobenius sont denses dans ce groupe.

\textbf{(4.4.8)} Plaçons-nous dans les hypothèses de (4.4.1), et supposons de plus \(n\) pair, \(n=2n_1\), \(\mathrm{M}\) fini et \(\operatorname{Ev}(\mathrm{Y})\) non nul (i.e. \(\mathrm{S}\neq\varnothing\)). La forme \((,)\) sur \(\operatorname{Ev}(\mathrm{Y})\) est à valeurs dans \(\mathbf{Q}_\ell(-n)\). Elle correspond à une forme, encore notée \((,)\), sur \(\operatorname{Ev}(\mathrm{Y})(n_1)\), à valeurs dans \(\mathbf{Q}_\ell\). Les cycles évanescents \(\delta\) sont des éléments de \(\operatorname{Ev}(\mathrm{Y})(n_1)\). Ils vérifient \((\delta,\delta)=(-1)^{n_1}.2\).

\textit{Théorème} \textbf{(4.4.9)}. — \textit{Sous les hypothèses de \textup{(4.4.8)} :}

a) \textit{Le \(\mathbf{Z}\)-module \(\mathrm{L}\) engendré par les cycles évanescents \(\delta\) est tel que \(\mathrm{L}\otimes\mathbf{Q}_\ell\xrightarrow{\sim}\operatorname{Ev}(\mathrm{Y})(n_1)\); la forme \((,)\) est à valeurs entières sur \(\mathrm{L}\).}

b) \textit{Les cycles évanescents \(\pm\delta\) forment dans \(\mathrm{L}\) un système de racines \(\mathrm{R}\) de type \(\mathrm{A}\), \(\mathrm{D}\) ou \(\mathrm{E}\). Le groupe \(\mathrm{M}\) est le groupe de Weyl correspondant.}

Dans \(\mathrm{M}\), les réflexions de la forme \(s_\delta\) sont caractérisées, indépendamment de \(\ell\), comme étant conjuguées à un élément non trivial d’un groupe d’inertie en un \(s\in\mathrm{S}\). Soient donc deux réflexions \(s_\delta\) et \(s_{\delta'}\). On a :
\[
\operatorname{Tr}(s_\delta s_{\delta'})
=(\dim\operatorname{Ev}(\mathrm{Y})-4)+(\delta,\delta')^2.
\]
\Page{231}
L’entier \(\operatorname{Tr}(s_\delta s_{\delta'})-\dim\operatorname{Ev}(\mathrm{Y})+4\) a donc une racine carrée dans \(\mathbf{Q}_\ell\) pour tout \(\ell\neq p\). C’est donc le carré d’un entier, et \((\delta,\delta')\in\mathbf{Z}\). Si les \(\delta_i\in\mathrm{R}\) forment une \(\mathbf{Q}_\ell\)-base de \(\operatorname{Ev}(\mathrm{Y})(n_1)\), on a donc \((\delta_i,\delta)\in\mathbf{Z}\) pour tout \(\delta\in\mathrm{R}\), et \(\delta\) est combinaison linéaire à coefficients rationnels des \(\delta_i\). L’assertion a) en résulte.

Les \(s_\delta\) forment une seule classe de conjugaison dans \(\mathrm{M}\), et \(\mathrm{M}\) est fini, engendré par les réflexions \(s_\delta\) de \(\mathrm{L}\). Il en résulte que \(\mathrm{R}\) est un système de racines irréductible dont toutes les racines sont de même longueur (donc de type \(\mathrm{A}\), \(\mathrm{D}\) ou \(\mathrm{E}\)), et \(\mathrm{M}\) est son groupe de Weyl.

\section*{(4.5) Application : le théorème du pgcd}

\textit{Les notations \textup{(0.7)} sont en vigueur.}

Le théorème suivant est utilisé par Katz et Messing [7] pour comparer les cohomologies \(\ell\)-adiques et cristallines.

\textit{Théorème} \textbf{(4.5.1)}. — \textit{Soient \(\mathrm{X}_0\) une variété projective lisse, absolument irréductible de dimension \(n+1\) sur \(\mathbf{F}_q\), et \((\mathrm{X}_t)_{t\in\mathbf{P}^1}\) un pinceau de Lefschetz de sections hyperplanes de \(\mathrm{X}\), défini sur \(\mathbf{F}_q\). Pour \(p=2\) et \(n\) pair, on le suppose assez général. Soit \(\mathrm{S}\subset\mathbf{P}^1\) l’ensemble des \(t\) pour lesquels \(\mathrm{X}_t\) est singulier. Alors, le polynôme \(\det(1-\mathrm{F}t,\mathrm{H}^n(\mathrm{X}))\) est le ppcm des polynômes \(f(\mathrm{T})=\prod(1-\alpha_i\mathrm{T})\) ayant la propriété suivante :}

\((*)\) \textit{Pour \(t\) un point de \(\mathbf{P}^1-\mathrm{S}\) défini sur \(\mathbf{F}_{q^k}\), et \(\mathrm{F}\) l’endomorphisme de Frobenius de \(\mathrm{X}_t\), relatif à sa \(\mathbf{F}_{q^k}\)-structure naturelle, on a :}
\[
\textstyle
\prod(1-\alpha_i^k\mathrm{T})
\quad\text{\textit{divise}}\quad
\det(1-\mathrm{F}\mathrm{T},\mathrm{H}^n(\mathrm{X}_t)).
\]

Écrivons \(\mathrm{H}^n(\mathrm{X}_t)=\mathrm{H}^n(\mathrm{X})+\operatorname{Ev}(\mathrm{X}_t)\) (4.3.9). Si les \(\beta_i\) sont les valeurs propres de Frobenius sur \(\mathrm{H}^n(\mathrm{X})\) (comptées avec leur multiplicité), cette décomposition donne, pour \(t\) défini sur \(\mathbf{F}_{q^k}\) :
\[
\textstyle
\det(1-\mathrm{F}\mathrm{T},\mathrm{H}^n(\mathrm{X}_t))
=\prod(1-\beta_i^k\mathrm{T}).\det(1-\mathrm{F}\mathrm{T},\operatorname{Ev}(\mathrm{X}_t)),
\]
de sorte que le polynôme \(\prod(1-\beta_i\mathrm{T})=\det(1-\mathrm{F}\mathrm{T},\mathrm{H}^n(\mathrm{X}))\) satisfait à \((*)\). Il reste à voir que tout polynôme satisfaisant à \((*)\) le divise. Nous prouverons plus précisément :

\textit{Lemme} \textbf{(4.5.2)}. — \textit{Quels que soient les polynômes :}
\[
\textstyle
\mathrm{P}(\mathrm{T})=\prod(1-\gamma_i\mathrm{T})
\qquad\text{\textit{et}}\qquad
\mathrm{Q}(\mathrm{T})=\prod(1-\delta_j\mathrm{T})
\]
\textit{(les \(\gamma_i\) et \(\delta_j\) sont supposés \(\neq0\)), si pour tout point \(t\in\mathbf{P}^1-\mathrm{S}\), défini sur \(\mathbf{F}_{q^k}\), on a :}
\[
\textstyle
\prod(1-\gamma_i^k\mathrm{T})
\mid
\prod(1-\delta_j^k\mathrm{T}).\det(1-\mathrm{F}\mathrm{T},\operatorname{Ev}(\mathrm{X}_t)),
\]
\textit{alors \(\mathrm{P}\mid\mathrm{Q}\).}

Déduisons le lemme (4.5.2) de son cas particulier où \(\mathrm{P}\) est de degré 1 : \(\mathrm{P}=(1-\gamma t)\). On procède par récurrence sur \(\deg\mathrm{P}\). Si \(\deg\mathrm{P}>0\) (seul cas non trivial), et que \(\gamma\) est
\Page{232}
l’un des \(\gamma_i\), le lemme pour \(\mathrm{P}=1-\gamma t\) montre que \(\gamma\) coïncide avec l’un des \(\delta_i\), et on applique l’hypothèse de récurrence à \(\mathrm{P}/(1-\gamma\mathrm{T})\) et \(\mathrm{Q}/(1-\gamma\mathrm{T})\).

Il reste à prouver (4.5.2) pour \(\mathrm{P}=(1-\gamma\mathrm{T})\). Raisonnons par l’absurde, et supposons \(\gamma\) distinct de tous les \(\delta_j\). Pour chaque \(j\), soit \(k_j\) le générateur de l’idéal des entiers \(k\) tels que \(\gamma^k=\delta_j^k\), et soit \(\mathrm{N}\) le produit des \(k_j\) non nuls. L’hypothèse assure que, pour \(k\equiv1(\mathrm{N})\), on a \(\gamma^k\neq\delta_j^k\), et donc :
\[
1-\gamma^k\mathrm{T}
\mid\det(1-\mathrm{F}\mathrm{T},\operatorname{Ev}(\mathrm{X}_t))
\qquad
(t\text{ sur }\mathbf{F}_{q^k},\ k\equiv1(\mathrm{N})).
\tag*{\textbf{(4.5.3)}}
\]
En particulier, \(\gamma\) est une unité \(\ell\)-adique.

Soit \(u\) un point géométrique de \(\mathbf{P}^1-\mathrm{S}\). Le groupe \(\pi_1((\mathbf{P}^1-\mathrm{S})_0,u)\) agit sur \(\mathrm{E}=\operatorname{Ev}(\mathrm{X}_u)\). Par ailleurs, il s’envoie dans \(\pi_1(\operatorname{Spec}(\mathbf{F}_q))=\operatorname{Gal}(\mathbf{F}/\mathbf{F}_q)=\widehat{\mathbf{Z}}\). Soit \(\mathrm{G}\) l’image de \(\pi_1((\mathbf{P}^1-\mathrm{S})_0,u)\) par l’application diagonale dans \(\operatorname{GL}(\mathrm{E})\times\widehat{\mathbf{Z}}\). C’est une extension de \(\widehat{\mathbf{Z}}\) par le groupe de monodromie géométrique \(\mathrm{M}\) déterminé au numéro précédent. Notons \(\rho\) sa projection dans \(\operatorname{GL}(\mathrm{E})\), et \(\deg\) sa projection dans \(\widehat{\mathbf{Z}}\). Pour \(t\in(\mathbf{P}^1-\mathrm{S})_0(\mathbf{F}_{q^k})\), et \(\mathrm{F}_t\) l’image dans \(\mathrm{G}\) du Frobenius correspondant, on a :
\[
\deg(\mathrm{F}_t)=k
\qquad\text{et}\qquad
\det(1-\mathrm{F}\mathrm{T},\operatorname{Ev}(\mathrm{X}_t))
=\det(1-\rho(\mathrm{F}_t)\mathrm{T}).
\]
Les \(\mathrm{F}_t\) étant denses dans \(\mathrm{G}\) (\v{C}ebotarev), on déduit par continuité de (4.5.3) que

\textbf{(4.5.4)} Pour \(g\in\mathrm{G}\), avec \(\deg(g)\equiv1(\mathrm{N})\), \(\gamma^{\deg(g)}\) est valeur propre de \(\rho(g)\). Nous allons voir que ceci est absurde, même si on se limite aux \(g\) tels que \(\deg(g)=1\). Puisque \(\rho(g)\) normalise \(\mathrm{M}\), et que les \(\rho(g)\) — pour \(\deg(g)=1\) — forment une classe latérale sous \(\mathrm{M}\), il suffit de vérifier le

\textit{Lemme} \textbf{(4.5.5)}. — \textit{Si \(g\in\operatorname{GL}(\mathrm{E})\) normalise le groupe de monodromie géométrique \(\mathrm{M}\), les \(gm\) pour \(m\in\mathrm{M}\) n’ont pas de valeur propre commune.}

Dans cet énoncé, on peut remplacer \(\mathrm{M}\) par son adhérence \(\mathrm{M}^{-}\) pour la topologie de Zariski. Celle-ci est soit le groupe symplectique d’une forme alternée non dégénérée sur \(\mathrm{E}\), soit le groupe orthogonal d’une forme quadratique non dégénérée sur \(\mathrm{E}\), soit un groupe de Weyl de type \(\mathrm{A}\), \(\mathrm{D}\) ou \(\mathrm{E}\). L’élément \(g\) est selon le cas une similitude symplectique, une similitude orthogonale, ou un multiple scalaire d’un automorphisme du système de racines. Il est loisible de remplacer \(g\) par un \(gm\) \((m\in\mathrm{M}^{-})\), et de le multiplier par un scalaire \(\lambda\) (ce qui multiplie toutes les valeurs propres par \(\lambda\)). Dans les deux premiers cas, ceci nous ramène à montrer que les éléments du groupe symplectique (resp. orthogonal) n’ont pas de valeur propre commune. On observe que \(\mathrm{Id}\) et \(-\mathrm{Id}\) déjà n’en ont pas. Dans le cas d’un groupe de Weyl, on se ramène à supposer que \(g\) est un automorphisme du système de racines, et seule compte sa classe mod \(\mathrm{W}\), i.e. l’automorphisme correspondant \(\bar{g}\) du diagramme de Dynkin.

Si \(\bar{g}\) est trivial, il s’agit de prouver que les éléments de \(\mathrm{W}\) n’ont pas de valeur propre commune, i.e. (puisque \(e\in\mathrm{W}\)) de trouver un élément n’ayant pas la valeur propre 1. On peut prendre une transformation de Coxeter. Le cas où \(\bar{g}\) est l’involution

% ---- 233-240 ----
\Page{233}
d’opposition (triviale pour \(\mathrm{A}_1,\mathrm{D}_{2k},\mathrm{E}_7,\mathrm{E}_8\), d’ordre \(2\) pour \(\mathrm{A}_k\) \((k>1)\), \(\mathrm{D}_{2k+1},\mathrm{E}_6\)) se ramène au précédent en remplaçant \(g\) par \(-g\). Il reste à traiter les cas \(\mathrm{D}_{2k}\), \(\bar g\) d’ordre \(2\), et \(\mathrm{D}_4\), \(\bar g\) d’ordre \(3\).

\textit{Cas \(\mathrm{D}_{2k}\), \(\bar g\) d’ordre \(2\)} : le système de racines est l’ensemble des vecteurs \(\pm e_i\pm e_j\) \((i\ne j)\) de \(\mathbf{Z}^{2k}\), \(\mathrm{W}\cup g\mathrm{W}\) est le groupe des permutations, avec changement de signe éventuel, des \(e_i\), et \(\mathrm{W}\) le sous-groupe des permutations avec un nombre pair de changements de signe. Les éléments :
\[
\left(
\begin{array}{l}
e_1,\ldots,e_{2k-1},e_{2k}\\
e_1,\ldots,e_{2k-1},-e_{2k}
\end{array}
\right)
\qquad\text{et}\qquad
\left(
\begin{array}{l}
e_1,\ldots,e_{2k-1},e_{2k}\\
e_2,\ldots,e_{2k},-e_1
\end{array}
\right)
\]
ont respectivement pour valeurs propres \(\pm1\), et les racines \((2k)^{\text{ième}}\) de \(-1\), donc n’ont aucune valeur propre en commun.

\textit{Cas \(\mathrm{D}_4\), \(\bar g\) d’ordre \(3\)} : le système de racines est l’ensemble des quaternions entiers de Hurwitz de norme \(2\), et on peut prendre pour \(g\) la multiplication par \(\frac{1}{2}(1+i+j+k)\). Ses valeurs propres sont les nombres \(\frac{1}{2}(1\pm\sqrt{-3})\), et \(g\) et \(-g\) n’ont pas de valeur propre commune — d’où le résultat puisque \(-1\) est dans le groupe de Weyl.

\Page{234}
\section*{\normalfont V. — APPLICATION AU \(\mathbf{Q}_{\ell}\)-TYPE D’HOMOTOPIE}

L’objet de ce paragraphe est de justifier certains des résultats annoncés dans [3]. Il s’agit d’attacher à chaque variété algébrique \(\mathrm{X}\) sur un corps algébriquement clos \(k\) une \(\mathbf{Q}_{\ell}\)-algèbre différentielle graduée \(\mathrm{A}(\mathrm{X})\), bien définie à quasi-isomorphisme près, fonctorielle en \(\mathrm{X}\), qui « soit » le \(\mathbf{Q}_{\ell}\)-type d’homotopie de \(\mathrm{X}\). En particulier, on veut que \(\mathrm{H}^*\mathrm{A}(\mathrm{X})=\mathrm{H}^*(\mathrm{X},\mathbf{Q}_{\ell})\). On veut aussi que, si \(k=\mathbf{C}\), le modèle minimal de \(\mathrm{A}(\mathrm{X})\) se déduise du modèle minimal de Sullivan [6] du type d’homotopie rationnelle de \(\mathrm{X}(\mathbf{C})\) en étendant les scalaires de \(\mathbf{Q}\) à \(\mathbf{Q}_{\ell}\).

La méthode à laquelle je pensais pour construire \(\mathrm{A}(\mathrm{X})\), lorsque je rédigeais [3], s’est évanouie. Celle que j’emploie ici utilise de façon essentielle une construction due indépendamment à Grothendieck et à Miller [8], décrite au n\textsuperscript{o} \(1\).

\section*{(5.1) Le \(\mathbf{Z}\{t\}\)-complexe de De Rham, d’après Grothendieck et Miller}

\textbf{(5.1.1)} \textit{L’algèbre des polynômes à puissances divisées} sur \(\mathbf{Z}\) est le sous-anneau \(\mathbf{Z}\{t\}\) de \(\mathbf{Q}[t]\) de \(\mathbf{Z}\)-base les \(\frac{t^n}{n!}\). La \textit{\(t\)-graduation} de \(\mathbf{Z}\{t\}\) est celle pour laquelle \(\frac{t^n}{n!}\) est de degré \(n\).

\textbf{(5.1.2)} Nous abrégerons \textit{différentiel gradué} en DG, et les algèbres DG seront toujours supposées commutatives, au sens de la règle de Koszul (exemple type d’algèbre DG : les formes différentielles extérieures sur une variété).

\textbf{(5.1.3)} Soit \(\mathrm{A}\) un anneau commutatif, et posons \(\mathrm{A}\{t\}=\mathrm{A}\otimes\mathbf{Z}\{t\}\). Pour \(\mathrm{A}=\mathbf{Q}\), on trouve simplement l’anneau des polynômes \(\mathbf{Q}[t]\). Une \(\mathrm{A}\{t\}\)\textit{-algèbre \(t\)-graduée} DG est une \(\mathrm{A}\{t\}\)-algèbre DG, dont chaque composante homogène est munie d’une graduation, la \textit{\(t\)-graduation}, qui en fait un \(\mathrm{A}\{t\}\)-module gradué, et telle que le produit et la différentielle \(d\) soient homogènes de \(t\)-degré \(0\).

\textbf{(5.1.4)} La construction « \(\mathbf{Z}\{t\}\)-complexe de De Rham », qui sera définie plus bas, est un foncteur contravariant qui à chaque ensemble simplicial \(\mathrm{X}\) associe une \(\mathbf{Z}\{t\}\)-algèbre \(t\)-graduée DG : \(\mathrm{X}\mapsto\Omega(\mathrm{X})\). On note \(\Omega^p(\mathrm{X})\) la partie de degré \(p\) pour la structure DG (degré extérieur).

Les \(t\)-degré et degré extérieur d’un élément bihomogène satisfont :
\[
0\leq\text{degré extérieur}\leq t\text{-degré}.
\tag*{\textbf{(5.1.4.1)}}
\]

\Page{235}
On dispose de plus d’un morphisme fonctoriel de \(\mathbf{Z}\{t\}\)-algèbres bigraduées :
\[
\mathrm{H}^*\Omega(\mathrm{X})
\longrightarrow
\mathrm{H}^*(\mathrm{X},\mathbf{Z})\otimes\mathbf{Z}\{t\}
\tag*{\textbf{(5.1.4.2)}}
\]
qui identifie la cohomologie de \(\Omega(\mathrm{X})\) à la somme des \(\mathrm{H}^p(\mathrm{X},\mathbf{Z})\otimes\frac{t^n}{n!}\) pour \(p\leq n\).

\textbf{(5.1.5)} Plus généralement, pour tout anneau commutatif \(\mathrm{A}\) (les cas qui nous importent sont \(\mathrm{A}=\mathbf{Z},\mathbf{Z}/\ell^n,\mathbf{Q}\)), on dispose d’un foncteur « \(\mathrm{A}\{t\}\)-complexe de De Rham » \(\Omega(\mathrm{X},\mathrm{A})\) obéissant au formalisme (5.1.4), avec \(\mathbf{Z}\) remplacé par \(\mathrm{A}\). Il est covariant en \(\mathrm{A}\) et l’isomorphisme \((5.1.4.2)_{\mathrm{A}}\) est fonctoriel en \(\mathrm{A}\). Pour \(\mathrm{A}=\mathbf{Q}\), on dispose d’un morphisme de \(\Omega(\mathrm{X},\mathbf{Q})\) dans la \(\mathbf{Q}[t]\)-algèbre \(t\)-graduée DG déduite du complexe de De Rham-Sullivan \(\Omega_{\mathrm{DRS}}(\mathrm{X})\) de \(\mathrm{X}\) par extension des scalaires de \(\mathbf{Q}\) à \(\mathbf{Q}[t]\). Cet isomorphisme est compatible avec (5.1.4.2). Il identifie la partie de \(t\)-degré \(n\) de \(\Omega(\mathrm{X},\mathbf{Q})\) à \(t^n\otimes\)(le sous-complexe de \(\Omega_{\mathrm{DRS}}(\mathrm{X})\) formé des formes de degré \(\leq n\) en les coordonnées barycentriques de chaque simplexe).

Si \(\mathrm{B}\) est un quotient de \(\mathrm{A}\), le morphisme naturel de \(\Omega^p(\mathrm{X},\mathrm{A})\otimes_{\mathrm{A}}\mathrm{B}\) dans \(\Omega^p(\mathrm{X},\mathrm{B})\) est un isomorphisme, sauf en la composante de \(t\)-degré \(p\).

\textbf{(5.1.6)} Quelques préliminaires, avant de donner l’idée de la construction, puis la construction elle-même.

\textbf{(5.1.6.1)} Si \(\sigma\) est un ensemble fini, on notera \(\mathbf{Z}\{(x_i)_{i\in\sigma}\}\) le sous-anneau de \(\mathbf{Q}[(x_i)_{i\in\sigma}]\) de \(\mathbf{Z}\)-base les monômes \(\prod\limits_{i\in\sigma}\frac{x_i^{n_i}}{n_i!}\). Soit \(\Omega'(\sigma)\) le complexe de De Rham correspondant : c’est le sous-complexe du complexe de De Rham de \(\mathbf{Q}[(x_i)_{i\in\sigma}]\) de \(\mathbf{Z}\{(x_i)_{i\in\sigma}\}\)-base les produits de \(dx_i\) (prendre \textit{un} produit pour chaque partie finie de \(\sigma\)). Posons \(t=\sum x_i\), et soit enfin \(\Omega(\sigma)\) le quotient de \(\Omega'(\sigma)\) par l’idéal engendré par \(dt\). C’est une \(\mathbf{Z}\{t\}\)-algèbre \(t\)-graduée DG, si on prend pour \(t\)-graduation la graduation pour laquelle \(x_i\) et \(dx_i\) sont de \(t\)-degré \(1\). On l’appelle la \(\mathbf{Z}\{t\}\)\textit{-algèbre \(t\)-graduée} DG \textit{des formes différentielles relatives (rel. à \(t\)), polynomiales, entières à puissances divisées}, sur \(\mathbf{R}^{\sigma}\).

\textbf{(5.1.6.2)} Quel que soit l’anneau \(\mathrm{A}\), on pose de même \(\Omega(\sigma,\mathrm{A})=\Omega(\sigma)\otimes\mathrm{A}\).

\textbf{(5.1.6.3)} Une construction plus intrinsèque est la suivante.

\textit{a)} Si \(\mathrm{L}\) est un \(\mathrm{A}\)-module libre, on considère \(\Omega(\mathrm{L})=\Gamma^*(\mathrm{L})\otimes\bigwedge\limits^{*}\mathrm{L}\); le degré extérieur est celui de la puissance extérieure, et le \(t\)-degré est la somme des degrés extérieurs et symétriques; la différentielle, compatible aux puissances divisées, est caractérisée par \(d(x\otimes1)=1\otimes x\).

\textit{b)} Si \(\mathrm{M}\) est un quotient libre de \(\mathrm{L}\), on considère le quotient \(\Omega(\mathrm{L},\mathrm{M}):=\Gamma^*(\mathrm{L})\otimes\bigwedge\limits^{*}\mathrm{M}\) de \(\Omega(\mathrm{L})\).

\textit{c)} On prend \(\Omega(\sigma,\mathrm{A})=\Omega(\mathrm{A}^{\sigma},\mathrm{A}^{\sigma}\text{ modulo }\mathrm{A}\text{ diagonal})\).

\Page{236}
\textbf{(5.1.6.4)} La \(\mathbf{Z}\{t\}\)-algèbre \(t\)-graduée DG \(\Omega(\sigma)\) est un foncteur contravariant de \(\sigma\). C’est clair sur (5.1.6.3), le couple \((\mathbf{Z}^{\sigma},\mathbf{Z}\text{ diagonal})\) étant un foncteur contravariant en \(\sigma\). Plus concrètement, à \(\varphi:\tau\to\sigma\) on associe \(\varphi^*\) qui envoie \(x_i\) sur \(\sum\limits_{\varphi j=i}x_j\). Par exemple, pour \(\tau\subset\sigma\) et \(\varphi\) le morphisme d’inclusion, \(\varphi^*\) s’obtient en faisant \(x_i=0\) pour \(i\notin\tau\).

\textbf{(5.1.7)} Voici l’idée de la construction dans le contexte parallèle des « schémas simpliciaux » de Godement : on part d’un ensemble \(\mathrm{S}\), et d’un ensemble \(\mathcal{S}\) de parties finies non vides de \(\mathrm{S}\), les simplexes. On suppose que chaque partie à \(1\) élément de \(\mathrm{S}\) est un simplexe, et que tout sous-ensemble non vide d’un simplexe est un simplexe. Pour \(\sigma\in\mathcal{S}\), on note \(|\sigma|\) le simplexe (usuel) de \(\mathbf{R}^{(\mathrm{S})}\) tendu par les vecteurs de base \(e_i\), \(i\in\sigma\), et l’on prend pour espace \(\mathrm{X}\) la réunion des \(|\sigma|\) pour \(\sigma\in\mathcal{S}\). Soit \(\widehat{\mathrm{X}}\) le cône sur \(\mathrm{X}\), la réunion des cônes de sommet \(\mathrm{O}\) tendus par les \(|\sigma|\).

Il nous sera plus commode de considérer les espaces homotopes \(\mathrm{X}'\), réunion des espaces affines \(\sigma'\) tendus par les \(|\sigma|\), et \(\widehat{\mathrm{X}}'\), réunion des espaces vectoriels \(\hat{\sigma}'\) tendus par les \(|\sigma|\). Soit \(t:\mathbf{R}^{(\mathrm{S})}\to\mathbf{R}\) la fonction « somme des coordonnées ». L’espace \(\mathrm{X}'\) est la fibre en \(t=1\) de la restriction \(\widehat{\mathrm{X}}'\to\mathbf{R}\) de \(t\) à \(\widehat{\mathrm{X}}'\).

Pour chaque \(\sigma\), soit \(\Omega(\sigma)\) la \(\mathbf{Z}\{t\}\)-algèbre DG des formes différentielles relatives, rel. à \(t\), polynomiales, entières à puissances divisées, sur \(\hat{\sigma}'\). Par définition, une \(p\)-forme \(\alpha\in\Omega^p(\mathrm{X})\) est une famille \(\alpha_{\sigma}\in\Omega^p(\sigma)\) \((\sigma\in\mathcal{S})\) de \(p\)-formes, telle que pour \(\tau\subset\sigma\), \(\alpha_{\tau}\) soit la restriction de \(\alpha_{\sigma}\in\Omega^p(\sigma)\) à \(\Omega^p(\tau)\).

\textbf{(5.1.8)} Voici la définition formelle. Restreignons le foncteur \(\Omega(\sigma)\) ((5.1.6.4)) aux ensembles finis \(\Delta_n=[0,n]\) et aux applications croissantes entre ceux-ci. On obtient une \(\mathbf{Z}\{t\}\)-algèbre \(t\)-graduée DG simpliciale \(\Omega\).

On pose :
\[
\Omega(\mathrm{X},\mathrm{A})=\operatorname{Hom}(\mathrm{X},\Omega\otimes\mathrm{A})
\]
(somme sur \(i,p\) des groupes des morphismes d’ensembles simpliciaux de \(\mathrm{X}\) dans la composante \((i,p)\)-bihomogène de \(\Omega\otimes\mathrm{A}\)).

\textbf{(5.1.9)} Pour la preuve des propriétés fondamentales de cette construction, je renvoie à [8]. Indiquons seulement que :

\textit{a)} Le morphisme (5.1.5) dans \(\Omega_{\mathrm{DRS}}(\mathrm{X})\otimes\mathbf{Q}[t]\) s’obtient en envoyant \(x_i\) sur (coordonnée barycentrique \(x_i\))\(\otimes t\).

\textit{b)} Le point clef est le calcul des groupes d’homologie (faut-il dire d’homotopie ?) du groupe simplicial \(\Omega\) : chaque \(\mathrm{H}_i\) est isomorphe à \(\mathbf{Z}\), et est placé en degré extérieur et \(t\)-degré \(i\). Ce calcul résulte facilement de (5.1.6.3), si on se rappelle que le groupe simplicial \(\sigma\mapsto\mathrm{A}^{\sigma}\) est homotope à \(0\) (son normalisé est réduit à \(\mathrm{A}\xrightarrow{\sim}\mathrm{A}\) en degrés homo-\Page{237}logiques \(0\) et \(1\)) : on a \(\Gamma^*(\mathrm{A}^{\sigma})\sim\Gamma^*(0)=(\sigma\mapsto\mathrm{A})\), et \(\bigwedge\limits^{*}(\mathrm{A}^{\sigma}\text{ modulo }\mathrm{A}\text{ diagonal})\) se calcule par les théorèmes de décalage de Quillen : son normalisé est l’algèbre symétrique du \(\mathrm{A}\)-module \(\mathrm{A}\), le \(\operatorname{Sym}^n\) étant placé en degré homologique \(n\).

\section*{(5.2) Le \(\mathbf{Q}_{\ell}\)-type d’homotopie}

\textbf{(5.2.1)} Soit \(\mathbf{X}\) un système projectif filtrant d’ensembles simpliciaux \(\mathrm{X}^{\alpha}\). Posons :
\[
\mathrm{H}^i(\mathbf{X},\mathbf{Z}/\ell^n)
=\mathop{\mathrm{lim}\,\mathrm{ind}}\mathrm{H}^i(\mathrm{X}^{\alpha},\mathbf{Z}/\ell^n).
\tag*{\textbf{(5.2.1.1)}}
\]
On a encore la relation (5.1.4.2) entre ces groupes et la \(\mathbf{Z}/\ell^n\{t\}\)-algèbre \(t\)-graduée DG :
\[
\Omega(\mathbf{X},\mathbf{Z}/\ell^n)
=\mathop{\mathrm{lim}\,\mathrm{ind}}\Omega(\mathrm{X}^{\alpha},\mathbf{Z}/\ell^n)
\tag*{\textbf{(5.2.1.2)}}
\]
(exactitude des limites inductives).

Faisons maintenant l’hypothèse :
\[
\text{\textit{les }}\mathrm{H}^i(\mathbf{X},\mathbf{Z}/\ell)
\text{\textit{ sont des groupes finis.}}
\tag*{(F)}
\]
Les arguments de SGA5 V (5.3.1) s’appliquent, montrent que chaque :
\[
\mathrm{H}^i(\mathbf{X},\mathbf{Z}_{\ell})
=\mathop{\mathrm{lim}\,\mathrm{proj}}\mathrm{H}^i(\mathrm{X},\mathbf{Z}/\ell^n)
\tag*{\textbf{(5.2.1.3)}}
\]
est un \(\mathbf{Z}_{\ell}\)-module de type fini, et qu’on a les suites exactes habituelles :
\[
0\longrightarrow
\mathrm{H}^i(\mathbf{X},\mathbf{Z}_{\ell})\otimes\mathbf{Z}/\ell^n
\longrightarrow
\mathrm{H}^i(\mathbf{X},\mathbf{Z}/\ell^n)
\longrightarrow
\operatorname{Tor}_{1}^{\mathbf{Z}_{\ell}}
(\mathrm{H}^{i+1}(\mathbf{X},\mathbf{Z}_{\ell}),\mathbf{Z}/\ell^n)
\longrightarrow0
\]
(coefficients universels). Posons
\[
\Omega(\mathbf{X},\mathbf{Z}_{\ell})
=\lim_{n}\mathrm{proj}\,\Omega(\mathbf{X},\mathbf{Z}/\ell^n)
\tag*{\textbf{(5.2.1.4)}}
\]
(limite projective bigraduée); c’est une \(\mathbf{Z}_{\ell}\{t\}\)-algèbre \(t\)-graduée DG. On a :

\textit{a)} les \(\mathrm{H}^i(\mathbf{X},\mathbf{Z}/\ell^n)\) sont finis; leur système projectif vérifie donc la condition de Mittag-Leffler ML;

\textit{b)} en chaque bidegré \((i,p)\), sauf pour \(i=p\), les morphismes de transition entre les \({}_i\Omega^p(\mathbf{X},\mathbf{Z}/\ell^n)\) sont surjectifs ((5.1.5)), et la condition ML est vérifiée; pour \(i=p\), on a une suite exacte :
\[
{}_p\Omega^{p-1}(\mathbf{X},\mathbf{Z}/\ell^n)
\xrightarrow{d}
{}_p\Omega^p(\mathbf{X},\mathbf{Z}/\ell^n)
\longrightarrow
\mathrm{H}^p(\mathbf{X},\mathbf{Z}/\ell^n)
\longrightarrow0,
\]
d’où ML puisque la cohomologie vérifie ML.

Grâce à \textit{a)}, \textit{b)}, le passage à la limite projective (5.2.1.4) est exact, d’où un morphisme :
\[
\mathrm{H}^*\Omega(\mathbf{X},\mathbf{Z}_{\ell})
\longrightarrow
\mathrm{H}^*(\mathbf{X},\mathbf{Z}_{\ell})
\otimes_{\mathbf{Z}_{\ell}}\mathbf{Z}_{\ell}\{t\}
\tag*{\textbf{(5.2.1.5)}}
\]
qui identifie la cohomologie de \(\Omega(\mathbf{X},\mathbf{Z}_{\ell})\) à la somme des \(\mathrm{H}^p(\mathbf{X},\mathbf{Z}_{\ell})\otimes\frac{t^n}{n!}\) pour \(p\leq n\).

Il est essentiel, dans cette construction, que la limite inductive, sur \(\alpha\), ait précédé la limite projective, sur \(n\).

\Page{238}
Pour tout module \(t\)-gradué sur \(\mathbf{Q}[t]\) (par exemple sur \(\mathbf{Q}_{\ell}[t]\)), notons \(\lim\nolimits_t\mathrm{M}\) la limite inductive du système inductif des composantes homogènes \(\mathrm{M}_k\) de \(\mathrm{M}\), avec pour morphisme de transition de \(\mathrm{M}_n\) à \(\mathrm{M}_{n+m}\) la multiplication par \(t^m\).

Posons :
\[
\mathrm{H}^*(\mathbf{X},\mathbf{Q}_{\ell})
=\mathrm{H}^*(\mathbf{X},\mathbf{Z}_{\ell})
\otimes_{\mathbf{Z}_{\ell}}\mathbf{Q}_{\ell}
\tag*{\textbf{(5.2.1.6)}}
\]
et
\[
\mathrm{A}(\mathbf{X})
=\lim\nolimits_t\mathbf{Q}_{\ell}\otimes\Omega(\mathbf{X},\mathbf{Z}_{\ell}).
\tag*{\textbf{(5.2.1.7)}}
\]
On déduit de (5.2.1.5) par passage à la limite un isomorphisme :
\[
\mathrm{H}^*\mathrm{A}(\mathbf{X})=\mathrm{H}^*(\mathbf{X},\mathbf{Q}_{\ell}).
\tag*{\textbf{(5.2.1.8)}}
\]

\textbf{(5.2.2)} Soit \(\mathrm{X}\) une variété algébrique sur un corps algébriquement clos \(k\). Elle peut être définie sur un sous-corps algébriquement clos dénombrable \(k_1\) de \(k\). Soit donc \(\mathrm{X}_1\) sur \(k_1\), dont \(\mathrm{X}\) se déduit par extension des scalaires.

Parce que \(k_1\) est dénombrable, on peut trouver un système projectif dénombrable, indexé par les entiers, de morphismes étales et surjectifs \(\varphi_{\alpha}:\mathrm{X}_1^{\alpha}\to\mathrm{X}_1\), tel que pour tout morphisme étale et surjectif \(\varphi:\mathrm{U}\to\mathrm{X}_1\), il existe \(\alpha\) et une factorisation de \(\varphi_{\alpha}\) par \(\varphi\). Pour chaque \(\alpha\), et chaque ensemble fini \(\sigma\), soit \((\mathrm{X}_1^{\alpha}/\mathrm{X}_1)^{\sigma}\) le produit fibré sur \(\mathrm{X}_1\) d’une famille indexée par \(\sigma\) de copies de \(\mathrm{X}_{\alpha}^{1}\). Lorsque \(\sigma\) parcourt les ensembles finis \(\Delta_n=[0,n]\), ces produits fibrés forment un schéma simplicial. On pose :
\[
\mathbf{X}=(\pi_0((\mathrm{X}_{\alpha}^{1}/\mathrm{X}_1)^{\Delta_n}))_{\alpha}
\tag*{\textbf{(5.2.2.1)}}
\]
et
\[
\mathrm{A}(\mathrm{X})=\mathrm{A}(\mathbf{X}).
\tag*{\textbf{(5.2.2.2)}}
\]

Il eût été plus naturel de prendre, plutôt que les \((\mathrm{X}_1^{\alpha}/\mathrm{X}_1)^{\Delta_n}\), un système projectif cofinal d’hyperrecouvrements étales de \(\mathrm{X}_1\); d’après Artin cela revient essentiellement au même.

Par définition de la cohomologie étale, si on avait considéré des hyperrecouvrements, et d’après Artin avec la définition choisie, on a :
\[
\mathrm{H}^*(\mathbf{X},\mathbf{Z}/\ell^n)
=\mathrm{H}^*(\mathrm{X}_1,\mathbf{Z}/\ell^n).
\]
On sait par ailleurs que \(\mathrm{H}^*(\mathrm{X}_1,\mathbf{Z}/\ell^n)\xrightarrow{\sim}\mathrm{H}^*(\mathrm{X},\mathbf{Z}/\ell^n)\). La condition (F) est vérifiée d’après SGA4\(\frac{1}{2}\) [Th. finitude] (1.1).

Par définition de la cohomologie \(\ell\)-adique, on a donc un isomorphisme :
\[
\mathrm{H}^*\mathrm{A}(\mathrm{X})
\xrightarrow{\sim}
\mathrm{H}^*(\mathrm{X},\mathbf{Q}_{\ell}).
\]

\textbf{(5.2.3)} Supposons \(\mathrm{X}\) défini sur un sous-corps dénombrable \(k_0\) de \(k\), et prenons pour \(k_1\) la clôture algébrique de \(k_0\) dans \(k\). On peut alors choisir le système projectif des \(\mathrm{X}_1^{\alpha}\) défini sur \(k_0\). Le groupe de Galois \(\operatorname{Gal}(k_1/k_0)\) agit alors sur \(\mathbf{X}\), donc, par transport de structure, sur \(\mathrm{A}(\mathrm{X})\). Cette action est compatible avec l’action de Galois sur la cohomologie \(\ell\)-adique de \(\mathrm{X}\).

\Page{239}
On prendra garde que \(\mathrm{A}(\mathrm{X})\) est énorme — c’était le prix à payer pour qu’il dépende canoniquement de \(\mathbf{X}\) — de sorte qu’on ne peut guère parler de continuité pour l’action de Galois. On peut seulement dire que Galois agit continûment sur les \(\Omega(\mathbf{X},\mathbf{Z}/\ell^n)\) (munis de la topologie discrète), pour tout \(n\).

\textbf{(5.2.4)} Faisons \(k=\mathbf{C}\), et soit un système projectif \(\mathrm{X}_1^{\alpha}\) comme en (5.2.2). On peut alors trouver un système projectif d’espaces \(\mathrm{X}^{\alpha}\) au-dessus de \(\mathrm{X}(\mathbf{C})\) et de morphismes \(a_{\alpha}:\mathrm{X}^{\alpha}\to\mathrm{X}_1^{\alpha}(\mathbf{C})\), tel que chaque \(\mathrm{X}^{\alpha}\) soit la somme disjointe des ouverts d’un recouvrement ouvert de \(\mathrm{X}\), les intersections \(p\) à \(p\) de ces ouverts étant vides ou contractiles (pour tout \(p\geq1\)).

Soit \(\mathbf{X}(\mathbf{C})\) le système projectif des nerfs \(\mathrm{N}_{\alpha}=\pi_0(\mathrm{X}^{\alpha}/\mathrm{X}(\mathbf{C})^{\Delta_n})\) des recouvrements \(\mathrm{X}^{\alpha}\). Chacun de ces nerfs a le type d’homotopie de \(\mathrm{X}\). Le système projectif des \(a_{\alpha}\) fournit un morphisme de \(\mathbf{X}(\mathbf{C})\) dans \(\mathbf{X}\).

On va maintenant déduire du théorème de comparaison SGA4 XVI (4.1).

\textit{Construction} \textbf{(5.2.5)}. — \textit{Construction d’isomorphismes entre les modèles minimaux de la \(\mathbf{Q}_{\ell}\)-algèbre \textup{DG} \(\mathrm{A}(\mathrm{X})\) et le tensorisé avec \(\mathbf{Q}_{\ell}\) des modèles minimaux \(\mathcal{M}(\mathrm{X}(\mathbf{C}))\) de Sullivan \textup{[6]} du type d’homotopie rationnelle de l’espace topologique \(\mathrm{X}(\mathbf{C})\).}

D’après SGA4 XVI (4.1), le morphisme :
\[
\Omega(\mathbf{X},\mathbf{Z}/\ell^n)
\xrightarrow{a}
\Omega(\mathbf{X}(\mathbf{C}),\mathbf{Z}/\ell^n)
\]
est un quasi-isomorphisme : il induit sur la cohomologie des isomorphismes déduits des isomorphismes de comparaison :
\[
\mathrm{H}^i(\mathrm{X},\mathbf{Z}/\ell^n)
\xrightarrow{\sim}
\mathrm{H}^i(\mathrm{X}(\mathbf{C}),\mathbf{Z}/\ell^n).
\]
On a aussi des quasi-isomorphismes :
\[
\Omega(\mathrm{N}_{\alpha},\mathbf{Z}/\ell^n)
\longrightarrow
\Omega(\mathbf{X}(\mathbf{C}),\mathbf{Z}/\ell^n).
\]
Par passage à la limite, on trouve des quasi-isomorphismes :
\[
\Omega(\mathbf{X},\mathbf{Z}_{\ell})
\longrightarrow
\Omega(\mathbf{X}(\mathbf{C}),\mathbf{Z}_{\ell})
\longleftarrow
\Omega(\mathrm{N}_{\alpha},\mathbf{Z}_{\ell}).
\]
On a de plus des quasi-isomorphismes :
\[
\Omega(\mathrm{N}_{\alpha},\mathbf{Z})\otimes\mathbf{Z}_{\ell}
\longrightarrow
\Omega(\mathrm{N}_{\alpha},\mathbf{Z}_{\ell}),
\]
et
\[
\lim\nolimits_t\Omega(\mathrm{N}_{\alpha},\mathbf{Z})\otimes\mathbf{Q}
\longrightarrow
\Omega_{\mathrm{DRS}}(\mathrm{N}_{\alpha})
\sim
\Omega_{\mathrm{DRS}}(\mathrm{X}(\mathbf{C})),
\]
d’où une chaîne de quasi-isomorphismes :
\[
\begin{aligned}
&\mathrm{A}(\mathbf{X})
\longrightarrow
\mathrm{A}(\mathbf{X}(\mathbf{C}))
\longleftarrow
(\varinjlim\nolimits_t\Omega(\mathrm{N}_{\alpha},\mathbf{Z})\otimes\mathbf{Q})
\otimes\mathbf{Q}_{\ell}\\
&\qquad\qquad\qquad
\longrightarrow
\Omega_{\mathrm{DRS}}(\mathrm{N}_{\alpha})\otimes\mathbf{Q}_{\ell}
\sim
\Omega_{\mathrm{DRS}}(\mathrm{X}_{\mathbf{C}})\otimes\mathbf{Q}_{\ell}.
\end{aligned}
\]
Passer au modèle minimal remplace ces quasi-isomorphismes par des isomorphismes.

\Page{240}
\section*{(5.3) Graduations par le poids.}

\textbf{(5.3.1)} Soit \(\mathrm{A}\) une algèbre DG sur un corps \(\mathrm{K}\) de caractéristique \(0\). Pour nous, \(\mathrm{K}\) sera \(\mathbf{Q}\) ou \(\mathbf{Q}_{\ell}\). On suppose que \(\mathrm{H}^0(\mathrm{A})=\mathrm{K}\) et que les \(\mathrm{H}^i(\mathrm{A})\) sont de dimension finie. Pour la définition des modèles minimaux \(\mathcal{M}\to\mathrm{A}\) de \(\mathrm{A}\), je renvoie à [6]. On déduit facilement de la construction inductive de \(\mathcal{M}\) qu’il existe une filtration croissante exhaustive de \(\mathcal{M}\) par des sous-algèbres de type fini \(\mathcal{M}_n\) \((n\geq0)\) telle que

\((*)\) Tout automorphisme de \(\mathcal{M}\) est homotope à un automorphisme de \(\mathcal{M}\) qui respecte chaque \(\mathcal{M}_n\).

Si l’automorphisme \(\sigma\) de \(\mathcal{M}\) respecte la filtration par les \(\mathcal{M}_n\), il y a un sens à parler des valeurs propres de \(\sigma\), et de la graduation de \(\mathcal{M}\) par les sous-espaces propres généralisés : après extension des scalaires à une clôture algébrique \(\bar{\mathrm{K}}\) de \(\mathrm{K}\), \(\mathcal{M}\otimes\bar{\mathrm{K}}\) est somme directe des \(\mathcal{M}(a)=\bigcup\limits_m\operatorname{Ker}(\sigma-a)^m\) \((a\in\bar{k}^*)\), et ceux-ci forment une graduation de \(\mathcal{M}\otimes\bar{k}\) de type \(\bar{\mathrm{K}}^*\).

On déduit facilement de la construction inductive de \(\mathcal{M}\) que toute valeur propre de \(\sigma\) sur \(\mathcal{M}\otimes\bar{\mathrm{K}}\) est un produit de valeurs propres de \(\sigma\) sur les \(\mathrm{H}^i(\mathrm{A})\otimes\bar{\mathrm{K}}\), voire même sur les quotients sphériques des \(\mathrm{H}^i(\mathrm{A})\otimes\bar{\mathrm{K}}\). Plus précisément, une valeur propre de \(\sigma\) sur \(\mathcal{M}^n\otimes\bar{\mathrm{K}}\) est un produit \(\prod\alpha_i\) de valeurs propres \(\alpha_i\) de \(\sigma\) sur \(\mathrm{H}^{n(i)}\), avec \(\sum n(i)\geq n\).

Soit \(\Gamma\) le sous-groupe du groupe multiplicatif de \(\bar{\mathrm{K}}^*\) engendré par les valeurs propres de \(\sigma\). Si \(v:\Gamma\to\mathbf{Z}\) est un homomorphisme, les \(\mathcal{M}(n)=\bigoplus\limits_{v(a)=n}\mathcal{M}(a)\) forment une graduation (de type \(\mathbf{Z}\)) de \(\mathcal{M}\otimes\bar{\mathrm{K}}\). Si pour toute valeur propre \(\alpha\) de \(\sigma\) sur \(\mathrm{H}^*(\mathrm{A})\), on a \(v(\alpha)\geq0\), cette graduation est à degrés \(\geq0\). Si pour toute valeur propre \(\alpha\) de \(\sigma\) sur \(\mathrm{H}^i(\mathrm{A})\), on a \(v(\alpha)\geq i\) (et ce pour tout \(i\)), alors \(\mathcal{M}^n\otimes\bar{\mathrm{K}}\) est entièrement en degrés \(\geq n\). Si \(\Gamma\) est contenu dans \(\Gamma'\), stable par \(\operatorname{Gal}(\bar{\mathrm{K}}/\mathrm{K})\), et que \(v\) se prolonge en \(v':\Gamma'\to\mathbf{Z}\) invariant par Galois, cette graduation de \(\mathcal{M}\otimes\bar{\mathrm{K}}\) est définie sur \(\mathrm{K}\), i.e. provient d’une graduation de \(\mathcal{M}\).

\textbf{(5.3.2)} Le cas qui nous intéresse est celui où \(\mathrm{K}=\mathbf{Q}_{\ell}\), et où, pour un entier \(q>1\) convenable, puissance d’un nombre premier, on peut prendre pour \(\Gamma'\) l’ensemble des nombres algébriques \(\alpha\) dont tous les conjugués complexes ont la même valeur absolue, celle-ci étant de la forme \(q^{n(\alpha)/2}\) (une puissance entière de \(q^{1/2}\)). On prend \(v'(\alpha)=n(\alpha)\).

Les graduations par le poids \(\ell\)-adiques mentionnées dans [3] sont obtenues par cette méthode, pour \(\sigma\) un Frobenius. Voici quelques exemples.

\textbf{(5.3.3)} Soit \(\mathrm{X}\) une variété algébrique complexe. Elle peut être définie sur un sous-corps \(k_0\) de \(\mathbf{C}\), de type fini sur \(\mathbf{Q}\), i.e. est déduite de \(\mathrm{X}_0/k_0\) par extension des scalaires. Ceci permet de faire agir \(\operatorname{Gal}(\bar{k}_0/k_0)\) sur \(\mathrm{A}(\mathrm{X})\), défini comme en (5.2). Soit \(\mathcal{M}\) le modèle minimal de Sullivan du type d’homotopie rationnelle de l’espace topologique \(\mathrm{X}(\mathbf{C})\). D’après (5.2.5), \(\mathcal{M}\otimes\mathbf{Q}_{\ell}\) est un modèle minimal de \(\mathrm{A}(\mathrm{X})\). Plus précisément, on dispose d’une classe d’homotopie naturelle de quasi-isomorphismes \(\varphi\) de \(\mathcal{M}\otimes\mathbf{Q}_{\ell}\) avec \(\mathrm{A}(\mathrm{X})\);

% ---- 241-248 ----
\Page{241}
ceux-ci induisent en cohomologie les isomorphismes naturels de
\(\mathrm{H}^*(\mathrm{X}(\mathbf{C}),\mathbf{Q})\otimes\mathbf{Q}_\ell\)
avec \(\mathrm{H}^*(\mathrm{X},\mathbf{Q}_\ell)\). Pour tout automorphisme
\(\sigma\) de \(\mathrm{A}(\mathrm{X})\), on dispose donc d’un automorphisme
\(\widetilde{\sigma}\) de \(\mathscr{M}\otimes\mathbf{Q}_\ell\), bien déterminé
à homotopie près, tel que \(\varphi\widetilde{\sigma}\) soit homotope à
\(\sigma\varphi\). On peut choisir \(\widetilde{\sigma}\) de façon à respecter
une filtration comme en (5.3.1). Ceci s’applique en particulier à
\(\sigma\in\operatorname{Gal}(\bar{k}_0/k_0)\). Si l’on prend pour \(\sigma\)
un Frobenius comme dans [3], les hypothèses de (5.3.2) sont vérifiées par
\(\widetilde{\sigma}\), et la graduation correspondante de
\(\mathscr{M}\otimes\mathbf{Q}_\ell\) induit sur la cohomologie une graduation
qui scinde le tensorisé avec \(\mathbf{Q}_\ell\) de la filtration par le
poids de la théorie de Hodge sur
\(\mathrm{H}^*(\mathrm{X}(\mathbf{C}),\mathbf{Q})\)
\textit{(loc. cit.)}. Un raisonnement de Sullivan montre alors l’existence
d’une graduation de \(\mathscr{M}\) ayant les mêmes propriétés. Plus
précisément :

\textit{Théorème} \textbf{(5.3.4)}. — \textit{Il existe des graduations
\(\mathscr{W}\) de \(\mathscr{M}\) compatibles à la filtration par le poids
des \(\mathrm{H}^i(\mathscr{M})=\mathrm{H}^i(\mathrm{X}(\mathbf{C}),\mathbf{Q})\),
en ce sens que
\(\mathrm{W}_m\mathrm{H}^i(\mathrm{X}(\mathbf{C}),\mathbf{Q})
=\bigoplus\limits_{j\leq m}\mathscr{W}_j\mathrm{H}^i(\mathscr{M})\).
Deux telles graduations sont conjuguées par un automorphisme de
\(\mathscr{M}\) induisant l’identité sur
\(\operatorname{Gr}^{\mathrm{W}}\mathrm{H}^*(\mathrm{X}(\mathbf{C}),\mathbf{Q})\).}

Choisissons une filtration \((\mathscr{M}_n)\) de \(\mathscr{M}\) comme en
(5.3.1). Les graduations de \(\mathscr{M}\), compatibles à la filtration
par les \((\mathscr{M}_n)\), s’interprètent alors comme les homomorphismes
de \(\mathbf{G}_m\) dans le groupe proalgébrique
\(\mathrm{G}_1=\operatorname{Aut}((\mathscr{M},\mathscr{M}_n))
=\operatorname{lim\,proj}\operatorname{Aut}((\mathscr{M}_k))_{k\leq n}\)
des automorphismes de \(\mathscr{M}\) qui respectent la filtration de
\(\mathscr{M}\) par les \(\mathscr{M}_n\). Soit \(\mathrm{G}_2\) le
sous-groupe de \(\mathrm{G}_1\times\mathbf{G}_m\) formé des couples
\((g,\lambda)\) tels que :

\((*)\) \(g\) respecte la filtration par le poids des
\(\mathrm{H}^i(\mathrm{X}(\mathbf{C}),\mathbf{Q})\); sur
\(\operatorname{Gr}^{\mathrm{W}}_j\mathrm{H}^i\), il agit par
multiplication par \(\lambda^j\).

On vérifie que le noyau de l’homomorphisme
\((g,\lambda)\mapsto\lambda\) est unipotent :
\[
\begin{tikzcd}[column sep=small]
0 \arrow[r] &
\mathrm{U} \arrow[r] &
\mathrm{G}_2 \arrow[r,"\alpha"] &
\mathbf{G}_m .
\end{tikzcd}
\]
Après extension des scalaires à \(\mathbf{Q}_\ell\), le morphisme
\(\alpha\) admet une section. C’est l’existence des graduations (5.3.3).
Il est donc surjectif, comme morphisme de schémas, sur \(\mathbf{Q}\).
Le groupe proalgébrique \(\mathrm{G}_2\) est donc, sur \(\mathbf{Q}\),
extension de \(\mathbf{G}_m\) par \(\mathrm{U}\). On sait qu’une telle
extension est toujours triviale, et que deux sections en sont toujours
\(\mathrm{U}\)-conjuguées. Ceci prouve l’existence, et l’unicité lorsqu’on
ne considère que des graduations induisant une graduation sur chaque
\(\mathscr{M}_n\).

Pour conclure, il resterait à prouver le lemme technique suivant, dont
la vérification un peu pénible est laissée au lecteur.

\textit{Lemme} \textbf{(5.3.5)}. — \textit{Pour tout ensemble fini
d’automorphismes de \(\mathscr{M}\), il existe une filtration
\textup{(5.3.1)} de \(\mathscr{M}\) stable par tous ces automorphismes.}

\textbf{(5.3.6)} Pour \(\mathrm{X}\) une variété algébrique sur un corps
algébriquement clos quelconque, on peut encore définir une filtration
par le poids sur \(\mathrm{H}^*(\mathrm{X},\mathbf{Q}_\ell)\), en terme
des sous-espaces propres généralisés d’automorphismes de Frobenius
convenables, et les constructions précédentes fournissent des
graduations du modèle minimal de \(\mathrm{A}(\mathrm{X})\), induisant
sur la cohomologie une graduation qui scinde la filtration par le poids.

\Page{242}
\textit{Corollaire} \textbf{(5.3.7)}. — \textit{Pour \(\mathrm{X}\)
propre et lisse, un modèle minimal de \(\mathrm{A}(\mathrm{X})\) est un
modèle minimal de \(\mathrm{H}^*(\mathrm{X},\mathbf{Q}_\ell)\), muni de
la structure \(\mathrm{DG}\) pour laquelle \(d=0\).}

Soit \(\mathrm{W}\) une graduation par le poids d’un modèle minimal
\(\mathscr{M}\) de \(\mathrm{A}(\mathrm{X})\). Le groupe
\(\mathrm{H}^i(\mathrm{X},\mathbf{Q}_\ell)\) étant purement de poids
\(i\), \(\mathscr{M}^n\) est entièrement de degrés \(\geq n\)
((5.3.1)). On a donc
\(d\mathrm{W}_n(\mathscr{M}^n)\subset
\mathrm{W}_n(\mathscr{M}^{n+1})=0\); puisque
\(\mathrm{W}_n(\mathscr{M}^n)\) est formé de cycles, on dispose d’une
projection naturelle de \(\mathrm{W}_n(\mathscr{M}^n)\) dans
\(\mathrm{H}^n(\mathrm{X},\mathbf{Q}_\ell)\). Elle est surjective,
puisque \(\mathrm{H}^n\) est purement de poids \(n\). Soit
\(\varphi_n:\mathscr{M}^n\to\mathrm{H}^n(\mathrm{X},\mathbf{Q}_\ell)\)
le composé de cette projection avec la projection de \(\mathscr{M}^n\)
sur \(\mathrm{W}_n(\mathscr{M}^n)\) qui annule les
\(\mathrm{W}_i(\mathscr{M}^n)\), pour \(i>n\). On vérifie facilement
que \(\varphi=(\varphi_n)\) est un quasi-isomorphisme d’algèbre
\(\mathrm{DG}\), d’où le corollaire.

\textbf{(5.3.8)} Pour \(\mathrm{X}\) normal, on obtient de même un
théorème analogue au théorème de Morgan sur le système des quotients
de \(\pi_1(\mathrm{X})\) qui sont des groupes de Lie \(\ell\)-adiques
unipotents, cf. [9].

Signalons quelques questions passées sous silence.

\textbf{(5.3.9)} Pour \(\mathscr{M}\) un modèle minimal, et
\((\mathscr{M}_n)\) une filtration comme en (5.3.1), soit \(\mathrm{G}\)
le groupe proalgébrique des automorphismes de \(\mathscr{M}\) qui
respectent les \(\mathscr{M}_n\), et \(\mathrm{H}\) le sous-groupe de
ceux qui sont homotopes à l’identité. Posons
\(\operatorname{Aut}h(\mathscr{M})=\mathrm{G}/\mathrm{H}\).
C’est un groupe proalgébrique.

Si \(\mathrm{X}\) est une variété algébrique sur \(k\) algébriquement
clos, provenant par extension des scalaires de \(\mathrm{X}_0/k_0\)
dénombrable, l’action naturelle de \(\operatorname{Gal}(\bar{k}_0/k_0)\)
sur \(\mathrm{A}(\mathrm{X})\) fournit un homomorphisme de groupes :
\[
\tag*{\textbf{(5.3.9.1)}}
\operatorname{Gal}(\bar{k}_0/k_0)
\to\operatorname{Aut}h(\mathscr{M})(\mathbf{Q}_\ell).
\]
Il faudrait vérifier que cet homomorphisme est \textit{continu}, i.e.
que pour tout quotient algébrique (\(=\)de type fini) \(\mathrm{Q}\)
de \(\operatorname{Aut}h(\mathscr{M})\), l’homomorphisme de
\(\operatorname{Gal}(\bar{k}_0/k_0)\) dans \(\mathrm{Q}(\mathbf{Q}_\ell)\)
déduit de (5.3.9.1) est continu, si \(\mathrm{Q}(\mathbf{Q}_\ell)\)
est muni de sa topologie usuelle.

On peut espérer des résultats plus précis, mais je crois qu’il n’y a
\textit{pas} lieu d’espérer une action raisonnable de
\(\operatorname{Gal}(\bar{k}_0/k_0)\) sur \(\mathscr{M}\).

\textbf{(5.3.10)} Il faudrait vérifier que les filtrations croissantes
attachées aux graduations par le poids de (5.3.3) ou de (5.3.6) sont
toutes conjuguées entre elles sous le groupe des automorphismes de
\(\mathscr{M}\) (le modèle minimal \(\ell\)-adique) homotopes à l’identité.

Dans le cadre (5.3.3), où \(\mathscr{M}\) est le \(\ell\)-adifié du
modèle minimal rationnel \(\mathscr{M}_{\mathbf{Q}}\) de Sullivan,
leur classe de conjugaison devrait être définie sur \(\mathbf{Q}\)
(ou, ce qui revient au même, être l’orbite d’une filtration déduite
par extension des scalaires d’une filtration de
\(\mathscr{M}_{\mathbf{Q}}\)).

\Page{243}
\section*{VI. — LE FORMALISME DES FAISCEAUX MIXTES}

\section*{(6.1) Stabilité}

\textbf{(6.1.1)} Dans ce numéro, nous établissons dans les contextes
suivants que la catégorie des faisceaux mixtes est stable par toutes
les opérations usuelles.

a) \textit{Égale caractéristique \(p\)} : on ne considère que des
schémas de type fini sur \(\mathbf{F}_p\); \textit{faisceau} signifie
faisceau de Weil.

b) \textit{Égale caractéristique \(0\)} : il s’agit plus précisément
de considérer des schémas de type fini sur \(\mathbf{Z}\), et des
faisceaux (au sens : \(\overline{\mathbf{Q}}_\ell\)-faisceau
constructible), en se donnant à chaque instant le droit de remplacer
\(\mathrm{X}\) par \(\mathrm{X}[1/n]\).

En ce qui concerne les schémas, cela revient à considérer des schémas
de type fini sur \(\mathbf{Q}\) : leur catégorie est la 2-limite
inductive des catégories des schémas de type fini sur les
\(\mathbf{Z}[1/n]\). Il n’en va pas de même pour les faisceaux :
si \(\mathrm{X}\) est de type fini sur \(\mathbf{Z}\), j’ignore si un
faisceau constructible sur \(\mathrm{X}_{\mathbf{Q}}\) provient toujours
d’un faisceau constructible sur un \(\mathrm{X}[1/n]\). Si
\(\mathscr{F}\) et \(\mathscr{G}\) sont deux faisceaux constructibles
sur \(\mathrm{X}\), on a toutefois :
\[
\operatorname{Hom}\text{ sur }\mathrm{X}_{\mathbf{Q}}(\mathscr{F},\mathscr{G})
=\operatorname{lim\,ind}\operatorname{Hom}\text{ sur }
\mathrm{X}[1/n](\mathscr{F},\mathscr{G}).
\]
En d’autres termes, la 2-limite inductive des catégories de faisceaux
sur les \(\mathrm{X}[1/n]\) s’identifie à une sous-catégorie pleine
de celle des faisceaux sur \(\mathrm{X}_{\mathbf{Q}}\).

La terminologie suivante permet de donner des énoncés uniformes dans
les contextes a) et b) : pour \(\mathrm{X}_{\mathbf{Q}}\) un schéma
de type fini sur \(\mathbf{Q}\), déduit par changement de base de
\(\mathrm{X}\) de type fini sur \(\mathbf{Z}\), et
\(\mathscr{F}_{\mathbf{Q}}\) un faisceau sur \(\mathrm{X}_{\mathbf{Q}}\),
nous dirons que \(\mathscr{F}_{\mathbf{Q}}\) est mixte, pur, ... si
\(\mathscr{F}_{\mathbf{Q}}\) se déduit d’un faisceau \(\mathscr{F}\)
sur un \(\mathrm{X}[1/n]\), et que la restriction de \(\mathscr{F}\)
à \(\mathrm{X}[1/nm]\) est mixte, pure..., pour \(m\) convenable.

Les résultats obtenus dans le contexte b) seront présentés sous la
forme de variantes de ceux obtenus dans le contexte a). On pourrait
aussi donner des variantes potentielles ((3.4.10)); nous ne les
expliciterons pas.

\textit{Théorème} \textbf{(6.1.2)}. — \textit{Soient
\(f:\mathrm{X}\to\mathrm{Y}\) un morphisme de schémas de type fini
sur \(\mathbf{F}_p\) (resp. \(\mathbf{Q}\)), et \(\mathscr{F}\) un
faisceau mixte sur \(\mathrm{X}\). Alors, les faisceaux
\(\mathrm{R}^if_*\mathscr{F}\) sont mixtes.}

Compte tenu du théorème de changement de base générique
(SGA4\(\frac{1}{2}\) [Th. finitude] (1.9)), cet énoncé résulte du
suivant :

\Page{244}
\textit{Lemme} \textbf{(6.1.3)}. — \textit{Soient \(\mathrm{S}\)
un schéma de type fini sur \(\mathbf{Z}[1/\ell]\),
\(f:\mathrm{X}\to\mathrm{Y}\) un morphisme de \(\mathrm{S}\)-schémas
de type fini, et \(\mathscr{F}\) un faisceau mixte sur \(\mathrm{X}\).
Alors, il existe un ouvert dense \(\mathrm{U}\) de \(\mathrm{S}\)
au-dessus duquel les faisceaux \(\mathrm{R}^if_*\mathscr{F}\) sont
mixtes.}

Reprenons la démonstration de SGA4\(\frac{1}{2}\) [Th. finitude]
(2.1) à (2.8).

\textbf{(6.1.4)} \textit{Cas où \(\mathrm{X}\) est lisse sur
\(\mathrm{Y}=\mathrm{S}\), où \(\mathscr{F}\) est lisse, et où les
\(\mathrm{R}^if_!\mathscr{F}\) sont lisses.} Dans ce cas, les
\(\mathrm{R}^if_*\mathscr{F}\) se déduisent par dualité de Poincaré
des \(\mathrm{R}^if_!\mathscr{F}\), auxquels (3.3.1) s’applique.

\textbf{(6.1.5)} \textit{Cas où \(\mathrm{X}\) est lisse sur
\(\mathrm{Y}=\mathrm{S}\) et où \(\mathscr{F}\) est lisse.}

On se ramène à (6.1.4) en remplaçant \(\mathrm{S}\) par un ouvert
dense convenable.

\textbf{(6.1.6)} Prouvons par récurrence sur \(n\) l’assertion :

\((*)_n\) La conclusion de (6.1.3) est vraie si \(\mathrm{S}\) est
intègre, de point générique \(\eta\), que \(f\) est un plongement
ouvert d’image dense, et que \(\dim\mathrm{X}_\eta\leq n\).

Pour \(n=0\), quitte à rétrécir \(\mathrm{S}\), on a
\(\mathrm{X}=\mathrm{Y}\) : \((*)_0\) est évident. Supposons
\((*)_{n-1}\), et prouvons \((*)_n\). Dans \((*)_{n-1}\), on peut
remplacer « plongement ouvert » par « plongement » comme on le voit
en factorisant en plongement ouvert et plongement fermé.

\textit{Lemme} \textbf{(6.1.7)}. — \textit{Quitte à rétrécir
\(\mathrm{S}\), la conclusion de \textup{(6.1.3)} vaut sur
\(\mathrm{Y}'\subset\mathrm{Y}\), le complément \(\mathrm{Y}_1\)
de \(\mathrm{Y}'\) étant fini sur \(\mathrm{S}\).}

L’assertion est locale sur \(\mathrm{Y}\), qu’on peut supposer
affine : \(\mathrm{Y}\subset\mathbf{A}_{\mathrm{S}}^n\).
L’hypothèse de récurrence \((*)_{n-1}\) s’applique à :
\[
\begin{tikzcd}[column sep=small,row sep=small]
\mathrm{X} \arrow[rr,"f"] \arrow[dr] &&
\mathrm{Y} \arrow[dl,"\mathrm{pr}_i"] \\
& \mathbf{A}_{\mathrm{S}}^1 &
\end{tikzcd}
\]
Il existe donc pour chaque \(i\) un ouvert dense \(\mathrm{U}_i\)
de \(\mathrm{A}_{\mathrm{S}}^1\) tel que la conclusion de (6.1.3)
vaille au-dessus de \(\mathit{pr}_i^{-1}(\mathrm{U}_i)\); elle vaut
au-dessus de la réunion des \(\mathit{pr}_i^{-1}(\mathrm{U}_i)\),
et (6.1.7) en résulte.

\textbf{(6.1.8)} Prouvons \((*)_n\) pour \(\mathrm{X}\) lisse sur
\(\mathrm{S}\) et \(\mathscr{F}\) lisse. Le problème étant local
sur \(\mathrm{Y}\), on peut supposer \(\mathrm{Y}\) affine, puis
projectif (remplacer \(\mathrm{Y}\) par son adhérence dans un espace
projectif). Soient \(i:\mathrm{Y}_1\hookrightarrow\mathrm{Y}\) et
\(j:\mathrm{Y}'\hookrightarrow\mathrm{Y}\) garantis par (6.1.7)
(après rétrécissement de \(\mathrm{S}\)).
\[
\begin{tikzcd}[column sep=small,row sep=small]
\mathrm{X} \arrow[r,hook,"f"] \arrow[dr] &
\mathrm{Y} \arrow[d,"b"] &
\mathrm{Y}_1 \arrow[l,"i"'] \arrow[dl] \\
& \mathrm{S} &
\end{tikzcd}
\]

\Page{245}
Le raisonnement suivant, où apparaît formellement une catégorie
dérivée de faisceaux \(\ell\)-adiques, sera justifié ci-dessous.
Appliquons \(\mathrm{R}b_*\) au triangle défini par la suite exacte :
\[
\tag*{(1)}
0\to j_!j^*\mathrm{R}f_*\mathscr{F}
\to\mathrm{R}f_*\mathscr{F}
\to i_!i^*\mathrm{R}f_*\mathscr{F}\to0
\]
on obtient un triangle :
\[
\tag*{(2)}
\to\mathrm{R}b_*j_!j^*\mathrm{R}f_*\mathscr{F}
\to\mathrm{R}a_*\mathscr{F}
\to b_{1*}i^*\mathrm{R}f_*\mathscr{F}\to
\]
dans lequel les faisceaux de cohomologie des deux premiers termes
sont mixtes sur un ouvert dense de \(\mathrm{S}\), d’après (3.3.1)
pour le premier (\(b\) est propre) et (6.1.5) pour le second.
Les faisceaux de cohomologie de \(b_{1*}i^*\mathrm{R}f_*\mathscr{F}\)
sont donc mixtes, et on en déduit que ceux de
\(i^*\mathrm{R}f_*\mathscr{F}\), puis ceux de
\(\mathrm{R}f_*\mathscr{F}\), le sont également.

Pour justifier cela, on écrit \(\mathscr{F}\) comme déduit d’une
limite projective d’un système projectif de faisceaux de
\(\mathrm{R}/m^n\)-modules libres localement constants, se déduisant
les uns des autres par réduction, comme en (1.1.1), on considère le
système projectif correspondant de triangles (1) et (2), et on
répète l’argument.

\textbf{(6.1.9)} Prouvons \((*)_n\) en général. L’usage qui sera fait
d’une catégorie dérivée se justifie comme ci-dessus.

On commence par se ramener au cas où dans \(\mathrm{X}\) existe
un ouvert dense \(\mathrm{V}\) lisse sur \(\mathrm{S}\). Pour
\(\mathrm{S}\) spectre d’un corps parfait, il suffit de remplacer
\(\mathrm{X}\) par \(\mathrm{X}_{\mathrm{red}}\) (et \(\mathrm{Y}\)
par \(\mathrm{Y}_{\mathrm{red}}\)). En général, il faut rapetisser
\(\mathrm{S}\), faire un changement de base fini radiciel et
surjectif \(\mathrm{S}'\to\mathrm{S}\), et remplacer \(\mathrm{X}\)
et \(\mathrm{Y}\) par \(\mathrm{X}_{\mathrm{red}}\) et
\(\mathrm{Y}_{\mathrm{red}}\). La topologie étale étant insensible
aux morphismes finis radiciels et surjectifs, ceci est innocent.
Quitte à rétrécir \(\mathrm{V}\), on peut supposer \(\mathscr{F}\)
lisse sur \(\mathrm{V}\) :
\[
\begin{tikzcd}[column sep=small]
\mathrm{V} \arrow[r,hook,"j"] &
\mathrm{X} \arrow[r,"f"] &
\mathrm{Y}.
\end{tikzcd}
\]
Ceci permet d’appliquer (6.1.8) à \(j\) et \(fj\). Définissons
\(\Delta\) par le triangle :
\[
\tag*{(1)}
\to\mathscr{F}\to\mathrm{R}j_*j^*\mathscr{F}\to\Delta\to.
\]
Les faisceaux de cohomologie de \(\Delta\) sont à support dans
\(\mathrm{X}-\mathrm{V}\), et
\(\dim(\mathrm{X}-\mathrm{V})_\eta<n\). L’hypothèse de récurrence
permet donc de supposer que \(\mathrm{R}f_*\Delta\) est mixte.
Appliquons \(\mathrm{R}f_*\) au triangle (1); on trouve un triangle :
\[
\to\mathrm{R}f_*\mathscr{F}
\to\mathrm{R}(fj)_*j^*\mathscr{F}
\to\mathrm{R}f_*\Delta\to
\]
dans lequel deux des sommets sont mixtes. Le troisième l’est donc
également.

\textbf{(6.1.10)} Prouvons (6.1.2). Le problème est local sur
\(\mathrm{Y}\), qu’on peut supposer affine. Prenant un recouvrement
affine de \(\mathrm{X}\) et invoquant sa suite spectrale de Leray,
on se ramène à avoir \(\mathrm{X}\) également affine. On peut alors
factoriser \(f\) en un plongement ouvert suivi d’un morphisme propre :
\(f=gj\), d’où \(\mathrm{R}f_*=\mathrm{R}g_*\mathrm{R}j_*\).
Le plongement ouvert est justiciable d’un \((*)_n\), le morphisme
propre de (3.3.1).

\Page{246}
On a, comme conséquence formelle de (6.1.2)
(cf. SGA4\(\frac{1}{2}\) [Th. finitude] (1.5)) :

\textit{Corollaire} \textbf{(6.1.11)}. — \textit{Les quatre opérations
\(\mathrm{R}^if_*\), \(\mathrm{R}^if_!\), \(f^*\) et
\(\mathrm{R}^if^!\) transforment faisceaux mixtes en faisceaux mixtes.}

La même stabilité vaut pour les foncteurs \(\operatorname{Ext}\) locaux.

\textbf{(6.1.12)} Nous allons maintenant déduire de (6.1.2) une
stabilité analogue pour les faisceaux de cycles évanescents.
Soient donc \(\mathrm{S}\) une courbe lisse sur \(\mathbf{F}_q\)
(resp. \(\mathbf{Q}\)), \(s\) un point fermé de \(\mathrm{S}\),
\(\mathrm{X}\) un \(\mathrm{S}\)-schéma de type fini et \(\mathscr{F}\)
un faisceau sur \(\mathrm{X}\). Faisant le changement de base
\(\mathrm{S}_{(s)}\to\mathrm{S}\), on obtient
\((\mathrm{X}',\mathscr{F}')\) au-dessus du trait hensélien
\(\mathrm{S}_{(s)}\). Soient \(\eta\) le point générique de
\(\mathrm{S}_{(s)}\), \(\bar{s}\) un point géométrique au-dessus de
\(s\) et \(\bar{\eta}\) un point géométrique générique du localisé
strict \(\mathrm{S}_{(\bar{s})}\). Les \textit{faisceaux de cycles
évanescents} \(\mathrm{R}^i\Psi(\mathscr{F}')\) sont des faisceaux
sur \(\mathrm{X}_{\bar{s}}\), munis d’une action continue de
\(\operatorname{Gal}(\bar{\eta},\eta)\), au-dessus de son action sur
\(\mathrm{X}_{\bar{s}}\) (via \(\operatorname{Gal}(\bar{s}/s)\)).
La restriction de cette action au groupe d’inertie est
automatiquement quasi unipotente.

Définissons, pour un faisceau \(\mathscr{G}\), muni d’une telle
action \(\rho\) de \(\operatorname{Gal}(\bar{\eta}/\eta)\), ce que
signifie être \textit{mixte} :

\textit{a)} Un faisceau sur \(\mathrm{X}_{\bar{s}}\) muni d’une action
continue de \(\operatorname{Gal}(\bar{s}/s)\) s’identifie à un
faisceau sur \(\mathrm{X}_s\) (SGA7 XIII (1.1.3)). Si l’action
\(\rho\) se factorise par \(\operatorname{Gal}(\bar{s}/s)\), on dit
que \((\mathscr{G},\rho)\) est mixte si le faisceau correspondant
sur \(\mathrm{X}_s\) l’est.

\textit{b)} Si \((\mathscr{G},\rho)\) admet une filtration finie
\(\mathrm{F}\) telle que \(\operatorname{Gr}_{\mathrm{F}}(\rho)\)
se factorise par \(\operatorname{Gal}(\bar{s}/s)\) (cas unipotent),
on dit que \((\mathscr{G},\rho)\) est mixte si
\(\operatorname{Gr}_{\mathrm{F}}(\mathscr{G})\) l’est, au sens
\textit{a)}.

\textit{c)} L’hypothèse \textit{b)} devient vraie après avoir remplacé
\(\operatorname{Gal}(\bar{\eta}/\eta)\) par un sous-groupe d’indice
fini (quasi-unipotence). Géométriquement, ceci revient à remplacer
\(\mathrm{S}_{(s)}\) par un trait fini sur \(\mathrm{S}_{(s)}\).
On dit que \((\mathscr{G},\rho)\) est mixte s’il le devient — au
sens \textit{b)} — après un tel changement de trait.

\textit{Théorème} \textbf{(6.1.13)}. — \textit{Avec les notations
ci-dessus, si \(\mathscr{F}\) est mixte, les faisceaux de cycles
évanescents sont également mixtes.}

Grâce au théorème de monodromie, quitte à remplacer
\(\mathrm{S}_{(s)}\) par un revêtement fini convenable, on peut
supposer que le groupe d’inertie \(\mathrm{I}\) agit de façon
unipotente, via son quotient \(\mathbf{Z}_\ell(1)\). On peut alors
appliquer les constructions des numéros (1.6.1), (1.6.14).
Les formules (1.6.14) montrent qu’il suffit de vérifier que les
faisceaux \((\mathrm{R}^i\Psi\mathscr{F}')^{\mathrm{I}}\) sont
mixtes. Soient \(u\) l’inclusion de \(\mathrm{X}_s\) dans
\(\mathrm{X}\), et \(v\) celle de l’ouvert complémentaire.
Procédant comme en (3.6.1), on voit que ces faisceaux sont quotients
des faisceaux \(u^*\mathrm{R}^iv_*(v^*\mathscr{F})\), et on applique
(6.1.11).

\Page{247}
\section*{(6.2) Complexes purs}

Dans ce numéro, nous nous plaçons dans le cadre (6.1.1) \textit{a)}.

\textbf{(6.2.1)} Pour tout schéma \(\mathrm{X}\) de type fini sur
\(\mathbf{F}_p\) : \(a:\mathrm{X}\to\operatorname{Spec}(\mathbf{F}_p)\),
nous posons \(\mathrm{K}_{\mathrm{X}}=\mathrm{R}a^!\overline{\mathbf{Q}}_\ell\).
C’est le \textit{complexe dualisant}.

\textit{Exemple} : si \(\mathrm{X}\) est lisse purement de dimension
\(n\), \(\mathrm{K}_{\mathrm{X}}\) est réduit au faisceau
\(\overline{\mathbf{Q}}_\ell(n)\) placé en degré \(-2n\) :
\(\mathrm{K}_{\mathrm{X}}=\overline{\mathbf{Q}}_\ell(n)[2n]\).

Nous notons \(\mathrm{D}\) le foncteur contravariant « dual à valeur
dans le complexe dualisant » :
\[
\mathrm{D}:\mathrm{D}_c^b(\mathrm{X})\to\mathrm{D}_c^b(\mathrm{X}):
\quad \mathrm{K}\to\mathrm{R}\,\mathit{Hom}(\mathrm{K},\mathrm{K}_{\mathrm{X}}).
\]
Ce foncteur est involutif. Pour tout morphisme
\(f:\mathrm{X}\to\mathrm{Y}\), il échange les foncteurs
\(\mathrm{R}f^!\) et \(f^*\), ainsi que les foncteurs
\(\mathrm{R}f_!\) et \(\mathrm{R}f_*\). On a aussi :
\[
\mathrm{D}(\mathrm{K}\otimes\mathrm{L})
=\mathrm{R}\,\mathit{Hom}(\mathrm{K},\mathrm{D}\mathrm{L}).
\]

\textit{Définition} \textbf{(6.2.2)}. — \textit{Un complexe
\(\mathrm{K}\)} est mixte de poids \(\leq n\) \textit{si pour tout
\(i\), le faisceau \(\mathscr{H}^i\mathrm{K}\) est mixte de poids
ponctuels \(\leq i+n\).}

Avec cette définition, le théorème fondamental (3.3.1) fournit la

\textit{Variante} \textbf{(6.2.3)}. — \textit{Soit
\(f:\mathrm{X}\to\mathrm{Y}\) un morphisme de schémas de type fini
sur \(\mathbf{F}_p\). Si
\(\mathrm{K}\in\operatorname{Ob}\mathrm{D}_c^b(\mathrm{X})\) est
mixte de poids \(\leq n\), alors \(\mathrm{R}f_!\mathrm{K}\) est
mixte de poids \(\leq n\).}

Soit en effet la suite spectrale :
\[
\mathrm{E}_1^{pq}=\mathrm{R}^pf_!\mathscr{H}^q\mathrm{K}
\Rightarrow\mathscr{H}^{p+q}\mathrm{R}f_!\mathrm{K}.
\]
D’après (3.3.1), le terme \(\mathrm{E}_1^{pq}\) est mixte de poids
ponctuels \(\leq p+(q+n)\).

L’aboutissement \(\mathscr{H}^{p+q}\mathrm{R}f_!\mathrm{K}\) est
donc mixte de poids ponctuels \(\leq(p+q)+n\).

\textit{Définition} \textbf{(6.2.4)}. — \textit{Un complexe
\(\mathrm{K}\)} est pur de poids \(n\) \textit{s’il est mixte de
poids \(\leq n\), et que \(\mathrm{D}\mathrm{K}\) est mixte de
poids \(\leq-n\). Un faisceau \(\mathscr{F}\) est pur de poids
\(n\) si le complexe réduit à \(\mathscr{F}\) en degré \(0\) l’est.}

\textit{Exemples} \textbf{(6.2.5)}. — \textit{a)} Quel que soit
\(\mathrm{N}\), un complexe \(\mathrm{K}\) est mixte de poids
\(\leq n\) si et seulement si \(\mathrm{K}(\mathrm{N})[2\mathrm{N}]\)
l’est. Dans la définition (6.2.4) de la pureté, on peut donc remplacer
le complexe dualisant \(\mathrm{K}_{\mathrm{X}}\) par un complexe
localement de la forme \(\mathrm{K}_{\mathrm{X}}(\mathrm{N})[2\mathrm{N}]\),
et \(\mathrm{D}\mathrm{K}\) par
\(\mathrm{D}\mathrm{K}(\mathrm{N})[2\mathrm{N}]\).

\textit{b)} En particulier, pour \(\mathrm{X}\) lisse, on peut
remplacer la condition sur \(\mathrm{D}\mathrm{K}\) par :
« \(\mathrm{R}\,\mathit{Hom}(\mathrm{K},\overline{\mathbf{Q}}_\ell)\)
est mixte de poids \(\leq-n\) ». Si les faisceaux de cohomologie de
\(\mathrm{K}\in\operatorname{Ob}\mathrm{D}_c^b(\mathrm{X})\) sont
lisses, le complexe \(\mathrm{K}\) est pur de poids \(n\) si et
seulement si chaque \(\mathscr{H}^i\mathrm{K}\) est ponctuellement
pur de poids \(i+n\).

\Page{248}
\textit{c)} Soient \(\mathrm{X}\) une courbe lisse, et
\(j:\mathrm{U}\hookrightarrow\mathrm{X}\) un ouvert dense de
\(\mathrm{X}\). Si \(\mathscr{F}\) est un faisceau pur de poids
\(n\) sur \(\mathrm{U}\), alors \(j_*\mathscr{F}\) est pur de poids
\(n\) (appliquer (1.8.8.1) et la formule
\(\mathrm{R}\,\mathit{Hom}(j_*\mathscr{F},\mathbf{Q}_\ell)
=j_*\,\mathit{Hom}(\mathscr{F},\mathbf{Q}_\ell)\)).

\textit{d)} Pour \(\mathrm{X}\) lisse,
\(j:\mathrm{U}\hookrightarrow\mathrm{X}\) le complément d’un
diviseur lisse \(\mathrm{D}\), et \(\mathscr{F}\) un faisceau lisse
ponctuellement pur de poids \(n\) sur \(\mathrm{U}\), modérément
ramifié le long de \(\mathrm{D}\), le faisceau \(j_*\mathscr{F}\)
est encore pur de poids \(n\) (même argument).

\textit{Proposition} \textbf{(6.2.6)}. — \textit{Si
\(f:\mathrm{X}\to\mathrm{Y}\) est propre, et que
\(\mathrm{K}\in\operatorname{Ob}\mathrm{D}_c^b(\mathrm{X})\) est
pur, alors \(\mathrm{R}f_*\mathrm{K}\) est pur du même poids.}

On applique (6.2.3) à \(\mathrm{K}\) et à \(\mathrm{D}\mathrm{K}\),
compte tenu de ce que
\(\mathrm{D}\mathrm{R}f_*\mathrm{K}
=\mathrm{R}f_*\mathrm{D}\mathrm{K}\) puisque \(f\) est propre.

\textit{Variante} \textbf{(6.2.7)}. — Pour tout schéma \(\mathrm{X}\)
de type fini sur \(\mathbf{Z}[1/\ell]\) :
\[
a:\mathrm{X}\to\operatorname{Spec}(\mathbf{Z}[1/\ell]),
\]
posons \(\mathrm{K}'_{\mathrm{X}}=\mathrm{R}a^!\overline{\mathbf{Q}}_\ell\).
On a encore un foncteur de passage au dual
\(\mathrm{D}'\mathrm{K}
=\mathrm{R}\operatorname{Hom}(\mathrm{K},\mathrm{K}'_{\mathrm{X}})\),
et une définition de la pureté parallèle à (6.2.4). Si
\(\mathrm{X}\) est de type fini sur \(\mathbf{F}_p\), on a
\(\mathrm{K}'_{\mathrm{X}}=\mathrm{K}_{\mathrm{X}}(-1)[-2]\) :
les foncteurs \(\mathrm{D}\) et \(\mathrm{D}'\) se déduisent l’un
de l’autre par décalage et twist et, d’après (6.2.5) \textit{a)},
la définition ci-dessus de « pur de poids \(n\) » généralise
(6.2.4). La proposition (6.2.6) reste vraie.

\textbf{(6.2.8)} Soient \(k\) un corps algébriquement clos,
\(\mathrm{S}\) le spectre de l’hensélisé de \(k[\mathrm{T}]\) en
\((\mathrm{T})\), \(s\) le point fermé de \(\mathrm{S}\), \(\eta\)
le point générique et \(\bar{\eta}\) un point générique géométrique.
Soit \(f:\mathrm{X}\to\mathrm{S}\) un morphisme propre.

Un complexe \(\mathrm{K}\in\operatorname{Ob}\mathrm{D}_c^b(\mathrm{X})\)
est dit \textit{potentiellement pur} si
\((\mathrm{X}/\mathrm{S},\mathrm{K})\) provient d’une situation
arithmétique du type suivant :

\textit{a)} \(\operatorname{Spec}(k)\) est un point géométrique
générique d’un schéma intègre \(\mathrm{A}_0\), de type fini sur
\(\mathbf{Z}[1/\ell]\);

\textit{b)} \(\mathrm{S}\) provient par extension des scalaires à
\(k\), et hensélisation, d’une courbe lisse \(\mathrm{S}_0\) sur
\(\mathrm{A}_0\) munie d’une section;

\textit{c)} \(f\) provient d’un morphisme propre
\(f_0:\mathrm{X}_0\to\mathrm{S}_0\);

\textit{d)} \(\mathrm{K}\) provient d’un complexe pur
\(\mathrm{K}_0\in\operatorname{Ob}\mathrm{D}_c^b(\mathrm{X}_0)\).

Le théorème suivant généralise (3.6.1) :

\textit{Théorème} \textbf{(6.2.9)}
\textit{(théorème local des cycles invariants, pour les complexes purs)}.
— \textit{Avec les notations ci-dessus, soit \(\mathrm{K}\) un
complexe potentiellement pur, sur \(\mathrm{X}\). Alors, le morphisme
de spécialisation :}
\[
\mathit{sp}^*:\mathrm{H}^*(\mathrm{X}_s,\mathrm{K})
\to\mathrm{H}^*(\mathrm{X}_{\bar{\eta}},\mathrm{K})^{
\operatorname{Gal}(\bar{\eta}/\eta)}
\]
\textit{est surjectif.}

% ---- 249-252 ----
\Page{249}
La démonstration est parallèle à celle de (3.6.1). On dispose d’une croix (3.6.1) (8)
\[
\begin{tikzcd}[column sep=1.3em, row sep=1.5em]
& & \mathrm{H}^{i+1}_{\mathrm{X}_s}(\mathrm{X},\mathrm{K}) & & \\
0 \arrow[r]
& \mathrm{H}^{i-1}(\mathrm{X}_{\bar{\eta}},\mathrm{K})_{\mathrm{I}}(-1) \arrow[r]
& \mathrm{H}^{i}(\mathrm{X}_{\eta},\mathrm{K}) \arrow[u] \arrow[r]
& \mathrm{H}^{i}(\mathrm{X}_{\bar{\eta}},\mathrm{K})^{\mathrm{I}} \arrow[r]
& 0 \\
& & \mathrm{H}^{i}(\mathrm{X}_s,\mathrm{K})
    \arrow[u]
    \arrow[ur, dashed, "\mathrm{sp}^{*}"']
& &
\end{tikzcd}
\]
et \(\mathrm{H}^{i+1}_{\mathrm{X}_s}(\mathrm{X},\mathrm{K})
=[\mathrm{H}^{-i-1}(\mathrm{X}_s,\mathrm{DK})]^{\vee}\).
Ceci exprime que \(\mathrm{D}\) échange les foncteurs \(\mathrm{R}i^!\) et \(i^*\), pour \(i\) l’inclusion de \(\mathrm{X}_s\) dans \(\mathrm{X}\), et commute à \(\mathrm{R}\Gamma(\mathrm{X}_s,\ )\). L’analogue de (3.6.2) (resp. (3.6.3)) se prouve en remplaçant la référence à (3.3.9) (resp. (3.3.4)) par une référence à (6.2.6) (resp. (6.2.3)) et on conclut comme en (3.6).

\textbf{(6.2.10)} Dans la fin de ce numéro, nos arguments seront géométriques, sauf une référence à (6.2.9) et une référence à (3.4).

Soit donc \(k\) un corps algébriquement clos. Pour tout schéma \(\mathrm{X}\) de type fini sur \(k\) : \(a:\mathrm{X}\to\operatorname{Spec}(k)\) on définit comme en (6.2.1) \(\mathrm{K}_{\mathrm{X}}=\mathrm{R}a^!\overline{\mathbf{Q}}_{\ell}\) et \(\mathrm{D}(\mathrm{K})=\mathrm{R}\,\mathit{Hom}(\mathrm{K},\mathrm{K}_{\mathrm{X}})\). On a le même formalisme qu’en (6.2.1).

Soit \(\mathrm{X}\subset\mathbf{P}\) une variété projective. Les sections hyperplanes de \(\mathrm{X}\) sont paramétrées par les points de l’espace projectif dual \(\check{\mathbf{P}}\), et chaque droite de \(\check{\mathbf{P}}\) paramétrise un pinceau de sections hyperplanes :
\[
\begin{tikzcd}[column sep=1.5em, row sep=1.5em]
\mathrm{X}
& \mathrm{Z} \arrow[l, "q"'] \arrow[d, "f"]
& \widetilde{\mathrm{X}} \arrow[l, hook', "i"'] \arrow[d, "f"] \\
& \check{\mathbf{P}}
& \mathrm{D} \arrow[l, hook', "i"']
\end{tikzcd}
\tag*{\textbf{(6.2.10.1)}}
\]
Pour \(\mathrm{K}\) dans \(\mathrm{D}_c^b(\mathrm{X})\), les \(\mathrm{R}^if_*q^*\mathrm{K}\) sont lisses sur un ouvert dense \(\mathrm{U}\) de \(\check{\mathbf{P}}\). On appellera \textit{générale} une section hyperplane \(\mathrm{Y}_u\) correspondant à \(u\in\mathrm{U}\). Le groupe fondamental \(\pi_1(\mathrm{U},u)\) agit sur \(\mathrm{H}^*(\mathrm{Y}_u,\mathrm{K})\). Si \(\mathrm{D}\) est assez générale, et que \(u\in\mathrm{V}=\mathrm{U}\cap\mathrm{D}\), l’image de \(\pi_1(\mathrm{V},u)\) dans \(\mathrm{GL}\,\mathrm{H}^*(\mathrm{Y}_u,\mathrm{K})\) coïncide avec celle de \(\pi_1(\mathrm{U},u)\). En particulier, on a :
\[
\mathrm{H}^*(\mathrm{Y}_u,\mathrm{K})^{\pi_1(\mathrm{U},u)}
=
\mathrm{H}^*(\mathrm{Y}_u,\mathrm{K})^{\pi_1(\mathrm{V},u)}
\qquad
\text{(\(\mathrm{D}\) assez général).}
\tag*{\textbf{(6.2.10.2)}}
\]

\medskip
\textit{Proposition} \textbf{(6.2.11)}. — \textit{Avec les notations précédentes, soit \(n\) un entier (dans les bons cas, la dimension de \(\mathrm{Y}_u\)), et supposons que pour chaque \(i\), le support de \(\mathscr{H}^i(\mathrm{DK}[-2n-2])\) est de dimension \(\leq n+1-i\). Alors :}

(i) \textit{Pour toute section hyperplane \(\mathrm{Y}\), on a :}
\begin{align*}
& \mathrm{H}^i(\mathrm{X},\mathrm{K})
  \xrightarrow{\sim}
  \mathrm{H}^i(\mathrm{Y},\mathrm{K})
  \qquad \text{\textit{pour}} \qquad i<n \\
\text{\textit{et}}\qquad
& \mathrm{H}^n(\mathrm{X},\mathrm{K})
  \hookrightarrow
  \mathrm{H}^n(\mathrm{Y},\mathrm{K}).
\end{align*}

(ii) \textit{Pour \(\mathrm{D}\) un pinceau général, et \(u\in\mathrm{D}\), l’image de \(\mathrm{H}^n(\mathrm{X},\mathrm{K})\) dans \(\mathrm{H}^n(\mathrm{Y}_u,\mathrm{K})\) coïncide avec l’image de \(\mathrm{H}^0(\mathrm{D},\mathrm{R}^nf_*(i^*q^*\mathrm{K}))\).}

\Page{250}
La preuve est parallèle à celle de (4.1.6). On vérifie d’abord le théorème de Lefschetz faible \(\mathrm{H}_c^i(\mathrm{X}-\mathrm{Y},\mathrm{K})=0\) pour \(i\leq n\). Par dualité, cette nullité se ramène à celle de \(\mathrm{H}^{-i}(\mathrm{X}-\mathrm{Y},\mathrm{DK})\) pour \(-i\geq -n\), qu’on a prise pour hypothèse. L’assertion (i) en résulte, par la suite exacte longue de cohomologie.

Pour \(\mathrm{Y}\) une section hyperplane assez générale, la restriction de \(\mathrm{K}\) à \(\mathrm{Y}\) vérifie l’hypothèse, avec \(n\) remplacé par \(n-1\). Cela résulte du théorème d’acyclicité générique SGA4\(\frac{1}{2}\) [Th. finitude] pour \(\mathrm{Z}/\check{\mathbf{P}}\) et \(q^*\mathrm{K}\) : on a \(\mathrm{D}(\mathrm{K}|\mathrm{Y})=((\mathrm{DK})|\mathrm{Y})(-1)[-2]\). Ceci permet d’appliquer (4.3.7) comme en (4.3.8), et de conclure comme en (4.3.9).

\medskip
\textit{Corollaire} \textbf{(6.2.12)} \textit{(théorème global des cycles invariants, pour les complexes purs).} — \textit{Sous les hypothèses de \textup{(6.2.11)}, si \(\mathrm{K}\) est potentiellement pur, alors, pour \(\mathrm{Y}_u\) une section hyperplane générale de \(\mathrm{X}\), \(\mathrm{H}^n(\mathrm{X},\mathrm{K})\) s’identifie au sous-groupe \(\mathrm{H}^n(\mathrm{Y}_u,\mathrm{K})^{\pi_1(\mathrm{U},u)}\) des invariants de la monodromie dans \(\mathrm{H}^n(\mathrm{Y}_u,\mathrm{K})\).}

\medskip
Il suffit de le vérifier pour une section hyperplane assez générale. Pour un pinceau \(\mathrm{D}\) assez général, \(i^*q^*\mathrm{K}\) sur \(\widetilde{\mathrm{X}}\) est potentiellement pur (on a \(\mathrm{D}i^*q^*\mathrm{K}=i^*q^*\mathrm{DK}\)). On peut donc appliquer (6.2.9) en chaque point de \(\mathrm{D}-\mathrm{V}\), et on trouve que :
\[
\mathrm{H}^0(\mathrm{D},\mathrm{R}^nf_*(i^*q^*\mathrm{K}))
\twoheadrightarrow
\mathrm{H}^n(\mathrm{Y}_u,\mathrm{K})^{\pi_1(\mathrm{V},u)}.
\]
D’après (6.2.11), \(\mathrm{H}^n(\mathrm{X},\mathrm{K})\) s’envoie donc sur \(\mathrm{H}^n(\mathrm{Y}_u,\mathrm{K})^{\pi_1(\mathrm{V},u)}
\supset
\mathrm{H}^n(\mathrm{Y}_u,\mathrm{K})^{\pi_1(\mathrm{U},u)}\).

\medskip
\textit{Théorème} \textbf{(6.2.13)}. — \textit{Soient \(\mathrm{X}\subset\mathbf{P}\) un schéma projectif, \(\eta\) la première classe de Chern de \(\mathcal{O}(1)\) et \(\mathrm{K}\in\operatorname{Ob}\mathrm{D}_c^b(\mathrm{X})\), un complexe potentiellement pur. On identifie \(\mathbf{Z}_{\ell}\) à \(\mathbf{Z}_{\ell}(1)\), sur \(k\), pour pouvoir regarder \(\eta\) comme étant dans \(\mathrm{H}^2(\mathrm{X},\mathbf{Z}_{\ell})\).}

\textit{Soit \(n\) un entier, et supposons que \(\mathrm{K}\) et \(\mathrm{DK}[-2n]\) vérifient la condition suivante : pour tout \(i\), la dimension du support du faisceau \(\mathscr{H}^i\) est \(\leq n-i\) (on pose par convention \(\dim\varnothing=-\infty\)).}

\textit{Alors, pour tout \(i\geq0\), le cup-produit itéré}
\[
\eta^i:\mathrm{H}^{n-i}(\mathrm{X},\mathrm{K})
\to
\mathrm{H}^{n+i}(\mathrm{X},\mathrm{K})
\]
\textit{est un isomorphisme.}

\medskip
La preuve est parallèle à celle de (4.1.1). Elle procède par récurrence sur la dimension de \(\mathrm{X}\).

Si \(\mathrm{X}\) est de dimension \(0\), et que \(\mathscr{H}^i\mathrm{K}\neq0\), l’hypothèse donne \(0\leq n-i\) et \(0\leq n-(-i+2n)=i-n\), soit \(i=n\). Le groupe \(\mathrm{H}^i\mathrm{K}\) est donc nul pour \(i\neq n\), et l’assertion en résulte.

Si \(\mathrm{Y}\) est une section hyperplane assez générale de \(\mathrm{X}\), la restriction de \(\mathrm{K}\) à \(\mathrm{Y}\) vérifie encore les hypothèses, avec \(n\) remplacé par \(n-1\). Ceci permet comme en (4.1.1) de ne considérer que l’application :
\[
\eta:\mathrm{H}^{n-1}(\mathrm{X},\mathrm{K})
\to
\mathrm{H}^{n-1}(\mathrm{Y},\mathrm{K})
\to
\mathrm{H}^{n+1}(\mathrm{X},\mathrm{K}).
\]
\Page{251}
La seconde flèche est transposée du morphisme de restriction
\[
\mathrm{H}^{n-1}(\mathrm{X},\mathrm{DK}[-2n])
\to
\mathrm{H}^{n-1}(\mathrm{Y},\mathrm{DK}[-2n]).
\]
Il s’agit donc de vérifier que la dualité parfaite entre \(\mathrm{H}^{n-1}(\mathrm{Y},\mathrm{K})\) et \(\mathrm{H}^{n-1}(\mathrm{Y},\mathrm{DK}[-2n])\) (où \(\mathrm{DK}[-2n]=(\mathrm{D}(\mathrm{K}|\mathrm{Y}))[-2(n-1)]\)) induit une dualité parfaite entre les parties fixes ((6.2.12)). Cela se vérifie par (4.1.4), la représentation de \(\pi_1(\mathrm{U},u)\) sur \(\mathrm{H}^*(\mathrm{Y}_u,\mathrm{K})\) étant semi-simple ((6.2.6) et (3.4.1) (iii)).

\medskip
\textit{Ajouté sur épreuves :} Soit \(\mathrm{X}\) une variété algébrique complexe, qu’on supposera irréductible et normale. On sait que la cohomologie intermédiaire de Mac Pherson et Goresky de \(\mathrm{X}\) peut se définir comme l’hypercohomologie de \(\mathrm{X}\) à coefficients dans un certain complexe de faisceaux \(\mathrm{I}_{\mathrm{X}}\). La construction de \(\mathrm{I}_{\mathrm{X}}\) admet un analogue en cohomologie étale, et ce sur un corps de base quelconque. O. Gabber vient de prouver la pureté de \(\mathrm{I}_{\mathrm{X}}\). Son théorème permet d’appliquer les résultats de ce numéro à la cohomologie de Mac Pherson - Goresky.

\section*{BIBLIOGRAPHIE}

{\small
\noindent\hangindent=2em
[1] P. \textsc{Deligne}, La conjecture de Weil, I, \textit{Publ. Math. IHES}, \textbf{43} (1974), 273-308.

\noindent\hangindent=2em
[2] P. \textsc{Deligne}, Théorie de Hodge, I, \textit{Actes ICM}, Nice, Gauthier-Villars, 1970, t. I, 425-430; II, \textit{Publ. Math. IHES}, \textbf{40} (1971), 5-58; III, \textit{Publ. Math. IHES}, \textbf{44} (1974), 5-77.

\noindent\hangindent=2em
[3] P. \textsc{Deligne}, Poids dans la cohomologie des variétés algébriques, \textit{Actes ICM}, Vancouver, 1974, 79-85.

\noindent\hangindent=2em
[4] P. \textsc{Deligne}, Formes modulaires et représentations \(\ell\)-adiques, \textit{Sém. Bourb.} 355 (févr. 1969), in \textit{LN}, \textbf{179} (Springer Verlag).

\noindent\hangindent=2em
[5] P. \textsc{Deligne}, Les constantes des équations fonctionnelles des fonctions L, in \textit{Proc. Antwerpen conf.}, vol. 2, \textit{LN}, \textbf{349} (Springer Verlag), 501-595.

\noindent\hangindent=2em
[6] P. \textsc{Deligne}, P. \textsc{Griffiths}, J. \textsc{Morgan} et D. \textsc{Sullivan}, Real homotopy theory of Kähler manifolds, \textit{Inv. Math.}, \textbf{29} (1975), 245-274.

\noindent\hangindent=2em
[7] N. \textsc{Katz} and W. \textsc{Messing}, Some consequences of the Riemann Hypothesis for Varieties over Finite Fields, \textit{Inv. Math.}, \textbf{23} (1974), 73-77.

\noindent\hangindent=2em
[8] E. \textsc{Miller}, De Rham cohomology with arbitrary coefficients, \textit{Topology}, \textbf{17} (2) (1978), 193-203.

\noindent\hangindent=2em
[9] J. \textsc{Morgan}, The algebraic topology of smooth algebraic manifolds, \textit{Publ. Math. IHES}, \textbf{48} (1978), 137-204.

\noindent\hangindent=2em
[10] J.-P. \textsc{Serre}, \textit{Corps locaux}, Publ. Math. Nancago, Hermann, 1962.

\noindent\hangindent=2em
[11] J.-P. \textsc{Serre}, \textit{Abelian \(\ell\)-adic representations and elliptic curves}, Benjamin, 1968.

\noindent\hangindent=2em
[12] J. \textsc{Steenbrink}, Limits of Hodge structures, \textit{Inv. Math.}, \textbf{31} (1976), 229-257.

\noindent\hangindent=2em
[13] D. \textsc{Sullivan}, Infinitesimal calculations in topology, \textit{Publ. Math. IHES}, \textbf{47} (1977), 269-331.
\par}

\section*{SIGLES}

{\small
\noindent\hangindent=4em
\makebox[4em][l]{EGA}Éléments de géométrie algébrique, par A. Grothendieck et J. Dieudonné, parus aux \textit{Publ. Math. IHES} (I = 4, II = 8, III = 11, 17, IV = 20, 24, 28, 32).

\noindent\hangindent=4em
\makebox[4em][l]{SGA}\textit{Séminaire de géométrie algébrique du Bois-Marie.} La liste complète des SGA figure dans SGA5; nous n’avons dû citer que :

\noindent\hangindent=4em
\makebox[4em][l]{SGA1}Revêtements étales et groupe fondamental, par A. Grothendieck, \textit{Lecture Notes in Mathematics}, \textbf{224}, Springer-Verlag, 1971.

\Page{252}
\noindent\hangindent=4em
\makebox[4em][l]{SGA4}Théorie des topos et cohomologie étale des schémas, par M. Artin, A. Grothendieck, J.-L. Verdier, \textit{Lecture Notes in Math.}, \textbf{269}, \textbf{270}, \textbf{305}, Springer-Verlag, 1972-1973.

\noindent\hangindent=4em
\makebox[4em][l]{SGA4\(\frac{1}{2}\)}Cohomologie étale, par P. Deligne, \textit{Lecture Notes in Mathematics}, \textbf{569}, Springer-Verlag, 1977.

\noindent\hangindent=4em
\makebox[4em][l]{SGA5}Cohomologie \(\ell\)-adique et fonctions L, par A. Grothendieck, \textit{Lecture Notes in Math.}, \textbf{589}, Springer-Verlag, 1977.

\noindent\hangindent=4em
\makebox[4em][l]{SGA7}Groupes de monodromie en géométrie algébrique, I par A. Grothendieck, II par P. Deligne et N. Katz, \textit{Lecture Notes in Mathematics}, \textbf{288}, \textbf{340}, Springer-Verlag, 1972, 1973.
\par}

\bigskip
{\leftskip=2em
\noindent Institut des Hautes Études Scientifiques,\\
35, route de Chartres,\\
91440 Bures-sur-Yvette,\\
France.
\par}

\bigskip
\begin{flushright}
\textit{Manuscrit reçu le 1\textsuperscript{\textup{er}} juin 1979.}
\end{flushright}

\end{document}
